
Skripsi
Selesai Skripsi 26 Mei



Apa yang bisa saya bantu hari ini?


Terbaru
Tabel perbaikan dokumen untuk pemeriksaan dosen
7 hari yang lalu
Memori
Hanya Anda
Purpose & context Ammar is a Statistics undergraduate at Universitas Indonesia (NPM 2206051582), supervised by Dr. Dra. Yekti Widyaningsih, M.Si. His thesis focuses on STATIS and Dynamic Biplot methods applied to spatio-temporal socio-economic analysis across Banten, DKI Jakarta, and West Java provinces. Current state Ammar is in the thesis revision stage, working to incorporate lecturer feedback into his thesis document. He has both an original and a revised version of his thesis, and has been reconciling handwritten revision notes from his supervisor with the actual changes made between versions. Claude helped produce a formatted revision tracking table (Word document) cataloguing substantive changes — including corrections to economic growth data, expanded methodology sections, new conceptual framework content, and style standardization — for submission or documentation purposes. Tools & resources Microsoft Word (for revision tracking deliverables) PDF/text-based thesis files for comparison Handwritten lecturer feedback notes (photo-based) Key learnings & principles The "Revisi yang Diminta" (requested revisions) column in the tracking table was reconstructed from observed textual changes and partially legible handwritten notes, and should be verified against the original handwritten feedback before formal use. Direct text comparison between document versions proved to be a reliable ground truth for identifying substantive revisions when source files are text-accessible.

Terakhir diperbarui 6 hari yang lalu

Petunjuk
Tambahkan instruksi untuk menyesuaikan respons Claude

File
17% kapasitas proyek terpakai
Mode pencarian

induk.pdf
7.001 baris

pdf



induk_rev.pdf
6.938 baris

pdf


IMG20260705WA0002.jpg

induk.pdf
Penyorotan sintaks telah dinonaktifkan karena ukuran kode.
UNIVERSITAS INDONESIA
PENDEKATAN STATIS DAN DYNAMIC BIPLOT UNTUK ANALISIS
SPASIO-TEMPORAL SOSIO-EKONOMI (STUDI KASUS: PROVINSI BANTEN,
DKI JAKARTA, DAN JAWA BARAT 2020–2024)
SKRIPSI
AMMAR HANAFI
2206051582
FAKULTAS MATEMATIKA DAN ILMU PENGETAHUAN ALAM
PROGRAM STUDI SARJANA STATISTIKA
DEPOK
JUNI 2026

UNIVERSITAS INDONESIA
PENDEKATAN STATIS DAN DYNAMIC BIPLOT UNTUK ANALISIS
SPASIO-TEMPORAL SOSIO-EKONOMI (STUDI KASUS: PROVINSI BANTEN,
DKI JAKARTA, DAN JAWA BARAT 2020–2024)
SKRIPSI
Diajukan sebagai salah satu syarat untuk memperoleh gelar
Sarjana Statistika
AMMAR HANAFI
2206051582
FAKULTAS MATEMATIKA DAN ILMU PENGETAHUAN ALAM
PROGRAM STUDI SARJANA STATISTIKA
DEPOK
JUNI 2026

HALAMAN PERSETUJUAN
Judul : Pendekatan STATIS dan Dynamic Biplot untuk Analisis SpasioTemporal Sosio-Ekonomi (Studi Kasus: Provinsi Banten, DKI
Jakarta, dan Jawa Barat 2020–2024)
Nama : Ammar Hanafi
NPM : 2206051582
Laporan Skripsi ini telah diperiksa dan disetujui.
... Juni 2026
Dr. Dra. Yekti Widyaningsih, M.Si.
Pembimbing Skripsi
iii

HALAMAN PERNYATAAN ORISINALITAS
Skripsi ini adalah hasil karya saya sendiri,
dan semua sumber baik yang dikutip maupun dirujuk
telah saya nyatakan dengan benar.
Nama : Ammar Hanafi
NPM : 2206051582
Tanda Tangan
Tanggal : ... Juni 2026
iv

HALAMAN PENGESAHAN
Skripsi ini diajukan oleh :
Nama : Ammar Hanafi
NPM : 2206051582
Program Studi : Sarjana Statistika
Judul Skripsi : Pendekatan STATIS dan Dynamic Biplot untuk Analisis Spasio-Temporal Sosio-Ekonomi (Studi Kasus:
Provinsi Banten, DKI Jakarta, dan Jawa Barat 2020–
2024)
Telah berhasil dipertahankan di hadapan Dewan Penguji dan diterima sebagai bagian
persyaratan yang diperlukan untuk memperoleh gelar Sarjana Statistika pada Program
Studi Sarjana Statistika, Fakultas Matematika dan Ilmu Pengetahuan Alam, Universitas
Indonesia.
DEWAN PENGUJI
Pembimbing I : Dr. Dra. Yekti Widyaningsih, M.Si. ( )
Penguji I : Sarini Abdullah, S.Si., M.Stats., Ph.D. ( )
Penguji II : Gianinna Ardaneswari, S.Si., M.Si. ( )
Ditetapkan di : Depok
Tanggal : ... Juni 2026
v

HALAMAN PERNYATAAN PENGGUNAAN KECERDASAN ARTIFISIAL
Saya menyatakan bahwa dalam penyusunan Tugas Akhir ini, saya menggunakan
teknologi Kecerdasan Buatan dengan platform Antigravity dan VS Code sebagai IDE
nya untuk tujuan membantu dalam penulisan Tugas Akhir ini dan saya bertanggung
jawab penuh atas keakuratan, orisinalitas, dan integritas seluruh isi Tugas Akhir ini.
Nama : Ammar Hanafi
NPM : 2206051582
Tanda Tangan :
Tanggal : ... Juni 2026
vi

KATA PENGANTAR
Assalamu’alaikum Warahmatullahi Wabarakatuh
Alhamdulillah. Puji syukur kehadirat Allah SWT karena atas segala rahmat dan
karunia-Nya, penulis dapat menjalani perkuliahan hingga menyelesaikan penulisan skripsi
ini. Penulisan skripsi ini dilakukan dalam rangka memenuhi salah satu syarat untuk memperoleh gelar Sarjana Statistika pada Program Studi S1 Statistika, Fakultas Matematika
dan Ilmu Pengetahuan Alam, Universitas Indonesia. Penulis menyadari bahwa selama
menjalani perkuliahan hingga penyusunan skripsi ini, tidak terlepas dari segala bantuan,
dukungan, bimbingan, serta doa dari orang-orang di sekitar. Oleh karena itu, penulis
mengucapkan terima kasih kepada:
1. Dr. Dra. Yekti Widyaningsih, M.Si. selaku Dosen Pembimbing skripsi penulis
yang telah memberikan waktu, tenaga, dan ilmu untuk membimbing penulis dalam
penyusunan skripsi ini sehingga selesai pada waktunya. Terima kasih atas segala
pengertian, dukungan, nasihat, dan doa yang diberikan sehingga penulis dapat
menyelesaikan skripsi ini.
2. Sarini Abdullah, S.Si., M.Stats., Ph.D. selaku Ketua Departemen Matematika FMIPA
UI sekaligus Dosen Penguji I pada sidang skripsi penulis. Terima kasih atas segala
bantuan, perhatian, bimbingan akademik, serta saran dan arahan yang diberikan
sehingga skripsi ini menjadi lebih baik.
3. Gianinna Ardaneswari, S.Si., M.Si. selaku Ketua Program Studi S1 Statistika sekaligus Dosen Penguji II pada sidang skripsi penulis. Terima kasih atas segala saran,
arahan, dan bimbingan akademik kepada penulis selama masa perkuliahan hingga
skripsi ini terselesaikan.
4. Seluruh keluarga Departemen Matematika FMIPA UI yang telah memberikan banyak
ilmu yang bermanfaat bagi penulis. Terima kasih juga kepada para staf yang telah
berkontribusi dalam mendukung kelancaran perkuliahan.
5. Kedua orang tua penulis yang sangat penulis sayangi. Terima kasih atas segala
pengertian, dukungan, dan doa yang telah diberikan sehingga penulis dapat menjalani
perkuliahan hingga selesai.
6. Teman-teman seperjuangan angkatan 2022 serta seluruh pihak yang pernah penulis
kenal. Terima kasih atas dukungan dan bantuan yang begitu besar kepada penulis.
Penulis sangat bersyukur bisa bertemu dan mengenal kalian. Kontribusi kalian membuat penulis dapat menjalani perkuliahan dengan lebih mudah dan menyenangkan.
7. Seluruh pihak yang tanpa mengurangi rasa hormat, tidak bisa disebutkan satu per
satu.
vii

Semoga Allah SWT senantiasa memberikan kesehatan, kebahagiaan, keberkahan, serta
umur yang panjang. Penulis berharap skripsi ini dapat memberi manfaat untuk penelitianpenelitian selanjutnya. Akhir kata, penulis mengucapkan terima kasih kepada semua pihak
yang turut membantu dalam penelitian ini.
Wassalamu’alaikum Warahmatullahi Wabarakatuh
Depok, ... Juni 2026
Ammar Hanafi
viii

HALAMAN PERNYATAAN PERSETUJUAN PUBLIKASI TUGAS AKHIR
UNTUK KEPENTINGAN AKADEMIS
Sebagai sivitas akademik Universitas Indonesia, saya yang bertanda tangan di bawah ini:
Nama : Ammar Hanafi
NPM : 2206051582
Program Studi : Sarjana Statistika
Fakultas : Matematika dan Ilmu Pengetahuan Alam
Jenis Karya : Skripsi
demi pengembangan ilmu pengetahuan, menyetujui untuk memberikan kepada Universitas
Indonesia Hak Bebas Royalti Noneksklusif (Non-exclusive Royalty Free Right) atas
karya ilmiah saya yang berjudul:
Pendekatan STATIS dan Dynamic Biplot untuk Analisis Spasio-Temporal Sosio-Ekonomi
(Studi Kasus: Provinsi Banten, DKI Jakarta, dan Jawa Barat 2020–2024)
beserta perangkat yang ada (jika diperlukan). Dengan Hak Bebas Royalti Noneksklusif ini,
Universitas Indonesia berhak menyimpan, mengalihmedia/formatkan, mengelola dalam
bentuk pangkalan data (data base), merawat, dan memublikasikan tugas akhir saya selama
tetap mencantumkan nama saya sebagai penulis/pencipta dan sebagai pemilik Hak Cipta.
Demikian pernyataan ini saya buat dengan sebenarnya.
Dibuat di : Kota Depok
Pada tanggal : ... Juni 2026
Yang menyatakan
(Ammar Hanafi)
ix

ABSTRAK
Nama : Ammar Hanafi
Program Studi : Sarjana Statistika
Judul : Pendekatan STATIS dan Dynamic Biplot untuk Analisis
Spasio-Temporal Sosio-Ekonomi (Studi Kasus: Provinsi
Banten, DKI Jakarta, dan Jawa Barat 2020–2024)
Wilayah strategis Banten, DKI Jakarta, dan Jawa Barat sebagai pusat ekonomi nasional
mengalami perubahan struktural kompleks selama 2020–2024 akibat pandemi COVID19 dan fase pemulihan yang tidak merata antar-kabupaten/kota. Pendekatan multivariat
konvensional memiliki keterbatasan dalam menangkap interaksi simultan antara dimensi
spasial, indikator sosial-ekonomi, dan pergeseran temporal secara terpadu. Penelitian ini
bertujuan mengevaluasi stabilitas longitudinal dan memetakan trajektori spasio-temporal
11 indikator pembangunan strategis yang meliputi aspek ekonomi (PDRB), kesejahteraan
(Kemiskinan, Garis Kemiskinan, Pangan), pembangunan manusia (Pendidikan, Kesehatan),
ketenagakerjaan, serta kesetaraan (Gender, Gini Ratio) di 41 kabupaten/kota pada ketiga
provinsi tersebut.
Metode STATIS (Structuration des Tableaux a Trois Indices de la Statistique ` atau Strukturisasi Tabel dengan Tiga Indeks dalam Statistik) diterapkan melalui tahapan: (1) standardisasi Z-score terbobot demografis, (2) evaluasi inter-struktur menggunakan koefisien
RV , (3) pembentukan struktur kompromi (Wc) melalui pembobotan spektral optimal sebagaimana didefinisikan oleh Lavit et al. (1994), (4) validasi signifikansi via uji permutasi
Monte Carlo (B = 1000), dan (5) visualisasi lintasan (trajectory) menggunakan Dynamic
Biplot. Pengelompokan lanjutan diterapkan melalui algoritma K-Means (k = 3) pada
koordinat kompromi yang direduksi via Principal Component Analysis (PCA).
Analisis inter-struktur menunjukkan konsistensi yang sangat tinggi dengan seluruh koefisien RV melampaui 0,90 (p < 0,001), membuktikan bahwa hierarki relatif antar-wilayah
tetap stabil meskipun terjadi fluktuasi absolut akibat pandemi. Struktur kompromi mampu
menjelaskan 71,88% total variansi melalui dua dimensi utama: Dimensi 1 merepresentasikan ”Kapasitas Sosio-Ekonomi” (berkorelasi kuat dengan PDRB dan IPM), sedangkan
Dimensi 2 mencerminkan ”Resiliensi Sosial dan Kesehatan” (didominasi Angka Harapan
Hidup). Dynamic Biplot mengidentifikasi tiga pola trajektori: pemulihan berbentuk V pada
pusat metropolitan, trajektori stagnan di kabupaten selatan Banten dan Jawa Barat, serta
pertumbuhan progresif di wilayah penyangga. Pengelompokan K-Means menghasilkan
klaster ”Maju,” ”Berkembang,” dan ”Tertinggal” yang memiliki keselarasan tinggi dengan
kategori normatif IPM UNDP, mengindikasikan bahwa kesenjangan regional bersumber
dari faktor struktural yang persisten.
Penelitian ini dibatasi pada tiga provinsi di Jawa bagian Barat dengan 11 indikator
terpilih dan rentang waktu lima tahun (2020–2024), sehingga generalisasi ke wilayah
lain memerlukan verifikasi lebih lanjut. Kebaruan penelitian ini terletak pada penerapan
pertama metode STATIS untuk analisis dinamika pembangunan regional pascapandemi di
konteks Jawa bagian Barat, yang diintegrasikan dengan sistem dashboard web interaktif
berbasis Python mencakup uji diagnostik (KMO, Bartlett), visualisasi klaster spasiotemporal, serta ekspor otomatis ke format LaTeX dan Excel sebagai instrumen analitik
bagi peneliti maupun pemangku kebijakan.
x Universitas Indonesia

Kata kunci: Analisis Data Tiga Arah, Data Longitudinal, K-Means Clustering, PCA,
Visualisasi Interaktif.
xi Universitas Indonesia

EXTENDED ABSTRACT (VERSI BAHASA INDONESIA)
Skripsi Sarjana
Pendekatan STATIS dan Dynamic Biplot untuk Analisis
Spasio-Temporal Sosio-Ekonomi (Studi Kasus: Provinsi Banten, DKI
Jakarta, dan Jawa Barat 2020–2024)
Nama Mahasiswa: Ammar Hanafi
NPM: 2206051582
Program Studi: Program Studi Sarjana Statistika
Wilayah strategis Banten, DKI Jakarta, dan Jawa Barat merupakan pusat aktivitas
ekonomi nasional di Indonesia yang mengalami variasi struktural signifikan selama periode 2020–2024 sebagai konsekuensi dari pandemi global dan fase pemulihan setelahnya.
Pendekatan multivariat tradisional sering kali menghadapi keterbatasan dalam menangkap
interaksi simultan antara unit spasial, indikator sosial, dan pergeseran temporal dalam
satu kerangka kerja yang terpadu. Penelitian ini menjawab tantangan tersebut dengan
menerapkan metode STATIS (Structuration des Tableaux a Trois Indices de la Statistique ` ),
sebuah kerangka analisis data tiga-arah yang kokoh, untuk menyintesis struktur kompromi
stabil yang merepresentasikan pola pembangunan bersama di 41 kabupaten dan kota.
Tujuan penelitian ini adalah untuk mengevaluasi stabilitas longitudinal dari 11 indikator
pembangunan strategis dan memetakan trajektori spasio-temporal untuk membedakan
resiliensi regional. Metodologi penelitian melibatkan standardisasi Z-score terbobot menggunakan bobot demografis, evaluasi inter-struktur melalui koefisien RV , dan pembentukan
representasi kompromi (Wc) melalui pembobotan spektral optimal. Uji permutasi Monte
Carlo (B = 1000) digunakan untuk memvalidasi signifikansi struktur yang ditemukan,
sebagaimana diilustrasikan dalam kerangka skematik pada Bab 3.
Analisis inter-struktur menunjukkan tingkat konsistensi yang sangat tinggi pada konfigurasi spasial selama periode lima tahun, dengan seluruh koefisien RV melebihi 0,90
(p < 0,05). Hal ini menunjukkan bahwa meskipun nilai aktual indikator pembangunan
mengalami fluktuasi, hierarki relatif dan hubungan struktural antar-wilayah, antar-indikator,
dan lintas tahun tetap bersifat persisten, sebagaimana dirangkum pada Tabel 1. Struktur
kompromi yang dihasilkan mampu menjelaskan sebesar 71,88% dari total variansi melalui
dua dimensi utama (Gambar 1). Dimensi 1 merepresentasikan sumbu utama “Taraf Kemajuan Pembangunan Manusia dan Ekonomi,” yang menunjukkan korelasi positif kuat
dengan PDRB ADHK, IPM, dan Rata-rata Lama Sekolah, serta korelasi negatif kuat dengan Persentase Penduduk Miskin. Dimensi 2 menangkap “Tekanan Pasar Tenaga Kerja,”
yang terutama dipengaruhi oleh variasi Tingkat Pengangguran Terbuka (loading +0,86)
dan Tingkat Partisipasi Angkatan Kerja (loading −0,49).
xii Universitas Indonesia

Tabel 1: Stabilitas Inter-struktur (Koefisien RV antar Tahun)
Tahun 2020 2021 2022 2023 2024
2020 1.0000 0.9555 0.9429 0.9533 0.9545
2021 0.9555 1.0000 0.9534 0.9487 0.9396
2022 0.9429 0.9534 1.0000 0.9608 0.9497
2023 0.9533 0.9487 0.9608 1.0000 0.9827
2024 0.9545 0.9396 0.9497 0.9827 1.0000
Proyeksi data tahunan ke dalam ruang kompromi melalui Biplot Dinamis mengungkapkan tiga pola trajektori regional yang berbeda (Gambar 2). Pusat-pusat metropolitan berbasis jasa, khususnya Kota Tangerang Selatan dan seluruh kota DKI Jakarta, menunjukkan
pola pemulihan berbentuk V yang khas, di mana penurunan posisi pada 2020–2021 diikuti
oleh peningkatan struktural yang kuat; Kota Tangerang Selatan mencatat kenaikan PC1
terbesar (+0,85 poin) sepanjang 2020–2024. Sebaliknya, kabupaten-kabupaten di bagian
selatan Banten dan Jawa Barat, termasuk Kabupaten Pangandaran (−1,37 poin) dan Garut
(−0,61 poin), menunjukkan trajektori stagnan bahkan menurun, yang mengindikasikan
bahwa kesenjangan pembangunan melebar selama periode pengamatan. Beberapa wilayah
penyangga seperti Kota Sukabumi menunjukkan pola pertumbuhan progresif dengan
keberhasilan naik klaster pada 2023.
Gambar 1: Biplot Kompromi Lintas Tahun: Hubungan Struktural antar Wilayah, antar
Indikator, dan antar Wilayah dan Indikator
Pengelompokan lanjut menggunakan algoritma K-Means (k = 3) pada koordinat
kompromi berhasil mempartisi 41 wilayah menjadi tiga kelompok: Klaster 1 “Tertinggal”
(12 wilayah, didominasi Banten Selatan dan Jawa Barat pedalaman selatan), Klaster 2
xiii Universitas Indonesia

“Berkembang” (16 wilayah, koridor industri dan penyangga), serta Klaster 3 “Maju” (13
wilayah metropolitan). Matriks transisi temporal mengonfirmasi stabilitas sempurna pada
2020–2021, diikuti restrukturisasi terbesar pada 2023–2024 dengan 14 wilayah berpindah
dari Klaster 2 ke Klaster 1 dan 4 wilayah dari Klaster 3 ke Klaster 2, sementara sepuluh
wilayah metropolitan inti tetap bertahan di Klaster 3 sepanjang seluruh periode. Validasi eksternal terhadap kategori normatif IPM UNDP mengonfirmasi keselarasan struktural yang
tinggi, di mana 11 dari 13 anggota Klaster 3 berkategori IPM “Sangat Tinggi”. Temuan
ini mengindikasikan bahwa kesenjangan regional di kawasan Banten–DKI Jakarta–Jawa
Barat didorong oleh faktor struktural yang mendalam, sehingga memerlukan intervensi
kebijakan yang terarah pada wilayah yang tertinggal, bukan strategi pembangunan yang
seragam. Gambar 2 memperlihatkan pergeseran kondisi sosial ekonomi kabupaten/kota
dari tahun ke tahun (2020–2024) yang mencerminkan perbedaan trajektori antar ketiga
kelompok tersebut.
Gambar 2: Trajektori Spasio-temporal: Pergeseran Kondisi Sosial Ekonomi
Kabupaten/Kota dari Tahun ke Tahun (2020–2024)
Sebagai kesimpulan, metode STATIS memberikan perspektif longitudinal yang komprehensif dan melampaui pendekatan analisis statis tradisional dalam memantau dinamika
regional. Stabilitas struktural dari representasi kompromi membuktikan bahwa penggerak
utama ketimpangan regional bersifat persisten meskipun terjadi pandemi global. Uji permutasi Monte Carlo (B = 1000) memvalidasi bahwa struktur yang ditemukan signifikan
secara statistik (p < 0,05) dan bukan merupakan hasil fluktuasi acak. Aplikasi dashboard
interaktif yang dikembangkan menggunakan framework Dash/Plotly menyediakan sarana
eksplorasi visual dinamis bagi peneliti dan pemangku kebijakan untuk memantau trajektori
xiv Universitas Indonesia

pembangunan wilayah secara mandiri dan berbasis data.
Kata Kunci: Analisis Data Tiga Arah, Data Longitudinal, K-Means Clustering, PCA,
Visualisasi Interaktif.
xv Universitas Indonesia

EXTENDED ABSTRACT
Undergraduate Thesis
STATIC and Dynamic Biplot Approaches for Socio-Economic
Spatio-Temporal Analysis (Case Studies: Province of Banten, DKI
Jakarta, and West Java 2020–2024)
Student Name : Ammar Hanafi
Student ID Number (NPM) : 2206051582
Study Program : Undergraduate Study Program of Statistics
The strategic regions of Banten, DKI Jakarta, and West Java constitute the core national
economic activity in Indonesia, which underwent significant structural variations during the
2020–2024 period as a consequence of the global pandemic and the subsequent recovery
phase. Traditional multivariate approaches often encounter limitations in capturing the
simultaneous interactions between spatial units, social indicators, and temporal shifts
within a unified framework. This research addresses this gap by applying the STATIS
(Structuration des Tableaux a Trois Indices de la Statistique ` ) method, a robust three-way
data analysis framework, to synthesize a stable compromise structure that represents
the common development patterns across 41 regencies and cities. The objective is to
evaluate the longitudinal stability of 11 strategic development indicators and map the
spatio-temporal trajectories to differentiate regional resilience. The methodology involves
weighted Z-score standardization using demographic weights, inter-structure evaluation
through RV coefficients, and the construction of a compromise representation (Wc) via
optimal spectral weighting. A Monte Carlo permutation test (B = 1000) is employed
to validate the significance of the discovered structure, as illustrated in the schematic
framework in Chapter 3.
Analysis of the inter-structure reveals an exceptionally high degree of consistency in
the spatial configurations across the five-year period, with all RV coefficients exceeding
0.90 (p < 0.05). This demonstrates that while the absolute magnitudes of development
indicators fluctuated, the underlying relative hierarchy and structural relationships between
regions, between indicators, and across years remained remarkably persistent, as summarized in Table 2. The resulting compromise structure accounts for 71.88% of the total
variance through two principal dimensions (Figure 3). Dimension 1 represents the primary
axis of “Human Development and Economic Advancement,” showing strong positive
correlations with GRDP, HDI, and Mean Years of Schooling, alongside strong negative
correlations with Poverty Rate. Dimension 2 captures “Labor Market Pressure,” primarily
driven by the Open Unemployment Rate (loading +0.86) and the Labor Force Participation
Rate (loading −0.49).
xvi Universitas Indonesia

Tabel 2: Inter-structure Stability (Inter-year RV Coefficients)
Year 2020 2021 2022 2023 2024
2020 1.0000 0.9555 0.9429 0.9533 0.9545
2021 0.9555 1.0000 0.9534 0.9487 0.9396
2022 0.9429 0.9534 1.0000 0.9608 0.9497
2023 0.9533 0.9487 0.9608 1.0000 0.9827
2024 0.9545 0.9396 0.9497 0.9827 1.0000
The projection of yearly data into the compromise space through a Dynamic Biplot
reveals three distinct regional trajectory patterns (Figure 4). Service-based metropolitan
centers, most notably Kota Tangerang Selatan and all DKI Jakarta cities, exhibited a
characteristic V-shape recovery pattern, where an initial decline in 2020–2021 was followed
by a robust structural rebound; Kota Tangerang Selatan recorded the largest PC1 increase
(+0.85 points) over the 2020–2024 period. In contrast, regencies in the southern parts of
Banten and West Java, including Kabupaten Pangandaran (−1.37 points) and Garut (−0.61
points), displayed stagnant or declining trajectories, indicating that regional development
gaps widened rather than narrowed. Several buffer-zone cities such as Kota Sukabumi
demonstrated progressive growth patterns, successfully advancing to a higher cluster in
2023.
Gambar 3: Compromise Biplot (Across Years): Structural Relationship between Regions,
between Indicators, and between Regions and Indicators
Advanced clustering using the K-Means algorithm (k = 3) on the compromise coordinates successfully partitioned 41 regions into three groups: Cluster 1 “Underdeveloped”
(12 regions, predominantly southern Banten and inland southern West Java), Cluster 2
xvii Universitas Indonesia

“Developing” (16 regions, industrial corridors and buffer zones), and Cluster 3 “Advanced”
(13 metropolitan regions). The temporal transition matrix confirmed perfect stability in
2020–2021, followed by the largest restructuring in 2023–2024, where 14 regions moved
from Cluster 2 to Cluster 1 and 4 regions from Cluster 3 to Cluster 2, while ten core
metropolitan regions remained in Cluster 3 throughout the entire period. External validation against normative UNDP HDI categories confirmed a high degree of structural
alignment, with 11 of 13 Cluster 3 members carrying the “Very High” HDI classification.
These findings indicate that regional disparities in the Banten–DKI Jakarta–West Java area
are driven by deep-seated structural factors, necessitating targeted policy interventions for
underdeveloped regions rather than uniform development strategies. Figure 4 illustrates the
movement of socio-economic conditions of regencies/cities from year to year (2020–2024),
reflecting the divergent trajectory patterns across these three clusters.
Gambar 4: Spatio-temporal Trajectories: Movement of Socio-economic Conditions of
Regencies/Cities from Year to Year (2020–2024)
In conclusion, the STATIS method provides a comprehensive longitudinal perspective that outperforms traditional static approaches in monitoring regional dynamics. The
structural stability of the compromise representation demonstrates that the core drivers of
regional inequality remained persistent despite the global pandemic. A Monte Carlo permutation test (B = 1000) validated that the discovered structure is statistically significant
(p < 0.05) and not attributable to random fluctuation. The interactive dashboard application developed using the Dash/Plotly framework provides a dynamic visual exploration
tool for researchers and policymakers to independently monitor regional development
xviii Universitas Indonesia

trajectories and design evidence-based strategies for equitable regional development.
Keywords: Interactive Visualization, K-Means Clustering, Longitudinal Data, PCA, ThreeWay Data Analysis.
xix Universitas Indonesia

DAFTAR ISI
LEMBAR PERSETUJUAN . . . . . . . . . . . . . . . . . . . . . . . . . . . . . iii
HALAMAN PERNYATAAN ORISINALITAS . . . . . . . . . . . . . . . . . . iv
HALAMAN PENGESAHAN . . . . . . . . . . . . . . . . . . . . . . . . . . . . v
HALAMAN PERNYATAAN PENGGUNAAN KECERDASAN ARTIFISIAL . vi
KATA PENGANTAR . . . . . . . . . . . . . . . . . . . . . . . . . . . . . . . . vii
HALAMAN PERNYATAAN PERSETUJUAN PUBLIKASI . . . . . . . . . . . ix
ABSTRAK . . . . . . . . . . . . . . . . . . . . . . . . . . . . . . . . . . . . . . x
EXTENDED ABSTRACT (BAHASA INDONESIA) . . . . . . . . . . . . . . . xii
EXTENDED ABSTRACT (INDONESIA) . . . . . . . . . . . . . . . . . . . . . xii
EXTENDED ABSTRACT . . . . . . . . . . . . . . . . . . . . . . . . . . . . . . xvi
DAFTAR ISI . . . . . . . . . . . . . . . . . . . . . . . . . . . . . . . . . . . . . xx
DAFTAR TABEL . . . . . . . . . . . . . . . . . . . . . . . . . . . . . . . . . . xxv
DAFTAR GAMBAR . . . . . . . . . . . . . . . . . . . . . . . . . . . . . . . . .xxviii
DAFTAR LAMPIRAN . . . . . . . . . . . . . . . . . . . . . . . . . . . . . . . xxx
BAB 1 PENDAHULUAN . . . . . . . . . . . . . . . . . . . . . . . . . . . . . 2
1.1 Latar Belakang . . . . . . . . . . . . . . . . . . . . . . . . . . . . . . . 2
1.2 Rumusan Masalah . . . . . . . . . . . . . . . . . . . . . . . . . . . . . . 5
1.3 Tujuan Penelitian . . . . . . . . . . . . . . . . . . . . . . . . . . . . . . 6
1.4 Manfaat Penelitian . . . . . . . . . . . . . . . . . . . . . . . . . . . . . 7
1.5 Ruang Lingkup Penelitian . . . . . . . . . . . . . . . . . . . . . . . . . . 7
xx Universitas Indonesia

BAB 2 LANDASAN TEORI . . . . . . . . . . . . . . . . . . . . . . . . . . . 8
2.1 Indikator Kesejahteraan Rakyat . . . . . . . . . . . . . . . . . . . . . . . 8
2.1.1 Indeks Pembangunan Manusia (IPM) . . . . . . . . . . . . . . . 8
2.1.2 Produk Domestik Regional Bruto (PDRB ADHK) . . . . . . . . . 9
2.1.3 Persentase Penduduk Miskin (PPM) . . . . . . . . . . . . . . . . 9
2.1.4 Garis Kemiskinan (GK) . . . . . . . . . . . . . . . . . . . . . . . 9
2.1.5 Prevalensi Ketidakcukupan Konsumsi Pangan (PKKP) . . . . . . 9
2.1.6 Tingkat Pengangguran Terbuka (TPT) . . . . . . . . . . . . . . . 10
2.1.7 Tingkat Partisipasi Angkatan Kerja (TPAK) . . . . . . . . . . . . 10
2.1.8 Rata-Rata Lama Sekolah (RLS) . . . . . . . . . . . . . . . . . . 10
2.1.9 Angka Harapan Hidup Laki-laki . . . . . . . . . . . . . . . . . . 10
2.1.10 Angka Harapan Hidup Perempuan . . . . . . . . . . . . . . . . . 11
2.1.11 Indeks Pembangunan Gender (IPG) . . . . . . . . . . . . . . . . 11
2.1.12 Gini Ratio (GR) . . . . . . . . . . . . . . . . . . . . . . . . . . . 11
2.1.13 Jumlah Penduduk (JP) . . . . . . . . . . . . . . . . . . . . . . . 11
2.2 Tinjauan Penelitian Terdahulu . . . . . . . . . . . . . . . . . . . . . . . . 12
2.3 Analisis Data Tiga-Arah (Three-Way Data) . . . . . . . . . . . . . . . . . 12
2.4 Konsep Spasio-Temporal dalam Analisis Multivariat . . . . . . . . . . . . 13
2.5 Standardisasi Data . . . . . . . . . . . . . . . . . . . . . . . . . . . . . . 14
2.6 Statistika Deskriptif . . . . . . . . . . . . . . . . . . . . . . . . . . . . . 15
2.6.1 Ukuran Pemusatan . . . . . . . . . . . . . . . . . . . . . . . . . 15
2.6.2 Ukuran Penyebaran . . . . . . . . . . . . . . . . . . . . . . . . . 16
2.6.3 Ukuran Bentuk Distribusi . . . . . . . . . . . . . . . . . . . . . . 17
2.7 Pemilihan Indikator dalam Analisis Multivariat . . . . . . . . . . . . . . 17
2.7.1 Uji Kaiser-Meyer-Olkin (KMO) . . . . . . . . . . . . . . . . . . 18
2.7.2 Uji Bartlett tentang Spherisitas (Bartlett’s Test of Sphericity) . . . 19
2.8 Dasar Aljabar Linier dan Dekomposisi Matriks . . . . . . . . . . . . . . 20
2.9 Korelasi Pearson . . . . . . . . . . . . . . . . . . . . . . . . . . . . . . . 22
2.10 Principal Component Analysis (PCA) dan Biplot . . . . . . . . . . . . . 23
2.11 Metode STATIS . . . . . . . . . . . . . . . . . . . . . . . . . . . . . . . 26
2.11.1 Tahap 1: Intra-struktur . . . . . . . . . . . . . . . . . . . . . . . 26
2.11.2 Tahap 2: Inter-struktur . . . . . . . . . . . . . . . . . . . . . . . 27
2.11.3 Tahap 3: Analisis Kompromi . . . . . . . . . . . . . . . . . . . . 29
2.11.4 Tahap 4: Ekstraksi Koordinat dan Relasi Transisi . . . . . . . . . 29
2.11.5 Klasifikasi Struktur: Analisis Transversal dan Longitudinal . . . . 32
2.11.6 Uji Signifikansi Stabilitas Struktur (Permutasi Monte Carlo) . . . 32
xxi Universitas Indonesia

2.11.7 Model Vektor Struktur dan Kerangka Inferensial STATIS . . . . . 34
2.12 Dynamic Biplot dan Visualisasi Trajektori . . . . . . . . . . . . . . . . . 37
2.12.1 Proyeksi Lintasan (Trajectory) . . . . . . . . . . . . . . . . . . . 37
2.12.2 Jarak Euclidean antar Trajektori . . . . . . . . . . . . . . . . . . 38
2.12.3 Rotasi Baseline . . . . . . . . . . . . . . . . . . . . . . . . . . . 38
2.13 Metode Pengelompokan (Clustering) . . . . . . . . . . . . . . . . . . . . 39
2.13.1 K-Means Clustering . . . . . . . . . . . . . . . . . . . . . . . . . 39
2.13.2 Penentuan Jumlah klaster Optimal (Metode Elbow) . . . . . . . . 39
2.13.3 Koefisien Siluet (Silhouette Score) . . . . . . . . . . . . . . . . . 40
2.13.4 Agglomerative Hierarchical Clustering (Dendrogram) . . . . . . . 40
2.13.5 Matriks Transisi Temporal . . . . . . . . . . . . . . . . . . . . . 41
2.13.6 Penyelarasan Label Klaster Lintas Waktu . . . . . . . . . . . . . 41
2.13.7 Kategorisasi dan Perankingan Cluster . . . . . . . . . . . . . . . 42
2.14 Uji Autokorelasi Spasial (Global Moran’s I) . . . . . . . . . . . . . . . . 43
2.14.1 Matriks Pembobot Spasial . . . . . . . . . . . . . . . . . . . . . 43
2.14.2 Analisis Autokorelasi Spasial (Global Moran’s I) . . . . . . . . . 43
2.15 Visualisasi Interaktif menggunakan Plotly dan Dash . . . . . . . . . . . . 44
2.16 Kerangka Berpikir . . . . . . . . . . . . . . . . . . . . . . . . . . . . . . 45
BAB 3 METODOLOGI PENELITIAN . . . . . . . . . . . . . . . . . . . . . 47
3.1 Statistika Deskriptif . . . . . . . . . . . . . . . . . . . . . . . . . . . . . 47
3.2 Jenis Penelitian . . . . . . . . . . . . . . . . . . . . . . . . . . . . . . . 48
3.3 Desain Penelitian . . . . . . . . . . . . . . . . . . . . . . . . . . . . . . 48
3.4 Populasi dan Sampel . . . . . . . . . . . . . . . . . . . . . . . . . . . . 50
3.5 Teknik Pengumpulan Data . . . . . . . . . . . . . . . . . . . . . . . . . 50
3.5.1 Pra-pemrosesan Data, Penskalaan Visual, dan Parameter Analisis . 51
3.6 Validasi Instrumen . . . . . . . . . . . . . . . . . . . . . . . . . . . . . . 53
3.7 Teknik Analisis Data . . . . . . . . . . . . . . . . . . . . . . . . . . . . 55
3.7.1 Prosedur Komputasi dan Ilustrasi Numerik STATIS . . . . . . . . 55
3.7.1.1 Tahap 0: Pra-pemrosesan dan Bobot Demografis . . . . 56
3.7.1.2 Percabangan Standardisasi: STATIS vs. Dynamic Biplot
(PCA per-Tahun) . . . . . . . . . . . . . . . . . . . . . 58
3.7.1.3 Tahap 1: Standardisasi Data STATIS (Zk) . . . . . . . . 59
3.7.1.4 Tahap 2: Matriks Intra-struktur (Wk) . . . . . . . . . . 61
3.7.1.5 Tahap 3: Analisis Inter-struktur (Koefisien RV) . . . . . 62
xxii Universitas Indonesia

3.7.1.6 Tahap 4: Penentuan Bobot Kompromi (αk) – Dua Skema
Pembobotan . . . . . . . . . . . . . . . . . . . . . . . 63
3.7.1.7 Tahap 5: Sintesis Matriks Kompromi – Skenario 1 vs
Skenario 2 . . . . . . . . . . . . . . . . . . . . . . . . 64
3.7.1.8 Tahap 6: Proyeksi Kompromi – Perbandingan Skenario
1 dan Skenario 2 . . . . . . . . . . . . . . . . . . . . . 67
3.7.1.9 Tahap 7: Uji Signifikansi dan Clustering . . . . . . . . 75
3.7.1.10 Tahap 7a: Komputasi Vektor Struktur dan Analisis
Transversal-Longitudinal . . . . . . . . . . . . . . . . . 77
3.7.1.11 Tahap 8: Analisis Perbandingan Klaster (9 Jenis) . . . . 79
3.7.2 Analisis Autokorelasi Spasial (Global Moran’s I) . . . . . . . . . 82
BAB 4 HASIL DAN PEMBAHASAN . . . . . . . . . . . . . . . . . . . . . . 83
4.1 Eksplorasi Data Awal . . . . . . . . . . . . . . . . . . . . . . . . . . . . 83
4.1.1 Statistika Deskriptif . . . . . . . . . . . . . . . . . . . . . . . . . 83
4.1.2 Distribusi Indikator (Boxplot) . . . . . . . . . . . . . . . . . . . 85
4.2 Uji Kelayakan Data Multivariat . . . . . . . . . . . . . . . . . . . . . . . 89
4.2.1 Hasil Uji KMO dan Bartlett . . . . . . . . . . . . . . . . . . . . 89
4.3 Analisis Inter-struktur . . . . . . . . . . . . . . . . . . . . . . . . . . . . 89
4.3.1 Koefisien RV . . . . . . . . . . . . . . . . . . . . . . . . . . . . 89
4.3.2 Vektor Bobot (Alpha) dan Representativitas Tahun . . . . . . . . 91
4.4 Analisis Struktur Kompromi . . . . . . . . . . . . . . . . . . . . . . . . 91
4.4.1 Dekomposisi Eigenvalue (Scree Plot) . . . . . . . . . . . . . . . 92
4.4.2 Peta Kompromi (Skor Objek) . . . . . . . . . . . . . . . . . . . . 93
4.4.3 Biplot Kompromi (Skor + Loading) . . . . . . . . . . . . . . . . 93
4.4.4 Kontribusi dan Kualitas Representasi Indikator . . . . . . . . . . 99
4.5 Uji Stabilitas Struktur Kompromi (Permutation Test) . . . . . . . . . . . 99
4.5.1 Prosedur dan Hasil Uji Permutasi . . . . . . . . . . . . . . . . . . 100
4.5.2 Interpretasi . . . . . . . . . . . . . . . . . . . . . . . . . . . . . 101
4.5.3 Analisis Vektor Struktur (β˜) dan Dominansi Interstruktur . . . . . 101
4.6 Analisis Dinamika Spasio-Temporal (Trajektori) . . . . . . . . . . . . . . 104
4.6.1 Gambaran Umum Trajektori 2020 hingga 2024 . . . . . . . . . . 104
4.6.2 Dinamika Biplot Tahunan . . . . . . . . . . . . . . . . . . . . . . 106
4.6.3 Dinamika Struktur Wilayah Tahunan . . . . . . . . . . . . . . . . 107
4.6.4 Pergeseran Centroid Tahunan . . . . . . . . . . . . . . . . . . . . 109
4.6.5 Jarak Pergeseran Per-Entitas Tahunan . . . . . . . . . . . . . . . 109
xxiii Universitas Indonesia

4.6.6 Dinamika Struktur Indikator Tahunan . . . . . . . . . . . . . . . 110
4.7 Analisis Klaster Spasio-Temporal . . . . . . . . . . . . . . . . . . . . . . 111
4.7.1 Penentuan Jumlah Klaster Optimal . . . . . . . . . . . . . . . . . 111
4.7.2 Pola Pengelompokan Tahunan . . . . . . . . . . . . . . . . . . . 113
4.7.3 Matriks Transisi Klaster STATIS . . . . . . . . . . . . . . . . . . 116
4.7.4 Klaster Trajektori Keseluruhan . . . . . . . . . . . . . . . . . . . 119
4.8 Analisis Komparatif Sembilan Jenis Perbandingan Klaster . . . . . . . . . 122
4.8.1 Perbandingan 1 — Matriks Transisi Klaster Biplot Antar-Tahun . 123
4.8.2 Perbandingan 3 — Kongruensi Klaster Biplot vs STATIS per Tahun124
4.8.3 Perbandingan 4 — Matriks Transisi Kategori IPM UNDP . . . . . 127
4.8.4 Perbandingan 5 — Validasi Eksternal Klaster Biplot terhadap Kategori IPM . . . . . . . . . . . . . . . . . . . . . . . . . . . . . . 129
4.8.5 Perbandingan 6 — Validasi Eksternal Klaster STATIS Tahunan
terhadap Kategori IPM . . . . . . . . . . . . . . . . . . . . . . . 132
4.8.6 Perbandingan 7 — Konsistensi Klaster Kompromi terhadap Klaster
Biplot per Tahun . . . . . . . . . . . . . . . . . . . . . . . . . . 134
4.8.7 Perbandingan 8 — Konsistensi Klaster Kompromi terhadap Klaster
STATIS Tahunan . . . . . . . . . . . . . . . . . . . . . . . . . . 136
4.8.8 Perbandingan 9 — Validasi Eksternal Final Klaster Kompromi
terhadap Kategori IPM . . . . . . . . . . . . . . . . . . . . . . . 138
4.9 Uji Autokorelasi Spasial (Global Moran’s I) . . . . . . . . . . . . . . . . 141
4.10 Analisis Hotspot (Local Moran’s I) . . . . . . . . . . . . . . . . . . . . . 141
4.11 Ringkasan Temuan dan Implikasi Kebijakan . . . . . . . . . . . . . . . . 142
4.11.1 Pembahasan Berdasarkan Tujuan Penelitian . . . . . . . . . . . . 142
4.11.2 Kontribusi terhadap Manfaat Penelitian . . . . . . . . . . . . . . 143
4.11.3 Implikasi Kebijakan . . . . . . . . . . . . . . . . . . . . . . . . . 144
BAB 5 KESIMPULAN DAN SARAN . . . . . . . . . . . . . . . . . . . . . . 145
5.1 Kesimpulan . . . . . . . . . . . . . . . . . . . . . . . . . . . . . . . . . 145
5.2 Saran . . . . . . . . . . . . . . . . . . . . . . . . . . . . . . . . . . . . . 147
DAFTAR PUSTAKA . . . . . . . . . . . . . . . . . . . . . . . . . . . . . . . . . 148
xxiv Universitas Indonesia

DAFTAR TABEL
Tabel 1 Stabilitas Inter-struktur (Koefisien RV antar Tahun) . . . . . . . . . xiii
Tabel 2 Inter-structure Stability (Inter-year RV Coefficients) . . . . . . . . . xvii
Tabel 2.1 Ringkasan Penelitian Terdahulu . . . . . . . . . . . . . . . . . . . . 12
Tabel 2.2 Kriteria Interpretasi Nilai KMO . . . . . . . . . . . . . . . . . . . . 19
Tabel 2.3 Rangkuman Dualitas Relasi Transisi STATIS . . . . . . . . . . . . 32
Tabel 3.1 Hasil Uji KMO dan Uji Bartlett per Tahun . . . . . . . . . . . . . . 54
Tabel 3.2 Data Mentah Sampel Ilustratif: 4 Wilayah × 3 Tahun × 5 Variabel . 56
Tabel 3.3 Statistika Deskriptif (tanpa Bobot) Tahun 2020 . . . . . . . . . . . . 56
Tabel 3.6 Statistika Deskriptif (dengan Bobot) Tahun 2020 (N = 4) . . . . . . 57
Tabel 3.9 Matriks Korelasi Terbobot Indikator Tahun 2020 (S2020) . . . . . . . 61
Tabel 3.12 Matriks Interstruktur H W (Identik Kedua Skenario) . . . . . . . . . 63
Tabel 3.13 Koefisien RV Wilayah (Identik Kedua Skenario) . . . . . . . . . . . 63
Tabel 3.14 Bobot Alpha (Wilayah) . . . . . . . . . . . . . . . . . . . . . . . . 64
Tabel 3.15 Bobot Alpha (Indikator) . . . . . . . . . . . . . . . . . . . . . . . . 64
Tabel 3.16 Dekomposisi Nilai Eigen (scenario1) . . . . . . . . . . . . . . . . . 66
Tabel 3.17 Dekomposisi Nilai Eigen (scenario2) . . . . . . . . . . . . . . . . . 66
Tabel 3.18 Diagnostik Dominansi Model Rank-1 (Areia et al., 2008) . . . . . . 66
Tabel 3.19 Koordinat Skor Kompromi (scenario1) . . . . . . . . . . . . . . . . 68
Tabel 3.26 Loadings Kompromi (scenario1) . . . . . . . . . . . . . . . . . . . 72
Tabel 3.27 Koordinat Skor Kompromi (scenario2) . . . . . . . . . . . . . . . . 72
Tabel 3.28 Loadings Kompromi (scenario2) . . . . . . . . . . . . . . . . . . . 73
Tabel 3.35 Ringkasan Anggota Klaster Kompromi STATIS (Skenario 1) . . . . 77
Tabel 3.36 Matriks Transisi Klaster 2020 → 2021 . . . . . . . . . . . . . . . . 77
Tabel 3.37 Matriks Transisi Klaster 2021 → 2022 . . . . . . . . . . . . . . . . 77
Tabel 3.38 Komponen Vektor Struktur β˜ per Situasi Waktu . . . . . . . . . . . 78
Tabel 3.39 Komponen Vektor Struktur β˜
gk per Provinsi (Analisis Transversal) . 79
Tabel 3.40 Matriks Transisi Klaster Biplot (Unweighted) 2020 → 2021 . . . . . 81
Tabel 3.41 Nama Daerah Transisi Klaster Biplot 2020 → 2021 . . . . . . . . . 81
xxv Universitas Indonesia

Tabel 3.42 Nama Daerah Perpindahan Klaster 2020 → 2021 . . . . . . . . . . 81
Tabel 3.43 Matriks Perbandingan Klaster Biplot vs STATIS (Tahunan) (2020) . 81
Tabel 3.44 Nama Daerah Perbandingan Klaster Biplot vs STATIS (Tahunan)
(2020) . . . . . . . . . . . . . . . . . . . . . . . . . . . . . . . . . 82
Tabel 4.1 Statistika Deskriptif Indikator Pembangunan Manusia Tahun 2020–
2024 (N = 41) . . . . . . . . . . . . . . . . . . . . . . . . . . . . . 84
Tabel 4.2 Hasil Uji KMO dan Uji Bartlett per Tahun . . . . . . . . . . . . . . 89
Tabel 4.3 Koefisien RV Wilayah (Identik Kedua Skenario) . . . . . . . . . . . 90
Tabel 4.4 Bobot Alpha (Wilayah) . . . . . . . . . . . . . . . . . . . . . . . . 91
Tabel 4.5 Bobot Alpha (Indikator) . . . . . . . . . . . . . . . . . . . . . . . . 91
Tabel 4.6 Dekomposisi Nilai Eigen Dua Dimensi Pertama Matriks Kompromi 92
Tabel 4.7 Koordinat Skor Kompromi (scenario1) . . . . . . . . . . . . . . . . 96
Tabel 4.8 Koordinat Skor Kompromi (scenario2) . . . . . . . . . . . . . . . . 97
Tabel 4.9 Loadings Kompromi (scenario1) . . . . . . . . . . . . . . . . . . . 98
Tabel 4.10 Loadings Kompromi (scenario2) . . . . . . . . . . . . . . . . . . . 98
Tabel 4.11 Rangkuman Hasil Uji Permutasi . . . . . . . . . . . . . . . . . . . 101
Tabel 4.12 Komponen Vektor Struktur β˜ per Situasi Waktu . . . . . . . . . . . 102
Tabel 4.13 Diagnostik Dominansi Model Rank-1 (Areia et al., 2008) . . . . . . 102
Tabel 4.14 Komponen Vektor Struktur β˜
gk per Provinsi (Analisis Transversal) . 103
Tabel 4.15 Skor Intra-Struktur Wilayah Tahun 2020 hingga 2024 . . . . . . . . 108
Tabel 4.19 Skor Intra-Struktur Wilayah Tahun 2023 hingga 2024 . . . . . . . . 108
Tabel 4.22 Pergeseran Centroid Tahunan — STATIS (Skor Intra Fk) . . . . . . 109
Tabel 4.23 Loading Intra-Struktur Indikator Tahun 2020 hingga 2022 . . . . . . 110
Tabel 4.27 Loading Intra-Struktur Indikator Tahun 2023 hingga 2024 . . . . . . 110
Tabel 4.30 Matriks Transisi Klaster 2023 → 2024 . . . . . . . . . . . . . . . . 117
Tabel 4.31 Nama Daerah Perpindahan Klaster 2023 → 2024 . . . . . . . . . . 117
Tabel 4.32 Matriks Transisi Klaster 2023 → 2024 . . . . . . . . . . . . . . . . 117
Tabel 4.33 Nama Daerah Perpindahan Klaster 2023 → 2024 . . . . . . . . . . 118
Tabel 4.34 Ringkasan Anggota Klaster Kompromi STATIS (Skenario 1) . . . . 122
Tabel 4.35 Ringkasan Anggota Klaster Kompromi STATIS (Skenario 2) . . . . 122
Tabel 4.36 Matriks Transisi Klaster Biplot (Unweighted) 2023 → 2024 . . . . . 123
Tabel 4.37 Nama Daerah Transisi Klaster Biplot 2023 → 2024 . . . . . . . . . 123
Tabel 4.38 Matriks Perbandingan Klaster Biplot vs STATIS (Tahunan) (2024) . 124
Tabel 4.39 Nama Daerah Perbandingan Klaster Biplot vs STATIS (Tahunan)
(2024) . . . . . . . . . . . . . . . . . . . . . . . . . . . . . . . . . 125
xxvi Universitas Indonesia

Tabel 4.40 Matriks Perbandingan Klaster Biplot vs STATIS (vs Kompromi) (2024)125
Tabel 4.41 Nama Daerah Perbandingan Klaster Biplot vs STATIS (vs Kompromi) (2024) . . . . . . . . . . . . . . . . . . . . . . . . . . . . . 125
Tabel 4.42 Matriks Perbandingan Klaster Biplot vs STATIS (Tahunan) (2024) . 126
Tabel 4.43 Nama Daerah Perbandingan Klaster Biplot vs STATIS (Tahunan)
(2024) . . . . . . . . . . . . . . . . . . . . . . . . . . . . . . . . . 126
Tabel 4.44 Matriks Perbandingan Klaster Biplot vs STATIS (vs Kompromi) (2024)126
Tabel 4.45 Nama Daerah Perbandingan Klaster Biplot vs STATIS (vs Kompromi) (2024) . . . . . . . . . . . . . . . . . . . . . . . . . . . . . 127
Tabel 4.46 Matriks Transisi Kategori IPM (UNDP) Tahun 2023 ke 2024 . . . . 128
Tabel 4.47 Nama Daerah Matriks Transisi Kategori IPM (UNDP) 2023 → 2024 128
Tabel 4.48 Matriks Transisi Kategori IPM (UNDP) Tahun 2023 ke 2024 . . . . 128
Tabel 4.49 Nama Daerah Matriks Transisi Kategori IPM (UNDP) 2023 → 2024 129
Tabel 4.50 Perbandingan Klaster Biplot vs Kategori IPM (UNDP) (2024) . . . 130
Tabel 4.51 Nama Daerah Perbandingan Klaster Biplot vs Kategori IPM (UNDP)
(2024) . . . . . . . . . . . . . . . . . . . . . . . . . . . . . . . . . 130
Tabel 4.52 Perbandingan Klaster Biplot vs Kategori IPM (UNDP) (2024) . . . 130
Tabel 4.53 Nama Daerah Perbandingan Klaster Biplot vs Kategori IPM (UNDP)
(2024) . . . . . . . . . . . . . . . . . . . . . . . . . . . . . . . . . 131
Tabel 4.54 Perbandingan Klaster STATIS Tahunan vs Kategori IPM (UNDP)
(2024) . . . . . . . . . . . . . . . . . . . . . . . . . . . . . . . . . 132
Tabel 4.55 Nama Daerah Perbandingan Klaster STATIS Tahunan vs IPM
(UNDP) (2024) . . . . . . . . . . . . . . . . . . . . . . . . . . . . 132
Tabel 4.56 Perbandingan Klaster STATIS Tahunan vs Kategori IPM (UNDP)
(2024) . . . . . . . . . . . . . . . . . . . . . . . . . . . . . . . . . 133
Tabel 4.57 Nama Daerah Perbandingan Klaster STATIS Tahunan vs IPM
(UNDP) (2024) . . . . . . . . . . . . . . . . . . . . . . . . . . . . 133
Tabel 4.58 Matriks Perbandingan Klaster Biplot vs Kompromi STATIS (2024) . 134
Tabel 4.59 Nama Daerah Perbandingan Klaster Biplot vs Kompromi STATIS
(2024) . . . . . . . . . . . . . . . . . . . . . . . . . . . . . . . . . 134
Tabel 4.60 Matriks Perbandingan Klaster Biplot vs Kompromi STATIS (2024) . 135
Tabel 4.61 Nama Daerah Perbandingan Klaster Biplot vs Kompromi STATIS
(2024) . . . . . . . . . . . . . . . . . . . . . . . . . . . . . . . . . 135
Tabel 4.62 Perbandingan Klaster Kompromi STATIS vs Klaster STATIS Tahunan (2024) . . . . . . . . . . . . . . . . . . . . . . . . . . . . . . . 136
xxvii Universitas Indonesia

Tabel 4.63 Nama Daerah Perbandingan Klaster Kompromi STATIS vs STATIS
Tahunan (2024) . . . . . . . . . . . . . . . . . . . . . . . . . . . . 136
Tabel 4.64 Perbandingan Klaster Kompromi STATIS vs Klaster STATIS Tahunan (2024) . . . . . . . . . . . . . . . . . . . . . . . . . . . . . . . 136
Tabel 4.65 Nama Daerah Perbandingan Klaster Kompromi STATIS vs STATIS
Tahunan (2024) . . . . . . . . . . . . . . . . . . . . . . . . . . . . 137
Tabel 4.66 Perbandingan Klaster Kompromi STATIS vs Kategori IPM (UNDP)
(2024) . . . . . . . . . . . . . . . . . . . . . . . . . . . . . . . . . 138
Tabel 4.67 Nama Daerah Perbandingan Klaster Kompromi STATIS vs IPM
(UNDP) (2024) . . . . . . . . . . . . . . . . . . . . . . . . . . . . 138
Tabel 4.68 Perbandingan Klaster Kompromi STATIS vs Kategori IPM (UNDP)
(2024) . . . . . . . . . . . . . . . . . . . . . . . . . . . . . . . . . 139
Tabel 4.69 Nama Daerah Perbandingan Klaster Kompromi STATIS vs IPM
(UNDP) (2024) . . . . . . . . . . . . . . . . . . . . . . . . . . . . 139
DAFTAR GAMBAR
Gambar 1 Biplot Kompromi Lintas Tahun: Hubungan Struktural antar
Wilayah, antar Indikator, dan antar Wilayah dan Indikator . . . . xiii
Gambar 2 Trajektori Spasio-temporal: Pergeseran Kondisi Sosial Ekonomi
Kabupaten/Kota dari Tahun ke Tahun (2020–2024) . . . . . . . . xiv
Gambar 3 Compromise Biplot (Across Years): Structural Relationship between Regions, between Indicators, and between Regions and
Indicators . . . . . . . . . . . . . . . . . . . . . . . . . . . . . . xvii
Gambar 4 Spatio-temporal Trajectories: Movement of Socio-economic Conditions of Regencies/Cities from Year to Year (2020–2024) . . . . xviii
Gambar 2.1 Ilustrasi Komponen Biplot (Skor dan Loading) . . . . . . . . . . 25
Gambar 2.2 Flowchart Mekanisme Interaktivitas Dash-Plotly . . . . . . . . . 45
Gambar 3.1 Flowchart Alur Metodologi Penelitian . . . . . . . . . . . . . . . 49
Gambar 3.2 Heatmap Korelasi Pearson Indikator Data Sampel Tahun 2020 . . 62
Gambar 3.3 Visualisasi Heatmap Matriks RV Data Sampel . . . . . . . . . . . 63
xxviii Universitas Indonesia

Gambar 3.4 Scree Plot Analisis Kompromi Data Sampel (Skenario 1) . . . . . 67
Gambar 3.5 Compromise Biplot Skenario 1 (Host = Wilayah): Skor Wilayah
dan Loading Indikator . . . . . . . . . . . . . . . . . . . . . . . 68
Gambar 3.6 Compromise Biplot Skenario 2 (Host = Indikator): Skor Wilayah
dan Loading Indikator . . . . . . . . . . . . . . . . . . . . . . . 72
Gambar 3.7 Trajektori Skenario 1 (Host = Wilayah) . . . . . . . . . . . . . . 75
Gambar 3.8 Trajektori Skenario 2 (Host = Indikator) . . . . . . . . . . . . . . 75
Gambar 4.1 Boxplot Tingkat Pengangguran Terbuka (TPT) Tahun 2020–2024 85
Gambar 4.2 Boxplot Tingkat Partisipasi Angkatan Kerja (TPAK) Tahun 2020–
2024 . . . . . . . . . . . . . . . . . . . . . . . . . . . . . . . . . 85
Gambar 4.3 Boxplot Angka Harapan Hidup Laki-Laki (AHH L) Tahun 2020–
2024 . . . . . . . . . . . . . . . . . . . . . . . . . . . . . . . . . 85
Gambar 4.4 Boxplot Angka Harapan Hidup Perempuan (AHH P) Tahun 2020–
2024 . . . . . . . . . . . . . . . . . . . . . . . . . . . . . . . . . 85
Gambar 4.5 Boxplot Gini Ratio (GR) Tahun 2020–2024 . . . . . . . . . . . . 86
Gambar 4.6 Boxplot Garis Kemiskinan (GK) Tahun 2020–2024 . . . . . . . . 86
Gambar 4.7 Boxplot Indeks Pembangunan Manusia (IPM) Tahun 2020–2024 . 86
Gambar 4.8 Boxplot Indeks Pemberdayaan Gender (IPG) Tahun 2020–2024 . 86
Gambar 4.9 Boxplot PDRB Atas Dasar Harga Konstan (PDRB ADHK) Tahun
2020–2024 . . . . . . . . . . . . . . . . . . . . . . . . . . . . . 87
Gambar 4.10 Boxplot Prevalensi Ketidakcukupan Konsumsi Pangan (PKKP)
Tahun 2020–2024 . . . . . . . . . . . . . . . . . . . . . . . . . . 87
Gambar 4.11 Boxplot Rata-Rata Lama Sekolah (RLS) Tahun 2020–2024 . . . . 87
Gambar 4.12 Boxplot Persentase Penduduk Miskin (PPM) Tahun 2020–2024 . 87
Gambar 4.13 Heatmap Korelasi Pearson Antar-Indikator Tahun 2020 (41 Kabupaten/Kota) . . . . . . . . . . . . . . . . . . . . . . . . . . . . . 88
Gambar 4.14 Heatmap Koefisien RV Antar-Tahun (2020 hingga 2024) . . . . . 90
Gambar 4.15 Scree Plot Nilai Eigen Matriks Kompromi . . . . . . . . . . . . . 92
Gambar 4.16 Peta Skor Kompromi 41 Kabupaten/Kota (2020 hingga 2024) . . 93
Gambar 4.17 Biplot Kompromi Skenario 1 (Host = Wilayah) . . . . . . . . . . 94
Gambar 4.18 Biplot Kompromi Skenario 2 (Host = Indikator) . . . . . . . . . . 94
Gambar 4.19 Kontribusi Indikator Dim 1 – Skenario 1 . . . . . . . . . . . . . . 99
Gambar 4.20 Kontribusi Indikator Dim 1 – Skenario 2 . . . . . . . . . . . . . . 99
Gambar 4.21 Kualitas Representasi (cos2) – Skenario 1 . . . . . . . . . . . . . 99
Gambar 4.22 Kualitas Representasi (cos2) – Skenario 2 . . . . . . . . . . . . . 99
xxix Universitas Indonesia

Gambar 4.23 Distribusi Null Uji Permutasi vs. Nilai Eigen Observasi (λ
obs
1
) . . 100
Gambar 4.24 Lintasan Temporal Komponen Vektor Struktur β˜
k (2020 hingga
2024) . . . . . . . . . . . . . . . . . . . . . . . . . . . . . . . . 104
Gambar 4.25 Trajektori Pergeseran 41 Kabupaten/Kota pada Ruang Kompromi
(2020 hingga 2024) . . . . . . . . . . . . . . . . . . . . . . . . . 105
Gambar 4.26 Biplot Tahunan 2020 (Masa Awal Pandemi) . . . . . . . . . . . . 106
Gambar 4.31 Kurva Elbow dan Dendrogram Tahun 2020 . . . . . . . . . . . . 111
Gambar 4.36 Pengelompokan K-Means (k = 3) Tahun 2020: Biplot (kiri) dan
Peta (kanan) . . . . . . . . . . . . . . . . . . . . . . . . . . . . . 113
Gambar 4.41 Perbandingan Peta Klaster STATIS Periode 2023 vs 2024 (Kiri:
Tahun 2023; Kanan: Tahun 2024) . . . . . . . . . . . . . . . . . 118
Gambar 4.42 Scatter Plot Klaster Trajektori pada Ruang Kompromi . . . . . . 120
Gambar 4.43 Peta Persebaran Klaster Trajektori Keseluruhan . . . . . . . . . . 120
Gambar 4.44 Distribusi Proporsi Klaster Trajektori Keseluruhan . . . . . . . . 121
Gambar 4.45 Perbandingan Peta Klaster Biplot (Unweighted): Tahun 2023 (kiri)
vs 2024 (kanan) . . . . . . . . . . . . . . . . . . . . . . . . . . . 124
Gambar 4.46 Perbandingan Peta Klaster Biplot (kiri) vs STATIS (kanan) Tahun
2024 . . . . . . . . . . . . . . . . . . . . . . . . . . . . . . . . . 127
Gambar 4.47 Perbandingan Peta Kategori IPM UNDP: Tahun 2023 (kiri) vs
2024 (kanan) . . . . . . . . . . . . . . . . . . . . . . . . . . . . 129
Gambar 4.48 Perbandingan Peta Klaster Biplot (kiri) vs Kategori IPM UNDP
(kanan) Tahun 2024 . . . . . . . . . . . . . . . . . . . . . . . . . 131
Gambar 4.49 Perbandingan Peta Klaster STATIS Tahunan (kiri) vs Kategori
IPM UNDP (kanan) Tahun 2024 . . . . . . . . . . . . . . . . . . 133
Gambar 4.50 Perbandingan Peta Klaster Kompromi (kiri) vs Klaster Biplot
Tahun 2024 (kanan) . . . . . . . . . . . . . . . . . . . . . . . . . 135
Gambar 4.51 Perbandingan Peta Klaster Kompromi (kiri) vs Klaster STATIS
Tahun 2024 (kanan) . . . . . . . . . . . . . . . . . . . . . . . . . 137
Gambar 4.52 Perbandingan Peta Klaster Kompromi (kiri) vs Kategori IPM
UNDP Tahun 2024 (kanan) . . . . . . . . . . . . . . . . . . . . . 139
Gambar 1 Tampilan Antarmuka Aplikasi Dashboard Analisis SpasioTemporal Sosio-Ekonomi . . . . . . . . . . . . . . . . . . . . . . 167
xxx Universitas Indonesia

1
DAFTAR LAMPIRAN
Lampiran 1 Data Penelitian . . . . . . . . . . . . . . . . . . . . . . . . . . . . . . 154
Lampiran 2 Kode Program Aplikasi Dashboard . . . . . . . . . . . . . . . . . . . 167
Universitas Indonesia

BAB 1
PENDAHULUAN
1.1 Latar Belakang
Kawasan strategis yang meliputi Provinsi Banten, DKI Jakarta, dan Jawa Barat tidak
hanya berperan sebagai mesin pertumbuhan ekonomi nasional, tetapi juga menjadi arena dinamika pembangunan manusia yang sangat kompleks dalam lima tahun terakhir. Meskipun
wilayah ini memberikan kontribusi yang sangat dominan terhadap perekonomian nasional
dengan rata-rata pangsa Produk Domestik Regional Bruto (PDRB) sekitar 33,8 persen
dari PDRB Indonesia (Badan Pusat Statistik, 2025b), di mana keberhasilan pembangunannya tidak dapat lagi diukur hanya melalui kacamata indikator makro ekonomi semata.
Terdapat interaksi multidimensi yang melibatkan kualitas hidup yang tecermin dalam
Indeks Pembangunan Manusia (IPM), Indeks Pembangunan Gender (IPG), serta tantangan
kesejahteraan sosial seperti kemiskinan dan ketahanan pangan yang terus mengalami perubahan. Kompleksitas ini menuntut pembuktian empiris yang lebih mendalam untuk melihat
bagaimana indikator-indikator tersebut terdistribusi secara nyata di seluruh kabupaten/kota.
Secara empiris, lanskap pembangunan di ketiga provinsi ini menunjukkan polarisasi
yang tajam antara pusat-pusat pertumbuhan utama dan wilayah penyangga. Meskipun data
PDRB di kota-kota besar mencatat performa ekonomi yang kuat dengan nilai mencapai
532.606,87 miliar rupiah (Badan Pusat Statistik, 2025b), capaian tersebut tidak serta-merta
diikuti oleh pemerataan kesejahteraan. Hal ini terlihat dari variasi spasial yang mencolok
(striking spatial variation) pada indikator kemiskinan, di mana Persentase Penduduk
Miskin (PPM) terentang lebar dari 2,29 persen hingga mencapai 15,06 persen di wilayah
lainnya. Kesenjangan ini juga terkonfirmasi oleh indikator daya beli, terlihat dari perbedaan
Garis Kemiskinan (GK) dan Prevalensi Ketidakcukupan Konsumsi Pangan (PKKP) yang
bervariasi lebar antarwilayah (Badan Pusat Statistik, 2024). Lebih jauh lagi, perbedaan
kapabilitas manusia juga terlihat mencolok; Rata-rata Lama Sekolah (RLS) dan Angka
Harapan Hidup (AHH) masih menunjukkan perbedaan yang nyata. Indeks Pembangunan
Gender (IPG) serta Tingkat Pengangguran Terbuka (TPT) masih berada di kisaran 6,55
hingga 10,12 persen pada rentang tahun 2020 hingga 2024. Namun, persoalan ketimpangan
antarwilayah ini menjadi semakin rumit ketika diletakkan dalam konteks waktu yang penuh
perubahan.
Selain perbedaan wilayah, tantangan analisis meningkat akibat adanya dinamika temporal yang sangat fluktuatif sepanjang periode 2020 hingga 2024. Periode ini diawali
dengan peristiwa pandemi COVID-19 yang menekan indikator makroekonomi dan kesejahteraan secara serentak terlihat dari perubahan pertumbuhan ekonomi yang sempat
2 Universitas Indonesia

3
mencapai angka 1,88 persen pada tahun 2020 (Badan Pusat Statistik, 2025b) namun diikuti oleh pola pemulihan yang lajunya tidak seragam di setiap daerah pada tahun-tahun
berikutnya. Perubahan drastis dari masa pandemi COVID-19 menuju masa pemulihan ini
diduga kuat telah mengubah struktur hubungan antar-indikator secara fundamental; sebagai
contoh, pola korelasi antara pertumbuhan ekonomi (PDRB) dengan tingkat pengangguran
(TPT) saat krisis sangat mungkin berbeda dibandingkan saat kondisi normal kembali pulih
pada tahun 2023 hingga 2024. Oleh karena itu, pendekatan analisis multivariat klasik
yang mengasumsikan struktur hubungan indikator bersifat statis (tetap sepanjang waktu)
menjadi tidak lagi relevan untuk menganalisis fenomena pada periode yang penuh fluktuasi
ini.
Untuk mengatasi keterbatasan asumsi stabilitas tersebut, diperlukan pendekatan analisis data tiga-arah (three-way data analysis) yang mampu menangani struktur data kompleks
yang terdiri dari objek (kabupaten/kota), indikator (indikator sosio-ekonomi), dan situasi (waktu). Penelitian ini mengusulkan penggunaan metode STATIS (Structuration des
Tableaux a Trois Indices de la Statistique ` ), sebuah metode eksploratif yang diperkenalkan
oleh Lavit et al. (1994); Oliveira et al. (2024). Berbeda dengan Analisis Komponen
Utama (PCA) konvensional yang hanya bekerja pada data dua dimensi, STATIS dirancang
khusus untuk mengekstraksi informasi dari sekumpulan tabel data yang berubah seiring
waktu. Metode ini memiliki keunggulan dalam membentuk sebuah struktur kompromi
(compromise structure) yang merepresentasikan pola hubungan umum antar-indikator
selama periode pengamatan, sekaligus memungkinkan identifikasi seberapa jauh struktur
ekonomi setiap tahunnya menyimpang dari pola umum tersebut.
Meskipun metode STATIS mampu menghasilkan struktur kompromi secara matematis,
kompleksitas hasil analisis dari matriks berdimensi tinggi sering kali sulit dipahami jika
hanya disajikan dalam bentuk tabel numerik. Oleh karena itu, penelitian ini mengintegrasikan pendekatan STATIS dengan teknik visualisasi Dynamic Biplot. Pendekatan ini
memungkinkan proyeksi posisi objek (kabupaten/kota) dan indikator (indikator ekonomi)
ke dalam ruang berdimensi rendah, sekaligus menggambarkan pergeseran atau lintasan
(trajectory) perubahan posisi tersebut dari tahun ke tahun. Dengan memanfaatkan konsep
trajectory, pergeseran karakteristik suatu wilayah misalnya, sebuah kabupaten yang pada
tahun 2020 memiliki tingkat pengangguran tinggi namun kemudian bergeser mendekati
karakteristik wilayah dengan IPM tinggi pada tahun 2024 dapat dilacak secara visual.
Visualisasi lintasan ini menjadi instrumen penting untuk membedakan antara fluktuasi
jangka pendek akibat pandemi dan perubahan struktural jangka panjang.
Penerapan metode eksploratif multivariat seperti ini mengisi celah penting dalam literatur analisis regional di Indonesia yang selama ini didominasi oleh pendekatan inferensial
Universitas Indonesia

4
klasik. Studi-studi terdahulu umumnya berfokus pada pemodelan regresi data panel atau
model spasial ekonometrik untuk mengestimasi besaran pengaruh indikator independen
terhadap indikator dependen. Sebagai contoh, penelitian yang dilakukan oleh Nuswantari
(2021) menggunakan pendekatan Dynamic Spatial Durbin Model (SDM) untuk memodelkan kemiskinan, sedangkan Ahmad (2024) menerapkan pendekatan spasial yang serupa
untuk menganalisis fenomena penyebaran penyakit kusta di Pulau Jawa. Fokus utama
dari penelitian-penelitian tersebut adalah pengukuran besaran pengaruh peubah-peubah
penjelas (explanatory variables) terhadap peubah respon yang terbukti nyata secara statistik. Meskipun pendekatan tersebut berhasil mengidentifikasi faktor-faktor determinan,
metode inferensial ini memiliki keterbatasan dalam memvisualisasikan struktur hubungan
antarindikator secara simultan (inter-variable structure) dan bagaimana struktur tersebut berevolusi dari waktu ke waktu (spatio-temporal). Hingga saat ini, belum terdapat
penelitian yang secara spesifik menggunakan pendekatan eksploratif untuk memetakan dinamika struktur multivariat indikator ekonomi secara komprehensif tanpa membedakannya
menjadi indikator bebas dan terikat.
Kompleksitas informasi yang dihasilkan dari analisis spasio-temporal ini menuntut
media penyajian yang lebih adaptif dibandingkan laporan statis konvensional. Untuk itu,
penelitian ini tidak hanya berhenti pada tataran analisis teoritis, tetapi juga mengembangkan
sebuah aplikasi visualisasi interaktif berbasis web menggunakan bahasa pemrograman
Python dengan kerangka kerja (framework) Dash dan Plotly. Pengembangan dashboard ini
bertujuan untuk memfasilitasi Eksplorasi Data (Exploratory Data Analysis) secara mandiri
oleh pengguna. Fitur-fitur interaktif yang disematkan seperti seleksi tahun, filter wilayah,
hingga rotasi baseline memungkinkan para pengambil kebijakan untuk mengamati pola
interaksi multivariat secara lebih mendalam. Melalui antarmuka ini, pola-pola tersembunyi
(latent patterns) yang mungkin terlewatkan dalam tabel statistik standar dapat terdeteksi
dengan lebih intuitif.
Secara keseluruhan, penelitian ini diarahkan untuk menjawab kebutuhan akan
metode pemantauan pembangunan daerah yang lebih menyeluruh dan dinamis. Dengan
memadukan kemampuan analitis metode STATIS dalam memahami struktur kompromi
(compromise structure) lintas tahun dan kemampuan naratif visual dari Dynamic Biplot.
Penelitian ini diharapkan dapat memberikan kontribusi ganda: secara akademis memperkaya khazanah penerapan statistik multivariat pada data ekonomi regional, dan secara
praktis menyediakan prototipe alat bantu visual yang relevan bagi pemerintah daerah dalam
mengevaluasi dampak kebijakan pasca-pandemi. Identifikasi struktur hubungan yang stabil
maupun yang berubah selama lima tahun terakhir akan menjadi landasan penting dalam
merumuskan strategi pembangunan yang lebih tangguh di masa depan.
Universitas Indonesia

5
Pemilihan lokasi penelitian yang mencakup seluruh kabupaten/kota di Provinsi Banten,
DKI Jakarta, dan Jawa Barat didasarkan pada karakteristik interkonektivitas ekonomi yang
kuat di antara ketiga wilayah ini. Sebagai kesatuan wilayah pembangunan utama di Pulau
Jawa, ketiga provinsi ini memiliki keterkaitan rantai pasok, mobilitas tenaga kerja, dan arus
investasi yang sangat intensif. Permasalahan ekonomi yang terjadi di pusat pertumbuhan
utama (DKI Jakarta) diproyeksikan akan merambat (spillover) ke wilayah penyangga di
Banten dan Jawa Barat dengan skala yang berbeda-beda. Oleh karena itu, analisis yang
hanya membatasi lingkup pada wilayah metropolitan Jabodetabek berpotensi mengabaikan
dampak sistemik yang dirasakan oleh wilayah kabupaten non-metropolitan yang secara
administratif berada dalam satu kesatuan provinsi namun memiliki karakteristik ekonomi
yang jauh berbeda.
Untuk memahami kompleksitas pembangunan tersebut, penelitian ini menggunakan
dua belas indikator strategis yang merepresentasikan dimensi pembangunan secara
menyeluruh. Pemilihan indikator ini tidak dilakukan tanpa alasan, melainkan dipertimbangkan untuk mencakup aspek makroekonomi (PDRB ADHK), kesejahteraan material
(PPM, GK, PKKP), kualitas modal manusia (IPM, RLS, AHH Laki-laki dan Perempuan),
kesetaraan (IPG, Gini Ratio), serta partisipasi pasar kerja (TPT, TPAK). Kombinasi indikator ini memungkinkan analisis untuk tidak hanya pada indikator pertumbuhan konvensional,
namun juga aspek inklusivitas pembangunan. Sebagai contoh, analisis simultan antara Gini
Ratio dan PDRB memungkinkan deteksi dini terhadap fenomena pertumbuhan yang tidak
merata (exclusive growth), sementara interaksi antara TPAK dan IPG dapat mengungkap
hambatan struktural dalam partisipasi angkatan kerja berbasis gender.
Ketiadaan studi yang menganalisis dinamika interaksi kedua belas indikator ini secara
spasio-temporal pada periode kritis 2020 hingga 2024 menciptakan sebuah ”kotak hitam”
(black box) dalam pemahaman mengenai resiliensi ekonomi regional. Tanpa analisis
struktur kompromi dan visualisasi lintasan yang jelas, para pemangku kebijakan kehilangan
panduan empiris mengenai bagaimana sebenarnya pola pemulihan berlangsung di tingkat
lokal apakah kembali ke struktur lama atau membentuk ekuilibrium baru. Berdasarkan
urgensi untuk mengungkap fenomena tersebut secara statistik dan visual, penelitian ini
disusun untuk menjawab pertanyaan-pertanyaan mendasar mengenai evolusi struktur
ekonomi dan pembangunan di ketiga provinsi tersebut.
1.2 Rumusan Masalah
Berdasarkan latar belakang yang telah diuraikan, permasalahan yang akan dikaji dalam
penelitian ini dirumuskan sebagai berikut:
Universitas Indonesia

6
1. Bagaimana karakteristik struktur kompromi (compromise structure) yang merepresentasikan pola hubungan multivariat sebelas indikator pembangunan di seluruh
kabupaten/kota Provinsi Banten, DKI Jakarta, dan Jawa Barat menggunakan metode
STATIS untuk memahami pola umum lintas tahun?
2. Bagaimana menguji struktur hubungan antar-indikator tersebut secara statistik selama periode 2020 hingga 2024 menggunakan uji permutasi matriks inter-struktur
untuk membuktikan keberadaan dinamika temporal?
3. Bagaimana pola lintasan (trajectory) perubahan karakteristik ekonomi kabupaten/kota dari masa krisis pandemi (2020) hingga masa pemulihan (2024)
berdasarkan visualisasi Dynamic Biplot?
4. Bagaimana pola pengelompokan spasial tiap tahunnya serta dinamika transisi atau
perubahan kabupaten/kota antar klaster dari tahun ke tahun?
5. Bagaimana merancang aplikasi berbasis web menggunakan framework Dash dan
Plotly sebagai alat bantu eksplorasi data yang memungkinkan pengguna melakukan
pemantauan indikator pembangunan secara interaktif?
1.3 Tujuan Penelitian
Sejalan dengan rumusan masalah di atas, tujuan dari penelitian ini adalah:
1. Menganalisis karakteristik struktur kompromi (compromise structure) yang merepresentasikan pola hubungan multivariat sebelas indikator pembangunan di seluruh
kabupaten/kota Provinsi Banten, DKI Jakarta, dan Jawa Barat menggunakan metode
STATIS untuk memahami pola umum lintas tahun.
2. Menganalisis struktur hubungan antar-indikator tersebut secara statistik selama
periode 2020 hingga 2024 menggunakan uji permutasi matriks inter-struktur untuk
membuktikan keberadaan dinamika temporal.
3. Menganalisis pola lintasan (trajectory) perubahan karakteristik ekonomi kabupaten/kota dari masa krisis pandemi (2020) hingga masa pemulihan (2024)
berdasarkan hasil visualisasi Dynamic Biplot.
4. Mengidentifikasi pola pengelompokan spasial tiap tahunnya serta dinamika transisi
atau perubahan kabupaten/kota antar klaster dari tahun ke tahun.
Universitas Indonesia

7
5. Merancang aplikasi berbasis web menggunakan framework Dash dan Plotly sebagai
alat bantu eksplorasi data yang memungkinkan pengguna melakukan pemantauan
indikator pembangunan secara interaktif.
1.4 Manfaat Penelitian
Penelitian ini bermanfaat untuk:
1. Secara akademis memperkaya khazanah penerapan statistik multivariat terutama
metode STATIS dan Dynamic Biplot pada data ekonomi regional di Indonesia.
2. Secara praktis menyediakan prototipe alat bantu visual berbasis web yang relevan
bagi pemerintah daerah dalam mengevaluasi perkembangan indikator pembangunan
dan dampak kebijakan pasca-pandemi.
1.5 Ruang Lingkup Penelitian
Agar pembahasan penelitian ini lebih terarah dan fokus, penulis menetapkan batasanbatasan masalah sebagai berikut:
1. Indikator Penelitian: Analisis menggunakan 11 indikator terpilih yang bersumber
dari Badan Pusat Statistik (BPS), meliputi: Indeks Pembangunan Gender (IPG),
Persentase Penduduk Miskin (PPM), PDRB ADHK, Rata-rata Lama Sekolah (RLS),
Angka Harapan Hidup Laki-laki (AHH L) dan Perempuan (AHH P), Garis Kemiskinan (GK), Prevalensi Ketidakcukupan Konsumsi Pangan (PKKP), Tingkat Pengangguran Terbuka (TPT), Tingkat Partisipasi Angkatan Kerja (TPAK), dan Gini Ratio
(GR) untuk periode 2020 hingga 2024.
2. Aspek Spasial: Istilah ”spasial” dalam penelitian ini merujuk pada unit analisis berbasis wilayah (kabupaten/kota), bukan pada penggunaan matriks pembobot
spasial (spatial weight matrix) seperti pada Ekonometrika Spasial.
3. Fungsionalitas Aplikasi: Aplikasi berbasis web yang dikembangkan berfungsi
sebagai alat bantu eksplorasi visual (Visual Exploratory Tool) bagi pengguna, bukan
sebagai Sistem Pendukung Keputusan (Decision Support System) yang menghasilkan
rekomendasi kebijakan berbasis data secara otomatis.
Universitas Indonesia

BAB 2
LANDASAN TEORI
Bab ini berisi uraian mengenai teori-teori dasar yang menjadi fondasi metodologi
penelitian. Pembahasan dimulai dari konsep indikator kesejahteraan rakyat sebagai latar
belakang substansi indikator penelitian, kemudian dilanjutkan dengan struktur data tiga
arah dan konsep spasio-temporal, statistika deskriptif, standardisasi data, pemilihan indikator (Uji KMO dan Bartlett), fondasi aljabar linier beserta dekomposisi matriks (nilai
eigen, dekomposisi spektral, dan Singular Value Decomposition), korelasi Pearson, Principal Component Analysis (PCA) beserta Biplot, algoritma metode STATIS beserta uji
permutasi, Dynamic Biplot dan visualisasi trajektori, metode pengelompokan (clustering),
uji autokorelasi spasial Global Moran’s I, dan terakhir adalah media visualisasi interaktif
Plotly dan Dash.
2.1 Indikator Kesejahteraan Rakyat
Tujuan utama pembangunan suatu wilayah adalah mewujudkan kesejahteraan rakyatnya. Untuk mengukur keberhasilan pembangunan, dibutuhkan indikator-indikator yang
mampu mencerminkan taraf kesejahteraan secara menyeluruh. Indikator kesejahteraan
rakyat cukup luas dan kompleks sehingga hanya dapat dinilai melalui indikator-indikator
terukur dari berbagai bidang pembangunan, di antaranya bidang kesehatan, pendidikan,
ketenagakerjaan, taraf dan pola konsumsi, serta kemiskinan (Badan Pusat Statistik, 2025a).
Pemilihan indikator yang tepat dan terukur sangat menentukan kualitas analisis multivariat
yang akan dilakukan.
2.1.1 Indeks Pembangunan Manusia (IPM)
Indeks Pembangunan Manusia (IPM) merupakan indeks komposit yang mengukur
capaian pembangunan manusia berbasis tiga dimensi dasar: (1) umur panjang dan hidup
sehat, tercermin dari angka harapan hidup saat lahir; (2) pengetahuan, tercermin dari
rata-rata lama sekolah dan harapan lama sekolah; serta (3) standar hidup layak, tercermin
dari pengeluaran per kapita yang disesuaikan. IPM dihitung sebagai rata-rata geometrik
dari ketiga indeks dimensi tersebut dengan skala nilai 0 hingga 100. Semakin tinggi nilai
IPM, semakin baik capaian pembangunan manusia di suatu wilayah (Badan Pusat Statistik,
2025a).
8 Universitas Indonesia

9
2.1.2 Produk Domestik Regional Bruto (PDRB ADHK)
Produk Domestik Regional Bruto Atas Dasar Harga Konstan (PDRB ADHK) mencerminkan nilai tambah bruto barang dan jasa yang dihasilkan oleh seluruh unit ekonomi di
suatu wilayah dalam kurun waktu tertentu, dihitung menggunakan harga pada tahun dasar
tertentu (2010) sehingga pengaruh inflasi dieliminasi. PDRB ADHK digunakan sebagai
ukuran pertumbuhan ekonomi riil suatu wilayah. Satuan yang digunakan umumnya adalah
juta atau miliar rupiah. Wilayah dengan PDRB ADHK yang tinggi menunjukkan kapasitas
produksi ekonomi yang kuat (Badan Pusat Statistik, 2025a).
2.1.3 Persentase Penduduk Miskin (PPM)
Persentase penduduk miskin adalah persentase penduduk yang memiliki rata-rata pengeluaran per kapita per bulan di bawah Garis Kemiskinan. Garis Kemiskinan (GK) mencerminkan nilai rupiah pengeluaran minimum yang diperlukan seseorang untuk memenuhi kebutuhan pokok hidupnya selama sebulan, baik kebutuhan makanan maupun non-makanan.
Daerah yang lebih sejahtera pada umumnya memiliki persentase penduduk miskin yang
relatif kecil (Badan Pusat Statistik, 2025a; Widyaningsih, 2024).
2.1.4 Garis Kemiskinan (GK)
Garis Kemiskinan (GK) adalah nilai rupiah pengeluaran minimum yang diperlukan
oleh seseorang untuk dapat memenuhi kebutuhan pokok makanan setara 2.100 kilo kalori
per orang per hari dan kebutuhan pokok non-makanan yang mencakup perumahan, sandang, pendidikan, dan kesehatan. GK digunakan sebagai batas untuk menentukan status
kemiskinan seseorang dan bervariasi antarwilayah bergantung pada pola konsumsi dan
harga komoditas setempat (Badan Pusat Statistik, 2025a).
2.1.5 Prevalensi Ketidakcukupan Konsumsi Pangan (PKKP)
Prevalensi Ketidakcukupan Konsumsi Pangan (PKKP) adalah estimasi proporsi populasi yang mengonsumsi energi (kalori) di bawah standar kecukupan minimum untuk
menjalani kehidupan yang sehat dan aktif. Indikator ini mencerminkan dimensi ketahanan
pangan suatu wilayah dan merupakan salah satu dari target Tujuan Pembangunan Berkelanjutan (SDGs) yang dipantau secara reguler oleh Badan Pusat Statistik (Badan Pusat
Statistik, 2025a).
Universitas Indonesia

10
2.1.6 Tingkat Pengangguran Terbuka (TPT)
Tingkat Pengangguran Terbuka (TPT) adalah persentase jumlah pengangguran terhadap jumlah angkatan kerja. Pengangguran didefinisikan sebagai penduduk usia kerja
yang sedang mencari pekerjaan, mempersiapkan usaha, atau tidak mencari pekerjaan
karena merasa tidak mungkin mendapatkannya, termasuk mereka yang baru mendapat pekerjaan tetapi belum mulai bekerja. Angkatan kerja adalah penduduk dengan
usia lebih dari 15 tahun yang bekerja atau sedang mencari pekerjaan. TPT yang tinggi
mengindikasikan pemanfaatan tenaga kerja yang belum optimal (Badan Pusat Statistik,
2025a; Widyaningsih, 2024).
2.1.7 Tingkat Partisipasi Angkatan Kerja (TPAK)
Tingkat Partisipasi Angkatan Kerja (TPAK) adalah persentase penduduk usia kerja
(lebih dari 15 tahun) yang termasuk dalam angkatan kerja, baik yang sudah bekerja maupun
yang sedang mencari pekerjaan. TPAK mencerminkan proporsi penduduk usia produktif
yang aktif secara ekonomi. Nilai TPAK yang tinggi menandakan tingginya partisipasi
masyarakat dalam kegiatan ekonomi (Badan Pusat Statistik, 2025a).
2.1.8 Rata-Rata Lama Sekolah (RLS)
Rata-rata lama sekolah adalah rata-rata jumlah tahun yang telah ditempuh oleh penduduk dengan usia lebih dari 25 tahun dalam menjalani pendidikan formal. Indikator
ini dihitung untuk penduduk dengan usia lebih dari 25 tahun dengan asumsi bahwa pada
usia tersebut proses pendidikan pada umumnya telah selesai. Semakin tinggi nilai RLS,
semakin tinggi tingkat pendidikan rata-rata penduduk suatu wilayah yang pada gilirannya
mendukung peningkatan produktivitas dan kesejahteraan (Badan Pusat Statistik, 2025a;
Widyaningsih, 2024).
2.1.9 Angka Harapan Hidup Laki-laki
Angka Harapan Hidup (AHH) laki-laki adalah perkiraan rata-rata jumlah tahun hidup
yang akan dijalani oleh seorang bayi laki-laki yang baru lahir pada suatu tahun tertentu,
dengan asumsi bahwa pola mortalitas pada saat kelahiran tetap berlaku sepanjang hidupnya.
Indikator ini mencerminkan derajat kesehatan penduduk laki-laki serta kualitas layanan
kesehatan yang tersedia bagi mereka (Badan Pusat Statistik, 2025a; Widyaningsih, 2024).
Universitas Indonesia

11
2.1.10 Angka Harapan Hidup Perempuan
Angka Harapan Hidup (AHH) perempuan adalah perkiraan rata-rata jumlah tahun
hidup yang akan dijalani oleh seorang bayi perempuan yang baru lahir pada suatu tahun
tertentu. Umumnya, nilai AHH perempuan lebih tinggi dibandingkan laki-laki karena
perbedaan biologis dan pola perilaku kesehatan. AHH perempuan merupakan komponen
penting dalam penghitungan Indeks Pembangunan Manusia (IPM) dan Indeks Pembangunan Gender (IPG) (Badan Pusat Statistik, 2025a; Widyaningsih, 2024).
2.1.11 Indeks Pembangunan Gender (IPG)
Indeks Pembangunan Gender (IPG) adalah rasio capaian IPM perempuan terhadap
IPM laki-laki. IPG mengukur kesenjangan gender dalam dimensi kesehatan, pendidikan,
dan standar hidup. Nilai IPG sebesar 100 menunjukkan kesetaraan sempurna antara lakilaki dan perempuan, sedangkan nilai yang lebih rendah menunjukkan adanya kesenjangan
yang merugikan perempuan (Badan Pusat Statistik, 2025a).
2.1.12 Gini Ratio (GR)
Gini Ratio adalah ukuran tingkat ketimpangan distribusi pengeluaran atau pendapatan
penduduk. Nilai Gini Ratio berada di antara 0 dan 1: nilai 0 berarti pemerataan sempurna
(setiap orang memiliki pengeluaran yang sama), sedangkan nilai 1 berarti ketimpangan
sempurna (satu orang menguasai seluruh pengeluaran). Gini Ratio yang tinggi menandakan
tingginya kesenjangan ekonomi antar kelompok penduduk dalam suatu wilayah (Badan
Pusat Statistik, 2025a).
2.1.13 Jumlah Penduduk (JP)
Jumlah Penduduk (JP) merupakan banyaknya jiwa yang mendiami suatu wilayah
administratif pada suatu titik waktu tertentu, sebagaimana dicatat melalui Sensus Penduduk
maupun estimasi antarsensus yang diselenggarakan oleh Badan Pusat Statistik (Badan
Pusat Statistik, 2025a). Dalam penelitian ini, JP tidak digunakan sebagai indikator
analisis dalam metode STATIS, melainkan berperan sebagai basis pembentukan bobot
demografis (mi). Bobot ini mengintegrasikan ukuran populasi setiap wilayah ke dalam
analisis multivariat sehingga wilayah dengan penduduk lebih besar memberikan kontribusi
yang proporsional terhadap struktur rata-rata (kontribusi ini dinamakan kompromi), dan
wilayah berpenduduk kecil tidak mengubah hasil analisis multivariat secara berlebihan
(Lavit et al., 1994).
Universitas Indonesia

12
2.2 Tinjauan Penelitian Terdahulu
Penelitian mengenai analisis statistik terhadap data regional di Pulau Jawa telah banyak
dilakukan dengan menggunakan berbagai pendekatan. Bagian ini merangkum beberapa
penelitian relevan yang menjadi rujukan dalam pengembangan metodologi penelitian ini,
sebagaimana disajikan pada Tabel 2.1.
Tabel 2.1: Ringkasan Penelitian Terdahulu
Peneliti
(Tahun) Wilayah Metode Hasil Utama
Nuswantari
(2021)
Pulau Jawa
(2012-2018)
Dynamic Spatial
Panel (Model
Durbin Spasial)
Model terbaik menggunakan pembobot
dua tetangga terdekat. Angka melek
huruf dan fasilitas kesehatan
berpengaruh signifikan terhadap
kemiskinan.
Ahmad
(2024)
Pulau Jawa
(2017-2021)
Dynamic Spatial
Durbin Model
(DSDM)
Menemukan adanya pengaruh lag
temporal dan spasiotemporal yang
signifikan terhadap prevalensi kusta.
Widyaningsih
(2024)
Pulau Jawa
(2022)
Biclustering Plaid
Model vs
Cluster-Biplot
Plaid model menghasilkan kelompok
wilayah yang lebih homogen
dibandingkan metode biplot
konvensional dalam mengelompokkan
indikator kesejahteraan.
Nuswantari (2021) menunjukkan bahwa pemodelan di Pulau Jawa harus mempertimbangkan efek dinamis (waktu) dan dependensi spasial secara simultan. Hal ini diperkuat
oleh Ahmad (2024) yang menggunakan DSDM untuk menangkap interaksi antar-wilayah
lintas waktu. Di sisi lain, Widyaningsih (2024) memberikan perspektif mengenai pentingnya metode pengelompokan yang mampu menangani objek dan indikator secara
bersamaan untuk melihat karakteristik regional yang unik. Penelitian ini menggabungkan
aspek temporal dari data tiga-arah menggunakan metode STATIS dan aspek keterkaitan
spasial untuk memberikan analisis yang lebih komprehensif.
2.3 Analisis Data Tiga-Arah (Three-Way Data)
Dalam statistika multivariat klasik, umumnya data direpresentasikan dalam struktur
dua-arah (two-way structure), yaitu sebuah matriks X berukuran n × p, dengan n merepresentasikan jumlah objek dan p merepresentasikan jumlah indikator. Struktur ini memadai
untuk menganalisis kondisi statis pada satu titik waktu tertentu. Namun, dalam studi yang
melibatkan pemantauan berkala, struktur data menjadi lebih kompleks dengan hadirnya
dimensi ketiga, yaitu waktu atau situasi (Kroonenberg, 2008).
Universitas Indonesia

13
Struktur data seperti ini dikenal sebagai data tiga-arah (three-way data) atau data
longitudinal multivariat. Secara formal, data tiga-arah dapat didefinisikan sebagai berikut:
Definisi 2.1 (Data Tiga-Arah). Misalkan terdapat n objek penelitian, p indikator, dan K
situasi (waktu). Data tiga-arah adalah sekumpulan K buah matriks data {X1, X2, . . . , XK},
dengan setiap matriks Xk (untuk k = 1, . . . , K) berukuran n × p.
Elemen skalar dari data ini dinotasikan sebagai xijk, yang merepresentasikan nilai dari
indikator ke-j pada objek ke-i pada situasi ke-k. Keseluruhan himpunan matriks ini sering
diilustrasikan sebagai sebuah balok data (data cube) dengan dimensi n × p × K.
Relevansi pendekatan tiga-arah ini semakin menguat dalam analisis data regional,
di mana interaksi kompleks antara wilayah, waktu, dan indikator sosial-ekonomi sangat
penting untuk dipahami (Biegert et al., 2023). Metode STATIS, yang akan dibahas
selanjutnya, dirancang secara spesifik untuk mendapatkan informasi dari struktur tiga-arah
{Xk} ini tanpa harus merata-ratakan data terlebih dahulu ke dalam satu matriks tunggal
yang menghilangkan karakteristik temporal (Lavit et al., 1994).
2.4 Konsep Spasio-Temporal dalam Analisis Multivariat
Analisis spasio-temporal bertujuan untuk mendeskripsikan atau memodelkan variabilitas fenomena yang terjadi dalam domain ruang dan waktu secara simultan. Tantangan
utama dalam analisis ini adalah adanya dua ketergantungan sekaligus: ketergantungan
antar-lokasi (spasial) dan ketergantungan antar-waktu (temporal).
Pendekatan konvensional dalam ekonometrika biasanya melakukan pendekatan terhadap masalah ini dengan menggunakan regresi data panel atau model dinamis spasial
(Spatial Dynamic Panel Data) untuk mengetahui pengaruh antar-lokasi, dengan struktur hubungan antar-lokasi didefinisikan menggunakan matriks pembobot spasial (spatial
weight matrix, Wspasial) berbasis ketetanggaan geografis yang memiliki asumsi bahwa
interaksi antar-wilayah hanya terjadi melalui batas geografis fisik yang bersifat konstan
sepanjang waktu (Baltagi, 2021). Dalam pendekatan multivariat eksploratif seperti STATIS,
definisi “spasial” dipandang dari perspektif yang berbeda.
Konsep Spasial dalam STATIS Dalam kerangka STATIS, aspek spasial direpresentasikan melalui konfigurasi kemiripan antar-objek dalam ruang multivariat, bukan melalui
matriks kontiguitas geografis. Struktur “spasial” pada situasi ke-k didefinisikan oleh matriks produk skalar Wk = XkXT
k
(berukuran n × n) yang merangkum “peta” kemiripan
antar-objek berdasarkan profil indikatornya.
Universitas Indonesia

14
Evolusi spasio-temporal dalam pendekatan ini difokuskan pada pengamatan perubahan
konfigurasi matriks Wk dari waktu ke waktu, untuk menjawab pertanyaan apakah pola
kemiripan antar-objek berubah selama periode pengamatan (Lavit et al., 1994).
2.5 Standardisasi Data
Perbedaan nilai yang besar antar-indikator akibat perbedaan satuan pengukuran (misalnya rupiah, persen, dan tahun) dapat menyebabkan bias dalam analisis multivariat sehingga
perlu menyamakan skala variabilitasnya. Karena 41 wilayah yang diamati merupakan populasi utuh dengan ukuran penduduk yang sangat bervariasi, standardisasi dalam kerangka
STATIS dilakukan menggunakan rata-rata dan simpangan baku terbobot (µ dan σ) agar
pusat data konsisten dengan D yang dijelaskan pada 2.3 (Lavit et al., 1994). Untuk itu,
terlebih dahulu didefinisikan bobot demografis dan matriks bobot berikut menggunakan
jumlah penduduk sebagai dasar pembobotan. Hal ini memastikan bahwa bobot yang digunakan dalam analisis STATIS mencerminkan proporsi populasi relatif dari setiap wilayah,
sehingga hasil analisis lebih representatif terhadap kondisi demografis yang sebenarnya.
Definisi 2.2 (Bobot Demografis Wilayah ke-i (mi)). Misalkan JPi =
1
K
PK
k=1 JPik adalah
rata-rata jumlah penduduk wilayah ke-i selama K tahun pengamatan. Bobot demografis
mi didefinisikan sebagai proporsi penduduk wilayah ke-i terhadap total penduduk seluruh
wilayah:
mi =
JPi
Xn
i
′=1
JPi
′
, i = 1, 2, . . . , n (2.1)
dengan syarat Pn
i=1 mi = 1 dan mi > 0. Penggunaan rata-rata lintas tahun memastikan
bahwa ruang terboboti bersifat konstan di semua situasi, sehingga jarak antar-konfigurasi
dapat dibandingkan secara valid (Lavit et al., 1994).
Definisi 2.3 (Matriks Bobot D). Matriks bobot D adalah matriks diagonal berukuran n×n
dengan elemen diagonal ke-i adalah bobot demografis mi:
D = diag(m1, m2, . . . , mn) (2.2)
Matriks D mendefinisikan suatu inner product terbobot pada ruang objek: ⟨x, y⟩D =
x
TDy. Seluruh penghitungan jarak dan produk skalar dalam STATIS digunakan
dalam ruang ini. Akar dan invers akarnya masing-masing dinotasikan sebagai D1/2 =
diag(√
m1, . . . , √mn) dan D−1/2 = diag(1/
√
m1, . . . , 1/
√
mn) (Lavit et al., 1994).
Universitas Indonesia

15
Definisi 2.4 (Standardisasi Z-Score Terbobot). Misalkan xijk adalah nilai indikator ke-j
pada wilayah ke-i tahun ke-k. Rata-rata populasi terbobot dan simpangan baku populasi
terbobot didefinisikan sebagai:
µjk =
Xn
i=1
mi xijk (2.3)
σjk =
vuutXn
i=1
mi (xijk − µjk)
2
(2.4)
Nilai terstandardisasi terbobot (weighted standardized value) dihitung sebagai:
zijk =
xijk − µjk
σjk
, i = 1, . . . , n; j = 1, . . . , p; k = 1, . . . , K (2.5)
Persamaan 2.5 memastikan bahwa Pn
i=1 mi zijk = 0 (pusat terbobot nol) dan
Pn
i=1 mi z
2
ijk = 1 (varians terbobot satu) (Lavit et al., 1994).
Hasil standardisasi menghasilkan matriks data standar terbobot Zk (berukuran n × p)
untuk setiap situasi ke-k. Prosedur ini memastikan bahwa tidak ada indikator yang
mendominasi analisis hanya karena skala pengukurannya lebih besar, sekaligus mempertahankan konsistensi ruang bobot D di seluruh situasi waktu.
2.6 Statistika Deskriptif
Sebelum menerapkan metode analisis multivariat, tahap awal yang penting adalah
melakukan eksplorasi data melalui Statistika Deskriptif. Statistika deskriptif bertujuan
meringkas dan mendeskripsikan karakteristik distribusi data secara numerik, sehingga
peneliti dapat memahami pola umum, keragaman, dan anomali pada data sebelum tahap
analisis data yang lebih mendalam (Johnson & Wichern, 2007).
2.6.1 Ukuran Pemusatan
Ukuran pemusatan mendeskripsikan nilai “tengah” dari data.
Definisi 2.5 (Rata-Rata (Mean)). Rata-rata dari n pengamatan x1, x2, . . . , xn suatu indikator adalah:
x¯ =
1
n
Xn
i=1
xi (2.6)
Rata-rata bersifat sensitif terhadap pencilan (outlier).
Universitas Indonesia

16
Definisi 2.6 (Median). Median (Q2) adalah nilai tengah dari suatu indikator setelah diurutkan dari yang terkecil hingga yang terbesar:
Q2 =



x(
n+1
2 )
, jika n ganjil
x(
n
2 )
+ x(
n
2 +1)
2
, jika n genap
(2.7)
Nilai median lebih robust terhadap pencilan dibandingkan nilai rata-rata.
2.6.2 Ukuran Penyebaran
Ukuran penyebaran mendeskripsikan seberapa jauh data bervariasi di sekitar pusat
distribusinya.
Definisi 2.7 (Simpangan Baku (Standar Deviasi)). Simpangan baku populasi dari n pengamatan adalah:
s =
vuut
1
n
Xn
i=1
(xi − x¯)
2
(2.8)
Simpangan baku merepresentasikan akar kuadrat dari ragam (variance) populasi s
2
. Karena
penelitian ini mencakup seluruh wilayah kabupaten/kota di Pulau Jawa (populasi utuh),
maka pembagi yang digunakan adalah n (Johnson & Wichern, 2007).
Definisi 2.8 (Nilai Minimum, Maksimum, dan Rentang Antarkuartil (IQR)). Nilai minimum dan maksimum adalah:
xmin = min(x1, . . . , xn), xmax = max(x1, . . . , xn) (2.9)
Kuartil pertama (Q1) dan ketiga (Q3) membagi distribusi data menjadi empat bagian,
masing-masing 25%. Rentang Antarkuartil (Interquartile Range, IQR) adalah:
IQR = Q3 − Q1 (2.10)
IQR merupakan ukuran penyebaran yang robust karena tidak dipengaruhi oleh pencilan.
Definisi 2.9 (Koefisien Variasi (CV)). Koefisien Variasi atau yang dalam bahasa Inggris
disebut sebagai Coefficient of Variation (CV) digunakan untuk mengukur variabilitas relatif
terhadap rata-rata, dengan rumus sebagai berikut:
CV =
s
x¯
× 100% (2.11)
Universitas Indonesia

17
CV berfungsi untuk membandingkan tingkat keragaman antar-indikator yang memiliki
satuan atau skala yang berbeda-beda, dan ukuran ini juga lazim digunakan dalam analisis
stabilitas (Departemen Pertanian (Deptan), 2003).
2.6.3 Ukuran Bentuk Distribusi
Ukuran bentuk memberikan informasi karakteristik bentuk simetris dan bentuk
keruncingan dari distribusi data.
Definisi 2.10 (Kemiringan Distribusi (Skewness)). Skewness mengukur tingkat asimetris
distribusi data:
g1 =
1
n
Pn
i=1(xi − x¯)
3
s
3
(2.12)
Interpretasi: g1 > 0 berarti distribusi data condong ke kanan (right-skewed, ekor panjang
di kanan); g1 < 0 berarti condong ke kiri (left-skewed); g1 ≈ 0 berarti distribusi data
hampir simetris. Skewness juga dapat diklasifikasikan dari rentang moderat hingga sangat
menceng yang biasa dijadikan sebagai acuan baku (Bulmer, 1979).
Definisi 2.11 (Kurtosis (Keruncingan)). Kurtosis mengukur ketebalan ekor dan keruncingan relatif distribusi data dibandingkan distribusi normal:
g2 =
1
n
Pn
i=1(xi − x¯)
4
s
4
− 3 (2.13)
Nilai negatif dari pengurangan 3 merupakan koreksi excess kurtosis (Weisstein, Eric W.,
2024): g2 > 0 berarti distribusi data lebih runcing (leptokurtik); g2 < 0 berarti lebih datar
(platikurtik) dibandingkan distribusi normal.
2.7 Pemilihan Indikator dalam Analisis Multivariat
Keberhasilan suatu analisis multivariat berbasis reduksi dimensi sangat bergantung
pada kualitas himpunan indikator yang digunakan. Indikator yang tidak memiliki hubungan
satu sama lain tidak akan menghasilkan komponen yang bermakna, sedangkan indikator
yang terlalu berkorelasi sempurna dapat menyebabkan matriks kovariansi menjadi singular
(Hair et al., 2019). Sebelum menerapkan metode analisis, perlu dilakukan serangkaian
uji statistik untuk memverifikasi bahwa himpunan indikator layak untuk dianalisis secara multivariat yang artinya terdapat struktur kovariansi yang cukup untuk diekstraksi
(factorable).
Dua kriteria utama yang lazim digunakan untuk mengetahui apakah data yang akan
digunakan memenuhi syarat untuk dilakukan analisis multivariat adalah: (1) Uji KaiserUniversitas Indonesia

18
Meyer-Olkin (KMO) yang menilai kecukupan (sampling adequacy) masing-masing indikator maupun keseluruhan himpunan indikator, dan (2) Uji Bartlett tentang Spherisitas yang
menguji apakah matriks korelasi antar-indikator secara keseluruhan berbeda signifikan
dari matriks identitas.
2.7.1 Uji Kaiser-Meyer-Olkin (KMO)
Statistik Kaiser-Meyer-Olkin (KMO) dikembangkan oleh Kaiser (1970, 1974) sebagai
ukuran untuk menilai kecukupan sampel (sampling adequacy) dalam analisis faktor dan
analisis komponen utama. Secara intuitif, KMO mengukur sejauh mana korelasi antarpasang indikator dapat “dijelaskan” oleh indikator-indikator lainnya, suatu kondisi yang
disebut korelasi parsial rendah. Apabila korelasi parsial antar-indikator rendah setelah
dikontrol indikator lain, maka ada faktor laten yang mendasari struktur korelasinya, kondisi
ini disyaratkan dalam analisis berbasis reduksi dimensi (Dziuban & Shirkey, 1974).
Definisi 2.12 (Statistik KMO). Misalkan rjk adalah koefisien korelasi Pearson antara
indikator ke-j dan ke-k, dan qjk adalah elemen dari invers matriks korelasi R−1. Korelasi
parsial antara indikator ke-j dan ke-k dihitung sebagai:
rjk.parsial = −
qjk
√qjj · qkk
(2.14)
Statistik KMO keseluruhan (overall KMO) didefinisikan sebagai:
KMO =
X
j̸=k
r
2
jk
X
j̸=k
r
2
jk +
X
j̸=k
r
2
jk.parsial
(2.15)
Nilai KMO berada di antara 0 dan 1. Semakin besar nilai KMO, semakin kecil jumlah
kuadrat korelasi parsial (penyebut) relatif terhadap jumlah kuadrat korelasi bivariate (pembilang). Hal ini mengindikasikan bahwa indikator-indikator saling berbagi faktor laten
yang kuat.
Selain statistik KMO keseluruhan, Kaiser (1974) mendefinisikan Measure of Sampling
Adequacy (MSAj ) untuk setiap indikator ke-j:
MSAj =
X
k̸=j
r
2
jk
X
k̸=j
r
2
jk +
X
k̸=j
r
2
jk.parsial
(2.16)
Universitas Indonesia

19
Interpretasi nilai KMO dan MSAj menurut Kaiser (1974) dan Hair et al. (2019)
disajikan pada Tabel 2.2 berikut:
Tabel 2.2: Kriteria Interpretasi Nilai KMO
Nilai KMO (x) Interpretasi
0,90 ≤ x ≤ 1,00 Marvelous (Luar Biasa)
0,80 ≤ x < 0,90 Meritorious (Sangat Baik)
0,70 ≤ x < 0,80 Middling (Cukup Baik)
0,60 ≤ x < 0,70 Mediocre (Biasa)
0,50 ≤ x < 0,60 Miserable (Buruk, namun dapat diterima)
x < 0,50 Unacceptable (Tidak diterima)
Apabila nilai MSAj suatu indikator lebih kecil dari 0,50; indikator tersebut dipertimbangkan untuk dikeluarkan dari analisis dan himpunan indikator diuji ulang hingga semua
indikator yang tersisa memenuhi syarat untuk dianalisis lebih lanjut (Hair et al., 2019;
Widyaningsih, 2024).
2.7.2 Uji Bartlett tentang Spherisitas (Bartlett’s Test of Sphericity)
Uji Bartlett tentang Spherisitas, dikembangkan oleh Bartlett (1954), bertujuan menguji
hipotesis nol bahwa matriks korelasi populasi adalah matriks identitas (R = I). Jika R = I,
tidak ada korelasi sama sekali antar-indikator sehingga tidak ada informasi bersama yang
dapat diekstraksi melalui analisis multivariat (Johnson & Wichern, 2007).
Definisi 2.13 (Uji Bartlett tentang Spherisitas). Diberikan matriks korelasi sampel R
(berukuran p × p, dengan p adalah banyaknya indikator) dari n observasi. Hipotesis uji
Bartlett adalah:
• H0: Σ = σ
2
I (matriks korelasi populasi adalah matriks identitas)
• H1: Σ ̸= σ
2
I (terdapat korelasi yang signifikan antar-indikator)
Statistik uji χ
2 dihitung dengan aproksimasi (Bartlett, 1954):
χ
2 = −

n − 1 −
2p + 5
6

ln |R| (2.17)
dengan |R| adalah determinan matriks korelasi. Statistik ini berdistribusi χ
2 dengan derajat
bebas:
df =
p(p − 1)
2
(2.18)
Universitas Indonesia

20
Logika Determinan Matriks Korelasi Apabila semua indikator saling tidak berkorelasi,
maka R = I sehingga |R| = 1 dan ln(1) = 0, sehingga nilai statistik uji χ
2 bernilai nol,
artinya korelasi antar indikator tidak signifikan. Sebaliknya, ketika indikator-indikator
saling berkorelasi erat, |R| → 0 sehingga ln |R| bernilai negatif dan statistik uji menjadi
besar, menghasilkan p-value yang kecil (Johnson & Wichern, 2007).
Jika p-value < α, maka H0 ditolak, yang berarti terdapat struktur korelasi yang
signifikan. Kondisi ini merupakan prasyarat yang diperlukan agar analisis reduksi dimensi
dapat mengekstraksi informasi yang bermakna.
Definisi 2.14 (Determinan Matriks). Determinan dari matriks bujur sangkar A berukuran
p × p, dinotasikan sebagai |A| atau det(A), adalah sebuah nilai skalar yang merangkum
properti-properti penting dari matriks tersebut, termasuk keterbalikannya (invertibility) dan
sifat singular. Secara formal, determinan didefinisikan sebagai (Anton & Rorres, 2013):
|A| =
X
σ∈Sp
sgn(σ)
Y
p
i=1
ai,σ(i) (2.19)
dengan:
• Sp adalah himpunan seluruh permutasi dari indeks {1, 2, . . . , p}.
• sgn(σ) adalah tanda (paritas) dari permutasi σ (bernilai +1 untuk permutasi genap,
dan −1 untuk permutasi ganjil).
• ai,σ(i) mewakili elemen matriks pada baris ke-i dan kolom ke-σ(i).
Untuk matriks korelasi populasi R, nilai determinan memberikan ukuran global atas
keberadaan multikolinieritas antar indikator. Apabila indikator-indikator saling bebas linier
(matriks korelasi membentuk matriks identitas), maka determinan bernilai maksimum,
yaitu |R| = 1. Sebaliknya, semakin tinggi korelasi linier antar indikator, nilai determinan
akan semakin mengecil mendekati 0. Jika determinan bernilai mutlak 0, matriks saling
berkorelasi sempurna sehingga tidak dapat diinvers (Johnson & Wichern, 2007).
2.8 Dasar Aljabar Linier dan Dekomposisi Matriks
Metode reduksi dimensi seperti Principal Component Analysis (PCA) dan STATIS
sangat bergantung pada operasi aljabar linier, khususnya dekomposisi matriks simetris.
Pemahaman mengenai nilai eigen, vektor eigen, dan sifat trace matriks menjadi prasyarat
untuk memahami bagaimana metode-metode ini mengekstraksi informasi struktural dari
data.
Universitas Indonesia

21
Definisi 2.15 (Nilai Eigen dan Vektor Eigen). Diberikan sebuah matriks persegi A berukuran p × p. Sebuah skalar λ disebut sebagai nilai eigen (eigenvalue) dari A jika terdapat
sebuah vektor tak-nol v (berukuran p × 1) yang disebut vektor eigen (eigenvector),
sedemikian sehingga berlaku:
Av = λv (2.20)
Vektor eigen menentukan “arah” dominan dari transformasi linier matriks A, sedangkan
nilai eigen menentukan besaran “kekuatan” atau variansi pada arah tersebut.
Definisi 2.16 (Invers Matriks). Sebuah matriks persegi A memiliki invers A−1jika
det(A) ̸= 0, sedemikian sehingga berlaku AA−1 = A−1A = I, dengan I adalah matriks
identitas. Invers dari matriks korelasi digunakan dalam penghitungan statistik KMO untuk
mendapatkan korelasi parsial melalui matriks anti-citra.
Karena matriks kovariansi selalu bersifat simetris (A = AT), berlaku teorema dekomposisi spektral berikut.
Teorema 2.1 (Dekomposisi Spektral). Setiap matriks simetris A berukuran p × p dapat
didekomposisi menjadi:
A = VΛVT =
X
p
j=1
λjvjv
T
j
(2.21)
dengan V adalah matriks ortogonal (p × p) yang kolom-kolomnya (vj ) adalah vektor
eigen yang saling tegak lurus, dan Λ adalah matriks diagonal berisi nilai eigen terurut
λ1 ≥ λ2 ≥ · · · ≥ λp ≥ 0 pada elemen diagonalnya.
Selain dekomposisi, konsep trace matriks juga menjadi fundamental, khususnya dalam
penghitungan koefisien korelasi antar-konfigurasi (Koefisien RV) pada metode STATIS.
Definisi 2.17 (Trace Matriks). Trace dari sebuah matriks persegi A berukuran n × n,
dinotasikan sebagai tr(A), adalah penjumlahan dari elemen-elemen diagonal utamanya:
tr(A) = Xn
i=1
aii (2.22)
Sifat penting dari trace yang digunakan dalam pengukuran jarak antar-matriks adalah
invariansi siklis: tr(AB) = tr(BA).
Definisi 2.18 (Dekomposisi Nilai Singular (SVD)). Singular Value Decomposition (SVD)
adalah faktorisasi yang dapat diterapkan pada matriks persegi panjang manapun A berukuran m × n:
A = USVT(2.23)
Universitas Indonesia

22
dengan:
• U (m × m): matriks ortogonal, kolom-kolomnya disebut vektor singular kiri.
• S (m × n): matriks diagonal yang berisi nilai singular σ1 ≥ σ2 ≥ · · · ≥ 0 pada
diagonal utamanya.
• VT(n × n): transpos dari matriks ortogonal V, kolom-kolomnya disebut vektor
singular kanan.
Untuk matriks simetris semi-definit positif Wc (seperti matriks kompromi STATIS), SVD
setara dengan dekomposisi spektral: nilai singular σi sama dengan nilai eigen λi dan
U = V.
Ekuivalensi SVD dan Eigendekomposisi untuk Matriks Simetris PSD Untuk matriks simetris semi-definit positif A = AT dengan v
TAv ≥ 0 untuk semua v, berlaku
identitas: jika A = QΛQTadalah dekomposisi spektral (dengan Λ = diag(λ1, . . . , λn),
λi ≥ 0, dan Q ortogonal), maka A = QΛQTsudah merupakan SVD dengan
U = V = Q dan Σ = Λ (nilai singular = nilai eigen). Ini karena nilai singular suatu
matriks didefinisikan sebagai akar kuadrat nilai eigen dari ATA = A2; untuk matriks PSD,
p
λi(A2) = pλ
2
i = λi = σi
. Akibatnya, dalam implementasi komputasional penelitian
ini, numpy.linalg.eigh (eigendekomposisi untuk matriks Hermitian/simetris) dan
numpy.linalg.svd menghasilkan hasil yang secara aljabar identik untuk Wc.
SVD memiliki kestabilan numerik yang lebih baik dibandingkan dekomposisi eigen
langsung untuk matriks nonsimetris. Untuk Wc yang simetris PSD, keduanya setara. Skor
kompromi STATIS (Fc) diperoleh dari:
Fc = U
√
S (2.24)
sehingga Wc = USUT = FcF
T
c
.
2.9 Korelasi Pearson
Korelasi Pearson adalah ukuran statistik yang mengkuantifikasi kekuatan dan arah
hubungan linier antara dua indikator. Koefisien ini digunakan secara luas dalam analisis
multivariat, khususnya pada pembentukan matriks korelasi yang merupakan fondasi Uji
KMO dan Bartlett, serta pada komputasi loadings (arah indikator dalam Biplot) (Johnson
& Wichern, 2007).
Universitas Indonesia

23
Definisi 2.19 (Koefisien Korelasi Pearson). Koefisien korelasi Pearson antara dua indikator
x dan y dari n pengamatan didefinisikan sebagai:
rxy =
Xn
i=1
(xi − x¯)(yi − y¯)
sXn
i=1
(xi − x¯)
2
·
sXn
i=1
(yi − y¯)
2
(2.25)
Nilai rxy berada di antara −1 dan 1:
• rxy = 1: korelasi positif sempurna (ketika x naik, y naik proporsional).
• rxy = −1: korelasi negatif sempurna (ketika x naik, y turun proporsional).
• rxy = 0: tidak ada hubungan linier.
Matriks Korelasi Untuk p indikator, koefisien Pearson antar setiap pasang indikator dapat
disusun menjadi matriks korelasi simetris R berukuran p × p, dengan elemen (j, k)-nya
adalah rjk. Diagonal utamanya selalu bernilai 1 (rjj = 1). Matriks korelasi inilah yang
menjadi dasar penghitungan statistik KMO dan Bartlett.
Dalam konteks Dynamic Biplot, loadings indikator ke-j pada komponen ke-ℓ dihitung
sebagai korelasi Pearson antara nilai rata-rata indikator tersebut lintas situasi (x¯j ) dengan
skor kompromi pada komponen ke-ℓ (F
(ℓ)
c ):
gjℓ = r(¯xj, F
(ℓ)
c
) (2.26)
Panjang vektor loading qg
2
j1 + g
2
j2 mencerminkan seberapa baik indikator tersebut direpresentasikan dalam ruang dua dimensi kompromi.
2.10 Principal Component Analysis (PCA) dan Biplot
Principal Component Analysis (PCA) atau dalam bahasa Indonesia dikenal sebagai
Analisis Komponen Utama merupakan metode yang pertama kali diinisiasi oleh Pearson
(1901) melalui gagasan penyesuaian bidang pada himpunan titik, kemudian berkembang
menjadi teknik reduksi dimensi data multivariat yang mentransformasi p indikator asli
yang saling berkorelasi menjadi q indikator baru (q < p) yang tidak berkorelasi, disebut
komponen utama (principal components) (Hotelling, 1933).
Misalkan X adalah matriks data berukuran n × p yang telah distandardisasi. PCA
mengekstraksi informasi dengan mendekomposisi matriks kovariansi S =
1
n−1XTX.
Universitas Indonesia

24
Berdasarkan Teorema Dekomposisi Spektral, S = VΛVT, menghasilkan dua elemen
utama:
1. Matriks Loadings (G) adalah matriks vektor eigen yang berisi bobot kontribusi
setiap indikator asli terhadap setiap komponen utama. Kolom ke-j dari V disebut
vektor loading ke-j.
2. Matriks Scores (F) adalah proyeksi data asli ke dalam ruang komponen utama,
dihitung sebagai:
F = XV (2.27)
Untuk memvisualisasikan hasil reduksi dimensi secara simultan antara objek dan
indikator, digunakan teknik Biplot.
Definisi 2.20 (Biplot). Biplot adalah representasi grafis dua dimensi yang memetakan baris
(objek) dan kolom (indikator) dari matriks data X secara bersamaan dalam satu grafik
(Gabriel, 1971). Biplot didasarkan pada aproksimasi matriks rank-rendah:
X ≈ FVT(2.28)
Dalam visualisasi ini:
• Objek ke-i digambarkan sebagai titik berdasarkan baris ke-i dari matriks Skor F.
Jarak Euclidean antar titik merepresentasikan kemiripan antar objek.
• Indikator ke-j digambarkan sebagai vektor (panah) berdasarkan baris ke-j dari
matriks Loadings G. Sudut antar vektor merepresentasikan korelasi antar indikator.
Sudut lancip (< 90◦) berarti korelasi positif, sedangkan sudut tumpul (> 90◦) berarti
korelasi negatif.
Gambar 2.1 mengilustrasikan hubungan visual antara komponen skor objek (F) dan
loading indikator (G) dalam representasi biplot.
Universitas Indonesia

25
Gambar 2.1: Ilustrasi Komponen Biplot (Skor dan Loading)
Skalasi Visual Biplot Dalam visualisasi biplot, seringkali terdapat perbedaan besaran
(magnitude) yang mencolok antara nilai skor objek dan loading indikator. Untuk menjamin keterbacaan grafik tanpa mengubah interpretasi statistik (sudut korelasi dan arah),
dilakukan penyesuaian skala (visual scaling) secara proporsional. Secara matematis, jika
F adalah matriks skor dan V adalah matriks loading, maka koordinat visual V∗ dihitung
sebagai:
V∗ = V ×

max(|F|)
max(|V|)
× c

(2.29)
dengan c adalah konstanta penyesuaian (umumnya bernilai 0,8). Prosedur ini memastikan
bahwa panah indikator tidak menumpuk di pusat koordinat dan dapat diobservasi secara
simultan bersama titik objek dalam satu ruang tampilan (Gabriel, 1971; Kroonenberg,
2008).
Catatan literatur: Selain skema penskalaan visual heuristik seperti Persamaan 2.29,
literatur aljabar linear modern pada konteks biplot merekomendasikan pendekatan yang
lebih konsisten secara matematis. Berdasarkan Abdi et al. (2012) – yang membahas
hubungan antara skor faktor dan loadings dalam konteks komponen utama dan analisis
struktur tiga-arah – salah satu prosedur yang direkomendasikan adalah mengalikan loading
Universitas Indonesia

26
pada komponen ke-ℓ dengan nilai singular (singular value) σℓ, di mana σℓ =
√
λℓ dengan
λℓ adalah nilai eigen komponen tersebut. Secara sederhana:
V∗
·ℓ = V·ℓ × σℓ (2.30)
Pendekatan ini membuat varians (kuadrat) dari koordinat loading pada dimensi ke-ℓ sepadan dengan nilai eigen dimensi tersebut, sehingga panah indikator menjadi commensurable
dengan sebaran skor (tanpa mengubah arah atau sudut yang merepresentasikan korelasi).
Perlu ditekankan bahwa prosedur ini bersifat representasional untuk tujuan visualisasi dan
tidak merubah penghitungan statistik dasar seperti nilai eigen, skor, atau uji signifikansi
(mis. permutasi). (Lihat Abdi et al. (2012) untuk pembahasan teoretis dan alasan aljabar
di balik transformasi ini.)
PCA dan Biplot bekerja dengan baik untuk satu tabel data pada satu titik waktu.
Namun, ketika data tersebar dalam K matriks yang berbeda-beda seiring waktu, PCA
konvensional tidak mampu menangkap dinamika evolusi struktural tersebut. Keterbatasan
inilah yang mendorong pengembangan metode STATIS.
2.11 Metode STATIS
STATIS merupakan singkatan dari Structuration des Tableaux a Trois Indices de la `
Statistique atau jika diartikan ke dalam bahasa Indonesia berbunyi ”Strukturisasi Tabel
dengan Tiga Indeks dalam Statistik”. STATIS adalah metode analisis data eksploratif yang
dirancang untuk menganalisis sekumpulan K tabel data dengan membandingkan struktur
internal antar-objek pada setiap situasi waktu (Lavit et al., 1994). Berbeda dengan metode
yang menggabungkan data secara naif (seperti Pooled PCA), STATIS mempertahankan
integritas struktur masing-masing situasi melalui tiga tahapan analisis data yang sistematis:
intra-struktur, inter-struktur, dan kompromi. Konsep inti STATIS dilandasi oleh teori
operator Escoufier (Escoufier, 1973) yang memperkenalkan cara mengukur “jarak” antara
dua konfigurasi data multivariat.
2.11.1 Tahap 1: Intra-struktur
Tahap pertama bertujuan mentransformasi tabel data terstandardisasi terbobot setiap
situasi menjadi representasi geometri yang dapat dibandingkan dalam ruang bobot D
yang sama. Misalkan Zk adalah matriks data terstandardisasi terbobot pada situasi ke-k
(Persamaan 2.5). STATIS mendefinisikan operator Escoufier Wk untuk setiap k =
Universitas Indonesia

27
1, . . . , K sebagai matriks produk skalar berukuran n × n (Lavit et al., 1994):
Wk = ZkZ
T
k
(2.31)
Elemen-elemen dalam Wk merepresentasikan kemiripan antar-objek pada situasi ke-k.
Karena analisis data dilakukan dalam ruang bobot D, maka penghitungan produk Hilbert–
Schmidt dan koefisien RV pada tahap berikutnya mengikutsertakan matriks bobot diagonal
D (Lavit et al., 1994).
Interpretasi Wk: Elemen Wk[i, j] yang bernilai besar menunjukkan bahwa objek ke-i
dan objek ke-j memiliki profil indikator yang sangat mirip pada situasi ke-k, sedangkan
nilai yang kecil menunjukkan perbedaan karakteristik yang mencolok.
Peran D dalam Wk dan Ekuivalensi Komputasional Matriks bobot D tidak muncul
secara eksplisit dalam Wk = ZkZ
T
k
karena D masuk melalui produk skalar Hilbert–
Schmidt pada tahap inter-struktur: Hkk′ = tr(WkDWk
′D) (Lavit et al., 1994). Secara
komputasional, ekspresi ini setara dengan mendefinisikan Wfk = D1/2ZkZ
T
k D1/2 =
(D1/2Zk)(D1/2Zk)
T
, sehingga:
tr(Wfk Wfk
′) = tr(D1/2ZkZ
T
k D Zk
′Z
T
k
′D1/2
) = tr(WkDWk
′D) = Hkk′
Kedua formulasi menghasilkan nilai Hkk′ yang identik. Dalam implementasi komputasional penelitian ini, digunakan bentuk Wfk karena lebih efisien secara numerik (perkalian
matriks biasa tanpa perlu menyertakan D secara eksplisit pada setiap langkah).
2.11.2 Tahap 2: Inter-struktur
Setelah representasi Wk diperoleh untuk setiap situasi, langkah selanjutnya adalah
mengukur seberapa mirip konfigurasi antar-situasi. Secara aljabar, matriks simetris Wk dan
Wk
′ dilengkapi dengan produk skalar Hilbert–Schmidt (Wk |Wk′) = tr(WkDWk′D)
(Lavit et al., 1994; Oliveira et al., 2024). Produk ini membentuk matriks interstruktur
berukuran K × K:
Hkk′ = tr(WkDWk
′D), k, k′ = 1, . . . , K. (2.32)
Matriks H = (Hkk′) merepresentasikan kemiripan geometris antar-situasi dalam kondisi
berbobot. Dekomposisi spektralnya memberikan gambaran grafis interstruktur (titik-titik
situasi di R
h
) sebagaimana dijelaskan oleh Lavit et al. (1994).
Universitas Indonesia

28
Untuk interpretasi data yang telah terstandardisasi (nilai diagonal menjadi 1), digunakan koefisien RV (Escoufier, 1973):
Definisi 2.21 (Koefisien RV). Koefisien RV antara dua matriks konfigurasi Wk dan Wk
′
didefinisikan sebagai:
RV (Wk,Wk
′) = tr(WkDWk
′D)
p
tr((WkD)
2
) tr((Wk
′D)
2
)
=
Hkk′
√
Hkk Hk
′k′
. (2.33)
Nilai RV berada di antara 0 dan 1. Nilai RV yang mendekati 1 mengindikasikan bahwa
konfigurasi antar-objek pada situasi ke-k sangat mirip dengan situasi ke-k
′
. Koefisien ini
secara geometris merupakan kosinus sudut antara dua matriks dalam ruang vektor matriks
terbobot.
Matriks RRV = (RV (Wk,Wk
′)) dapat digunakan untuk membuat heatmap dan untuk
mengetahui jarak antar-situasi dengan rumus jarak d
2
(k, k′) = 2(1 − RV ) antar-situasi,
namun pembobotan kompromi mengacu pada H: menurut Lavit et al. (1994). Koefisien
αk diambil sebagai komponen vektor eigen utama yang terkait dengan nilai-nilai eigen
terbesar dari dekomposisi spektral matriks interstruktur. Secara eksplisit,
Hu = θu, (2.34)
dengan H adalah matriks interstruktur, u adalah vektor eigen, dan θ adalah nilai eigen.
Vektor eigen pertama u1 (berkorespondensi dengan θ1 terbesar) dipilih dengan elemen
non-negatif lalu dinormalisasi:
αk =
u1k
X
K
k
′=1
u1k
′
. (2.35)
Kendala PK
k=1 αk = 1 memastikan matriks kompromi (yang akan dibahas pada sub-bab
2.11.3) sebanding dengan skala Wk (Lavit et al., 1994). Sebagai catatan, RRV dan H
memiliki nilai eigen berbeda. Hal ini sesuai dengan implementasi ACT/STATIS yang
menggunakan Persamaan (2.34)–(2.35) untuk α.
Contoh 2.1 (Logika Bobot). Apabila suatu situasi memiliki profil yang sangat anomali
dibandingkan situasi lainnya (misalnya akibat krisis), metode STATIS akan secara otomatis
memberikan bobot αk yang lebih kecil pada situasi tersebut. Hal ini memastikan bahwa
struktur kompromi tidak “terdistorsi” oleh situasi yang tidak representatif.
Universitas Indonesia

29
2.11.3 Tahap 3: Analisis Kompromi
Tahap ini membentuk struktur “rata-rata” yang disebut Kompromi. Matriks kompromi
Wc dihitung sebagai kombinasi linier terbobot (menggunakan bobot αk) dari matriks
intra-struktur Wk (Lavit et al., 1994):
Wc =
X
K
k=1
αkWk (2.36)
Matriks Wc (berukuran n × n) merepresentasikan pola hubungan antar-objek yang paling
stabil selama periode pengamatan. Karena setiap Wk adalah matriks simetris positif
semi-definit dan αk ≥ 0, maka Wc juga bersifat simetris positif semi-definit. Penurunan
sifat ini ditunjukkan melalui bentuk kuadratik untuk sebarang vektor v ∈ R
n
:
v
TWcv = vT
 X
K
k=1
αkWk
!
v =
X
K
k=1
αk(v
TWkv) ≥ 0 (2.37)
Hal ini menjamin bahwa dekomposisi spektral pada Wc akan selalu menghasilkan nilai
eigen yang non-negatif. Berdasarkan teorema 2.1, Wc dapat didekomposisi menjadi:
Wc = PΛPT(2.38)
dengan P adalah matriks vektor eigen tergeneralisasi yang memenuhi syarat ortogonalitas
terbobot PTDP = I, dan Λ adalah matriks diagonal berisi nilai eigen terurut λ1 ≥ λ2 ≥
· · · ≥ 0. Hal ini menunjukkan bahwa vektor eigen P dihitung dalam ruang terboboti D
untuk mengakomodasi heterogenitas objek (Lavit et al., 1994).
2.11.4 Tahap 4: Ekstraksi Koordinat dan Relasi Transisi
Tahap akhir dari metode STATIS adalah penentuan koordinat untuk objek dan indikator
dalam ruang bersama (biplot). Berdasarkan Abdi et al. (2012), terdapat dua skenario utama
yang menentukan arah proyeksi koordinat (transition formulas) berdasarkan fokus analisis
(host space):
1. Skenario 1: STATIS Klasik (Host = Wilayah/Objek)
Pada skenario ini, matriks produk skalar wilayah (Wc) dijadikan fondasi utama.
Karena bermula dari ruang observasi, skor kompromi wilayah ditarik langsung
secara alami, sedangkan komponen lainnya diturunkan melalui proyeksi.
Universitas Indonesia

30
• Skor Kompromi Wilayah (Fc): Diekstraksi langsung dari dekomposisi Wc:
Fc = D−1/2PΛ1/2
c
(2.39)
• Skor Parsial Wilayah (Fk): Merujuk pada Persamaan (40) Abdi et al. (2012),
skor parsial situasi k dihitung dengan memproyeksikan matriks produk skalar
situasi tersebut ke ruang kompromi:
Fk = WkDPΛ−1/2
c
(2.40)
• Loading Parsial Indikator (Gk): Merujuk pada Persamaan (39) Abdi et al.
(2012), ini merupakan relasi transisi yang memproyeksikan data ke ruang
observasi:
Gk = Z
T
k DFcΛ
−1/2
c
(2.41)
• Loading Kompromi Indikator (Gc): Pada skenario ini, loading kompromi
dihitung sebagai rata-rata terbobot dari loading parsialnya (Lavit et al., 1994):
Gc =
X
K
k=1
αkGk (2.42)
2. Skenario 2: Dual STATIS / CovSTATIS (Host = Indikator/Indikator)
Pada skenario ini, matriks korelasi indikator (Sc) menjadi fondasi utama. Karena
bermula dari ruang indikator, loading indikator ditarik secara alami, sedangkan
posisi wilayah diturunkan melalui relasi transisi.
• Loading Kompromi Indikator (Gc): Diekstraksi langsung dari dekomposisi
Sc:
Gc = VcΛ
1/2
s
(2.43)
• Loading Parsial Indikator (Gk): Melalui hukum dualitas, matriks kovariansi
parsial (Sk) diproyeksikan ke sumbu indikator (Lavit et al., 1994):
Gk = SkVcΛ
−1/2
s
(2.44)
• Skor Parsial Wilayah (Fk): Merujuk pada Persamaan (25) Abdi et al. (2012),
ini memproyeksikan data situasi k ke ruang indikator:
Fk = ZkGcΛ
−1/2
s
(2.45)
Universitas Indonesia

31
• Skor Kompromi Wilayah (Fc): Berbeda dengan Skenario 1, pada skenario ini
skor kompromi objek dihitung sebagai rata-rata terbobot dari skor parsialnya
(Lavit et al., 1994):
Fc =
X
K
k=1
αkFk (2.46)
3. Transformasi Penyesuaian Visual Biplot: Untuk kepentingan interpretasi praktis dan estetika grafik biplot, koordinat matematis hasil SVD murni (Fsvd, Gsvd)
seringkali memerlukan transformasi tambahan tanpa mengubah hubungan korelasi
antar-objek maupun indikator. Transformasi ini digunakan untuk menghasilkan
kolom koordinat visual pada sistem:
a. Penyelarasan PC1 (Mirroring): Agar dimensi pertama (PC1) memiliki makna
ekonomi yang konsisten (misalnya: nilai positif berarti wilayah lebih sejahtera),
dilakukan operasi pembalikan tanda jika ditemukan korelasi negatif antara PC1
dengan indikator performa utama:
Xvis = Xsvd × sign(ρ(P C1, Indikator Utama)) (2.47)
di mana X merepresentasikan matriks F atau G.
b. Skalasi Berbasis Nilai Singular (Scaling): Agar loading indikator sebanding dengan skor objek secara geometris, koordinat loading tiap komponen ℓ
dikalikan dengan nilai singular komponen tersebut (Abdi et al., 2012):
G
(ℓ)
vis = G(ℓ) × δℓ
, δℓ =
p
λℓ (2.48)
di mana λℓ adalah nilai eigen komponen ke-ℓ dari dekomposisi kompromi
STATIS dan δℓ adalah nilai singularnya. Penskalaan ini menjamin varians
koordinat loading visual pada tiap dimensi setara dengan nilai singular dimensi
tersebut, sehingga panah indikator dan titik objek berada dalam ruang geometris
yang dapat dibandingkan.
Rangkuman Dualitas Relasi Transisi: Sebagai ringkasan teoretis, tabel berikut menyajikan simetri aljabar antara kedua skenario STATIS yang menjamin sifat rata-rata terbobot
dalam visualisasi biplot:
Universitas Indonesia

32
Tabel 2.3: Rangkuman Dualitas Relasi Transisi STATIS
Elemen yang Dicari Skenario 1 (Host Objek) Skenario 2 (Host Indikator)
Ekstraksi Utama Fc = D−1/2PΛ1/2
c Gc = VcΛ
1/2
s
Relasi Proyeksi Sejenis Fk = WkD(PΛ−1/2
c ) Gk = Sk(VcΛ
−1/2
s )
Relasi Transisi Menyilang Gk = Z
T
k D(FcΛ
−1/2
c ) Fk = Zk(GcΛ
−1/2
s )
2.11.5 Klasifikasi Struktur: Analisis Transversal dan Longitudinal
Dalam pengembangan terbaru metodologi STATIS, Oliveira et al. (2024) mengklasifikasikan analisis data pada vektor struktur inter-struktur menjadi dua pendekatan
berdasarkan keberadaan urutan pada periode studi:
1. Analisis Transversal (Transversal Analysis): Analisis Transversal dilakukan untuk
membandingkan komponen yang setara namun tidak memiliki urutan waktu yang
mengikat (misalnya perbandingan antar kelompok atau juri).
2. Analisis Longitudinal (Longitudinal Analysis): Analisis Longitudinal dilakukan
ketika terdapat urutan yang jelas dalam periode studi, seperti data yang dikumpulkan
pada periode waktu yang berurutan. Analisis ini sangat memadai untuk memodelkan
garis tren lintasan (trajectories) dari waktu ke waktu.
Dalam konteks penelitian ini, karena data yang dianalisis mencakup periode tahun
2020 hingga 2024 secara berurutan, maka pendekatan yang digunakan adalah analisis
longitudinal, dengan arah dan besaran pergeseran lintasan objek menjadi fokus utama
interpretasi.
2.11.6 Uji Signifikansi Stabilitas Struktur (Permutasi Monte Carlo)
Metode STATIS secara deskriptif menghasilkan struktur kompromi, namun tidak
secara otomatis menjamin bahwa struktur tersebut signifikan secara statistik. Dapat
terjadi struktur kompromi yang terbentuk hanya merupakan konstruk semu dari data yang
sebenarnya acak dan tidak memiliki pola hubungan yang stabil antar-situasi.
Untuk memvalidasi hal ini, digunakan pengujian berbasis resampling, yaitu uji permutasi non-parametrik yang diadopsi dari berbagai literatur analisis multivariat (Kazi-Aoual
et al., 1995; Abdi et al., 2012; Legendre & Legendre, 2012). Uji ini tidak mensyaratkan
asumsi distribusi normal pada data. Mengingat jumlah seluruh kemungkinan permutasi
baris dari himpunan matriks sangat masif ((n!)K), menghitung permutasi eksak secara
menyeluruh sangat tidak efisien secara komputasional. Oleh karena itu, pengujian dieksekusi menggunakan pendekatan Simulasi Monte Carlo, dengan sampel permutasi acak
Universitas Indonesia

33
dalam jumlah besar (B kali iterasi) diambil untuk membangun aproksimasi numerik dari
distribusi eksak.
1. Hipotesis Uji:
• H0: θ1 = E [θ
π
1
] (Konfigurasi objek antar situasi waktu saling independen
secara linier, sehingga data bersifat acak dan tidak ada struktur bersama yang
stabil).
• H1: θ1 > E [θ
π
1
] (Konfigurasi objek antar situasi waktu memiliki ketergantungan linier parsial atau total yang signifikan, membentuk struktur kompromi
yang bermakna).
Notasi superskrip π (pi) pada rumusan hipotesis di atas merepresentasikan kondisi
permutasi. Dengan demikian, θ
π
1 melambangkan statistik nilai eigen dominan yang
dihasilkan dari data dummy (buatan) yang telah diacak (dipermutasi) untuk mensimulasikan kondisi distribusi H0, sedangkan E [θ
π
1
] merupakan nilai harapan atau nilai
rata-rata dari statistik uji yang dihasilkan dari data acak tersebut.
2. Statistik Uji dan Distribusi: Statistik yang digunakan adalah θ1, yaitu nilai eigen
terbesar dari matriks interstruktur H pada Persamaan (2.32). Nilai θ1 mengkuantifikasi proporsi informasi (variansi) matriks produk skalar yang dapat diringkas
secara konsisten dalam satu dimensi interstruktur. Semakin besar θ1, semakin kuat
bukti empiris adanya struktur spasio-temporal bersama yang stabil (Abdi et al.,
2012).
Distribusi teoretis (parametrik) dari θ1 dalam analisis multivariat dapat menjadi
semakin kompleks sehingga semakin sulit untukdidefinisikan. Sebagai solusinya,
distribusi probabilitas pada kondisi H0 (null distribution) diasumsikan benar. Kondisi
ini dibangun secara empiris melalui simulasi Monte Carlo terhadap data yang dipermutasi sehingga menghasilkan sekumpulan statistik uji sintetis yang dinotasikan
sebagai θ
π
1 yang jumlahnya sama dengan jumlah iterasi yang dilakukan (B).
3. Prosedur Simulasi Monte Carlo:
a. Menghitung nilai uji empiris θ
obs
1
(observasi asli) dari matriks interstruktur H
berdasar kumpulan matriks Zk.
b. Melakukan pengacakan baris (permutasi label objek) secara independen pada
setiap matriks situasi, menghasilkan sekumpulan matriks dummy acak Z
π
k
. Pengacakan ini mengubah struktur keterkaitan asli antar-situasi, mensimulasikan
kondisi di mana H0 berlaku.
Universitas Indonesia

34
c. Menghitung operator Escoufier terbobot untuk himpunan data tersimulasi:
Wπ
k = (D1/2Z
π
k
)(D1/2Z
π
k
)
T
. Notasi eksponen π mensyaratkan pembedaan
dari conjugate transpose atau rotasi matriks asliWk maupun matriks kompromi
Wc.
d. Membentuk matriks interstruktur baru Hπ yang elemennya bernilai Hπ
kk′ =
tr(Wπ
kWπ
k
′). Selanjutnya, mengekstraksi nilai eigen terbesarnya yang ditetapkan sebagai θ
π
1
.
e. Mengulangi iterasi pengacakan matriks pada langkah (b) hingga (d) secara
terus-menerus sebanyak B kali (dalam penelitian ini B = 1000) untuk menghasilkan histogram distribusi frekuensi null dari θ
π
1
.
4. Pengambilan Keputusan: Berbeda dengan uji statistik parametrik konvensional
yang mengandalkan tabel distribusi teoritis, distribusi sampling dari θ1 dalam konteks
STATIS tidak diketahui bentuk pastinya secara analitis. Oleh karena itu, pengujian
ini bersifat non-parametrik (distribution-free), di mana distribusi null dibangun
secara empiris dari data itu sendiri melalui proses permutasi (Kazi-Aoual et al.,
1995).
Nilai p-value dihitung sebagai proporsi statistik uji hasil permutasi (θ
π
1
) yang nilainya
lebih besar atau sama dengan statistik observasi (θ
obs
1
):
p-value =
X
B
b=1
I

θ
π(b)
1 ≥ θ
obs
1

+ 1
B + 1
(2.49)
dengan I(·) sebagai fungsi indikator Bernoulli yang bernilai 1 jika argumen perbandingan benar dan 0 jika sebaliknya. Elemen penambah +1 pada pembilang dan
penyebut merupakan operasi pseudo-count correction yang lazim digunakan dalam
algoritma Monte Carlo untuk mereduksi bias penaksir titik, sehingga p-value eksak
tidak bernilai nol absolut (0%) (Phipson & Smyth, 2010).
Apabila p-value < α (misalnya 0,05), maka H0 ditolak. Hal ini membuktikan bahwa
kompromi STATIS mencerminkan konstelasi evolusi sistem sosial-ekonomi yang
signifikan, bukan produk distorsi stokastik (Legendre & Legendre, 2012).
2.11.7 Model Vektor Struktur dan Kerangka Inferensial STATIS
Sebagaimana diuraikan di atas, STATIS klasik (Lavit et al., 1994) bersifat eksploratif:
kesimpulan mengenai ada-tidaknya struktur bersama didasarkan semata-mata pada inspeksi
Universitas Indonesia

35
grafis terhadap dekomposisi interstruktur. Areia et al. (2008) memperluas kerangka ini
dengan memperkenalkan formulasi model probabilistik yang memungkinkan inferensi
statistik formal atas parameter interstruktur. Pengembangan ini dirangkum kembali dalam
tinjauan metodologis terkini oleh Oliveira et al. (2024).
Model Rank-1 dan Estimator Vektor Struktur
Apabila titik-titik representasi situasi pada plot interstruktur cenderung terletak sepanjang sumbu eigen pertama struktur bersama derajat 1 (common structure of degree
1) maka Areia et al. (2008); Mexia & Oliveira (2007) mengajukan model berikut untuk
matriks interstruktur H:
H = λααT + E¯ , E¯ =
1
2

E + E
T

, (2.50)
dengan asumsi bahwa Vec(E) ∼ N (0, σ2IK2 ), dengan λ > 0 adalah skalar kegentingan
(magnitude), α adalah vektor satuan yang menentukan arah struktur bersama, dan E¯ adalah
komponen galat simetris.
Dalam model ini, informasi yang terkandung dalam H dipadatkan ke dalam vektor
struktur β:
β˜ = λ1α1, (2.51)
dengan λ1 dan α1 berturut-turut adalah nilai eigen dan vektor eigen pertama (terbesar)
dari H. Komponen ke-k dari β˜, yaitu β˜
k = λ1u1k (sebelum normalisasi α), merupakan
estimator tidak bias yang konsisten dari parameter βk = λuk (Areia et al., 2008; Oliveira
et al., 2024).
Besaran Sum of Squares residual dihitung sebagai:
V = ∥H∥
2
F − ∥β˜∥
2 =
X
K
k=1
X
K
k
′=1
H
2
kk′ − λ
2
1
, (2.52)
dengan ∥H∥
2
F = tr(HTH) = tr(H2
) (karena H simetris). Nilai V yang kecil relatif
terhadap ∥H∥
2
F mengindikasikan bahwa model rank-1 memberikan aproksimasi yang baik
terhadap H, sehingga interpretasi melalui satu vektor struktur β˜ adalah sahih.
Universitas Indonesia

36
Validasi Kegentingan Eigenvalue (Dominansi Struktur)
Validitas model bergantung pada dominansi eigenvalue pertama. Kuantitas yang
mengukur hal ini didefinisikan sebagai rasio:
τ =
λ
2
σ
2
, (2.53)
dengan σ
2 merupakan varians galat dari model pada Persamaan (2.50). Nilai τ yang
besar (yaitu ketika λ1 ≫ λj untuk j ≥ 2) mengonfirmasi bahwa satu komponen dominan
mampu menjelaskan hampir seluruh variansi interstruktur. Areia et al. (2008) menunjukkan
bahwa estimator λ1 dan α1 memiliki perilaku yang baik (good behaviors) melalui simulasi.
Secara empiris, dominansi ini dinilai melalui proportion of inertia pada sumbu pertama
interstruktur:
Inertia1 =
λ1
X
K
j=1
λj
, (2.54)
yang setara dengan kontribusi variansi dimensi pertama pada plot interstruktur.
Analisis Transversal dan Longitudinal (Keluarga Seri Studi)
Perluasan lebih lanjut oleh Areia et al. (2008), yang dirangkum kembali dalam Oliveira
et al. (2024), mempertimbangkan situasi di mana tersedia beberapa seri studi yang masingmasing berkorespondensi dengan satu perlakuan atau kelompok dalam suatu model eksperimental. Misalkan terdapat G kelompok (misalnya G provinsi), masing-masing memiliki
seri K matriks data {Zgk}
K
k=1 untuk g = 1, . . . , G. Untuk setiap kelompok g diperoleh
vektor struktur yang diestimasi β˜
g = λ1gα1g. Dua jenis analisis inferensial kemudian
dapat diterapkan (Oliveira et al., 2024):
1. Analisis Transversal: Menguji apakah vektor-vektor struktur dari kelompok yang
berbeda (β˜
1
, β˜
2
, . . . , β˜
G) berbeda secara signifikan satu sama lain, atau dengan kata
lain menguji ada-tidaknya pengaruh faktor kelompok terhadap struktur bersama.
Analisis ini umumnya dilakukan pada komponen-komponen homologus dari vektor
struktur (yaitu pada indeks situasi yang sama lintas kelompok) (Oliveira et al., 2024).
2. Analisis Longitudinal: Menguji apakah terdapat progresi atau tren yang signifikan
dalam komponen vektor struktur seiring berjalannya waktu. Analisis ini relevan
ketika seri studi disusun secara kronologis, dan pertanyaannya adalah apakah perubahan temporal yang teramati pada struktur bersama merupakan tren yang sistematis
atau sekadar fluktuasi acak (Areia et al., 2008; Oliveira et al., 2024).
Universitas Indonesia

37
Hubungan dengan Uji Permutasi Uji permutasi yang dijelaskan pada subbab sebelumnya merupakan pendekatan non-parametrik yang menguji hipotesis global mengenai
keberadaan struktur bersama yang setara dengan menguji H0 : λ = 0 dalam model
pada Persamaan (2.50). Analisis transversal dan longitudinal, di sisi lain, beroperasi
setelah struktur bersama dikonfirmasi ada, dan mengajukan pertanyaan inferensial yang
lebih spesifik mengenai perbedaan antar-kelompok dan tren temporal. Keduanya bersifat
komplementer.
2.12 Dynamic Biplot dan Visualisasi Trajektori
Hasil dekomposisi kompromi Fc hanya memberikan gambaran rata-rata posisi objek
selama periode pengamatan. Untuk memahami dinamika perubahan dari situasi ke situasi,
diperlukan perluasan yang disebut Dynamic Biplot atau Trajectory Biplot (Egido & Galindo,
2015).
Konsep dasar dari visualisasi ini adalah memproyeksikan data tabel parsial (situasi
tertentu) ke dalam ruang referensi yang stabil (ruang kompromi), sehingga posisi objek
dapat dibandingkan secara apple-to-apple lintas waktu.
2.12.1 Proyeksi Lintasan (Trajectory)
Misalkan Zk adalah matriks data terstandardisasi pada situasi ke-k (n × p) dan Vc
(berukuran p × 2) adalah matriks vektor eigen dari matriks kompromi indikator Sc. Koordinat objek pada situasi ke-k, dinotasikan Fk, dihitung melalui proyeksi linier (Lavit et al.,
1994):
Fk = ZkVc (2.55)
Persamaan ini memetakan kondisi setiap situasi ke dalam sistem koordinat global yang
sama. Trajektori atau lintasan objek ke-i didefinisikan sebagai urutan vektor koordinat
dari waktu ke waktu:
Ti = {fi,1,fi,2, . . . ,fi,K} (2.56)
Visualisasi trajektori menghubungkan titik-titik koordinat ini dengan garis berarah,
sehingga perubahan karakteristik objek, misalnya perpindahan dari kuadran dengan nilai
indikator tertentu yang rendah menuju kuadran dengan nilai yang tinggi dapat teramati
secara visual.
Universitas Indonesia

38
2.12.2 Jarak Euclidean antar Trajektori
Untuk mengelompokkan objek berdasarkan dinamika lintasannya (bukan hanya posisi
rata-ratanya), digunakan ukuran jarak yang memperhitungkan pergerakan lintas seluruh
situasi waktu. Jarak Euclidean antar-trajektori objek ke-i dan objek ke-j (masing-masing
memiliki K titik koordinat) didefinisikan sebagai:
d(i, j) =
vuutX
K
k=1
∥fi,k − fj,k∥
2
(2.57)
di mana fi,k ∈ R
2
adalah koordinat proyeksi objek ke-i pada situasi ke-k. Matriks
jarak ini kemudian digunakan sebagai masukan pada algoritma K-Means berbasis trajektori, sehingga objek-objek yang memiliki pola pergerakan serupa sepanjang waktu akan
dikelompokkan bersama (Egido & Galindo, 2015).
2.12.3 Rotasi Baseline
Solusi Principal Component Analysis (PCA) dan Biplot bersifat invarian terhadap
rotasi ortogonal, artinya rotasi sumbu koordinat tidak mengubah jarak antar-objek maupun
sudut antar-indikator, namun dapat memudahkan interpretasi (Kroonenberg, 2008). Dalam
implementasi dashboard, rotasi digunakan untuk menyelaraskan suatu vektor referensi
(indikator atau objek yang dipilih sebagai baseline) dengan sumbu horizontal (PC1),
sehingga perbandingan posisi relatif objek lain terhadap baseline tersebut lebih intuitif.
Sudut rotasi θ dihitung dari posisi vektor referensi v = (v1, v2)
T menggunakan fungsi
atan2:
θ = − atan2(v2, v1) (2.58)
Fungsi atan2 mengembalikan sudut dalam radian antara sumbu positif-x dan vektor
(v1, v2), dengan rentang nilai (−π, π]. Tanda negatif memastikan rotasi berlawanan jarum
jam agar vektor referensi disejajarkan ke sumbu kanan (positif PC1).
Matriks rotasi ortogonal 2 × 2 yang merotasi koordinat sebesar sudut θ adalah:
Q(θ) = 
cos θ − sin θ
sin θ cos θ
!
(2.59)
Dengan demikian, koordinat terotasi dihitung sebagai:
F
∗ = F Q(θ) (2.60)
Universitas Indonesia

39
Prosedur ini diterapkan secara seragam ke semua plot (skor, loading, dan trajektori seluruh
situasi waktu) agar konsistensi referensi spasial tetap terjaga. Hal ini memungkinkan
peneliti mensejajarkan indikator tertentu dengan sumbu horizontal untuk memudahkan
interpretasi posisi relatif antar-objek.
2.13 Metode Pengelompokan (Clustering)
Analisis pengelompokan (clustering) adalah metode statistik yang bertujuan mengelompokkan objek ke dalam beberapa kelompok sedemikian rupa sehingga objek yang
berada dalam satu kelompok memiliki tingkat kemiripan yang tinggi, sementara objek
antar kelompok bersifat heterogen.
2.13.1 K-Means Clustering
K-Means adalah algoritma partitioning clustering non-hierarkis yang mempartisi n
objek ke dalam k klaster sedemikian rupa sehingga variansi dalam klaster (within-klaster
variation) menjadi sekecil mungkin (MacQueen, 1967). Misalkan xi ∈ R
p
adalah vektor
observasi objek ke-i. Algoritma K-Means meminimumkan fungsi objektif Sum of Squared
Errors (SSE):
SSE =
X
k
j=1
X
xi∈Cj
∥xi − µj∥
2
(2.61)
dengan Cj adalah himpunan objek pada klaster ke-j dan µjadalah sentroid klaster ke-j.
Algoritma ini bekerja secara iteratif dengan memindahkan objek ke klaster bersentriod
terdekat hingga tidak ada objek yang berpindah tempat.
2.13.2 Penentuan Jumlah klaster Optimal (Metode Elbow)
Tantangan utama K-Means adalah penentuan jumlah klaster optimal k yang bersifat
apriori. Untuk menentukan k optimal secara empiris, metode Elbow digunakan sebagai
panduan heuristik (Bholowalia & Kumar, 2014). Metode Elbow mengevaluasi efektivitas
partisi dengan merencanakan kurva nilai inersia (SSE) pada sumbu Y terhadap nilai k
yang bervariasi (misalnya 1 hingga 10) pada sumbu X. Semakin banyak cluster, nilai
SSE secara matematis akan selalu menurun karena jarak objek ke sentroidnya terdekat
akan semakin memendek. Namun, pada suatu titik tertentu (k), penurunan drastis SSE
akan mulai melandai sehingga membentuk sudut (patahan/siku). Titik patahan inilah yang
menunjukkan diminishing returns terhadap penambahan jumlah klaster yang direkomendasikan sebagai jumlah klaster optimal agar model tidak mengalami overfitting (Syakur
Universitas Indonesia

40
et al., 2018).
2.13.3 Koefisien Siluet (Silhouette Score)
Koefisien Siluet diperkenalkan oleh Rousseeuw (1987) sebagai ukuran kuantitatif
yang mengevaluasi kualitas partisi klaster dari dua aspek sekaligus: (1) kepadatan dalam
klaster (cohesion) dan (2) keterpisahan antar-klaster (separation). Untuk setiap objek i
yang berada pada klaster Cq, didefinisikan rata-rata jarak intra-klaster:
a(i) = 1
|Cq| − 1
X
j∈Cq, j̸=i
d(i, j) (2.62)
dan rata-rata jarak ke klaster tetangga terdekat:
b(i) = min
Cr̸=Cq
1
|Cr|
X
j∈Cr
d(i, j) (2.63)
Nilai siluet objek ke-i adalah:
s(i) = b(i) − a(i)
maxa(i), b(i)
, s(i) ∈ [−1, 1] (2.64)
Nilai s(i) → +1 menunjukkan objek terklaster dengan baik; s(i) ≈ 0 berarti objek berada
di batas antar-klaster; dan s(i) < 0 mengindikasikan kemungkinan salah penempatan.
Koefisien Siluet rata-rata keseluruhan adalah:
SC =
1
n
Xn
i=1
s(i) (2.65)
Nilai SC yang lebih besar (mendekati +1) mengindikasikan struktur klaster yang lebih
koheren dan terpisah (Rousseeuw, 1987). Dalam penelitian ini, SC digunakan bersama
kurva Elbow sebagai konfirmasi kuantitatif pemilihan jumlah klaster optimal.
2.13.4 Agglomerative Hierarchical Clustering (Dendrogram)
Sebagai alat bantu komplementer bagi metode Elbow, pendekatan hierarkis Agglomerative Clustering dengan metode Ward’s Linkage digunakan untuk memvisualisasikan
struktur pengelompokan secara hierarkis melalui Dendrogram (Ward Jr, 1963). Metode
Ward menggunakan pendekatan Analysis of Variance (ANOVA) untuk mengevaluasi jarak
antar-cluster. Pada setiap tahap aglomerasi, metode ini menggabungkan sepasang klaster
yang meminimalkan peningkatan Error Sum of Squares (ESS) total. Peningkatan ESS
Universitas Indonesia

41
agregat saat menggabungkan klaster A dan B untuk membentuk A ∪ B dihitung sebagai:
∆(A, B) = nAnB
nA + nB
∥µA − µB∥
2
(2.66)
Jarak vertikal cabang dendrogram secara visual mengindikasikan ketidaksamaan antarkelompok.
2.13.5 Matriks Transisi Temporal
Setelah hasil clustering ditetapkan per situasi waktu, diperlukan alat untuk menganalisis dinamika perubahan keanggotaan objek antar-cluster. Matriks transisi M berdimensi
k ×k (dengan k adalah jumlah cluster) disusun untuk memetakan kondisi dari situasi k1 ke
situasi k2. Elemen Mxy mendeskripsikan jumlah objek yang awalnya berada pada klaster
x di situasi k1 namun berpindah ke klaster y di situasi k2:
T =


M1,1 M1,2 · · · M1,k
M2,1 M2,2 · · · M2,k
.
.
.
.
.
.
.
.
.
.
.
.
Mk,1 Mk,2 · · · Mk,k


(2.67)
Diagonal utama (Mj,j ) merepresentasikan objek yang tidak berpindah klaster (stabil), sementara elemen off-diagonal merepresentasikan objek yang mengalami perubahan keanggotaan klaster antar-situasi.
2.13.6 Penyelarasan Label Klaster Lintas Waktu
Sifat K-Means yang non-deterministik pada proses inisialisasi awal centroid menyebabkan pelabelan klaster bersifat arbitrer atau acak. Akibatnya, kelompok objek yang
secara spasial berada pada kuadran yang sama dapat memperoleh label bilangan yang
berbeda pada pengujian di tahun yang berbeda (misalnya, identitas klaster 1 pada tahun
2020 bergeser menjadi klaster 2 pada tahun 2021). Masalah ini dikenal dengan istilah
Label Switching atau Klaster Correspondence (Kuhn, 1955).
Untuk menganalisis transisi kondisi (transition matrix) pada data longitudinal, pencocokan label antar-periode waktu merupakan prosedur prapemrosesan yang wajib dilakukan.
Pencocokan ini menggunakan optimasi kombinatorial berlandaskan jarak centroid. Misalkan c1i adalah centroid klaster ke-i pada tahun dasar, dan c2j adalah centroid klaster ke-j
Universitas Indonesia

42
pada tahun uji. Jarak Euclidean antar centroid dapat diformulasikan sebagai:
D(c1i, c2j ) =
vuutX
p
h=1
(c1ih − c2jh)
2
(2.68)
Matriks jarak ini kemudian diselesaikan menggunakan Hungarian Algorithm (Kuhn,
1955) untuk menemukan kombinasi pemetaan (minimum weight bipartite matching)
sedemikian sehingga total biaya jarak atau fungsi objektif terminimisasi:
minX
k
i=1
X
k
j=1
D(c1i, c2j ) · xij (2.69)
dengan indikator biner xij ∈ {0, 1} mengindikasikan apakah centroid i dari kondisi awal
dipasangkan dengan centroid j pada kondisi tujuan. Dengan metode ini, keselarasan penamaan Klaster dapat dijamin tetap konsisten untuk merepresentasikan profil kesejahteraan
yang sama.
2.13.7 Kategorisasi dan Perankingan Cluster
Setelah proses penyelarasan label selesai dilakukan, cluster-klaster tersebut masih
memiliki label angka yang bersifat nominal (1, 2, . . . , K). Untuk memberikan interpretasi
yang bermakna secara ekonomi (misalnya: Rendah, Sedang, Tinggi), diperlukan kriteria
pengurutan yang objektif. Dalam penelitian ini, kriteria yang digunakan adalah nilai
rata-rata skor komponen utama pertama (PC1) pada ruang kompromi STATIS.
Misalkan fi,1 adalah skor objek ke-i pada komponen utama pertama dari matriks
kompromi Fc. Rata-rata skor untuk klaster ke-k didefinisikan sebagai:
PC1,k =
1
nk
X
i∈Ck
fi,1 (2.70)
dengan Ck adalah himpunan indeks objek dalam Klaster ke-k dan nk adalah jumlah
anggotanya. Label peringkat (1 untuk Rendah, hingga K untuk Tinggi) diberikan
berdasarkan urutan menaik dari nilai PC1,k. Hal ini didasari oleh sifat PC1 dalam analisis
ini yang merepresentasikan dimensi utama kesejahteraan rakyat, di mana nilai skor yang
lebih besar berkorespondensi dengan capaian indikator yang lebih baik secara rata-rata.
Universitas Indonesia

43
2.14 Uji Autokorelasi Spasial (Global Moran’s I)
Autokorelasi spasial mengacu pada korelasi antara nilai suatu atribut di satu lokasi
dengan nilai atribut yang sama di lokasi lain yang berdekatan, suatu manifestasi dari
Hukum Pertama Geografi Tobler yang menyatakan bahwa “segalanya berhubungan dengan
segalanya yang lain, tetapi hal-hal yang berdekatan lebih berhubungan daripada yang jauh”
(Johnson & Wichern, 2007).
2.14.1 Matriks Pembobot Spasial
Matriks Pembobot Spasial (Spatial Weight Matrix) Wc berukuran N × N merepresentasikan struktur ketetanggaan antar-wilayah. Elemen wij mendefinisikan seberapa besar
pengaruh wilayah j terhadap wilayah i (Anselin, 1988). Salah satu definisi ketetanggaan
yang umum digunakan adalah Queen Contiguity: wilayah i dan j dianggap bertetangga
jika mereka berbagi setidaknya satu titik batas (sisi atau sudut):
wij =



1 jika wilayah i dan j bertetangga
0 selainnya
, wii = 0 (2.71)
Untuk menghilangkan pengaruh perbedaan jumlah tetangga antar-wilayah, dilakukan
standardisasi baris (row standardization):
w
∗
ij =
wij
X
N
k=1
wik
(2.72)
Bentuk matriks pembobot terstandardisasi-baris adalah W∗ = (w
∗
ij ) di mana setiap baris
berjumlah 1. Bentuk ini memastikan bahwa rata-rata berbobot spasial merupakan rata-rata
sederhana dari nilai-nilai tetangga.
2.14.2 Analisis Autokorelasi Spasial (Global Moran’s I)
Sejalan dengan hasil pengelompokan pada skor kompromi, dilakukan pengujian untuk
mendeteksi apakah terdapat ketergantungan spasial secara keseluruhan. Pengujian ini
merujuk pada koefisien Global Moran’s I yang divalidasi menggunakan pendekatan nonparametrik (permutasi) melalui simulasi Monte Carlo. Pendekatan ini dilakukan dengan
mengacak posisi nilai atribut secara spasial berkali-kali untuk membangun distribusi null,
sehingga p-value diperoleh secara empiris tanpa bergantung pada asumsi normalitas data.
Universitas Indonesia

44
Hal ini menjamin validitas pengujian pada data regional yang memiliki keterbatasan jumlah
amatan (n = 41).
Definisi 2.22 (Global Moran’s I). Misalkan xi adalah nilai suatu atribut pada wilayah ke-i.
Global Moran’s I menggunakan matriks pembobot awal (belum terstandardisasi baris)
Wc dengan elemen wij , dan didefinisikan sebagai:
I = X
N
i
X
j
wij
·
X
i
X
j
wij (xi − x¯)(xj − x¯)
X
i
(xi − x¯)
2
(2.73)
dengan N adalah jumlah wilayah dan x¯ adalah rata-rata atribut. Ekspektasi di bawah
hipotesis nol (tidak ada autokorelasi spasial) adalah E[I] = −1/(N − 1).
• Nilai I signifikan positif (I ≫ E[I]): spatial clustering, yaitu wilayah bernilai tinggi
bertetangga dengan wilayah bernilai tinggi.
• Nilai I signifikan negatif (I ≪ E[I]): spatial dispersion, yaitu pola menyebar.
• Nilai I ≈ E[I]: pola acak, tidak ada autokorelasi spasial.
Signifikansi statistik diuji menggunakan pendekatan permutasi (conditional randomization) yang tidak memerlukan asumsi normalitas. Hal ini serupa dengan logika uji
permutasi pada stabilitas struktur STATIS.
2.15 Visualisasi Interaktif menggunakan Plotly dan Dash
Kompleksitas informasi yang dihasilkan oleh metode-metode yang telah dijelaskan
di subbab sebelumnya menuntut media penyajian yang lebih adaptif dibandingkan grafik
statis. Untuk menjawab kebutuhan ini, tersedia pustaka grafis Plotly (Plotly Technologies
Inc., 2015) dalam kerangka kerja aplikasi web Dash (Plotly Technologies Inc., 2024).
Plotly adalah pustaka visualisasi data deklaratif yang memungkinkan pembuatan grafik
interaktif saat diakses, termasuk grafik tiga dimensi (3D) yang dapat diputar dan diperbesar
secara langsung. Sementara itu, Dash menerapkan paradigma pemrograman reaktif (reactive programming) yang memungkinkan grafik Plotly merespons secara dinamis terhadap
masukan (input) pengguna. Secara konseptual, integrasi keduanya membentuk alur kerja
analitik interaktif yang digambarkan melalui flowchart pada Gambar 2.2 berikut.
Universitas Indonesia

45
Mulai
Input (Filter, Parameter & Dropdown)
Trigger Callback
(Dash Reactive Mechanism)
Komputasi & Pemrosesan Data
(Z-Score, STATIS Core Calculation)
Analisis Statistik Lanjut
(Clustering, Moran’s I, Permutation)
Generasi Objek Visualisasi
(Plotly Graph Objects)
Update Dashboard
(Output Rendering)
Dashboard Teraktualisasi
Interaksi
Berlanjut
Gambar 2.2: Flowchart Mekanisme Interaktivitas Dash-Plotly
Pendekatan ini memfasilitasi Analisis Data Eksploratori (Exploratory Data Analysis/EDA) yang mendalam dan interaktif yang memungkinkan peneliti melakukan eksplorasi
data secara real-time.
2.16 Kerangka Berpikir
Kerangka berpikir dalam penelitian ini disusun untuk mengintegrasikan berbagai
metode analisis multivariat dan analisis spasial dalam memahami dinamika kesejahteraan
rakyat di sebagian Pulau Jawa. Alur logika tersebut dimulai dari identifikasi fenomena
ketimpangan spasio-temporal pada indikator sosial-ekonomi, yang kemudian diproses
melalui tahapan-tahapan sistematis sebagai berikut:
1. Input Data: Pengumpulan data indikator kesejahteraan rakyat dari 41 kabupaten/kota di sebagian Pulau Jawa dalam rentang waktu tertentu, membentuk sebuah
struktur data tiga-arah (cube data).
2. Pra-pemrosesan: Melakukan standardisasi data menjadi Z-score untuk menghilangkan bias satuan dan skala, selanjutnya dilakukan uji kelayakan multivariat (KMO
dan Bartlett) untuk memastikan adanya struktur korelasi yang signifikan.
Universitas Indonesia

46
3. Analisis STATIS: Menerapkan tiga tahapan utama (intra-struktur, inter-struktur, dan
kompromi) untuk mendapatkan representasi gabungan yang paling stabil pada lintas
waktu tanpa mengabaikan variansi temporal dari masing-masing tahun.
4. Visualisasi Dinamis dan Pengelompokan: Mengekstraksi lintasan (trajectory)
setiap objek untuk melihat pergerakan profil kesejahteraan secara individual. Hasil
ini kemudian diklusterisasi menggunakan metode K-Means berbasis jarak trajektori
untuk membangun tipologi wilayah yang memiliki karakteristik perubahan serupa.
5. Validasi Spasial: Melakukan pengujian autokorelasi spasial Global Moran’s I untuk
memverifikasi apakah kelompok-kelompok yang terbentuk secara statistik memiliki
pola dependensi geografis yang nyata atau bersifat acak.
6. Output Dashboard: Mengintegrasikan seluruh hasil komputasi ke dalam platform
visualisasi berbasis web menggunakan Dash Plotly, yang memungkinkan interpretasi
data secara interaktif dan adaptif bagi pemangku kebijakan.
Dengan kerangka ini, analisis tidak hanya berhenti pada deskripsi angka rata-rata,
melainkan mampu mengungkap pola evolusi dan interdependensi antar-wilayah secara
mendalam.
Universitas Indonesia

BAB 3
METODOLOGI PENELITIAN
Bab ini berisi uraian mengenai tahapan-tahapan metodologi yang dilakukan dalam
penelitian di Provinsi Banten, DKI Jakarta, dan Jawa Barat. Uraian tersebut mencakup
penentuan jenis dan desain penelitian, populasi dan sampel penelitian, teknik pengumpulan data, serta teknik analisis data spasio-temporal secara komputasional menggunakan
algoritma STATIS dan analisis spasial.
3.1 Statistika Deskriptif
Analisis diawali dengan ringkasan statistik untuk memberikan gambaran umum mengenai distribusi data yang menggambarkan kesejahteraan rakyat pada seluruh kabupaten/kota
di Provinsi Banten, DKI Jakarta, dan Jawa Barat selama periode 2020 hingga 2024. Indikator yang dimaksud dalam hal ini adalah 12 indikator kesejahteraan rakyat sebagaimana
diidentifikasi pada Bagian 1.5, yang mencakup IPM, IPG, PPM, PDRB ADHK, RLS, AHH
Laki-laki, AHH Perempuan, GK, PKKP, TPT, TPAK, dan Gini Ratio. Prosedur ini melibatkan penghitungan ukuran pemusatan dan penyebaran data sebagaimana didefinisikan
pada Bab 2 yang meliputi:
1. Rata-rata (x¯): Digunakan untuk melihat level umum pencapaian indikator (merujuk
Persamaan 2.6).
2. Simpangan Baku (s): Untuk mengukur tingkat kesenjangan atau variabilitas data
antar wilayah (merujuk Persamaan 2.8).
3. Nilai Minimum dan Maksimum: Untuk mengidentifikasi rentang sebaran nilai
indikator dengan performa ekstrim (merujuk Persamaan 2.9).
4. Kuartil (Q1 dan Q3) serta Median (Q2): Untuk melihat sebaran data berdasarkan
pemisahan persentase kelompok amatan (merujuk Persamaan 2.7).
5. Kemiringan (Skewness): Untuk mengukur tingkat asimetri distribusi data (merujuk
Persamaan 2.12).
6. Keruncingan (Kurtosis): Untuk mengukur ketebalan ekor dan keruncingan relatif
distribusi data dibandingkan distribusi normal (merujuk Persamaan 2.13).
7. Koefisien Variasi (CV ): Digunakan untuk membandingkan tingkat keragaman
antar-indikator yang memiliki skala berbeda (merujuk Persamaan 2.11).
47 Universitas Indonesia

48
3.2 Jenis Penelitian
Penelitian ini menggunakan pendekatan kuantitatif dengan jenis penelitian deskriptifeksploratif. Secara khusus, penelitian ini memanfaatkan teknik analisis multivariat lanjut
(analisis tiga-arah) untuk mengeksplorasi struktur data longitudinal indikator pembangunan strategis di Provinsi Banten, DKI Jakarta, dan Jawa Barat. Pendekatan ini dipilih
karena data yang diolah berupa data statistik sekunder dan akan dianalisis secara objektif
menggunakan algoritma STATIS, K-Means clustering, dan analisis autokorelasi spasial
untuk menarik kesimpulan mengenai stabilitas dan dinamika wilayah.
3.3 Desain Penelitian
Desain penelitian ini dirancang secara sistematis untuk menjawab tantangan data
spasio-temporal. Tahapan penelitian dimulai dari input data, pra-pemrosesan, hingga
pembentukan hasil akhir berupa visualisasi interaktif. Rancangan jalannya penelitian
secara sistematis ditunjukkan pada Gambar 3.1 berikut.
Universitas Indonesia

49
Mulai
Input Data
(Upload dataset, pilih
indikator & tahun)
Eksplorasi Data Deskriptif
Hitung Bobot Demografis mi & Matriks D
Percabangan Zk:
STATIS: z
S = (x −
µjk)/σjk (terbobot mi)
PCA: z
P = (x − x¯)/s (biasa)
Dynamic Biplot
Z
PCA
k → SVD per-Tahun
→ Fk, Gk (per-tahun)
Uji KMO ≥ 0.5?
Uji Bartlett
p < 0.05?
Analisis STATIS
Wk → H, RV →
αk → Wc → SVD
Uji Permutasi
(Signifikansi Stabilitas Struktur)
Struktur Signifikan
(p < 0.05)?
Compromise Biplot &
Compromise Loading
Penentuan Jumlah Kluster Optimal
(Elbow + Silhouette)
K-Means Clustering
pada Skor Kompromi
Penyelarasan Label
(Hungarian Algorithm)
Perankingan Level
(Skor Rata-rata PC1)
Uji Autokorelasi Spasial
(Global Moran’s I)
Visualisasi Peta Kluster & Trajektori Wilayah
Selesai
Ganti/Tambah
Indikator
Struktur Tidak
Stabil – Selesai
PCA
STATIS
Ya
Ya
Ya
Tidak
Tidak
Tidak
Gambar 3.1: Flowchart Alur Metodologi Penelitian
Universitas Indonesia

50
3.4 Populasi dan Sampel
Populasi dalam penelitian ini adalah seluruh wilayah administratif setingkat kabupaten
dan kota di Pulau Jawa. Mengingat ketersediaan data dan fokus analisis pada wilayah yang
secara fungsional terhubung dalam koridor metropolitan Jawa bagian barat, penelitian ini
mencakup seluruh kabupaten/kota di tiga provinsi berikut:
1. Unit Observasi (n): 41 wilayah yang terdiri dari seluruh kabupaten/kota di Provinsi
Banten, DKI Jakarta, dan Jawa Barat. Pemilihan ketiga provinsi ini didasarkan
pada kelengkapan data yang dipublikasikan secara konsisten oleh BPS (purposive
sampling pada tingkat provinsi), sehingga seluruh wilayah di dalamnya tercakup
secara penuh (complete enumeration) tanpa pengambilan sampel lebih lanjut.
2. Jumlah Situasi (K): 5 tahun pengamatan (2020 hingga 2024).
3. Jumlah Indikator (p): 11 indikator kesejahteraan rakyat sebagai input analisis
STATIS, ditambah 2 variabel penunjang (IPM dan JP) yang tidak dimasukkan
sebagai input.
4. Sampel Ilustratif (Bab 3): Untuk keperluan rekonstruksi algoritmik dan validasi
numerik pada Bagian 3.7.1, dipilih 4 wilayah secara acak sederhana (simple random
sampling): Kota Jakarta Pusat, Kabupaten Bandung Barat, Kota Tasikmalaya, dan
Kota Cilegon. Sampel ini tidak digunakan untuk membangun model, melainkan
semata-mata sebagai wahana ilustrasi langkah-langkah perhitungan manual STATIS
dengan n = 4 wilayah, K = 3 tahun (2020–2022), dan p = 3 indikator (PPM, RLS,
TPT).
3.5 Teknik Pengumpulan Data
Data dalam penelitian ini merupakan data sekunder yang dikumpulkan melalui studi
dokumentasi pada publikasi resmi Badan Pusat Statistik (BPS).
1. Sumber Data: Website resmi BPS Pusat dan BPS Provinsi yang diperoleh melalui
portal Open Data atau publikasi Provinsi Dalam Angka. Data untuk Provinsi Banten,
DKI Jakarta, dan Jawa Barat periode 2020 hingga 2024 diperoleh dari publikasi
resmi BPS masing-masing provinsi (Badan Pusat Statistik Provinsi Banten, 2020,
2021, 2022, 2023, 2024; Badan Pusat Statistik Provinsi DKI Jakarta, 2020, 2021,
2022, 2023, 2024; Badan Pusat Statistik Provinsi Jawa Barat, 2020, 2021, 2022,
2023, 2024).
Universitas Indonesia

51
2. Metode Pengumpulan Data: Pengumpulan data dilakukan dengan mengunduh
berkas metadata dan tabel statistik dalam format spreadsheet (.xlsx dan .csv).
3. Indikator Penelitian: Indikator penelitian yang dikumpulkan terdiri dari 11 indikator sebagai berikut:
• IPG: Indeks Pembangunan Gender (0 hingga 100).
• PPM: Persentase Penduduk Miskin (%).
• PDRB ADHK: Produk Domestik Regional Bruto Atas Dasar Harga Konstan
tahun 2010 (Juta Rupiah).
• RLS: Rata-rata Lama Sekolah (Tahun).
• AHH L: Angka Harapan Hidup Laki-laki (Tahun).
• AHH P: Angka Harapan Hidup Perempuan (Tahun).
• GK: Garis Kemiskinan (Rupiah/Kapita/Bulan).
• PKKP: Prevalensi Ketidakcukupan Konsumsi Pangan (%).
• TPT: Tingkat Pengangguran Terbuka (%).
• TPAK: Tingkat Partisipasi Angkatan Kerja (%).
• GR: Gini Ratio (0 hingga 1).
Selain itu, dikumpulkan pula data tambahan berikut yang tidak dimasukkan sebagai
input analisis STATIS dan Dynamic Biplot:
• IPM: Indeks Pembangunan Manusia (0 hingga 100), khusus untuk pembanding
eksternal (benchmark) untuk mengevaluasi posisi klaster tersebut terhadap
standar UNDP (United Nations Development Programme). Hal ini dilakukan
agar klaster yang terbentuk murni merepresentasikan kemiripan struktur dari
11 indikator yang telah dijelaskan sebelumnya.
• JP (Jumlah Penduduk): Banyaknya jiwa per wilayah (Orang), khusus untuk pembentukan bobot demografis (mi) sesuai dengan Persamaan 2.3 dan
Persamaan 2.4.
3.5.1 Pra-pemrosesan Data, Penskalaan Visual, dan Parameter Analisis
Dalam rangka menjamin reproduksibilitas dan ketajaman interpretasi hasil, tahapan
pra-pemrosesan dan parameter analisis dilakukan sebagai berikut:
Universitas Indonesia

52
1. Standardisasi per-tahun (Z-score): Setiap matriks tahunan distandardisasi secara
terpisah sehingga komparasi antarvariabel di dalam tahun menjadi setara. Terdapat
dua skema standardisasi yang digunakan secara paralel sesuai jalur analisis (lihat
Subbagian 3.7.1.2): (1) z
PCA = (x − x¯)/s untuk Dynamic Biplot — menggunakan
mean aritmetika biasa dengan semua wilayah berbobot setara; (2) z
STATIS = (x −
µjk)/σjk untuk STATIS — menggunakan mean dan simpangan baku terbobot mi
agar wilayah berpenduduk besar berkontribusi proporsional.
2. Bobot demografis (mi): Untuk beberapa analisis (mis. pembentukan matriks bobot
objek), dihitung bobot demografis berdasarkan jumlah penduduk tiap wilayah sesuai
Persamaan 2.3. Bobot ini dipakai untuk memperhitungkan kontribusi relatif populasi
dalam pembentukan kompromi.
3. Uji Kelayakan Multivariat: Sebelum menerapkan STATIS, dilakukan Uji KMO
dan Uji Bartlett untuk memastikan kecukupan korelasi antarvariabel; ambang batas
yang digunakan adalah KMO ≥ 0,5 dan Bartlett p < 0,05.
4. Penskalaan Loadings untuk Visualisasi: Untuk tujuan representasi grafis (biplot/trajectory), panjang vektor loading disesuaikan mengikuti Abdi et al. (2012)
(lihat Persamaan 2.48): setiap elemen loading intra g
(asli)
jℓ pada dimensi ℓ dikalikan
dengan nilai singular kompromi σ
(c)
ℓ =
q
λ
(c)
ℓ
, menghasilkan nilai visual:
g
(visual)
jℓ = g
(asli)
jℓ ×
q
λ
(c)
ℓ
(3.1)
Pada data sampel, λ
(c)
1 = 2,3008 dan λ
(c)
2 = 0,6992, sehingga δ1 =
√
2,3008 =
1,5168 dan δ2 =
√
0,6992 = 0,8362. Sebagai ilustrasi, loading intra PPM tahun
2020 sebelum penskalaan adalah g
(asli)
PPM,1 = 0,8306 (PC1) dan g
(asli)
PPM,2 = 0,5569 (PC2),
sehingga:
g
(visual)
PPM,1,2020 = 0,8306 × 1,5168 = 1,2599
g
(visual)
PPM,2,2020 = 0,5569 × 0,8362 = 0,4657
Apabila konfigurasi mirror PC1 aktif (seperti pada Skenario 1 studi kasus ini),
nilai PC1 Visual selanjutnya dinegatifkan: g
(visual)
PPM,1 = −1,2599. Skor intra Fk tidak
mengalami perubahan amplitudo — hanya tanda yang dapat berubah bila mirror
aktif. Penyesuaian ini hanya berlaku untuk plotting dan tidak mengubah nilai analitis
atau kesimpulan inferensial.
Universitas Indonesia

53
5. Uji Permutasi Stabilitas STATIS: Untuk mengevaluasi signifikansi kestabilan
struktur kompromi digunakan uji permutasi Monte Carlo terhadap matriks interstruktur sebanyak Nperm percobaan (default 1000). Nilai p dihitung sebagai proporsi
permutasi yang menghasilkan statistik lebih ekstrim daripada observasi.
6. Pemilihan Jumlah Klaster: Penentuan jumlah klaster (k) untuk K-Means dilakukan
melalui kombinasi metode extitelbow dan skor extitsilhouette; nilai default yang
diuji di skrip adalah k = 2–5 dan hasil akhir dipilih berdasarkan keseimbangan
antara compactness dan interpretabilitas.
7. Uji Autokorelasi Spasial: Untuk menguji keberadaan clustering spasial pada skor
kompromi, dihitung Global Moran’s I dengan prosedur permutasi Monte Carlo pada
bobot spasial yang dibentuk dari file shapefile administratif (lihat lampiran teknis
untuk detail pembentukan matriks bobot).
3.6 Validasi Instrumen
Karena data yang digunakan berasal dari lembaga resmi negara (BPS) yang telah
memiliki protokol sensus dan survei (seperti Susenas dan Sakernas) yang terstandardisasi,
validitas dan reliabilitas instrumen pengumpulan data primer dianggap sudah terpenuhi.
Namun, pada tahap penelitian ini dilakukan validasi data sekunder melalui:
1. Penyelarasan Metadata: Memastikan definisi operasional indikator tetap konsisten
di seluruh provinsi dan tahun pengamatan.
2. Penataan Data (Data Wrangling): Data mentah yang diperoleh dari berbagai publikasi BPS direstrukturisasi ke dalam format tiga-arah yang diwajibkan oleh sistem
komputasi: setiap tahun disimpan sebagai lembar (sheet) tersendiri dalam berkas
.xlsx, dengan baris mewakili wilayah dan kolom mewakili indikator. Seluruh
41 wilayah pada ketiga provinsi tersedia lengkap untuk semua indikator dan tahun
yang diamati, sehingga tidak diperlukan penghapusan baris atau kolom. Data hasil
penataan disimpan dalam berkas Data Final.xlsx (lihat Lampiran 5.2) yang
selanjutnya diunggah langsung ke dalam sistem dashboard untuk diproses.
3. Standardisasi Z-score Terbobot: Karena indikator yang digunakan memiliki satuan
yang beragam (persen, rupiah, tahun), dilakukan transformasi Z-score terbobot agar
setiap indikator dapat dibandingkan dalam ruang terboboti D yang konsisten. Prosedur ini mengikuti Persamaan 2.5 yang mengggunakan rata-rata populasi terbobot
µjk (Persamaan 2.3) dan simpangan baku populasi terbobot σjk (Persamaan 2.4)
Universitas Indonesia

54
yang diterapkan pada setiap tabel situasi k, sehingga setiap indikator pada setiap
tahun memiliki rata-rata terbobot nol dan varians terbobot satu.
Untuk memastikan data layak dianalisis secara multivariat, dilakukan serangkaian
pengujian kecukupan sampling dan korelasi pada setiap periode pengamatan:
1. Indeks Kaiser-Meyer-Olkin (KMO): Pengujian ini dilakukan untuk menilai kecukupan sampel (sampling adequacy) dari 11 indikator yang digunakan untuk mengukur kesejahteraan rakyat. Penghitungan KMO merujuk pada rasio korelasi sederhana dan korelasi parsial sesuai Persamaan 2.15. Nilai KMO ≥ 0, 50 menjadi syarat
mutlak agar indikator-indikator tersebut dianggap memiliki informasi bersama yang
cukup untuk diringkas melalui analisis dimensi tiga-arah.
2. Uji Bartlett’s Test of Sphericity: Uji ini digunakan untuk memverifikasi adanya
struktur korelasi yang signifikan antar indikator (menguji hipotesis bahwa matriks korelasi bukan matriks identitas). Statistik uji χ
2 dihitung berdasarkan determinan matriks korelasi sebagaimana didefinisikan pada Persamaan 2.17. Jika p-value < 0, 05,
maka indikator-indikator tersebut dikonfirmasi memiliki hubungan interdependensi
yang kuat, sehingga valid untuk diproses lebih lanjut menggunakan metode STATIS
guna mengekstraksi struktur kompromi lintas waktu.
Penerapan kedua uji di atas pada dataset sampel (n = 4 wilayah, p = 3 indikator,
K = 3 tahun) menghasilkan nilai sebagaimana dirangkum pada Tabel 3.1.
Tabel 3.1: Hasil Uji KMO dan Uji Bartlett per Tahun
Tahun Nilai KMO Interpretasi Nilai KMO χ
2 Bartlett df P-Value Bartlett Keputusan
2020 0,2502 Unacceptable (Tidak diterima) 3,6300 3 0,3038 Periksa
2021 0,2620 Unacceptable (Tidak diterima) 5 3 0,1721 Periksa
2022 0,2917 Unacceptable (Tidak diterima) 7,2000 3 0,0659 Periksa
Catatan: Nilai KMO yang rendah (≈ 0,25–0,29) dan p-value Bartlett yang tidak signifikan
pada dataset sampel merupakan konsekuensi langsung dari jumlah observasi (n = 4)
yang sangat terbatas relatif terhadap jumlah indikator (p = 3). Rasio n/p ≈ 1,3 ini jauh
di bawah anjuran minimal n/p ≥ 5 dalam literatur analisis faktor. Pada dataset aktual
yang terdiri dari 41 kabupaten/kota dengan 11 indikator, kedua uji menghasilkan nilai
yang memenuhi ambang batas yang disyaratkan, sebagaimana disajikan secara lengkap
dalam Bab 4.
Contoh penghitungan: Statistik Bartlett untuk tahun 2020 menggunakan Persamaan 2.17
dengan n = 4, p = 3. Determinan matriks korelasi |R2020| = 0,0441 (dihitung dari
Universitas Indonesia

55
matriks S2020 pada Tabel 3.9). Substitusi ke rumus:
χ
2 = −

n − 1 −
2p + 5
6

ln |R| = −

4 − 1 −
2(3) + 5
6

ln(0,0441)
= −

3 −
11
6

(−3,1179) = −(1,1667)(−3,1179) = 3,6300
dengan df = p(p − 1)/2 = 3(2)/2 = 3 dan p-value = 0,3038. Statistik KMO diperoleh
dari Persamaan 2.15 yang membandingkan korelasi bivariate dengan korelasi parsial
antar-indikator; untuk tahun 2020 diperoleh KMO2020 = 0,2502.
3.7 Teknik Analisis Data
Teknik analisis data mengacu pada algoritma STATIS (Structuration des Tableaux a`
Trois Indices de la Statistique) untuk mengelola kompleksitas data tiga-arah serta analisis
autokorelasi spasial. Agar tidak terjadi kerancuan persepsi, penelitian ini membedakan
dua konsep ”tiga-arah” yang saling berkaitan:
• Struktur Data Tiga-Arah (Input Cube): Merupakan struktur data mentah yang
dianalisis, yang terdiri dari 41 Kabupaten/Kota sebagai objek (n), 11 Indikator
Kesejahteraan Rakyat sebagai indikator (p), dan 5 tahun pengamatan sebagai situasi
waktu (K). Seluruh pengujian awal (KMO dan Bartlett) dilakukan pada dimensi
indikator (p) ini.
• Visualisasi Dinamis (Output Trajectory): Merupakan hasil reduksi dimensi oleh
STATIS, di mana 11 indikator tersebut diringkas menjadi 2 Komponen Utama (PC1
dan PC2), sementara dimensi situasi (K) direpresentasikan melalui pergeseran atau
lintasan wilayah di dalam ruang PC1 dan PC2 tersebut.
3.7.1 Prosedur Komputasi dan Ilustrasi Numerik STATIS
Bagian ini menyajikan rekonstruksi algoritmik secara sistematis dengan menyajikan
ilustrasi penghitungan manual (step-by-step) untuk memvalidasi fungsionalitas sistem
komputasi yang dibangun. Ilustrasi dilakukan menggunakan dataset sampel yang terdiri
dari 4 wilayah (n = 4) selama 3 periode waktu (K = 3) dengan 3 indikator utama
(p = 3): Persentase Penduduk Miskin (PPM), Rata-rata Lama Sekolah (RLS), dan Tingkat
Pengangguran Terbuka (TPT). Wilayah sampel meliputi: (1) Kota Jakarta Pusat, (2)
Kabupaten Bandung Barat, (3) Kota Tasikmalaya, dan (4) Kota Cilegon. Data lengkap
sampel yang digunakan dalam ilustrasi ini disajikan pada Tabel 3.2.
Universitas Indonesia

56
Tabel 3.2: Data Mentah Sampel Ilustratif: 4 Wilayah × 3 Tahun × 5 Variabel
No Provinsi Kabupaten/Kota 2020 2021 2022
PPM RLS TPT IPM JP PPM RLS TPT IPM JP PPM RLS TPT IPM JP
1 DKI
JAKARTA
Kota Jakarta Pusat 4,51 11,38 10,97 81,39 1.056.896 4,94 11,39 7,75 81,79 1.057.465 4,90 11,53 5,88 82,13 1.053.482
2 JAWA
BARAT
Bandung Barat 10,49 8,19 12,25 68,08 1.782.190 11,30 8,20 11,65 68,30 1.808.420 10,82 8,22 9,63 69,04 1.834.230
3 JAWA
BARAT
Kota Tasikmalaya 12,97 9,33 7,99 72,97 713.670 13,13 9,52 7,66 73,20 723.100 12,72 9,53 6,62 73,83 732.480
4 BANTEN Kota Cilegon 3,69 9,87 12,69 73,05 433,76 4,24 10,08 10,13 73,30 439,48 3,64 10,34 8,10 73,95 445,06
Catatan: PPM = Persentase Penduduk Miskin (%); RLS = Rata-rata Lama Sekolah (Tahun); TPT = Tingkat Pengangguran Terbuka
(%); IPM = Indeks Pembangunan Manusia; JP = Jumlah Penduduk. Analisis STATIS menggunakan tiga indikator utama: PPM, RLS,
TPT.
Catatan penyajian: Seluruh objek matematis pada bagian ini, termasuk di dalamnya Wk,
Sk, H, Fc, Gc, Fk, dan Gk merupakan sebuah bentuk matriks. Penyajian dalam format
tabel pada penelitian ini semata-mata untuk memudahkan pembacaan nilai numerik; notasi
dan sifat matriks sepenuhnya berlaku dalam konteks komputasi dan derivasi analitis.
3.7.1.1 Tahap 0: Pra-pemrosesan dan Bobot Demografis
Analisis diawali dengan ringkasan statistik untuk periode 2020 hingga 2022 (Tabel 3.3–
3.5). Statistika deskriptif tanpa bobot disajikan pada tabel berikut.
Tabel 3.3: Statistika Deskriptif (tanpa Bobot) Tahun 2020
Indikator Mean (µ) Std Dev (σ) Min Q1 Median Q3 Max Skew (γ1) Kurt (γ2) CV (%)
TPT 10,9750 1,8356 7,9900 10,2250 11,6100 12,3600 12,6900 -0,7875 -1,0029 16,7200
PPM 7,9150 3,9252 3,6900 4,3050 7,5000 11,1100 12,9700 0,1296 -1,7889 49,5900
IPM 73,8725 4,7843 68,0800 71,7475 73,0100 75,1350 81,3900 0,5232 -0,9383 6,4800
RLS 9,6925 1,1476 8,1900 9,0450 9,6000 10,2475 11,3800 0,2269 -1,0941 11,8400
JP 888297,44 641534,50 433,76 535360,94 885283,00 1238219,50 1782190,00 0,0131 -1,1379 72,2200
Tabel 3.4: Statistika Deskriptif (tanpa Bobot) Tahun 2021
Indikator Mean (µ) Std Dev (σ) Min Q1 Median Q3 Max Skew (γ1) Kurt (γ2) CV (%)
TPT 9,2975 1,6810 7,6600 7,7275 8,9400 10,5100 11,6500 0,2894 -1,6215 18,0800
PPM 8,4025 3,8749 4,2400 4,7650 8,1200 11,7575 13,1300 0,0703 -1,8757 46,1200
IPM 74,1475 4,8533 68,3000 71,9750 73,2500 75,4225 81,7900 0,5357 -0,9353 6,5500
RLS 9,7975 1,1451 8,2000 9,1900 9,8000 10,4075 11,3900 -0,0062 -1,1160 11,6900
JP 897356,12 650095,55 439,48 542434,87 890282,50 1245203,75 1808420,00 0,0305 -1,1276 72,4500
Tabel 3.5: Statistika Deskriptif (tanpa Bobot) Tahun 2022
Indikator Mean (µ) Std Dev (σ) Min Q1 Median Q3 Max Skew (γ1) Kurt (γ2) CV (%)
TPT 7,5575 1,4390 5,8800 6,4350 7,3600 8,4825 9,6300 0,2951 -1,4124 19,0400
PPM 8,0200 3,8356 3,6400 4,5850 7,8600 11,2950 12,7200 0,0504 -1,8309 47,8300
IPM 74,7375 4,7052 69,0400 72,6325 73,8900 75,9950 82,1300 0,5227 -0,9386 6,3000
RLS 9,9050 1,2052 8,2200 9,2025 9,9350 10,6375 11,5300 -0,0662 -1,2117 12,1700
JP 905159,27 658311,85 445,06 549471,27 892981,00 1248669,00 1834230,00 0,0522 -1,1146 72,7300
Catatan: Koefisien Variasi (CV) menunjukkan rasio simpangan baku terhadap rata-rata,
dinyatakan dalam persen; CV yang lebih tinggi mengindikasikan variasi relatif yang lebih
besar antar-wilayah untuk indikator tersebut.
Universitas Indonesia

57
Selain statistika deskriptif biasa, digunakan pula statistika deskriptif terbobot yang
memperhitungkan jumlah penduduk (JP) di wilayah ke-i (mi) sebagai bobot demografis
(Persamaan 2.3 dan 2.4). Hal ini penting karena STATIS bekerja pada ruang terboboti D.
Tabel 3.6: Statistika Deskriptif (dengan Bobot) Tahun 2020 (N = 4)
Indikator Mean (µ) Std Dev (σ) Min Q1 Median Q3 Max Skewness (γ1) Kurtosis (γ2) CV (%)
TPT 11,0148 1,6168 7,9900 10,2250 11,6100 12,3600 12,6900 -1,0949 -0,3590 14,6781
PPM 9,2290 3,1906 3,6900 4,3050 7,5000 11,1100 12,9700 -0,5969 -1,1973 34,5718
IPM 72,9836 5,7374 68,0800 71,7475 73,0100 75,1350 81,3900 0,6111 -1,3746 7,8613
RLS 9,3588 1,3752 8,1900 9,0450 9,6000 10,2475 11,3800 0,6249 -1,3636 14,6938
Tabel 3.7: Statistika Deskriptif (dengan Bobot) Tahun 2021 (N = 4)
Indikator Mean (µ) Std Dev (σ) Min Q1 Median Q3 Max Skewness (γ1) Kurtosis (γ2) CV (%)
TPT 9,6978 1,9684 7,6600 7,7275 8,9400 10,5100 11,6500 -0,0171 -1,9986 20,2969
PPM 9,7961 3,2132 4,2400 4,7650 8,1200 11,7575 13,1300 -0,7398 -1,2057 32,8005
IPM 73,2585 5,8151 68,3000 71,9750 73,2500 75,4225 81,7900 0,6169 -1,3699 7,9378
RLS 9,4051 1,3763 8,2000 9,1900 9,8000 10,4075 11,3900 0,5445 -1,4306 14,6337
Tabel 3.8: Statistika Deskriptif (dengan Bobot) Tahun 2022 (N = 4)
Indikator Mean (µ) Std Dev (σ) Min Q1 Median Q3 Max Skewness (γ1) Kurtosis (γ2) CV (%)
TPT 7,9195 1,7433 5,8800 6,4350 7,3600 8,4825 9,6300 -0,0786 -1,9194 22,0134
PPM 9,4597 3,0329 3,6400 4,5850 7,8600 11,2950 12,7200 -0,7052 -1,2051 32,0610
IPM 73,8586 5,6426 69,0400 72,6325 73,8900 75,9950 82,1300 0,6132 -1,3730 7,6397
RLS 9,4585 1,4273 8,2200 9,2025 9,9350 10,6375 11,5300 0,5706 -1,4083 15,0902
Matriks data mentah Xk untuk setiap situasi tahun k = 2020, 2021, 2022 hanya
mencakup tiga indikator utama (PPM, RLS, TPT) yang digunakan dalam analisis STATIS,
sehingga kolom IPM dan JP tidak termasuk dalam matriks input. Matriks Xk disusun
dengan baris mewakili wilayah dan kolom mewakili indikator, sebagai berikut:
X2020 =


4, 51 11, 38 10, 97
10, 49 8, 19 12, 25
12, 97 9, 33 7, 99
3, 69 9, 87 12, 69


X2021 =


4, 94 11, 39 7, 75
11, 30 8, 20 11, 65
13, 13 9, 52 7, 66
4, 24 10, 08 10, 13


X2022 =


4, 90 11, 53 5, 88
10, 82 8, 22 9, 63
12, 72 9, 53 6, 62
3, 64 10, 34 8, 10


Penghitungan bobot demografis mi untuk setiap wilayah berdasarkan data rata-rata
jumlah penduduk (JP) dari dataset sampel (Tabel 3.3–3.5) dihitung satu per satu menggunakan Persamaan 2.1 sebagai berikut:
m1 =
JP1
P4
i=1 JPi
=
1.055.947, 67
3.587.750, 43
= 0,294320
m2 =
JP2
P4
i=1 JPi
=
1.808.280, 00
3.587.750, 43
= 0,504015
Universitas Indonesia

58
m3 =
JP3
P4
i=1 JPi
=
723.083, 33
3.587.750, 43
= 0,201542
m4 =
JP4
P4
i=1 JPi
=
439, 43
3.587.750, 43
= 0,000122
Verifikasi total bobot: m1+m2+m3+m4 = 0,294320+0,504015+0,201542+0,000122 =
0,999999 ≈ 1 dengan JPi merupakan rata-rata jumlah penduduk untuk wilayah ke-i
selama tiga tahun pengamatan. Verifikasi total bobot memastikan bahwa jumlahnya
mendekati 1, sehingga memenuhi syarat sebagai distribusi probabilitas yang valid untuk
pembentukan matriks bobot objek D dalam ruang terboboti STATIS.
Matriks diagonal D yang dihasilkan adalah:
D = diag(0,294320, 0,504015, 0,201542, 0,000122)
3.7.1.2 Percabangan Standardisasi: STATIS vs. Dynamic Biplot (PCA per-Tahun)
Sebelum melanjutkan ke tahap standardisasi, perlu ditegaskan bahwa penelitian ini
menggunakan dua jalur analisis paralel yang menghasilkan matriks Zk dengan nilai
numerik yang berbeda, meskipun notasinya sama:
1. Dynamic Biplot (PCA per-Tahun) menggunakan standardisasi biasa tanpa bobot.
Rata-rata dan simpangan baku dihitung secara aritmetika biasa, di mana setiap
wilayah diperlakukan setara tanpa memperhatikan besar populasi:
x¯jk =
1
n
Xn
i=1
xijk, sjk =
vuut
1
n
Xn
i=1

xijk − x¯jk2, zPCA
ijk =
xijk − x¯jk
sjk
(3.2)
Implementasinya dalam dashboard menggunakan
sklearn.preprocessing.StandardScaler yang menerapkan formula di atas. Akibatnya, Kota Depok (JP ≈ 2,3 juta) dan Kabupaten Kepulauan
Seribu (JP ≈ 0,02 juta) mendapat bobot yang setara dalam menentukan pusat
distribusi.
2. Analisis STATIS menggunakan standardisasi terbobot berbasis bobot demografis
mi (dihitung pada Tahap 0 di atas). Rata-rata dan simpangan baku dihitung secara
terbobot sehingga wilayah dengan populasi besar berkontribusi proporsional lebih
Universitas Indonesia

59
besar:
µjk =
Xn
i=1
mi xijk, σjk =
vuutXn
i=1
mi

xijk − µjk2
, zSTATIS
ijk =
xijk − µjk
σjk
(3.3)
Standardisasi ini menjamin Pmizijk = 0 dan Pmiz
2
ijk = 1, sehingga skala
konsisten lintas indikator dan lintas tahun dalam ruang terboboti D.
Perbedaan mendasar ini mengakibatkan z
PCA
ijk ̸= z
STATIS
ijk secara numerik, sehingga
matriks kemiripan intra-struktur Wk = ZkZ
T
k
yang digunakan STATIS juga berbeda dari
Zk yang masuk ke SVD pada PCA per-tahun. Bagian berikut menjelaskan standardisasi
terbobot secara detail dengan data sampel.
3.7.1.3 Tahap 1: Standardisasi Data STATIS (Zk)
Proses standardisasi di bawah ini dilakukan untuk setiap indikator pada setiap tahun
secara terpisah, menggunakan rata-rata terbobot µjk dan simpangan baku terbobot σjk
yang dihitung berdasarkan bobot demografis mi. Sebagai contoh, untuk indikator pertama
(PPM) pada tahun 2020 X2020[:, 1] = (4,51; 10,49; 12,97; 3,69)T. Langkah-langkahnya
adalah sebagai berikut:
1. Rata-rata terbobot µjk:
Nilai rata-rata terbobot untuk indikator pertama (PPM) dihitung menggunakan
Persamaan 2.3 sebagai berikut:
µ1,2020 =
X
4
i=1
mi· xi,1,2020
= 0,294320(4,51) + 0,504015(10,49)
+ 0,201542(12,97) + 0,000122(3,69)
= 1,327383 + 5,287117 + 2,613999 + 0,000450 = 9,228957
2. Simpangan baku terbobot σjk:
Nilai simpangan baku terbobot untuk indikator pertama (PPM) dihitung mengguUniversitas Indonesia

60
nakan Persamaan 2.4 sebagai berikut:
σ1,2020 =
vuutX
4
i=1
mi(xi,1,2020 − µ1,2020)
2
=

0,294320(4,51 − 9,228957)2 + 0,504015(10,49 − 9,228957)2
+ 0,201542(12,97 − 9,228957)2 + 0,000122(3,69 − 9,228957)2
1/2
=

0,294320(−4,718957)2 + 0,504015(1,261043)2
+ 0,201542(3,741043)2 + 0,000122(−5,538957)2
1/2
=

0,294320(22,268555) + 0,504015(1,590229)
+ 0,201542(13,995402) + 0,000122(30,680045)1/2
=
p
6,553881 + 0,801500 + 2,820653 + 0,003743
=
p
10,179777 = 3,190576
3. Z-score terbobot zijk:
Nilai Z-score terbobot masing-masing wilayah untuk indikator pertama (PPM)
dihitung menggunakan Persamaan 2.5 sebagai berikut:
z1,1,2020 =
4,51 − 9,228957
3,190576
=
−4,718957
3,190576
= −1,479024 ≈ −1,48
z2,1,2020 =
10,49 − 9,228957
3,190576
=
1,261043
3,190576
= 0,395237 ≈ 0,40
z3,1,2020 =
12,97 − 9,228957
3,190576
=
3,741043
3,190576
= 1,172529 ≈ 1,17
z4,1,2020 =
3,69 − 9,228957
3,190576
=
−5,538957
3,190576
= −1,736036 ≈ −1,74
Prosedur di atas diulang untuk setiap indikator dan tahun, menghasilkan matriks terstandar:
Z2020 =


−1, 48 1, 47 −0, 03
0, 40 −0, 85 0, 76
1, 17 −0, 02 −1, 87
−1, 74 0, 37 1, 04


Z2021 =


−1, 51 1, 44 −0, 99
0, 47 −0, 88 0, 99
1, 04 0, 08 −1, 04
−1, 73 0, 49 0, 22


Z2022 =


−1, 50 1, 45 −1, 17
0, 45 −0, 87 0, 98
1, 07 0, 05 −0, 75
−1, 92 0, 62 0, 10


Universitas Indonesia

61
3.7.1.4 Tahap 2: Matriks Intra-struktur (Wk)
Dari matriks terstandar Zk, dibentuk matriks intra-struktur Wk = ZkZ
T
k
. Sebagai
contoh pada tahun 2020:
W2020[1, 1] = X
3
j=1
z
2
1,j,2020
= (−1,479012)2 + (1,469744)2 + (−0,027681)2
= 2,187478 + 2,160150 + 0,000766
= 4,348394
W2020[1, 2] = X
3
j=1
z1,j,2020 · z2,j,2020
= (−1,479012 × 0,395235) + (1,469744 × −0,849962)
+ (−0,027681 × 0,764025)
= −0,584558 − 1,249227 − 0,021149 = −1,854935
Matriks korelasi terbobot antar-indikator Sk = Z
T
k DZk disajikan pada Tabel 3.9,
Tabel 3.10, dan Tabel 3.11.
Tabel 3.9: Matriks Korelasi Terbobot Indikator Tahun 2020 (S2020)
Indikator PPM RLS TPT
PPM 1 -0,8141 -0,2781
RLS -0,8141 1 -0,3313
TPT -0,2781 -0,3313 1
Tabel 3.10: Matriks Korelasi Terbobot Indikator Tahun 2021 (S2021)
Indikator PPM RLS TPT
PPM 1 -0,8307 0,4576
RLS -0,8307 1 -0,8751
TPT 0,4576 -0,8751 1
Tabel 3.11: Matriks Korelasi Terbobot Indikator Tahun 2022 (S2022)
Indikator PPM RLS TPT
PPM 1 -0,8277 0,5779
RLS -0,8277 1 -0,9363
TPT 0,5779 -0,9363 1
Universitas Indonesia

62
Sebagai pelengkap, Gambar 3.2 menyajikan heatmap korelasi Pearson antar-indikator
untuk data sampel tahun 2020. Nilai pada diagonal menunjukkan nama variabel (korelasi
sempurna = 1,00), sedangkan sel di luar diagonal memuat nilai koefisien Pearson r dan
bintang signifikansi (*** p < 0,001; ** p < 0,01; * p < 0,05; ns tidak signifikan). Warna
biru tua menandakan korelasi positif kuat dan merah tua menandakan korelasi negatif kuat.
Gambar 3.2: Heatmap Korelasi Pearson Indikator Data Sampel Tahun 2020
3.7.1.5 Tahap 3: Analisis Inter-struktur (Koefisien RV)
Kemiripan struktural antar-situasi dikuantifikasi melalui matriks interstruktur H (Persamaan 2.32) dan koefisien RV (Persamaan 2.33). Elemen Hkk′ = tr(WkDWk
′D) yang
secara aljabar setara dengan tr(SkSk
′) di mana Sk = Z
T
k DZk. Karena setelah standardisasi
terbobot diagonal Sk bernilai 1, elemen diagonal H11 diturunkan sebagai berikut:
H11 = tr(S
2
2020) = 12 + 12 + 12 + 2
(−0,8141)2 + (−0,2781)2 + (−0,3313)2

= 3 + 2(0,6628 + 0,0773 + 0,1098) = 4,6998
Matriks H lengkap disajikan pada Tabel 3.12 berikut.
Universitas Indonesia

63
Tabel 3.12: Matriks Interstruktur H W (Identik Kedua Skenario)
Tahun 2020 2021 2022
2020 4,6998 5,0009 5,1018
2021 5,0009 6,3306 6,5334
2022 5,1018 6,5334 6,7915
Substitusi nilai H11 = 4,6998, H12 = 5,0009, H22 = 6,3306 ke Persamaan 2.33:
RV (W2020,W2021) = H12
√
H11 · H22
=
5,0009
√
4,6998 × 6,3306
=
5,0009
√
29,7525
=
5,0009
5,4546
= 0,9168
Koefisien RV lengkap disajikan pada Tabel 3.13 dan divisualisasikan pada Gambar 3.3.
Tabel 3.13: Koefisien RV Wilayah (Identik Kedua Skenario)
Tahun 2020 2021 2022
2020 1 0,9168 0,9030
2021 0,9168 1 0,9964
2022 0,9030 0,9964 1
Gambar 3.3: Visualisasi Heatmap Matriks RV Data Sampel
3.7.1.6 Tahap 4: Penentuan Bobot Kompromi (αk) – Dua Skema Pembobotan
Langkah kritis dalam STATIS adalah penentuan bobot kompromi αk, yang mencerminkan kontribusi setiap situasi (tahun) terhadap struktur bersama. Sesuai dengan kerangka
Universitas Indonesia

64
kerja Lavit et al. (1994), penelitian ini menyediakan dua skema pembobotan yang menghasilkan dua skenario analisis:
• Skenario 1 – ACT Klasik (Host = Wilayah): Matriks intra-struktur Wk = Z˜
kZ˜T
k
berukuran n × n menggambarkan kemiripan antar-wilayah. Matriks interstruktur
HW dihitung dari tr(D1/2WkD1/2· D1/2Wk
′D1/2
).
• Skenario 2 – Dual STATIS (Host = Indikator): Matriks intra-struktur Sk =
Z
T
k DZk berukuran p × p menggambarkan kemiripan antar-indikator. Matriks interstruktur HS dihitung dari tr(SkSk
′).
Dekomposisi spektral HW u = θu menghasilkan vektor eigen pertama u1 dan bobot
αW,k; demikian pula HSv = ϑv menghasilkan αS,k. Kedua skema bobot disajikan pada
Tabel 3.14 dan Tabel 3.15.
Tabel 3.14: Bobot Alpha (Wilayah)
Tahun Alpha (Wilayah)
2020 0,2881
2021 0,3503
2022 0,3616
Tabel 3.15: Bobot Alpha (Indikator)
Tahun Alpha (Indikator)
2020 0,2794
2021 0,3556
2022 0,3650
Contoh penghitungan (Skenario 1): Matriks HW (Tabel 3.12) diselesaikan dengan
eigendekomposisi Persamaan 2.34. Nilai eigen pertama adalah λ1 = 17,1859 dengan vektor eigen ternormalisasi satuan u1 = (0,4967; 0,6038; 0,6233)T
, sehingga
∥u1∥1 = 0,4967 + 0,6038 + 0,6233 = 1,7238. Substitusi ke Persamaan 2.35:
α2020 =
0,4967
1,7238
= 0,2881, α2021 =
0,6038
1,7238
= 0,3503, α2022 =
0,6233
1,7238
= 0,3616
sebagaimana dikonfirmasi pada Tabel 3.14.
Dari kedua tabel di atas, perhatikan bahwa nilai αW,k dan αS,k dapat berbeda.
Skenario 1 menekankan kestabilan posisi relatif antar-wilayah, sedangkan Skenario 2
menekankan konsistensi struktur hubungan antar-indikator.
3.7.1.7 Tahap 5: Sintesis Matriks Kompromi – Skenario 1 vs Skenario 2
Matriks kompromi disintesis melalui kombinasi linear tertimbang. Kedua skenario
menghasilkan matriks kompromi yang berbeda:
Universitas Indonesia

65
Skenario 1 (ACT Klasik – Host = Wilayah): Matriks kompromi Wc =
P
k αW,k Wk
berukuran n × n, merupakan rataan terbobot matriks intra-struktur wilayah:
W(1)
c = αW,2020 W2020 + αW,2021 W2021 + αW,2022 W2022
Dekomposisi spektral W(1)
c = PΛPT memberikan eigenvector P dan eigenvalue Λ. Skor
kompromi wilayah diperoleh melalui:
F
(1)
c = D−1/2PΛ1/2
, G(1)
c =
X
k
αW,k Z
T
k DF(1)c Λ
−1/2

(3.4)
Contoh penghitungan: Sebagai contoh elemen matriks kompromi Wc (Persamaan 2.36),
elemen [1, 2] mencerminkan kemiripan antara Kota Jakarta Pusat dan Kabupaten Bandung
Barat. Terlebih dahulu dihitung Wk[1, 2] = z
T
1,kz2,k dari baris pertama dan kedua Zk:
W2020[1, 2] = (−1,48)(0,40) + (1,47)(−0,85) + (−0,03)(0,76)
= −0,592 − 1,250 − 0,023 = −1,864
W2021[1, 2] = (−1,51)(0,47) + (1,44)(−0,88) + (−0,99)(0,99)
= −0,710 − 1,267 − 0,980 = −2,957
W2022[1, 2] = (−1,50)(0,45) + (1,45)(−0,87) + (−1,17)(0,98)
= −0,675 − 1,262 − 1,147 = −3,083
Kombinasi linear tertimbang menggunakan αk dari Tabel 3.14:
Wc[1, 2] = 0,2881 × (−1,864) + 0,3503 × (−2,957) + 0,3616 × (−3,083)
= −0,537 − 1,036 − 1,115 = −2,688
Nilai eigen Wc disajikan pada Tabel 3.16; nilai eigen pertama λ
(c)
1 = 2,3008 digunakan
sebagai Λ1/2 pada persamaan di atas.
Skenario 2 (Dual STATIS – Host = Indikator): Matriks kompromi Sc =
P
k αS,k Sk
berukuran p × p, merupakan rataan terbobot matriks korelasi antar-indikator:
S
(2)
c = αS,2020 S2020 + αS,2021 S2021 + αS,2022 S2022
Universitas Indonesia

66
Dekomposisi SVD S
(2)
c = UΣUT memberikan eigenvector indikator U dan nilai singular
Σ. Loading (koordinat indikator) dan skor kompromi objek diperoleh melalui:
G(2)
c = U:,1:2 Σ
1/2
1:2 , F
(2)
c =
X
k
αS,k (Zk U:,1:2) (3.5)
Hasil dekomposisi nilai eigen untuk masing-masing skenario disajikan pada Tabel 3.16
dan Tabel 3.17.
Tabel 3.16: Dekomposisi Nilai Eigen (scenario1)
Dimensi Eigenvalue Varians (%) Kumulatif (%)
Dim 1 2,3008 76,6900 76,6900
Dim 2 0,6992 23,3100 100
Dim 3 0 0 100
Dim 4 0 0 100
Tabel 3.17: Dekomposisi Nilai Eigen (scenario2)
Dimensi Eigenvalue Varians (%) Kumulatif (%)
Dim 1 2,2692 75,6400 75,6400
Dim 2 0,7058 23,5300 99,1700
Dim 3 0,0250 0,8300 100
Tabel 3.18: Diagnostik Dominansi Model Rank-1 (Areia et al., 2008)
Kuantitas Nilai
λ1 (Nilai Eigen Dominan) 17,1859
∥H∥
2
F
(Frobenius2) 295,7348
V (SS Residual) 0,3803
V /∥H∥
2
F
(Rasio Sisa) 0,0013
Inertia1 (Dominansi %) 96,43%
λ1/λ2 (Rasio Kegentingan) 27,8833
Gambar 3.4 menyajikan scree plot yang memperlihatkan dekomposisi nilai eigen
matriks kompromi beserta kurva kumulatif variansi yang dijelaskan oleh setiap dimensi
(Skenario 1).
Universitas Indonesia

67
Gambar 3.4: Scree Plot Analisis Kompromi Data Sampel (Skenario 1)
3.7.1.8 Tahap 6: Proyeksi Kompromi – Perbandingan Skenario 1 dan Skenario 2
Koordinat kompromi dari kedua skenario disajikan berdampingan agar perbedaan
pendekatan dapat dipahami secara komprehensif.
Skenario 1 – Skor Wilayah (F
(1)
c ) dan Loading Indikator (G
(1)
c ): Skor wilayah mencerminkan posisi relatif tiap wilayah pada ruang kompromi, sedangkan loading indikator
mencerminkan arah dan kekuatan kontribusi tiap indikator. Gambar 3.5 menyajikan compromise biplot Skenario 1 yang menampilkan posisi wilayah dan arah loading indikator
secara simultan dalam satu bidang.
Universitas Indonesia

68
Gambar 3.5: Compromise Biplot Skenario 1 (Host = Wilayah): Skor Wilayah dan Loading
Indikator
dengan skor kompromi wilayah (F
(1)
c , Persamaan 2.39) dan loading kompromi indikator (G
(1)
c , Persamaan 2.42) yang disajikan pada Tabel 3.19 dan Tabel 3.26.
Tabel 3.19: Koordinat Skor Kompromi (scenario1)
Kabupaten/Kota PC1 (Asli) PC2 (Asli) PC1 (Visual) PC2 (Visual)
Kota Jakarta Pusat -2,2496 -0,3712 2,2496 -0,3712
Bandung Barat 1,2664 -0,4478 -1,2664 -0,4478
Kota Tasikmalaya 0,1188 1,6629 -0,1188 1,6629
Kota Cilegon -1,2227 -1,4783 1,2227 -1,4783
Skor parsial tahunan Fk (Persamaan 2.40) memproyeksikan matriks intra-struktur
tiap tahun ke ruang kompromi. Tabel menyajikan dua pasang kolom: PC1/PC2 (Asli)
adalah koordinat mentah langsung dari perhitungan STATIS sebelum penyesuaian apapun,
sedangkan PC1/PC2 (Visual) adalah koordinat setelah transformasi penyesuaian visual
(penyelarasan arah sumbu dan skalasi proporsional) sebagaimana diuraikan pada Poin 3..
Hasil per tahun disajikan pada tabel berikut:
Universitas Indonesia

69
Tabel 3.20: Skor Intra Tahun 2020 (F2020)
Kabupaten/Kota PC1 (Asli) PC2 (Asli) PC1 (Visual) PC2 (Visual)
Kota Jakarta Pusat -1,7840 -0,9027 1,7840 -0,9027
Bandung Barat 0,9296 -0,6366 -0,9296 -0,6366
Kota Tasikmalaya 0,2810 2,9117 -0,2810 2,9117
Kota Cilegon -0,9879 -2,3239 0,9879 -2,3239
Tabel 3.21: Skor Intra Tahun 2021 (F2021)
Kabupaten/Kota PC1 (Asli) PC2 (Asli) PC1 (Visual) PC2 (Visual)
Kota Jakarta Pusat -2,3610 -0,0858 2,3610 -0,0858
Bandung Barat 1,3894 -0,4687 -1,3894 -0,4687
Kota Tasikmalaya -0,0260 1,2981 0,0260 1,2981
Kota Cilegon -1,1970 -1,1055 1,1970 -1,1055
Tabel 3.22: Skor Intra Tahun 2022 (F2022)
Kabupaten/Kota PC1 (Asli) PC2 (Asli) PC1 (Visual) PC2 (Visual)
Kota Jakarta Pusat -2,5125 -0,2243 2,5125 -0,2243
Bandung Barat 1,4156 -0,2772 -1,4156 -0,2772
Kota Tasikmalaya 0,1300 1,0214 -0,1300 1,0214
Kota Cilegon -1,4347 -1,1659 1,4347 -1,1659
Dari koordinat skor intra Fk, pergeseran posisi rata-rata (centroid) seluruh wilayah
antar-tahun dapat dikuantifikasi menggunakan jarak Euclidean (Persamaan 2.68). Misalkan
ck =

fk,1, fk,2

adalah centroid kolom skor intra tahun k, maka D(ck, ck+1) mengukur
besarnya pergeseran kolektif seluruh wilayah dalam ruang kompromi dari satu tahun ke
tahun berikutnya.
Contoh penghitungan jarak pergeseran kolektif: Dengan n = 4 wilayah dari Tabel 3.20
dan 3.21, centroid masing-masing tahun diperoleh sebagai rata-rata aritmetika koordinat
seluruh wilayah:
c2020 =
 −1,7840+0,9296+0,2810−0,9879
4
,
−0,9027−0,6366+2,9117−2,3239
4

= (−0,3903; −0,2379)
c2021 =
 −2,3610+1,3894−0,0260−1,1970
4
,
−0,0858−0,4687+1,2981−1,1055
4

= (−0,5487; −0,0905)
Universitas Indonesia

70
Substitusi ke Persamaan 2.68 dengan p = 2:
D(c2020, c2021) = q−0,3903 − (−0,5487)2+

−0,2379 − (−0,0905)2
=
p
(0,1584)2 + (−0,1474)2
=
p
0,0251 + 0,0217 = p0,0468 ≈ 0,2163
Untuk tahun 2022, c2022 = (−0,6004; −0,1615), sehingga D(c2021, c2022) =
p
(0,0517)2 + (0,0710)2 ≈ 0,0879. Nilai D(c2020, c2021) > D(c2021, c2022)
mengindikasikan bahwa pergeseran kolektif terbesar terjadi pada transisi 2020 ke 2021.
Contoh penghitungan jarak pergeseran per-entitas: Selain centroid kolektif, pergeseran posisi setiap entitas (wilayah atau indikator) secara individual antar-tahun dihitung menggunakan rumus yang sama (Persamaan 2.68). Misalkan posisi Kota Jakarta
Pusat pada ruang skor intra (Fk, Skenario 1) adalah p2020 = (−1,7840; −0,9027) dan
p2021 = (−2,3610; −0,0858) (dari Tabel 3.20–3.21), maka:
D(p2020, p2021) = q−2,3610 − (−1,7840)2
+

−0,0858 − (−0,9027)2
=
p
(−0,5770)2 + (0,8169)2 =
p
0,3329 + 0,6673 ≈ 1,0001
Untuk periode 2021–2022, diperoleh D(p2021, p2022) ≈ 0,2053, jauh lebih kecil,
mengindikasikan posisi wilayah ini cenderung stabil setelah transisi 2020–2021. Nilai per-entitas untuk seluruh wilayah dan indikator tersedia pada folder jarak/ hasil
ekspor.
Loading parsial tahunan Gk (Persamaan 2.41) diperoleh dengan memproyeksikan Z
T
k
ke ruang kompromi. Contoh penghitungan: elemen G2020[PPM, PC1] dihitung sebagai
berikut (dengan qλ
(c)
1 =
√
2,3008 = 1,5168):
G2020[PPM, 1] =
X
4
i=1
mi zi,PPM,2020 Fc[i, 1]
q
λ
(c)
1
Universitas Indonesia

71
Numerator = (0,2943)(−1,48)(−2,25) + (0,5040)(0,40)(1,27)
+ (0,2015)(1,17)(0,12) + (0,0001)(−1,74)(−1,22)
= 0,980 + 0,256 + 0,028 + 0,000 = 1,264
∴ G2020[PPM, 1] = 1,264
1,5168
≈ 0,8306
Loading Gk untuk semua tahun disajikan pada tabel berikut. Serupa dengan tabel skor
intra, kolom PC1/PC2 (Asli) adalah koordinat loading mentah hasil STATIS, sedangkan
PC1/PC2 (Visual) telah melalui skalasi proporsional dan penyelarasan arah sumbu sebagaimana diuraikan pada Poin 3., agar panah indikator tetap sebanding dengan sebaran skor
wilayah pada biplot. Loading kompromi Gc =
P
k αW,kGk (Persamaan 2.42) diperoleh
dari rata-rata tertimbangnya:
Gc[PPM, 1] = 0,2881 × 0,8306 + 0,3503 × 0,8732 + 0,3616 × 0,8621
= 0,2393 + 0,3059 + 0,3117 = 0,8569
Tabel 3.23: Loading Intra Tahun 2020 (G2020)
Indikator PC1 (Asli) PC2 (Asli) PC1 (Visual) PC2 (Visual)
PPM 0,8306 0,5569 -1,2599 0,4657
RLS -0,9996 0,0289 1,5162 0,0242
TPT 0,3039 -0,9527 -0,4610 -0,7966
Tabel 3.24: Loading Intra Tahun 2021 (G2021)
Indikator PC1 (Asli) PC2 (Asli) PC1 (Visual) PC2 (Visual)
PPM 0,8732 0,4874 -1,3245 0,4075
RLS -0,9967 0,0813 1,5118 0,0679
TPT 0,8329 -0,5534 -1,2634 -0,4627
Tabel 3.25: Loading Intra Tahun 2022 (G2022)
Indikator PC1 (Asli) PC2 (Asli) PC1 (Visual) PC2 (Visual)
PPM 0,8621 0,5067 -1,3077 0,4236
RLS -0,9979 0,0645 1,5137 0,0539
TPT 0,9117 -0,4108 -1,3830 -0,3435
Universitas Indonesia

72
Tabel 3.26: Loadings Kompromi (scenario1)
Indikator PC1 (Asli) PC2 (Asli) PC1 (Visual) PC2 (Visual)
PPM 0,8569 0,5144 -1,2998 0,4301
RLS -0,9980 0,0601 1,5138 0,0503
TPT 0,7090 -0,6168 -1,0755 -0,5158
Skenario 2 – Skor Wilayah (F
(2)
c ) dan Loading Indikator (G
(2)
c ): Pada Skenario 2,
kompromi terbentuk pada ruang indikator, sehingga biplot lebih menekankan hubungan
antar-indikator. Proyeksi wilayah ke ruang ini menghasilkan skor yang mencerminkan
kemiripan profil indikator. Gambar 3.6 menyajikan compromise biplot Skenario 2 pada
ruang tersebut.
Gambar 3.6: Compromise Biplot Skenario 2 (Host = Indikator): Skor Wilayah dan
Loading Indikator
dengan skor kompromi wilayah (F
(2)
c , Persamaan 2.46) dan loading kompromi indikator (G
(2)
c , Persamaan 2.43) yang disajikan pada Tabel 3.27 dan Tabel 3.28.
Tabel 3.27: Koordinat Skor Kompromi (scenario2)
Kabupaten/Kota PC1 (Asli) PC2 (Asli) PC1 (Visual) PC2 (Visual)
Kota Jakarta Pusat 2,1848 0,4533 2,1848 0,4533
Bandung Barat -1,2879 0,3709 -1,2879 0,3709
Kota Tasikmalaya 0,0294 -1,5906 0,0294 -1,5906
Kota Cilegon 1,1103 1,5140 1,1103 1,5140
Universitas Indonesia

73
Tabel 3.28: Loadings Kompromi (scenario2)
Indikator PC1 (Asli) PC2 (Asli) PC1 (Visual) PC2 (Visual)
PPM -0,8260 -0,5578 -1,2443 -0,4687
RLS 0,9926 0,0262 1,4952 0,0220
TPT -0,7757 0,6276 -1,1684 0,5273
Contoh penghitungan (Skenario 2): Loading kompromi diekstraksi langsung dari
dekomposisi Sc (Persamaan 2.43: Gc = VcΛ
1/2
s ). Dengan vektor eigen pertama Vc[:
, 1] = (−0,5483; 0,6590; −0,5149)T dan √λs1 =
√
2,2692 = 1,5064:
Gc[PPM, 1] = (−0,5483) × 1,5064 = −0,8260
Gc[RLS, 1] = 0,6590 × 1,5064 = 0,9926
Gc[TPT, 1] = (−0,5149) × 1,5064 = −0,7757
Loading parsial tahunan Gk (Persamaan 2.44: Gk = SkVcΛ
−1/2
s ). Kolom PC1/PC2 (Asli)
adalah koordinat mentah STATIS, sedangkan PC1/PC2 (Visual) telah melalui skalasi proporsional (Poin 3.) agar sesuai rentang tampilan biplot. Contoh elemen G2020[PPM, PC1]
menggunakan S2020 dari Tabel 3.9:
Numerator = 1 × (−0,5483) + (−0,8141) × 0,6590 + (−0,2781) × (−0,5149)
= −0,5483 − 0,5365 + 0,1432 = −0,9416
G2020[PPM, 1] = −0,9416
1,5064
= −0,6251
Tabel 3.29: Loading Intra Tahun 2020 (G2020)
Indikator PC1 (Asli) PC2 (Asli) PC1 (Visual) PC2 (Visual)
PPM -0,6251 -1,0679 -0,9416 -0,8972
RLS 0,8470 0,3860 1,2760 0,3243
TPT -0,3855 1,0967 -0,5807 0,9214
Tabel 3.30: Loading Intra Tahun 2021 (G2021)
Indikator PC1 (Asli) PC2 (Asli) PC1 (Visual) PC2 (Visual)
PPM -0,8838 -0,4143 -1,3313 -0,3481
RLS 1,0390 -0,0845 1,5650 -0,0710
TPT -0,8912 0,4951 -1,3425 0,4159
Universitas Indonesia

74
Tabel 3.31: Loading Intra Tahun 2022 (G2022)
Indikator PC1 (Asli) PC2 (Asli) PC1 (Visual) PC2 (Visual)
PPM -0,9236 -0,3072 -1,3913 -0,2581
RLS 1,0588 -0,1413 1,5949 -0,1187
TPT -0,9618 0,3977 -1,4488 0,3341
Skor parsial Fk (Persamaan 2.45: Fk = ZkGcΛ
−1/2
s ). Contoh elemen
F2020[Kota Jakarta Pusat, PC1] dengan Z2020 baris ke-1 = (−1,48; 1,47; −0,03) dan
Gc[:, 1] = (−0,8260; 0,9926; −0,7757):
F2020[JakPus, 1] = (−1,48)(−0,8260) + (1,47)(0,9926) + (−0,03)(−0,7757)
1,5064
=
1,2225 + 1,4591 + 0,0233
1,5064
=
2,7049
1,5064
= 1,7937
Tabel 3.32: Skor Intra Tahun 2020 (F2020)
Kabupaten/Kota PC1 (Asli) PC2 (Asli) PC1 (Visual) PC2 (Visual)
Kota Jakarta Pusat 1,7937 1,0073 1,7937 1,0073
Bandung Barat -1,1702 0,2818 -1,1702 0,2818
Kota Tasikmalaya 0,3066 -2,1769 0,3066 -2,1769
Kota Cilegon 0,6633 1,9384 0,6633 1,9384
Tabel 3.33: Skor Intra Tahun 2021 (F2021)
Kabupaten/Kota PC1 (Asli) PC2 (Asli) PC1 (Visual) PC2 (Visual)
Kota Jakarta Pusat 2,2885 0,3093 2,2885 0,3093
Bandung Barat -1,3443 0,4028 -1,3443 0,4028
Kota Tasikmalaya 0,0191 -1,4598 0,0191 -1,4598
Kota Cilegon 1,1582 1,3275 1,1582 1,3275
Tabel 3.34: Skor Intra Tahun 2022 (F2022)
Kabupaten/Kota PC1 (Asli) PC2 (Asli) PC1 (Visual) PC2 (Visual)
Kota Jakarta Pusat 2,3831 0,1696 2,3831 0,1696
Bandung Barat -1,3229 0,4081 -1,3229 0,4081
Kota Tasikmalaya -0,1727 -1,2691 -0,1727 -1,2691
Kota Cilegon 1,4059 1,3708 1,4059 1,3708
Skor kompromi F
(2)
c diperoleh sebagai rata-rata terbobot (Persamaan 2.46: Fc =
P
k αS,kFk) menggunakan αS,k dari Tabel 3.15.
Universitas Indonesia

75
Dari perbandingan kedua biplot, dapat dilihat bahwa orientasi loading berbeda karena
Skenario 1 memaksimalkan varians antar-wilayah sementara Skenario 2 memaksimalkan
konsistensi struktur antar-indikator. Dalam praktiknya, Skenario 1 lebih cocok untuk
analisis hubungan antar-wilayah lintas waktu, sedangkan Skenario 2 lebih cocok untuk
analisis hubungan antar-indikator lintas waktu.
Lintasan pergeseran wilayah (Gambar 3.7 dan 3.8) mengilustrasikan dinamika temporal dari masing-masing skenario.
Gambar 3.7: Trajektori Skenario 1 (Host =
Wilayah)
Gambar 3.8: Trajektori Skenario 2 (Host =
Indikator)
3.7.1.9 Tahap 7: Uji Signifikansi dan Clustering
Sebelum melanjutkan ke tahap clustering, signifikansi struktur kompromi divalidasi
menggunakan uji permutasi Monte Carlo. Uji ini mengevaluasi apakah nilai eigen pertama
λobs dari matriks inter-struktur H secara statistik lebih besar daripada yang diharapkan
secara acak. Pada setiap permutasi b = 1, . . . , Nperm, baris-baris matriks data Zk diacak
secara terpisah untuk menghancurkan korespondensi antar-tahun, kemudian λ
(b) dihitung
ulang. Nilai p diperoleh menggunakan koreksi pseudo-count (Persamaan 2.49):
p-value =
X
B
b=1
I

θ
π(b)
1 ≥ θ
obs
1

+ 1
B + 1
Pada dataset sampel dengan B = 1000 permutasi, nilai eigen terobservasi adalah
θ
obs
1 = λ1 = 17,1859 (dari matriks interstruktur H, bukan dari Wc). Dari 1000 permutasi, terdapat 571 permutasi yang menghasilkan θ
π
1 ≥ 17,1859, sehingga substitusi ke
Universitas Indonesia

76
Persamaan 2.49:
p-value =
571 + 1
1000 + 1
=
572
1001
= 0,5714 (tidak signifikan)
Hasil ini tidak mengejutkan mengingat dataset sampel hanya terdiri dari 4 observasi; dengan
derajat kebebasan yang sangat terbatas, distribusi null hampir sepenuhnya menutupi nilai
observasi. Pada dataset aktual (41 wilayah), uji permutasi menghasilkan p < 0,05 yang
mengonfirmasi bahwa struktur kompromi STATIS signifikan secara statistik.
Setelah signifikansi dikonfirmasi, jumlah klaster optimal dipilih melalui kombinasi
kurva Elbow (Subbab 2.13.2) dan koefisien siluet (Subbab 2.13.3). Berikut adalah contoh
penghitungan koefisien siluet menggunakan k = 2 pada skor kompromi data sampel:
Contoh penghitungan koefisien siluet (k = 2, data sampel, Skenario 1): Skor kompromi empat wilayah sampel dari Tabel 3.19 adalah: Kota Jakarta Pusat (JP) =
(−2,2496; −0,3712), Kab. Bandung Barat (BB) = (1,2664; −0,4478), Kota Tasikmalaya
(TK) = (0,1188; 1,6629), dan Kota Cilegon (CL) = (−1,2227; −1,4783). Dengan
k = 2, K-Means menghasilkan dua klaster: C1 = {JP, CL} dan C2 = {TK, BB}. Jarak
Euclidean antar-pasang dihitung menggunakan Persamaan 2.68:
d(JP, CL) = p
(−1,0269)2 + (1,1071)2 ≈ 1,510
d(TK, BB) = p
(−1,1476)2 + (2,1107)2 ≈ 2,403
Untuk JP (di C1 dengan satu anggota lain yaitu CL):
a(JP) = d(JP, CL) = 1,510
b(JP) = d(JP, TK) + d(JP, BB)
2
=
3,122 + 3,517
2
= 3,320
s(JP) = 3,320 − 1,510
max(3,320; 1,510) =
1,810
3,320
≈ 0,545
Dengan cara yang sama diperoleh s(CL) ≈ 0,506, s(TK) ≈ 0,265, dan s(BB) ≈ 0,226
(lihat Persamaan 2.64). Koefisien siluet rata-rata (Persamaan 2.65):
SC(k = 2) = 0,545 + 0,506 + 0,265 + 0,226
4
=
1,542
4
≈ 0,386
Nilai ini menunjukkan pemisahan klaster yang cukup pada data sampel. Langkah yang
sama diterapkan pada data aktual (41 wilayah) untuk nilai k = 2 hingga 10, dan hasilnya
dikombinasikan dengan kurva Elbow untuk menetapkan k = 3 sebagai jumlah klaster
optimal (lihat Subbab 4.7.1).
Universitas Indonesia

77
Setelah validasi signifikansi, langkah terakhir adalah pengelompokan wilayah ke
dalam K = 3 klaster menggunakan algoritma K-Means pada skor kompromi Fc. Label
klaster diurutkan berdasarkan rata-rata skor PC1 (klaster dengan skor terendah = Klaster 1,
tertinggi = Klaster 3). Keanggotaan wilayah pada dataset sampel disajikan pada Tabel 3.35.
Tabel 3.35: Ringkasan Anggota Klaster Kompromi STATIS (Skenario 1)
Kelompok (Cluster)Jumlah Wilayah Daftar Kabupaten/Kota
Cluster 1 (Rendah)2 Bandung Barat, Kota Tasikmalaya
Cluster 2 (Tinggi) 2 Kota Jakarta Pusat, Kota Cilegon
Stabilitas pengelompokan lintas tahun dipantau melalui matriks transisi temporal
(Tabel 3.36 dan Tabel 3.37).
Tabel 3.36: Matriks Transisi Klaster 2020 → 2021
2020
2021 Cluster 1 (Rendah) (2021) Cluster 2 (Tinggi) (2021)
Cluster 1 (Rendah) (2020) 1 0
Cluster 2 (Tinggi) (2020) 1 2
Tabel 3.37: Matriks Transisi Klaster 2021 → 2022
2021
2022 Cluster 1 (Rendah) (2022) Cluster 2 (Tinggi) (2022)
Cluster 1 (Rendah) (2021) 2 0
Cluster 2 (Tinggi) (2021) 0 2
3.7.1.10 Tahap 7a: Komputasi Vektor Struktur dan Analisis TransversalLongitudinal
Setelah memperoleh matriks interstruktur HW (Skenario 1) dan memastikan signifikansinya melalui uji permutasi, parameter interstruktur dipadatkan ke dalam vektor
struktur β˜ = λ1α1 (Persamaan 2.51).
1. Vektor Struktur Gabungan (Global): Untuk data sampel, nilai eigen pertama dari
matriks HW adalah λ1 = 17,1859 (Tabel 3.18). Vektor eigen pertama ternormalisasi
satuan adalah u1 = (0,4967; 0,6038; 0,6233)T. Substitusi ke Persamaan 2.51
Universitas Indonesia

78
menghasilkan β˜
k = λ1 · u1k:
β˜
2020 = 17,1859 × 0,4967 = 8,5347
β˜
2021 = 17,1859 × 0,6038 = 10,3798
β˜
2022 = 17,1859 × 0,6233 = 10,7133
Kegentingan model rank-1 divalidasi melalui V (Persamaan 2.52) dan Inertia1 (Persamaan 2.54). Dengan ∥H∥
2
F = 295,7348 dan tr(H) = H11 + H22 + H33 =
4,6998 + 6,3306 + 6,7915 = 17,8219:
V = ∥H∥
2
F − λ
2
1 = 295,7348 − 17,18592 = 0,3803
Inertia1 =
λ1
tr(H)
=
17,1859
17,8219
= 96,43%
Nilai V/∥H∥
2
F = 0,3803/295,7348 = 0,0013 dan λ1/λ2 = 27,8833 mengonfirmasi
model rank-1 memadai. Hasil komputasi disajikan secara lengkap pada Tabel 3.38
berikut.
Tabel 3.38: Komponen Vektor Struktur β˜ per Situasi Waktu
Tahun (k) β˜
k = λ1u1k
2020 8,5347
2021 10,3798
2022 10,7133
2. Vektor Struktur Sub-populasi (Provinsi Jawa Barat): Untuk analisis transversal
provinsi pada data sampel, Jawa Barat memiliki 2 wilayah (nsub = 2) yaitu Kabupaten Bandung Barat dan Kota Tasikmalaya. Bobot lokal Dsub dibentuk dengan
menormalisasi bobot kedua wilayah tersebut sehingga berjumlah satu. Matriks interstruktur lokal Jawa Barat HJabar dibentuk dari produk skalar wilayah sub-populasi.
Eigendekomposisi HJabar menghasilkan nilai eigen pertama λ1,Jabar = 12,8858 dan
vektor eigen pertama u1,Jabar = (0,3712; 0,3305; 0,2982)T. Substitusi menghasilkan komponen β˜
Jabar,k:
β˜
Jabar,2020 = 12,8858 × 0,3712 = 4,7834
β˜
Jabar,2021 = 12,8858 × 0,3305 = 4,2591
β˜
Jabar,2022 = 12,8858 × 0,2982 = 3,8433
Hasil ini disajikan secara lengkap pada Tabel 3.39 berikut.
Universitas Indonesia

79
Tabel 3.39: Komponen Vektor Struktur β˜
gk per Provinsi (Analisis Transversal)
Tahun (k) JAWA BARAT
2020 4,7834
2021 4,2591
2022 3,8433
Perbedaan pola temporal antara β˜ gabungan yang cenderung meningkat (menunjukkan
kestabilan struktur yang menguat) dengan β˜
Jabar yang menurun mencerminkan dinamika
lokal sub-populasi yang berbeda dari tren global.
3.7.1.11 Tahap 8: Analisis Perbandingan Klaster (9 Jenis)
Untuk mengevaluasi konsistensi dan dinamika hasil pengelompokan secara komprehensif, penelitian ini melaksanakan sembilan jenis analisis perbandingan yang terstruktur
dalam empat kategori, sesuai dengan kerangka spasio-temporal STATIS. Penyelarasan
label klaster antar-tahun dilakukan menggunakan Hungarian Algorithm (Kuhn, 1955)
berbasis matriks jarak centroid (Persamaan 2.68), sehingga perbandingan antar-tahun
konsisten secara kombinatorial.
Contoh penghitungan (dataset sampel, Skenario 1, 2020 vs 2021): Hasil K-Means STATIS
pada F2020 menghasilkan tiga klaster dengan keanggotaan: Klaster 1 = {Kota Jakarta Pusat,
Kota Cilegon}, Klaster 2 = {Kota Tasikmalaya}, Klaster 3 = {Bandung Barat}. Centroid
masing-masing klaster dihitung dari rata-rata koordinat anggotanya (Tabel 3.20):
c1,2020 =

−1,7840+(−0,9879)
2
,
−0,9027+(−2,3239)
2

= (−1,3860; −1,6133)
c2,2020 = (0,2810; 2,9117), c3,2020 = (0,9296; −0,6366)
Demikian pula untuk tahun 2021 (Tabel 3.21): c1,2021 = (−1,7790; −0,5957), c2,2021 =
(−0,0260; 1,2981), c3,2021 = (1,3894; −0,4687). Matriks jarak 3 × 3 menggunakan
Persamaan 2.68:
D =


D(c1,2020, c1,2021) D(c1,2020, c2,2021) D(c1,2020, c3,2021)
D(c2,2020, c1,2021) D(c2,2020, c2,2021) D(c2,2020, c3,2021)
D(c3,2020, c1,2021) D(c3,2020, c2,2021) D(c3,2020, c3,2021)

 =


1,0909 3,2134 3,0021
4,0676 1,6425 3,5575
2,7089 2,1578 0,4895


Hungarian Algorithm memilih penugasan {1 → 1, 2 → 2, 3 → 3} dengan total biaya
1,0909 + 1,6425 + 0,4895 = 3,2229 (minimum di antara 3! = 6 permutasi), sehingga label
klaster konsisten antar-tahun dan tidak ada relabeling yang diperlukan pada dataset sampel
ini.
Universitas Indonesia

80
1. Perbandingan 1 – Transisi Biplot k → k+1: Matriks transisi klaster hasil K-Means
Biplot PCA (Unweighted) antar-tahun berturutan, untuk memonitor pergeseran
wilayah pada ruang PCA individu tanpa koreksi temporal.
2. Perbandingan 2 – Transisi STATIS k → k + 1: Matriks transisi klaster hasil KMeans pada koordinat STATIS tahunan (Fk) antar-tahun berturutan, untuk memantau
pergeseran berbasis struktur kompromi terbobot.
3. Perbandingan 3 – Biplot k vs STATIS k: Tabulasi silang antara label klaster Biplot
PCA (Unweighted) dengan label klaster STATIS Tahunan (Fk) pada tahun yang
sama, untuk mengukur perbedaan struktural akibat bobot temporal.
4. Perbandingan 4 – Transisi IPM k → k + 1: Matriks transisi klasifikasi kategori
IPM (UNDP) antar-tahun berurutan, sebagai baseline pergeseran normatif berbasis
standar internasional.
5. Perbandingan 5 – Biplot k vs IPM k: Tabulasi silang antara label klaster Biplot
PCA (Unweighted) dengan klasifikasi kategori IPM (UNDP) pada tahun yang sama,
sebagai validasi eksternal metode Biplot.
6. Perbandingan 6 – STATIS k vs IPM k: Tabulasi silang antara label klaster STATIS
Tahunan (Fk) dengan klasifikasi kategori IPM (UNDP) pada tahun yang sama,
sebagai validasi eksternal metode STATIS.
7. Perbandingan 7 – Kompromi STATIS vs Biplot k: Tabulasi silang antara label
klaster Kompromi STATIS (hasil clustering trajektori keseluruhan) dengan label
klaster Biplot PCA pada setiap tahun, untuk mengukur konsistensi posisi relatif
lintas metode.
8. Perbandingan 8 – Kompromi STATIS vs STATIS k: Tabulasi silang antara label
klaster Kompromi STATIS (clustering trajektori keseluruhan) dengan label klaster
STATIS Tahunan pada setiap tahun, untuk mengevaluasi seberapa representatif
kompromi terhadap gambaran tahunan.
9. Perbandingan 9 – Kompromi STATIS vs IPM k: Tabulasi silang antara label
klaster Kompromi STATIS dengan klasifikasi Kategori IPM (UNDP) pada tahun
yang sama, sebagai validasi eksternal final dari pengelompokan berbasis trajektori.
Tiga contoh representatif disajikan berikut untuk mendokumentasikan format keluaran
masing-masing kategori; hasil lengkap seluruh tahun tersedia pada dashboard interaktif.
Universitas Indonesia

81
Tabel 3.40: Matriks Transisi Klaster Biplot (Unweighted) 2020 → 2021
2020
2021 Cluster 1 (Rendah) (2021) Cluster 2 (Tinggi) (2021)
Cluster 1 (Rendah) (2020) 2 0
Cluster 2 (Tinggi) (2020) 0 2
Tabel 3.41: Nama Daerah Transisi Klaster Biplot 2020 → 2021
Kategori Daftar Wilayah
Cluster 1 (Rendah)
(2020) → Cluster 1
(Rendah) (2021)
Bandung Barat, Kota Tasikmalaya
Cluster 2 (Tinggi)
(2020) → Cluster 2
(Tinggi) (2021)
Kota Jakarta Pusat, Kota Cilegon
Contoh: Perbandingan 1 – Transisi Biplot (Tahun 2020 vs 2021):
Tabel 3.42: Nama Daerah Perpindahan Klaster 2020 → 2021
Kategori Daftar Wilayah
Cluster 1 (Rendah)
(2020) → Cluster 1
(Rendah) (2021)
Kota Tasikmalaya
Cluster 2 (Tinggi)
(2020) → Cluster 1
(Rendah) (2021)
Bandung Barat
Cluster 2 (Tinggi)
(2020) → Cluster 2
(Tinggi) (2021)
Kota Jakarta Pusat, Kota Cilegon
Contoh: Perbandingan 2 – Transisi STATIS (Skenario 1, Tahun 2020 vs 2021):
Tabel 3.43: Matriks Perbandingan Klaster Biplot vs STATIS (Tahunan) (2020)
Biplot (Unweighted) Statis 1 (Rendah) Statis 2 (Tinggi)
Biplot 1 (Rendah) 1 1
Biplot 2 (Tinggi) 0 2
Universitas Indonesia

82
Tabel 3.44: Nama Daerah Perbandingan Klaster Biplot vs STATIS (Tahunan) (2020)
Kategori Daftar Wilayah
Biplot 1 (Rendah) &
Statis 1 (Rendah)
Kota Tasikmalaya
Biplot 1 (Rendah) &
Statis 2 (Tinggi)
Bandung Barat
Biplot 2 (Tinggi) &
Statis 2 (Tinggi)
Kota Jakarta Pusat, Kota Cilegon
3.7.2 Analisis Autokorelasi Spasial (Global Moran’s I)
Contoh: Perbandingan 3 – Biplot vs STATIS (Skenario 1, Tahun 2020): Sebagai
tahap akhir metodologi, Global Moran’s I (Persamaan 2.73) diterapkan pada skor PC1
kompromi fi,1 menggunakan matriks bobot spasial wij berbasis ketetanggaan queen contiguity. Inferensi dilakukan via permutasi Monte Carlo (Nperm permutasi) tanpa bergantung
pada asumsi normalitas distribusi skor. Dataset sampel (n = 4) tidak memenuhi syarat
minimum matriks ketetanggaan yang bermakna; hasil pengujian pada dataset aktual (41
kabupaten/kota) dilaporkan dalam Bab 4.
Universitas Indonesia

BAB 4
HASIL DAN PEMBAHASAN
Bab ini menyajikan hasil analisis data longitudinal 41 kabupaten/kota di Provinsi
Banten, DKI Jakarta, dan Jawa Barat periode 2020 hingga 2024 menggunakan metode
STATIS (Structuration des Tableaux a Trois Indices de la Statistique ` ). Analisis dimulai
dari eksplorasi data awal dan uji kelayakan, dilanjutkan dengan analisis inter-struktur,
pembentukan struktur kompromi, uji stabilitas, analisis trajektori spasio-temporal, hingga
analisis klaster dan evaluasinya. Seluruh proses komputasi dilakukan menggunakan
dashboard interaktif berbasis Python (Dash/Plotly).
Data sekunder yang dianalisis pada bab ini bersumber dari publikasi tahunan referensi
statistik yang diterbitkan oleh Badan Pusat Statistik (BPS) dari masing-masing provinsi
pada tahun 2020 hingga 2024. Secara rinci, data Provinsi Banten diamati dari publikasi
rutin kelima tahunnya (Badan Pusat Statistik Provinsi Banten, 2020, 2021, 2022, 2023,
2024). Untuk Provinsi DKI Jakarta, data direkam dari dokumen “Provinsi DKI Jakarta
Dalam Angka” (Badan Pusat Statistik Provinsi DKI Jakarta, 2020, 2021, 2022, 2023, 2024).
Sementara itu, konstelasi data seluruh kabupaten/kota di provinsi Jawa Barat direferensikan
dari publikasi yang ekuivalen dalam periode yang sama (Badan Pusat Statistik Provinsi
Jawa Barat, 2020, 2021, 2022, 2023, 2024). Data lengkap seluruh 41 kabupaten/kota untuk
semua indikator dan tahun pengamatan tersedia pada Lampiran 5.2.
4.1 Eksplorasi Data Awal
4.1.1 Statistika Deskriptif
Langkah pertama sebelum menerapkan analisis multivariat adalah melakukan eksplorasi terhadap karakteristik distribusi masing-masing indikator pembangunan manusia.
Tabel 4.1 menyajikan ringkasan statistika deskriptif, yang meliputi nilai rata-rata (mean),
simpangan baku, nilai minimum, kuartil pertama (Q1), median, kuartil ketiga (Q3), nilai
maksimum, kemiringan (skewness), kurtosis, dan koefisien variasi (CV%), untuk setiap
indikator pada tiap tahun pengamatan.
83 Universitas Indonesia

84
Tabel 4.1: Statistika Deskriptif Indikator Pembangunan Manusia Tahun 2020–2024 (N =
41)
2020
Indikator Mean (µ) Std Dev (σ) Min Q1 Median Q3 Max Skew (γ1) Kurt (γ2) CV (%) GR 0,3592 0,0357 0,2850 0,3360 0,3520 0,3810 0,4270 0,0668 -0,7740 9,9500 TPT 10,1232 2,1596 5,0800 8,9500 10,6800 11,5400 14,2900 -0,4590 -0,3509 21,3300 TPAK 64,4700 3,8591 55,5700 62,0450 63,9100 66,5300 76,9650 0,7812 1,3209 5,9900 AHH P 73,4278 2,4619 66,5000 72,8700 74,0200 75,4200 76,7100 -1,2048 1,1055 3,3500
JP 1432347,14 1227766,25 433,76 332580,00 1226050,00 2434090,00 5409370,00 0,7879 0,6833 85,7200 PPM 7,6356 3,1055 2,2900 4,9400 7,6200 9,9800 14,8700 0,2597 -0,7499 40,6700
IPM 72,8971 5,7191 63,9100 68,0800 71,7000 77,8300 84,7200 0,4383 -0,9747 7,8500
IPG 89,9871 4,3626 79,0600 86,8100 90,9800 93,5300 96,1400 -0,6283 -0,2917 4,8500 PKKP 3,3263 1,3056 1,2600 2,1800 3,6000 4,2800 5,6400 -0,0754 -1,2654 39,2500 GK 470618,29 122257,45 310437,00 362212,00 457954,00 547415,00 782731,00 0,6325 -0,6175 25,9800 AHH L 69,6368 2,5445 62,6800 69,0300 70,2500 71,7200 73,2200 -1,1209 0,8502 3,6500 RLS 8,9244 1,6318 6,3000 7,5200 8,6300 10,6300 11,8100 0,2849 -1,2277 18,2800 PDRB ADHK 90668,15 114358,95 3251,76 21970,41 38598,43 93480,39 449518,14 1,8010 2,1881 126,13
2021
Indikator Mean (µ) Std Dev (σ) Min Q1 Median Q3 Max Skew (γ1) Kurt (γ2) CV (%) GR 0,3647 0,0529 0,2450 0,3400 0,3630 0,3880 0,5150 0,3782 1,0149 14,5000 TPT 9,1959 2,0346 3,2500 8,2300 9,3200 10,5800 13,0700 -0,6333 0,4640 22,1300 TPAK 64,3738 3,1827 56,7250 62,0800 63,9000 65,7450 74,9150 0,7701 1,6434 4,9400 AHH P 73,5822 2,4672 66,6100 73,0800 74,2100 75,5200 76,8900 -1,2313 1,1623 3,3500
JP 1447714,39 1242116,97 439,48 335810,00 1234830,00 2446687,00 5484150,00 0,7956 0,7128 85,8000 PPM 8,1639 3,1930 2,5700 5,3500 7,9700 10,7100 15,0600 0,1144 -0,9312 39,1100
IPM 73,1805 5,7523 64,0300 68,2900 72,1000 78,0600 84,9000 0,4213 -0,9879 7,8600
IPG 90,0432 4,3543 79,0700 86,9700 90,8700 93,6000 96,2800 -0,6450 -0,2436 4,8400 PKKP 3,8829 1,5201 1,5500 2,4600 3,8700 5,0000 7,1200 0,1816 -1,1459 39,1500 GK 486788,61 124407,79 320050,00 378819,00 467248,00 570330,00 792684,00 0,5875 -0,7418 25,5600 AHH L 69,8390 2,5705 62,8200 69,2700 70,5000 71,9400 73,4800 -1,1514 0,8882 3,6800 RLS 9,0273 1,6447 6,4100 7,5300 8,7700 10,7800 11,8200 0,2434 -1,3084 18,2200 PDRB ADHK 94041,02 118287,92 3371,78 22787,90 39981,19 97869,38 459971,27 1,7731 2,0407 125,78
2022
Indikator Mean (µ) Std Dev (σ) Min Q1 Median Q3 Max Skew (γ1) Kurt (γ2) CV (%) GR 0,3658 0,0610 0,2110 0,3180 0,3670 0,3940 0,5430 0,3696 0,8754 16,6800 TPT 7,8156 1,9985 1,5600 6,9800 8,1000 8,8300 10,7800 -0,9149 0,9909 25,5700 TPAK 65,6523 3,3273 58,6250 64,0300 65,0600 67,1600 80,0350 1,7204 6,4205 5,0700 AHH P 73,8912 2,4406 67,0300 73,3700 74,5000 75,8100 77,1800 -1,2024 1,1051 3,3000
JP 1462414,00 1255992,59 445,06 338940,00 1243320,00 2458707,00 5556310,00 0,8032 0,7413 85,8800 PPM 7,8027 3,0617 2,5000 5,0100 7,7200 10,1400 14,1100 0,1383 -0,9442 39,2400
IPM 73,8166 5,6575 64,7100 69,0400 72,7900 78,7700 85,2100 0,4127 -0,9878 7,6600
IPG 90,3415 4,2583 79,6900 87,1700 90,9700 93,8400 96,4800 -0,6361 -0,2386 4,7100 PKKP 5,5502 2,4654 1,5000 3,4500 5,7700 7,4300 10,4500 0,0747 -1,1068 44,4200 GK 512333,02 133660,54 335134,00 398884,00 485613,00 599514,00 842630,00 0,5993 -0,7304 26,0900 AHH L 70,0607 2,5400 63,1700 69,5000 70,6900 72,1400 73,7000 -1,1122 0,8040 3,6300 RLS 9,1566 1,6221 6,5900 7,8300 8,9000 10,8200 11,8400 0,2314 -1,3538 17,7200 PDRB ADHK 98968,25 124511,50 3516,57 24298,10 42012,96 103219,44 481954,45 1,7659 2,0038 125,81
2023
Indikator Mean (µ) Std Dev (σ) Min Q1 Median Q3 Max Skew (γ1) Kurt (γ2) CV (%) GR 0,3739 0,0545 0,2430 0,3420 0,3690 0,4280 0,4680 -0,2663 -0,4264 14,5700 TPT 7,2041 1,7473 1,5200 6,5200 7,3300 8,1100 10,5200 -0,9574 1,5426 24,2500 TPAK 66,2128 3,6045 59,8500 63,8250 66,0550 67,6950 80,2150 1,4114 3,8014 5,4400 AHH P 74,3205 2,4068 67,5600 73,7800 74,9700 76,2000 77,4500 -1,2365 1,1126 3,2400
JP 1476700,92 1269538,84 450,51 341980,00 1251540,00 2470054,00 5627020,00 0,8107 0,7696 85,9700 PPM 7,4420 2,8508 2,3800 4,8500 7,2700 9,5200 13,1300 0,0826 -0,9921 38,3100
IPM 74,4183 5,6277 65,2100 69,6100 73,3900 79,4600 86,0700 0,4201 -0,9571 7,5600
IPG 90,5937 4,1577 80,3700 87,3100 91,2900 94,0300 97,1600 -0,6001 -0,2853 4,5900 PKKP 4,6354 1,8772 1,5800 2,9100 5,0700 5,9800 8,6500 0,0616 -1,0247 40,5000 GK 553569,61 141787,63 367681,00 426355,00 520713,00 657247,00 891580,00 0,5879 -0,7939 25,6100 AHH L 70,4171 2,4823 63,7000 69,7900 70,9300 72,5000 74,1000 -1,0287 0,7650 3,5300 RLS 9,2263 1,6267 6,6000 7,8900 8,9100 10,8400 11,8800 0,2655 -1,3451 17,6300 PDRB ADHK 104051,93 130847,78 3305,13 25655,25 44087,22 108569,83 506560,74 1,7652 2,0079 125,75
2024
Indikator Mean (µ) Std Dev (σ) Min Q1 Median Q3 Max Skew (γ1) Kurt (γ2) CV (%) GR 0,3667 0,0561 0,2450 0,3250 0,3670 0,4240 0,4770 -0,3164 -0,6547 15,3000 TPT 6,5522 1,5032 1,5800 6,0800 6,4900 7,4000 9,1800 -0,9517 1,6482 22,9400 TPAK 67,2955 3,6266 58,2950 65,5350 67,2000 68,7950 80,2100 0,7021 2,7090 5,3900 AHH P 74,5900 2,4199 67,7800 74,0400 75,2900 76,4200 77,7500 -1,2313 1,0489 3,2400
JP 1504586,38 1293092,62 455,62 344850,00 1259230,00 2528160,00 5682300,00 0,7875 0,6512 85,9400 PPM 7,2541 2,8500 2,3400 4,6300 7,0500 9,4900 13,0300 0,1015 -1,0378 39,2900
IPM 74,9951 5,6228 65,8600 70,0300 73,9900 80,0200 86,9400 0,4290 -0,9383 7,5000
IPG 90,7871 4,0874 81,0200 87,5500 91,5500 94,2300 97,3000 -0,5647 -0,3873 4,5000 PKKP 4,6907 1,4802 1,8900 3,5500 5,2200 5,7800 7,9800 -0,4708 -0,6142 31,5500 GK 583233,61 146630,84 393464,00 455636,00 547912,00 685071,00 928278,00 0,5936 -0,7922 25,1400 AHH L 70,6402 2,4873 63,9000 70,0400 71,1000 72,7200 74,2900 -0,9979 0,7233 3,5200 RLS 9,2837 1,6529 6,6100 7,9000 8,9300 10,8500 11,9900 0,2552 -1,3503 17,8000 PDRB ADHK 109134,55 137364,99 3268,50 26863,08 46270,92 113998,21 532606,87 1,7693 2,0301 125,87
Catatan: Koefisien Variasi (CV) menunjukkan rasio simpangan baku terhadap rata-rata,
dinyatakan dalam persen; CV yang lebih tinggi mengindikasikan variasi relatif yang lebih
besar antar-wilayah untuk indikator tersebut.
Berdasarkan Tabel 4.1, terlihat bahwa indikator IPM memiliki nilai CV yang relatif
rendah (di bawah 8%) di seluruh tahun, yang mengindikasikan bahwa tingkat pembangunan
manusia antar-wilayah tidak terlalu bervariasi secara proporsional, meskipun selisih nilai
IPM antar wilayah tetap cukup besar secara aktual. Sebaliknya, indikator PDRB ADHK
menunjukkan CV terbesar (di atas 100%), mencerminkan kesenjangan kapasitas ekonomi
yang sangat lebar antara wilayah berbasis industri di Jabodetabek dengan wilayah agraris
di pedalaman Jawa Barat dan Banten.
Kedua, nilai skewness positif pada indikator PPM (Persentase Penduduk Miskin)
dan Garis Kemiskinan mengindikasikan bahwa distribusinya menceng ke kanan, artinya
sebagian besar wilayah memiliki tingkat kemiskinan moderat, namun terdapat sejumlah
kecil wilayah dengan tingkat kemiskinan yang sangat tinggi sehingga nilai rata-rata menjadi
lebih tinggi dari nilai yang umum dijumpai. Sebaliknya, indikator RLS (Rata-Rata Lama
Sekolah) memiliki distribusi yang lebih simetris, menandakan bahwa pemerataan akses
pendidikan sudah lebih baik.
Secara keseluruhan, data menunjukkan bahwa periode 2020 hingga 2024 diwarnai
oleh fluktuasi yang dapat dikaitkan dengan dampak pandemi COVID-19 pada 2020 hingga
2021 dan fase pemulihan pada 2022 hingga 2024.
Universitas Indonesia

85
4.1.2 Distribusi Indikator (Boxplot)
Untuk memvisualisasikan sebaran data secara lebih intuitif, termasuk mendeteksi
keberadaan outlier, berikut disajikan diagram kotak-garis (boxplot) setiap indikator sebagai
perbandingan tahunan 2020–2024 secara berpasangan.
Gambar 4.1: Boxplot Tingkat
Pengangguran Terbuka (TPT) Tahun
2020–2024
Gambar 4.2: Boxplot Tingkat Partisipasi
Angkatan Kerja (TPAK) Tahun 2020–2024
Gambar 4.1 menunjukkan lonjakan nilai median TPT pada tahun 2020 akibat pandemi
COVID-19, diikuti tren penurunan yang konsisten hingga 2024; penyempitan distribusi
yang terlihat mengindikasikan perbaikan kondisi ketenagakerjaan yang semakin merata
antar-wilayah. Gambar 4.2 memperlihatkan distribusi yang lebar dengan beberapa pencilan
rendah yang mencirikan wilayah dengan sektor informal yang dominan; median yang relatif
stabil antartahun menunjukkan bahwa struktur penyerapan tenaga kerja tidak mengalami
perubahan fundamental selama periode pengamatan.
Gambar 4.3: Boxplot Angka Harapan
Hidup Laki-Laki (AHH L) Tahun
2020–2024
Gambar 4.4: Boxplot Angka Harapan
Hidup Perempuan (AHH P) Tahun
2020–2024
Gambar 4.3 memperlihatkan peningkatan bertahap pada nilai median dan penyempitan
rentang distribusi, mengindikasikan konvergensi kualitas kesehatan di berbagai wilayah
Universitas Indonesia

86
meskipun masih terdapat beberapa pencilan pada nilai yang lebih rendah. Gambar 4.4
menunjukkan pola yang serupa namun secara konsisten berada pada nilai yang lebih tinggi
di setiap tahun, mencerminkan keunggulan harapan hidup perempuan yang merupakan
fenomena universal; tren peningkatan pada kedua indikator mengonfirmasi perbaikan
layanan kesehatan yang berlangsung merata selama periode pengamatan.
Gambar 4.5: Boxplot Gini Ratio (GR)
Tahun 2020–2024
Gambar 4.6: Boxplot Garis Kemiskinan
(GK) Tahun 2020–2024
Gambar 4.5 menampilkan distribusi yang cukup terkonsentrasi dengan median yang
berfluktuasi ringan, mengindikasikan bahwa tingkat ketimpangan distribusi pendapatan
antar-wilayah cenderung stagnan dan belum mengalami perbaikan yang signifikan selama
2020–2024. Gambar 4.6 sebaliknya menampilkan kenaikan nilai median dan pelebaran
rentang distribusi yang konsisten dari tahun ke tahun, yang sebagian besar disebabkan oleh
tekanan inflasi serta peningkatan standar kebutuhan hidup minimum yang ditetapkan BPS
sebagai acuan penghitungan.
Gambar 4.7: Boxplot Indeks Pembangunan
Manusia (IPM) Tahun 2020–2024
Gambar 4.8: Boxplot Indeks
Pemberdayaan Gender (IPG) Tahun
2020–2024
Gambar 4.7 menunjukkan bahwa nilai IPM mengalami peningkatan yang konsisten
dari 2020 hingga 2024, ditandai oleh kenaikan median setiap tahunnya yang mencerUniversitas Indonesia

87
minkan pemulihan dan pertumbuhan pembangunan manusia pascapandemi; rentang interkuartil yang relatif stabil mengindikasikan perbaikan yang terdistribusi cukup merata
antar-wilayah. Gambar 4.8 memperlihatkan distribusi IPG yang cenderung stagnan dengan
variasi antartahun yang kecil, mengindikasikan bahwa perkembangan kesetaraan gender
berlangsung lambat dan belum terjadi perubahan struktural yang signifikan selama periode
pengamatan.
Gambar 4.9: Boxplot PDRB Atas Dasar
Harga Konstan (PDRB ADHK) Tahun
2020–2024
Gambar 4.10: Boxplot Prevalensi
Ketidakcukupan Konsumsi Pangan (PKKP)
Tahun 2020–2024
Gambar 4.9 menampilkan outlier ekstrem yang sangat mencolok pada PDRB ADHK
yang secara konsisten mewakili wilayah DKI Jakarta—kapasitas ekonominya jauh melampaui seluruh kabupaten/kota lainnya di setiap tahun pengamatan sehingga distribusi bersifat
sangat menceng ke kanan. Gambar 4.10 memperlihatkan bahwa Prevalensi Ketidakcukupan Konsumsi Pangan cenderung menurun dari 2020 hingga 2022, namun mengalami
sedikit kenaikan pada 2023–2024 yang kemungkinan mencerminkan dampak tekanan
harga pangan global; nilai median yang masih cukup tinggi menunjukkan masih adanya
tantangan ketahanan pangan di sebagian wilayah.
Gambar 4.11: Boxplot Rata-Rata Lama
Sekolah (RLS) Tahun 2020–2024
Gambar 4.12: Boxplot Persentase
Penduduk Miskin (PPM) Tahun 2020–2024
Gambar 4.11 menunjukkan bahwa rata-rata lama sekolah memiliki distribusi yang
Universitas Indonesia

88
relatif simetris dengan median yang terus meningkat, mencerminkan perbaikan akses dan
capaian pendidikan yang cukup merata antar-wilayah dari tahun ke tahun. Gambar 4.12
memperlihatkan distribusi Persentase Penduduk Miskin yang menceng ke kanan (positively
skewed), di mana sebagian besar wilayah memiliki tingkat kemiskinan relatif rendah
namun terdapat beberapa wilayah dengan PPM yang jauh melebihi nilai sentral; terlihat
pula penurunan median secara bertahap dari 2020 ke 2024 yang mencerminkan tren
pengentasan kemiskinan pascapandemi.
Sebagai pelengkap eksplorasi, Gambar 4.13 menyajikan heatmap korelasi Pearson
antar-indikator untuk tahun 2020 (representatif dari seluruh periode). Diagonal utama
menunjukkan nama indikator, sedangkan warna di luar diagonal utama menunjukkan
besarnya nilai koefisien korelasi Pearson r beserta bintang signifikansi: *** p<0,001,
** p<0,01, * p<0,05, ns (tidak signifikan). Warna biru tua = korelasi positif kuat; merah
tua = korelasi negatif kuat; putih ≈ 0 (tidak berkorelasi).
Gambar 4.13: Heatmap Korelasi Pearson Antar-Indikator Tahun 2020 (41
Kabupaten/Kota)
Universitas Indonesia

89
4.2 Uji Kelayakan Data Multivariat
Sebelum data diproses lebih lanjut menggunakan metode STATIS, dilakukan dua
pengujian kelayakan untuk memastikan bahwa analisis multivariat valid untuk diterapkan:
Uji Kaiser-Meyer-Olkin (KMO) dan Uji Bartlett.
4.2.1 Hasil Uji KMO dan Bartlett
Mengacu pada formula yang disajikan pada Persamaan 2.15 (KMO) dan 2.17 (Bartlett)
di Bab 2, hasil pengujian dirangkum sebagai berikut.
Tabel 4.2: Hasil Uji KMO dan Uji Bartlett per Tahun
Tahun Nilai KMO Interpretasi Nilai KMO χ
2 Bartlett df P-Value Bartlett Keputusan
2020 0,7153 Middling (Cukup Baik) 501,2800 55 < 0, 001 Layak
2021 0,6732 Mediocre (Biasa) 487,5100 55 < 0, 001 Layak
2022 0,7436 Middling (Cukup Baik) 455,9800 55 < 0, 001 Layak
2023 0,7201 Middling (Cukup Baik) 469,6300 55 < 0, 001 Layak
2024 0,6555 Mediocre (Biasa) 459,8500 55 < 0, 001 Layak
Hasil uji KMO pada Tabel 4.2 menunjukkan bahwa nilai KMO berada pada rentang
0, 6555–0, 7436 dengan rata-rata 0, 7015 (kategori “M iddling(CukupBaik)”), yang
mengindikasikan bahwa data cukup memadai untuk dianalisis secara multivariat. Selanjutnya, hasil Uji Bartlett dengan p-value < 0,05 pada seluruh tahun memberikan cukup
bukti untuk menolak hipotesis nol bahwa matriks korelasi adalah matriks identitas. Hal ini
berarti terdapat interkorelasi signifikan antar-indikator, sehingga analisis dengan metode
multivariat (seperti STATIS) dapat dilakukan.
4.3 Analisis Inter-struktur
Tahap pertama dalam algoritma STATIS adalah analisis interstruktur, yang bertujuan
mengukur tingkat kemiripan atau konsistensi struktur antarpasangan tahun. Sesuai dengan
Persamaan 2.33 dan 2.32 (Bab 2), koefisien RV (RRV) dan matriks H dibentuk dari matriks
intra-struktur Wt dan Wt
′.
4.3.1 Koefisien RV
Mengikuti prosedur komputasi yang telah diilustrasikan pada Subbab 3.7.1.5 (Bab 3),
dihitung koefisien RV untuk seluruh pasangan tahun (2020–2024) menggunakan matriks
intra-struktur Wt yang dibentuk dari data standar (Z-score) untuk setiap tahun. Koefisien
RV memberikan ukuran kuantitatif tentang seberapa mirip struktur relatif antar-wilayah
Universitas Indonesia

90
pada dua tahun tertentu, dengan nilai berada di antara 0 (tidak ada kemiripan) hingga
1 (kemiripan sempurna). Hasilnya disajikan pada Tabel 4.3 dan divisualisasikan dalam
bentuk heatmap pada Gambar 4.14.
Tabel 4.3: Koefisien RV Wilayah (Identik Kedua Skenario)
Tahun 2020 2021 2022 2023 2024
2020 1 0,9760 0,9636 0,9543 0,9550
2021 0,9760 1 0,9776 0,9664 0,9626
2022 0,9636 0,9776 1 0,9772 0,9734
2023 0,9543 0,9664 0,9772 1 0,9824
2024 0,9550 0,9626 0,9734 0,9824 1
Gambar 4.14: Heatmap Koefisien RV Antar-Tahun (2020 hingga 2024)
Berdasarkan Gambar 4.14, seluruh nilai koefisien RV berada di atas 0, 9543, yang
mengindikasikan bahwa konfigurasi relatif antar-wilayah tidak berbeda signifikan dari
tahun ke tahun. Hal ini memiliki implikasi penting: meskipun terjadi perubahan akibat
pandemi COVID-19 pada 2020 hingga 2021 yang mengubah nilai aktual (seperti PDRB
yang turun), hierarki atau struktur kemiripan antar-wilayah tidak berubah secara drastis.
Wilayah yang secara relatif lebih maju tetap berada di posisi yang lebih baik dibandingkan
wilayah yang tertinggal, dan demikian pula sebaliknya.
Nilai RV antara tahun 2020 hingga 2021 yang sedikit lebih rendah dibandingkan
pasangan tahun lainnya mengindikasikan bahwa masa awal pandemi sempat menciptakan
Universitas Indonesia

91
sedikit pergeseran struktur, yang kemudian pulih dan kembali konsisten pada 2022 hingga
2024.
4.3.2 Vektor Bobot (Alpha) dan Representativitas Tahun
Berdasarkan nilai eigen dominan dari matriks interstruktur HW dan HS (Persamaan 2.34–2.35), diperoleh dua vektor bobot yang masing-masing merepresentasikan
skenario analisis yang berbeda:
• αW (Skenario 1 – Host Wilayah): Bobot berdasarkan kemiripan intra-struktur
antar-wilayah. Tabel 4.4 menyajikan hasilnya.
• αS (Skenario 2 – Host Indikator): Bobot berdasarkan kemiripan struktur korelasi
antar-indikator. Tabel 4.5 menyajikan hasilnya.
Tabel 4.4: Bobot Alpha (Wilayah)
Tahun Alpha (Wilayah)
2020 0,1939
2021 0,1938
2022 0,2010
2023 0,2076
2024 0,2038
Tabel 4.5: Bobot Alpha (Indikator)
Tahun Alpha (Indikator)
2020 0,1944
2021 0,1940
2022 0,2017
2023 0,2079
2024 0,2021
Dari Tabel 4.4 dan Tabel 4.5, terlihat bahwa sebagian besar bobot α untuk kelima
tahun relatif seimbang (berada di antara 0, 1938 hingga 0, 2076). Hal ini menegaskan
bahwa tidak ada satu tahun pun yang mendominasi atau mendistorsi gambaran rata-rata
(kompromi). Tahun dengan bobot terbesar di tahun 2023 dapat diinterpretasikan sebagai
tahun paling representatif terhadap kondisi rata-rata selama 2020 hingga 2024.
4.4 Analisis Struktur Kompromi
Setelah diperoleh bobot α, matriks kompromi dibentuk dan didekomposisi secara
spektral. Sesuai dengan dua skenario pembobotan, diperoleh dua representasi kompromi
yang akan dianalisis secara paralel:
• Skenario 1 (ACT Klasik, Host = Wilayah): W(1)
c =
P
t αW,tWt berukuran n × n;
hasil dekomposisi langsung memberikan koordinat skor wilayah.
• Skenario 2 (Dual STATIS, Host = Indikator): S
(2)
c =
P
t αS,tSt berukuran p × p;
hasil dekomposisi memberikan koordinat loading indikator secara langsung.
Universitas Indonesia

92
4.4.1 Dekomposisi Eigenvalue (Scree Plot)
Gambar 4.15 menyajikan scree plot, yaitu grafik nilai eigen beserta kurva kumulatif
variansi yang dijelaskan oleh setiap dimensi kompromi.
Gambar 4.15: Scree Plot Nilai Eigen Matriks Kompromi
Tabel 4.6: Dekomposisi Nilai Eigen Dua Dimensi Pertama Matriks Kompromi
Dimensi Eigenvalue Varians (%) Kumulatif (%)
Dim 1 6,6934 60,8500 60,8500
Dim 2 1,2132 11,0300 71,8800
Berdasarkan Gambar 4.15 dan Tabel 4.6, dua dimensi pertama mampu menjelaskan
71, 88% kumulatif variansi total. Secara spesifik, Dimensi 1 menjelaskan 60, 85% dan
Dimensi 2 menjelaskan 11, 03%. Nilai ini telah melampaui ambang batas 70% yang lazim
digunakan dalam analisis reduksi dimensi, sehingga representasi dua dimensi dianggap
sudah cukup informatif.
Tampak pada grafik bahwa terjadi “patahan” (elbow) yang jelas setelah Dimensi 2, di
mana peningkatan variansi yang dijelaskan oleh dimensi selanjutnya menjadi sangat kecil.
Hal ini memvalidasi keputusan untuk mereduksi data menjadi dua dimensi, yaitu menjadi
dua komponen (kompromi) utama dalam seluruh analisis berikutnya.
Universitas Indonesia

93
4.4.2 Peta Kompromi (Skor Objek)
Gambar 4.16 menampilkan posisi 41 kabupaten/kota pada ruang dua dimensi kompromi. Posisi ini merupakan representasi “rata-rata tertimbang” dari kondisi tiap wilayah
selama 5 tahun pengamatan.
Gambar 4.16: Peta Skor Kompromi 41 Kabupaten/Kota (2020 hingga 2024)
Dari Gambar 4.16, terlihat adanya pola pengelompokan yang jelas. Wilayah-wilayah
DKI Jakarta (Jakarta Pusat, Jakarta Selatan, dll.) secara konsisten berada jauh di sisi kanan
sumbu Dimensi 1, yang merepresentasikan tingginya PDRB dan kesejahteraan umum.
Sebaliknya, wilayah-wilayah seperti Kabupaten Lebak, Pandeglang, dan Garut berada di
sisi sebaliknya, yang mencerminkan masih tingginya tingkat kemiskinan dan rendahnya
kapasitas ekonomi relatif. Pola ini mengonfirmasikan bahwa Dimensi 1 menangkap
variansi utama kemajuan wilayah.
4.4.3 Biplot Kompromi (Skor + Loading)
Untuk mendapatkan pemahaman yang lebih menyeluruh mengenai hubungan antara
wilayah dan indikator, disajikan biplot kompromi dari kedua skenario. Gambar 4.17 dan
4.18 mengintegrasikan skor objek (Fc) dan loading indikator (Gc) dalam satu bidang.
Universitas Indonesia

94
Gambar 4.17: Biplot Kompromi Skenario 1
(Host = Wilayah)
Gambar 4.18: Biplot Kompromi Skenario
2 (Host = Indikator)
Catatan visualisasi: Untuk keperluan keterbacaan gambar Biplot yang disajikan pada
Gambar 4.17 dan Gambar 4.18, panjang panah indikator (loading) telah disesuaikan
secara representasional mengikuti rekomendasi Abdi et al. (2012). Secara praktis, setiap
koordinat loading pada komponen ke-ℓ dikalikan dengan nilai singular σℓ =
√
λℓ sehingga
panjang panah menjadi sepadan (commensurable) dengan variansi skor pada dimensi yang
sama. Penyesuaian ini bersifat visual semata — nilai skor kompromi, nilai eigen, serta
hasil uji statistik tidak berubah — dan skrip yang menghasilkan gambar ini menerapkan
transformasi tersebut hanya pada tahap plotting.
Interpretasi biplot ini dapat dilakukan dengan melihat arah dan panjang vektor indikator
serta posisi wilayah:
1. Dimensi 1: Menjelaskan 60,85% variansi. Berdasarkan nilai loading pada Tabel 4.9
dan Tabel 4.10, indikator dengan loading positif besar pada Dimensi 1 mencirikan
sumbu ini sebagai representasi kemajuan pembangunan secara umum. Indikator
dengan loading positif terkuat adalah Rata-rata Lama Sekolah (RLS, 0,9538), Angka
Harapan Hidup Laki-laki (AHH L, 0,8891), Angka Harapan Hidup Perempuan
(AHH P, 0,8831), Garis Kemiskinan (GK, 0,8830), Indeks Pemberdayaan Gender
(IPG, 0,8248), PDRB ADHK (0,7390), dan Gini Ratio (0,6229). Sebaliknya, Persentase Penduduk Miskin (PPM, −0,8816) dan PKKP (−0,8497) memiliki loading
negatif yang kuat. Dimensi 1 diinterpretasikan sebagai “Taraf Kemajuan Pembangunan Manusia dan Ekonomi”: wilayah dengan skor PC1 tinggi unggul pada
capaian pendidikan, kesehatan, dan produktivitas ekonomi, sekaligus memiliki angka
kemiskinan yang rendah.
2. Dimensi 2: Menjelaskan 11,03% variansi. Indikator dengan loading dominan
Universitas Indonesia

95
pada Dimensi 2 adalah Tingkat Pengangguran Terbuka (TPT, +0,8620), dengan
kontribusi negatif moderat dari Tingkat Partisipasi Angkatan Kerja (TPAK, −0,4946).
Dimensi 2 diinterpretasikan sebagai “Tekanan Pasar Tenaga Kerja”: skor PC2
yang tinggi menandakan wilayah dengan pengangguran terbuka tinggi dan partisipasi
kerja yang lebih rendah, sedangkan skor PC2 yang sangat negatif khas untuk wilayah
agraris atau berbasis sektor informal di mana hampir seluruh angkatan kerja sudah
terserap.
Koordinat numerik dari skor kompromi dariseluruh 41 kabupaten/kota pada ruang
kompromi disajikan pada Tabel 4.7 dan Tabel 4.8 berikut.
Universitas Indonesia

96
Tabel 4.7: Koordinat Skor Kompromi (scenario1)
Kabupaten/Kota PC1 (Asli) PC2 (Asli) PC1 (Visual) PC2 (Visual)
Kep. Seribu -2,3219 -1,1030 -2,3219 -1,1030
Kota Jakarta Selatan 4,5753 -1,6893 4,5753 -1,6893
Kota Jakarta Timur 3,5454 -0,3048 3,5454 -0,3048
Kota Jakarta Pusat 3,8726 -1,1918 3,8726 -1,1918
Kota Jakarta Barat 3,0485 -0,6863 3,0485 -0,6863
Kota Jakarta Utara 2,5146 -0,5856 2,5146 -0,5856
Bogor -0,8017 1,2607 -0,8017 1,2607
Sukabumi -2,4008 0,0804 -2,4008 0,0804
Cianjur -3,7243 -0,4450 -3,7243 -0,4450
Bandung 0,1358 -0,2011 0,1358 -0,2011
Garut -2,9254 0,2024 -2,9254 0,2024
Tasikmalaya -3,9480 -1,9026 -3,9480 -1,9026
Ciamis -2,2314 -2,2905 -2,2314 -2,2905
Kuningan -1,2415 1,8291 -1,2415 1,8291
Cirebon -2,5676 0,8021 -2,5676 0,8021
Majalengka -3,3402 -1,9813 -3,3402 -1,9813
Sumedang -1,0382 -0,6354 -1,0382 -0,6354
Indramayu -2,3926 -0,7923 -2,3926 -0,7923
Subang -1,9429 -0,2986 -1,9429 -0,2986
Purwakarta -1,6531 0,9821 -1,6531 0,9821
Karawang -0,3873 0,8745 -0,3873 0,8745
Bekasi 1,5854 0,6456 1,5854 0,6456
Bandung Barat -2,2252 1,8258 -2,2252 1,8258
Pangandaran -3,1715 -4,8306 -3,1715 -4,8306
Kota Bogor 1,5613 1,7859 1,5613 1,7859
Kota Sukabumi 0,7966 1,7032 0,7966 1,7032
Kota Bandung 2,6632 0,3163 2,6632 0,3163
Kota Cirebon < 0, 001 0,2750 < 0, 001 0,2750
Kota Bekasi 3,3477 0,3331 3,3477 0,3331
Kota Depok 3,3564 -0,1203 3,3564 -0,1203
Kota Cimahi 1,8737 1,5406 1,8737 1,5406
Kota Tasikmalaya -0,9956 -0,6269 -0,9956 -0,6269
Kota Banjar -2,0346 -1,1936 -2,0346 -1,1936
Pandeglang -5,3232 -0,1484 -5,3232 -0,1484
Lebak -5,0109 -0,8613 -5,0109 -0,8613
Tangerang -1,1628 -1,0035 -1,1628 -1,0035
Serang -3,9883 -0,1850 -3,9883 -0,1850
Kota Tangerang 1,3274 -1,3251 1,3274 -1,3251
Kota Cilegon -1,6039 -0,4886 -1,6039 -0,4886
Kota Serang -1,9140 -0,8495 -1,9140 -0,8495
Kota Tangerang Selatan 2,5591 -0,7669 2,5591 -0,7669
Universitas Indonesia

97
Tabel 4.8: Koordinat Skor Kompromi (scenario2)
Kabupaten/Kota PC1 (Asli) PC2 (Asli) PC1 (Visual) PC2 (Visual)
Kep. Seribu 2,3816 -1,1762 -2,3816 -1,1762
Kota Jakarta Selatan -4,5772 -1,4937 4,5772 -1,4937
Kota Jakarta Timur -3,5148 -0,3344 3,5148 -0,3344
Kota Jakarta Pusat -3,8494 -1,0104 3,8494 -1,0104
Kota Jakarta Barat -3,0121 -0,5570 3,0121 -0,5570
Kota Jakarta Utara -2,5775 -0,3136 2,5775 -0,3136
Bogor 0,7729 1,4663 -0,7729 1,4663
Sukabumi 2,3818 -0,0213 -2,3818 -0,0213
Cianjur 3,6949 -0,4256 -3,6949 -0,4256
Bandung -0,1756 -0,2382 0,1756 -0,2382
Garut 2,8810 0,0361 -2,8810 0,0361
Tasikmalaya 3,9765 -1,4783 -3,9765 -1,4783
Ciamis 2,2885 -2,2533 -2,2885 -2,2533
Kuningan 1,2332 1,6819 -1,2332 1,6819
Cirebon 2,5551 0,7920 -2,5551 0,7920
Majalengka 3,3673 -1,7852 -3,3673 -1,7852
Sumedang 1,0448 -0,6059 -1,0448 -0,6059
Indramayu 2,4535 -0,7554 -2,4535 -0,7554
Subang 1,9378 -0,5705 -1,9378 -0,5705
Purwakarta 1,6094 1,0552 -1,6094 1,0552
Karawang 0,4096 1,0082 -0,4096 1,0082
Bekasi -1,5566 0,5465 1,5566 0,5465
Bandung Barat 2,1812 1,7283 -2,1812 1,7283
Pangandaran 3,2029 -5,2378 -3,2029 -5,2378
Kota Bogor -1,5474 1,5695 1,5474 1,5695
Kota Sukabumi -0,8481 1,9201 0,8481 1,9201
Kota Bandung -2,6632 0,1180 2,6632 0,1180
Kota Cirebon -0,0478 0,4286 0,0478 0,4286
Kota Bekasi -3,3064 0,0100 3,3064 0,0100
Kota Depok -3,3049 -0,3471 3,3049 -0,3471
Kota Cimahi -1,8913 1,1807 1,8913 1,1807
Kota Tasikmalaya 0,9906 -0,4835 -0,9906 -0,4835
Kota Banjar 2,0604 -1,1151 -2,0604 -1,1151
Pandeglang 5,2527 0,4928 -5,2527 0,4928
Lebak 4,9545 -0,6147 -4,9545 -0,6147
Tangerang 1,1646 -0,8936 -1,1646 -0,8936
Serang 3,9528 0,3438 -3,9528 0,3438
Kota Tangerang -1,3445 -1,2020 1,3445 -1,2020
Kota Cilegon 1,5705 0,1352 -1,5705 0,1352
Kota Serang 1,8482 -0,6380 -1,8482 -0,6380
Kota Tangerang Selatan -2,5328 -0,5548 2,5328 -0,5548
Skor positif pada Dimensi 1 menunjukkan wilayah yang secara relatif lebih maju,
Universitas Indonesia

98
sedangkan skor negatif menunjukkan wilayah yang tertinggal (berdasarkan arah loading
pada Tabel 4.9). Pada Dimensi 2, skor negatif ekstrem seperti yang dimiliki Kabupaten
Pangandaran (skor intra PC2 hingga −4,76, Tabel 4.19) yang mencerminkan wilayah
dengan TPAK sangat tinggi dan TPT sangat rendah, khas wilayah pertanian dan ekonomi
informal. Sebaliknya, skor PC2 yang positif seperti pada Kabupaten Kuningan, mencerminkan tekanan pengangguran terbuka yang lebih besar, umum ditemui pada wilayah
semi-urban yang pertumbuhan angkatan kerjanya melampaui kapasitas serapan sektor
formal.
Tabel 4.9 dan Tabel 4.10 menyajikan nilai loading (Gc) masing-masing indikator
secara numerik untuk setiap skenario.
Tabel 4.9: Loadings Kompromi (scenario1)
Indikator PC1 (Asli) PC2 (Asli) PC1 (Visual) PC2 (Visual)
IPG 0,8248 -0,2499 2,1340 -0,2753
PPM -0,8816 0,0509 -2,2809 0,0560
PDRB ADHK 0,7390 -0,1576 1,9120 -0,1736
RLS 0,9538 -0,0801 2,4677 -0,0883
AHH L 0,8891 0,0982 2,3003 0,1081
AHH P 0,8831 0,1055 2,2847 0,1162
GK 0,8830 -0,1500 2,2846 -0,1653
PKKP -0,8497 0,1062 -2,1984 0,1169
TPT 0,2007 0,8620 0,5193 0,9495
TPAK -0,4764 -0,4946 -1,2324 -0,5448
GR 0,6229 0,1391 1,6115 0,1532
Tabel 4.10: Loadings Kompromi (scenario2)
Indikator PC1 (Asli) PC2 (Asli) PC1 (Visual) PC2 (Visual)
IPG -0,8222 -0,2330 2,1251 -0,2689
PPM 0,8789 0,0795 -2,2719 0,0917
PDRB ADHK -0,7381 -0,0973 1,9079 -0,1123
RLS -0,9528 -0,0922 2,4629 -0,1064
AHH L -0,8865 0,0217 2,2915 0,0251
AHH P -0,8795 0,0261 2,2732 0,0301
GK -0,8778 -0,1488 2,2690 -0,1717
PKKP 0,8560 0,1692 -2,2125 0,1953
TPT -0,2087 0,8643 0,5396 0,9976
TPAK 0,5005 -0,6283 -1,2938 -0,7252
GR -0,6505 0,2449 1,6815 0,2827
Loading yang lebih besar menunjukkan bahwa indikator tersebut memiliki kontribusi
yang lebih kuat dalam menentukan posisi wilayah pada dimensi tersebut. Misalnya, PDRB
Universitas Indonesia

99
ADHK memiliki loading yang sangat tinggi pada Dimensi 1, menegaskan perannya sebagai
indikator utama dalam membedakan tingkat kesejahteraan ekonomi antar-wilayah.
4.4.4 Kontribusi dan Kualitas Representasi Indikator
Gambar 4.19–4.22 secara berturut-turut menyajikan kontribusi relatif setiap indikator
terhadap Dimensi 1 dan kualitas representasinya (cos2) untuk masing-masing skenario.
Gambar 4.19: Kontribusi Indikator Dim 1 –
Skenario 1
Gambar 4.20: Kontribusi Indikator Dim 1 –
Skenario 2
Gambar 4.21: Kualitas Representasi (cos2)
– Skenario 1
Gambar 4.22: Kualitas Representasi (cos2)
– Skenario 2
Dari perbandingan kedua skenario, terlihat bahwa Skenario 1 lebih menekankan
separasi wilayah (loading indikator ekonomi mendominasi Dimensi 1), sedangkan Skenario
2 lebih menonjolkan hubungan antar-indikator sosial. Kedua skenario menghasilkan
informasi yang saling melengkapi.
4.5 Uji Stabilitas Struktur Kompromi (Permutation Test)
Untuk memvalidasi bahwa struktur kompromi yang diperoleh bukan terbentuk secara
kebetulan dari data, dilakukan uji permutasi sebanyak B = 1000 iterasi, sesuai dengan
Universitas Indonesia

100
prosedur yang telah diilustrasikan pada Subbab 3.7.1 (Bab 3) dan Persamaan 2.49 (Bab 2),
dengan statistik θ1 dari matriks H.
4.5.1 Prosedur dan Hasil Uji Permutasi
Pada setiap iterasi, entri-entri di dalam masing-masing kolom (indikator) dari matriks
data Xt diacak secara independen untuk setiap tahun (permutasi per kolom) sebagaimana
direkomendasikan oleh Abdi et al. (2012), guna menghancurkan korelasi antar-indikator
sekaligus mempertahankan sifat normalisasi asli (skala Z-score) dari masing-masing
kolom. Setelah pengacakan kolom tersebut, dihitung nilai eigen terbesar θ
(b)
1 dari matriks
interstruktur H yang baru terbentuk. Distribusi dari 1000 nilai θ
(b)
1
ini membentuk distribusi
null yang dibandingkan dengan nilai observasi θ
obs
1
, sebagaimana divisualisasikan pada
Gambar 4.23.
Gambar 4.23: Distribusi Null Uji Permutasi vs. Nilai Eigen Observasi (λ
obs
1
)
Hasil pengujian yang divisualisasikan pada Gambar 4.23 menunjukkan bahwa nilai
λ
obs
1 = 6, 6934 jauh melampaui seluruh nilai distribusi null. Secara formal, p-value =
0, 0010 yang jauh lebih kecil dari taraf signifikansi α = 0,05.
Universitas Indonesia

101
Tabel 4.11: Rangkuman Hasil Uji Permutasi
Parameter Nilai
Jumlah Permutasi (B) 1000
λ
obs
1
(Nilai Eigen Observasi) 6, 6934
p-value 0,0010
Taraf Signifikansi (α) 0,05
Keputusan Tolak H0
Interpretasi Struktur Signifikan (Stabil)
4.5.2 Interpretasi
Penolakan hipotesis nol H0 pada Tabel 4.11 berarti: struktur kompromi yang dihasilkan
metode STATIS secara statistik signifikan dan bukan terbentuk secara kebetulan. Dengan
kata lain, terdapat struktur multivariat yang nyata dan stabil yang mengatur hubungan
antar-wilayah di Provinsi Banten, DKI Jakarta, dan Jawa Barat selama periode 2020 hingga
2024. Hasil ini menjadi landasan yang kuat untuk melanjutkan interpretasi biplot dan
analisis trajektori.
4.5.3 Analisis Vektor Struktur (β˜) dan Dominansi Interstruktur
Mengacu pada model probabilistik yang dikembangkan oleh Areia et al. (2008) dan
dirangkum kembali dalam Oliveira et al. (2024) (Subbab 2.11.7), struktur bersama derajat-1
diasumsikan berlaku apabila matriks interstruktur H dapat diaproksimasi secara memadai
oleh model rank-1:
H ≈ λααT,
dengan estimator vektor struktur β˜ = λ1α1 (Persamaan 2.51). Pada penelitian ini, λ1 =
θ1 = 6, 6934 adalah nilai eigen pertama matriks H (sama dengan statistik observasi
pada uji permutasi), dan α1 = (u11, u12, . . . , u1K)
T
adalah vektor eigen pertama yang
dinormalisasi menjadi α.
Komponen Vektor Struktur β˜
Komponen β˜
k = λ1 ·u1k (di mana u1k adalah elemen ke-k dari vektor eigen pertama H
sebelum normalisasi) merupakan estimator dari “kontribusi terbobot” setiap situasi waktu
terhadap struktur bersama. Tabel 4.12 menyajikan hasil komputasi ini.
Universitas Indonesia

102
Tabel 4.12: Komponen Vektor Struktur β˜ per Situasi Waktu
Tahun (k) β˜
k = λ1u1k
2020 103,0345
2021 102,9818
2022 106,8355
2023 110,3180
2024 108,3078
Nilai β˜
k yang besar pada situasi ke-k menunjukkan bahwa konfigurasi antar-wilayah
pada tahun tersebut memberikan kontribusi yang dominan dalam membentuk struktur
bersama. Nilai yang kecil (relatif) mengindikasikan bahwa tahun tersebut memiliki deviasi
yang lebih besar dari pola struktur konsensus.
Validasi Dominansi: Inertia Dimensi Pertama dan Sum of Squares Residual
Validitas pendekatan rank-1 dinilai melalui dua kuantitas (Persamaan 2.54 dan 2.52):
1. Inertia Dimensi Pertama Interstruktur (Inertia1): Proporsi total nilai eigen matriks H yang dijelaskan oleh dimensi pertama. Nilai yang tinggi (mendekati 1)
mengonfirmasi bahwa satu komponen dominan cukup untuk merangkum keseluruhan informasi interstruktur.
2. Sum of Squares Residual (V ): Besaran V = ∥H∥
2
F − λ
2
1 mengukur varians
interstruktur yang tidak dapat dijelaskan oleh model rank-1. Nilai V yang kecil
relatif terhadap ∥H∥
2
F mengonfirmasi keselarasan model.
Tabel 4.13: Diagnostik Dominansi Model Rank-1 (Areia et al., 2008)
Kuantitas Nilai
λ1 (Nilai Eigen Dominan) 237,7726
∥H∥
2
F
(Frobenius2) 56548,0585
V (SS Residual) 12,2276
V /∥H∥
2
F
(Rasio Sisa) 0,0002
Inertia1 (Dominansi %) 97,52%
λ1/λ2 (Rasio Kegentingan) 79,5619
Berdasarkan Tabel 4.13, nilai Inertia1 yang mendekati 1,00 mengonfirmasi secara
kuantitatif bahwa struktur bersama derajat-1 yang diasumsikan dalam model Areia et al.
(2008) terpenuhi dengan baik pada data penelitian ini. Dengan demikian, estimasi β˜ dan
α yang diperoleh memiliki validitas statistik yang memadai.
Universitas Indonesia

103
Analisis Transversal: Perbandingan Vektor Struktur Antar-Provinsi
Mengacu pada kerangka Analisis Transversal dalam Oliveira et al. (2024), data penelitian dipartisi menjadi tiga seri studi yang berkorespondensi dengan tiga provinsi yang
diamati: Banten (8 wilayah), DKI Jakarta (6 wilayah), dan Jawa Barat (27 wilayah).
Untuk masing-masing provinsi g (g = 1, 2, 3), matriks interstruktur Hg dihitung menggunakan hanya wilayah-wilayah yang termasuk dalam provinsi tersebut, dan vektor struktur
β˜
g = λ1gα1g diestimasi secara independen.
Tabel 4.14: Komponen Vektor Struktur β˜
gk per Provinsi (Analisis Transversal)
Tahun (k) DKI JAKARTA JAWA BARAT BANTEN
2020 449,3325 72,5113 1579,2734
2021 447,8066 74,5115 1671,9900
2022 456,5275 77,2854 1600,9055
2023 468,8870 79,5852 1708,3063
2024 516,6398 75,6554 2030,7593
Tabel 4.14 menyajikan komponen β˜
gk untuk masing-masing provinsi g dan tahun k.
Perbedaan pola komponen antar-provinsi mencerminkan heterogenitas respons temporal
terhadap perubahan dari masing-masing kelompok wilayah. Ketiga provinsi menunjukkan
heterogenitas yang signifikan: DKI Jakarta mencatat kenaikan konsisten dari 449,33
(2020) ke 516,64 (2024), dengan akselerasi terbesar +47,75 poin (+10,2%) pada transisi
2023–2024, mencerminkan pemulihan ekonomi berbasis jasa yang kuat. Jawa Barat menunjukkan pola paling stabil (rentang 72,51–79,59) namun mengalami penurunan ringan
pada 2024 (75,66), mengindikasikan perlambatan struktural di akhir periode. Banten memperlihatkan fluktuasi terbesar (rentang 1.579,27–2.030,76): sempat turun pada 2021–2022,
kemudian melonjak drastis +322,45 poin (+18,9%) pada 2023–2024 yang mencerminkan
dinamika struktural yang paling tidak stabil, kemungkinan didorong oleh pergeseran
komposisi ekonomi manufaktur dan jasa di kawasan Tangerang Raya.
Analisis Longitudinal: Progresi Temporal Komponen Vektor Struktur
Analisis Longitudinal (Areia et al., 2008; Oliveira et al., 2024) mengevaluasi apakah
komponen β˜
k menunjukkan tren yang sistematis seiring berjalannya waktu. Gambar 4.24
menampilkan nilai β˜
k (untuk seri gabungan seluruh wilayah) pada sumbu-y terhadap tahun
pengamatan pada sumbu-x.
Universitas Indonesia

104
Gambar 4.24: Lintasan Temporal Komponen Vektor Struktur β˜
k (2020 hingga 2024)
Nilai β˜
k pada Tabel 4.12 mencerminkan bobot temporal terstruktur dari masing-masing
tahun dalam membentuk konsensus inter-struktur. Komponen β˜
k menampilkan tiga fase
temporal yang berbeda: (i) Stagnasi pandemi (2020–2021): nilai β˜
2020 = 103,03 dan
β˜
2021 = 102,98 hampir identik (selisih hanya 0,05 poin), mengindikasikan bahwa pandemi
tidak menyebabkan gangguan struktural yang berarti pada hierarki relatif antar-wilayah,
meskipun nilai absolut indikator mungkin mengalami kontraksi, tatanan inter-spasial tetap
terjaga; (ii) Pemulihan progresif (2022–2023): kenaikan bertahap ke β˜
2022 = 106,84 dan
puncak β˜
2023 = 110,32 mengonfirmasi pemulihan struktural yang nyata pasca-pandemi;
(iii) Stabilisasi (2024): moderasi ringan ke β˜
2024 = 108,31 mengindikasikan konsolidasi
struktur inter-relasi antar-wilayah, bukan kemunduran.
4.6 Analisis Dinamika Spasio-Temporal (Trajektori)
Analisis trajektori memanfaatkan proyeksi data setiap tahun ke dalam ruang kompromi
yang telah ditetapkan. Berbeda dengan peta kompromi yang hanya menampilkan posisi
“rata-rata”, analisis trajektori memperlihatkan lintasan pergeseran setiap wilayah dari tahun
ke tahun, sehingga dinamika temporalnya dapat dipahami secara komprehensif.
4.6.1 Gambaran Umum Trajektori 2020 hingga 2024
Gambar 4.25 menampilkan keseluruhan lintasan pergeseran 41 kabupaten/kota pada
bidang kompromi selama 2020 hingga 2024.
Universitas Indonesia

105
Gambar 4.25: Trajektori Pergeseran 41 Kabupaten/Kota pada Ruang Kompromi (2020
hingga 2024)
Secara umum, terdapat tiga pola utama trajektori yang dapat diidentifikasi dari Gambar 4.25:
1. Pola Pemulihan Berbentuk V (V-Shape Rebound): Wilayah metropolitan berbasis
jasa seperti kota-kota di DKI Jakarta menunjukkan penurunan posisi pada 2020
hingga 2021 (akibat kontraksi pandemi) namun berhasil kembali ke posisi semula
atau bahkan lebih tinggi pada 2022 hingga 2024.
2. Pola Progresif Bertahap (Monotone Improvement): Sejumlah wilayah menunjukkan peningkatan posisi yang konsisten dari tahun ke tahun tanpa penurunan,
mencerminkan pertumbuhan yang berkelanjutan.
3. Pola Stagnasi (Inertia): Wilayah-wilayah tertinggal menunjukkan lintasan yang sangat pendek dan hampir stagnan, mengindikasikan bahwa perubahan makroekonomi
tidak banyak mengubah posisi relatif mereka.
Universitas Indonesia

106
4.6.2 Dinamika Biplot Tahunan
Untuk memahami posisi spesifik setiap wilayah dan indikator pada tiap tahun, Gambar 4.26 hingga 4.30 menyajikan biplot tahunan secara individual.
Gambar 4.26: Biplot Tahunan 2020 (Masa
Awal Pandemi) Gambar 4.27: Biplot Tahunan 2021
Gambar 4.28: Biplot Tahunan 2022 (Awal
Masa Pemulihan) Gambar 4.29: Biplot Tahunan 2023
Universitas Indonesia

107
Gambar 4.30: Biplot Tahunan 2024
Membandingkan biplot pada Gambar 4.26 dan 4.30, terlihat pergeseran struktural yang
signifikan pada beberapa wilayah antara tahun 2020 dan 2024. Berdasarkan koordinat
skor intra pada Tabel 4.15 dan 4.19, pergeseran Dimensi 1 (PC1) terbesar ke arah positif
dialami oleh Kota Tangerang Selatan (+0,85 poin: dari 2,23 ke 3,09), diikuti Kota
Bekasi (+0,39) dan Kota Jakarta Pusat (+0,26), mengonfirmasi pemulihan kuat pada
kota-kota metropolitan berbasis jasa. Sebaliknya, pergeseran negatif terbesar dialami oleh
Kabupaten Pangandaran (−1,37 poin: dari −2,60 ke −3,97), disusul Kabupaten Garut
(−0,61) dan Kabupaten Cianjur (−0,51), mengindikasikan bahwa kabupaten pedalaman
Jawa Barat bagian selatan justru mengalami penurunan posisi relatif selama periode yang
sama.
4.6.3 Dinamika Struktur Wilayah Tahunan
analisis trajektori juga memungkinkan untuk mengamati perubahan posisi relatif antarwilayah pada tiap tahun. Tabel 4.15 dan 4.19 menyajikan skor kompromi intra-struktur
untuk masing-masing wilayah pada setiap tahun.
Universitas Indonesia

108
Tabel 4.15: Skor Intra-Struktur Wilayah Tahun 2020 hingga 2024
Tabel 4.16: Skor Intra Tahun
2020 (F2020)
Kabupaten/Kota PC1 (Asli) PC2 (Asli) PC1 (Visual) PC2 (Visual) Kep. Seribu -2,4946 -1,0804 -2,4946 -1,0804 Kota Jakarta Selatan 4,7264 -0,7688 4,7264 -0,7688 Kota Jakarta Timur 3,2871 -0,5831 3,2871 -0,5831 Kota Jakarta Pusat 3,9200 -0,6921 3,9200 -0,6921 Kota Jakarta Barat 2,9053 -0,0797 2,9053 -0,0797 Kota Jakarta Utara 2,4964 -0,4901 2,4964 -0,4901 Bogor -0,7714 1,6599 -0,7714 1,6599 Sukabumi -2,3806 0,5193 -2,3806 0,5193 Cianjur -3,2953 -0,4563 -3,2953 -0,4563 Bandung 0,3044 0,1156 0,3044 0,1156 Garut -2,5349 0,7781 -2,5349 0,7781 Tasikmalaya -4,0773 -2,0804 -4,0773 -2,0804 Ciamis -2,4376 -2,8807 -2,4376 -2,8807 Kuningan -1,1792 1,3542 -1,1792 1,3542 Cirebon -2,6304 1,0335 -2,6304 1,0335 Majalengka -3,2450 -2,5920 -3,2450 -2,5920 Sumedang -0,8621 -1,1387 -0,8621 -1,1387
Indramayu -2,6099 -1,3663 -2,6099 -1,3663 Subang -1,7863 -0,5727 -1,7863 -0,5727 Purwakarta -1,3807 1,3730 -1,3807 1,3730 Karawang -0,7398 0,1171 -0,7398 0,1171 Bekasi 1,6014 0,3242 1,6014 0,3242 Bandung Barat -2,0204 2,4503 -2,0204 2,4503 Pangandaran -2,5982 -4,3771 -2,5982 -4,3771 Kota Bogor 1,6180 1,6080 1,6180 1,6080 Kota Sukabumi 0,6049 2,4025 0,6049 2,4025 Kota Bandung 2,5924 -0,0698 2,5924 -0,0698 Kota Cirebon -0,0055 0,2968 -0,0055 0,2968 Kota Bekasi 3,0282 -0,3399 3,0282 -0,3399 Kota Depok 3,0639 -0,5731 3,0639 -0,5731 Kota Cimahi 1,9040 1,2714 1,9040 1,2714 Kota Tasikmalaya -1,2520 -1,2404 -1,2520 -1,2404 Kota Banjar -2,2202 -1,8799 -2,2202 -1,8799 Pandeglang -5,5059 -0,1055 -5,5059 -0,1055 Lebak -5,0571 -1,3031 -5,0571 -1,3031 Tangerang -1,0937 -0,0609 -1,0937 -0,0609 Serang -4,1050 -0,2479 -4,1050 -0,2479 Kota Tangerang 1,1965 -1,5058 1,1965 -1,5058 Kota Cilegon -1,7010 0,6711 -1,7010 0,6711 Kota Serang -1,8183 -0,7868 -1,8183 -0,7868 Kota Tangerang Selatan 2,2341 -0,8650 2,2341 -0,8650
Tabel 4.17: Skor Intra Tahun
2021 (F2021)
Kabupaten/Kota PC1 (Asli) PC2 (Asli) PC1 (Visual) PC2 (Visual) Kep. Seribu -2,4960 -1,6297 -2,4960 -1,6297 Kota Jakarta Selatan 4,3682 -1,4010 4,3682 -1,4010 Kota Jakarta Timur 3,4535 -0,4598 3,4535 -0,4598 Kota Jakarta Pusat 3,5633 -1,4125 3,5633 -1,4125 Kota Jakarta Barat 2,7284 -0,6437 2,7284 -0,6437 Kota Jakarta Utara 2,6166 -0,4332 2,6166 -0,4332 Bogor -0,8125 1,6170 -0,8125 1,6170 Sukabumi -2,1895 0,0797 -2,1895 0,0797 Cianjur -3,4367 -0,4044 -3,4367 -0,4044 Bandung 0,1201 -0,5178 0,1201 -0,5178 Garut -2,7261 0,4600 -2,7261 0,4600 Tasikmalaya -4,0037 -2,1133 -4,0037 -2,1133 Ciamis -2,3130 -2,8805 -2,3130 -2,8805 Kuningan -1,1283 1,6457 -1,1283 1,6457 Cirebon -2,6382 0,9318 -2,6382 0,9318 Majalengka -3,0717 -2,2535 -3,0717 -2,2535 Sumedang -1,2360 -1,0313 -1,2360 -1,0313
Indramayu -2,4043 -1,4763 -2,4043 -1,4763 Subang -1,6616 -0,3964 -1,6616 -0,3964 Purwakarta -1,2824 1,4685 -1,2824 1,4685 Karawang -0,5038 1,0774 -0,5038 1,0774 Bekasi 1,3779 -0,2026 1,3779 -0,2026 Bandung Barat -2,1674 2,5301 -2,1674 2,5301 Pangandaran -2,5940 -4,8963 -2,5940 -4,8963 Kota Bogor 1,7186 1,7562 1,7186 1,7562 Kota Sukabumi 0,7084 2,1920 0,7084 2,1920 Kota Bandung 2,6704 0,5184 2,6704 0,5184 Kota Cirebon 0,1510 0,9113 0,1510 0,9113 Kota Bekasi 3,2542 0,4183 3,2542 0,4183 Kota Depok 3,2971 0,1791 3,2971 0,1791 Kota Cimahi 2,0216 2,2224 2,0216 2,2224 Kota Tasikmalaya -1,0302 -0,9299 -1,0302 -0,9299 Kota Banjar -2,1637 -1,3541 -2,1637 -1,3541 Pandeglang -5,6607 -1,1583 -5,6607 -1,1583 Lebak -5,0151 -1,0759 -5,0151 -1,0759 Tangerang -1,1965 -0,7427 -1,1965 -0,7427 Serang -4,3897 -0,6520 -4,3897 -0,6520 Kota Tangerang 1,0581 -0,9518 1,0581 -0,9518 Kota Cilegon -1,8767 0,1254 -1,8767 0,1254 Kota Serang -1,9002 -0,2303 -1,9002 -0,2303 Kota Tangerang Selatan 2,1946 -0,6324 2,1946 -0,6324
Tabel 4.18: Skor Intra Tahun
2022 (F2022)
Kabupaten/Kota PC1 (Asli) PC2 (Asli) PC1 (Visual) PC2 (Visual) Kep. Seribu -2,6609 -1,1827 -2,6609 -1,1827 Kota Jakarta Selatan 4,3713 -1,8098 4,3713 -1,8098 Kota Jakarta Timur 3,6141 0,0661 3,6141 0,0661 Kota Jakarta Pusat 3,5474 -1,7398 3,5474 -1,7398 Kota Jakarta Barat 3,0528 -0,7996 3,0528 -0,7996 Kota Jakarta Utara 2,7174 -0,0817 2,7174 -0,0817 Bogor -0,7104 1,4232 -0,7104 1,4232 Sukabumi -2,6279 -0,4591 -2,6279 -0,4591 Cianjur -3,8567 -0,1697 -3,8567 -0,1697 Bandung 0,1906 -0,1063 0,1906 -0,1063 Garut -3,2219 -0,2056 -3,2219 -0,2056 Tasikmalaya -3,9956 -1,7949 -3,9956 -1,7949 Ciamis -1,9979 -2,0803 -1,9979 -2,0803 Kuningan -1,0980 1,8524 -1,0980 1,8524 Cirebon -2,5623 0,6521 -2,5623 0,6521 Majalengka -3,3045 -1,6111 -3,3045 -1,6111 Sumedang -0,7491 0,0458 -0,7491 0,0458
Indramayu -2,4469 -1,0389 -2,4469 -1,0389 Subang -1,7012 -0,2735 -1,7012 -0,2735 Purwakarta -1,8359 0,5985 -1,8359 0,5985 Karawang -0,4971 0,7499 -0,4971 0,7499 Bekasi 1,6353 0,9123 1,6353 0,9123 Bandung Barat -2,0948 1,7086 -2,0948 1,7086 Pangandaran -3,1161 -4,8291 -3,1161 -4,8291 Kota Bogor 1,5605 1,5472 1,5605 1,5472 Kota Sukabumi 0,5276 0,9738 0,5276 0,9738 Kota Bandung 2,5613 0,0584 2,5613 0,0584 Kota Cirebon 0,0481 0,2711 0,0481 0,2711 Kota Bekasi 3,4399 0,3121 3,4399 0,3121 Kota Depok 3,5290 -0,0834 3,5290 -0,0834 Kota Cimahi 1,7658 1,0687 1,7658 1,0687 Kota Tasikmalaya -1,0720 -0,5302 -1,0720 -0,5302 Kota Banjar -1,5992 -0,7466 -1,5992 -0,7466 Pandeglang -4,7428 0,3704 -4,7428 0,3704 Lebak -4,3814 0,1989 -4,3814 0,1989 Tangerang -1,0133 -0,9765 -1,0133 -0,9765 Serang -3,7159 0,1468 -3,7159 0,1468 Kota Tangerang 1,4533 -1,1260 1,4533 -1,1260 Kota Cilegon -1,9912 -1,1310 -1,9912 -1,1310 Kota Serang -1,6623 -0,4855 -1,6623 -0,4855 Kota Tangerang Selatan 2,7917 -0,4120 2,7917 -0,4120
Tabel 4.19: Skor Intra-Struktur Wilayah Tahun 2023 hingga 2024
Tabel 4.20: Skor Intra Tahun 2023 (F2023)
Kabupaten/Kota PC1 (Asli) PC2 (Asli) PC1 (Visual) PC2 (Visual)
Kep. Seribu -1,9880 -1,3104 -1,9880 -1,3104
Kota Jakarta Selatan 4,5312 -2,0797 4,5312 -2,0797
Kota Jakarta Timur 3,6892 -0,4308 3,6892 -0,4308
Kota Jakarta Pusat 4,1273 -0,9805 4,1273 -0,9805
Kota Jakarta Barat 3,2798 -0,9439 3,2798 -0,9439
Kota Jakarta Utara 2,4303 -0,6275 2,4303 -0,6275
Bogor -0,9257 0,9331 -0,9257 0,9331
Sukabumi -2,3371 0,0355 -2,3371 0,0355
Cianjur -4,1835 -0,3752 -4,1835 -0,3752
Bandung -0,0252 -0,4841 -0,0252 -0,4841
Garut -2,9745 0,0384 -2,9745 0,0384
Tasikmalaya -3,8870 -1,9319 -3,8870 -1,9319
Ciamis -2,1086 -1,7919 -2,1086 -1,7919
Kuningan -1,4326 2,6055 -1,4326 2,6055
Cirebon -2,6412 0,8547 -2,6412 0,8547
Majalengka -3,7095 -1,9189 -3,7095 -1,9189
Sumedang -1,0560 -0,4168 -1,0560 -0,4168
Indramayu -2,1127 0,1155 -2,1127 0,1155
Subang -2,1633 -0,1369 -2,1633 -0,1369
Purwakarta -1,9555 0,4661 -1,9555 0,4661
Karawang -0,2772 1,3379 -0,2772 1,3379
Bekasi 1,5891 0,9977 1,5891 0,9977
Bandung Barat -2,5939 1,2744 -2,5939 1,2744
Pangandaran -3,5152 -5,2657 -3,5152 -5,2657
Kota Bogor 1,8203 1,5281 1,8203 1,5281
Kota Sukabumi 1,0232 1,2056 1,0232 1,2056
Kota Bandung 3,0113 0,4985 3,0113 0,4985
Kota Cirebon 0,0484 -0,2400 0,0484 -0,2400
Kota Bekasi 3,5779 0,3525 3,5779 0,3525
Kota Depok 3,5872 0,0539 3,5872 0,0539
Kota Cimahi 1,8819 1,4421 1,8819 1,4421
Kota Tasikmalaya -0,4887 -0,1768 -0,4887 -0,1768
Kota Banjar -2,1534 -1,0657 -2,1534 -1,0657
Pandeglang -5,5305 1,3130 -5,5305 1,3130
Lebak -5,3797 0,0638 -5,3797 0,0638
Tangerang -1,4583 -0,8489 -1,4583 -0,8489
Serang -4,0692 0,9221 -4,0692 0,9221
Kota Tangerang 1,5774 -0,5859 1,5774 -0,5859
Kota Cilegon -1,5668 -0,3185 -1,5668 -0,3185
Kota Serang -2,3482 -0,9085 -2,3482 -0,9085
Kota Tangerang Selatan 2,4595 -0,6336 2,4595 -0,6336
Tabel 4.21: Skor Intra Tahun 2024 (F2024)
Kabupaten/Kota PC1 (Asli) PC2 (Asli) PC1 (Visual) PC2 (Visual)
Kep. Seribu -1,9977 -0,3337 -1,9977 -0,3337
Kota Jakarta Selatan 4,8746 -2,3228 4,8746 -2,3228
Kota Jakarta Timur 3,6643 -0,1300 3,6643 -0,1300
Kota Jakarta Pusat 4,1829 -1,1321 4,1829 -1,1321
Kota Jakarta Barat 3,2496 -0,9298 3,2496 -0,9298
Kota Jakarta Utara 2,3208 -1,2758 2,3208 -1,2758
Bogor -0,7842 0,7156 -0,7842 0,7156
Sukabumi -2,4619 0,2413 -2,4619 0,2413
Cianjur -3,8076 -0,8156 -3,8076 -0,8156
Bandung 0,1002 -0,0063 0,1002 -0,0063
Garut -3,1438 -0,0206 -3,1438 -0,0206
Tasikmalaya -3,7870 -1,6095 -3,7870 -1,6095
Ciamis -2,3132 -1,8830 -2,3132 -1,8830
Kuningan -1,3552 1,6416 -1,3552 1,6416
Cirebon -2,3712 0,5531 -2,3712 0,5531
Majalengka -3,3454 -1,5704 -3,3454 -1,5704
Sumedang -1,2850 -0,6748 -1,2850 -0,6748
Indramayu -2,4064 -0,2771 -2,4064 -0,2771
Subang -2,3731 -0,1345 -2,3731 -0,1345
Purwakarta -1,7764 1,0518 -1,7764 1,0518
Karawang 0,0551 1,0530 0,0551 1,0530
Bekasi 1,7143 1,1362 1,7143 1,1362
Bandung Barat -2,2283 1,2391 -2,2283 1,2391
Pangandaran -3,9706 -4,7580 -3,9706 -4,7580
Kota Bogor 1,0948 2,4814 1,0948 2,4814
Kota Sukabumi 1,0974 1,7995 1,0974 1,7995
Kota Bandung 2,4695 0,5602 2,4695 0,5602
Kota Cirebon -0,2397 0,1778 -0,2397 0,1778
Kota Bekasi 3,4150 0,8933 3,4150 0,8933
Kota Depok 3,2856 -0,1882 3,2856 -0,1882
Kota Cimahi 1,8021 1,7143 1,8021 1,7143
Kota Tasikmalaya -1,1595 -0,3090 -1,1595 -0,3090
Kota Banjar -2,0440 -0,9591 -2,0440 -0,9591
Pandeglang -5,1898 -1,2292 -5,1898 -1,2292
Lebak -5,2083 -2,2251 -5,2083 -2,2251
Tangerang -1,0430 -2,3321 -1,0430 -2,3321
Serang -3,6820 -1,1360 -3,6820 -1,1360
Kota Tangerang 1,3292 -2,4577 1,3292 -2,4577
Kota Cilegon -0,9079 -1,7153 -0,9079 -1,7153
Kota Serang -1,8243 -1,7970 -1,8243 -1,7970
Kota Tangerang Selatan 3,0868 -1,2871 3,0868 -1,2871
Universitas Indonesia

109
4.6.4 Pergeseran Centroid Tahunan
Selain mengamati lintasan individual setiap wilayah, dinamika kolektif seluruh
wilayah dapat dikuantifikasi melalui pergeseran centroid tahunan, yaitu posisi rata-rata
ck =

fk,1, fk,2

dari seluruh skor intra Fk pada tahun k. Jarak Euclidean antar-centroid
berturutan dihitung menggunakan Persamaan 2.68 (lihat penurunan contoh pada dataset
sampel di Subbab 3.7.1.8 Bab 3):
D(ck, ck+1) = q
(
¯fk,1 − ¯fk+1,1)
2 + ( ¯fk,2 − ¯fk+1,2)2
Nilai ini mengukur seberapa besar posisi “pusat berat” seluruh wilayah bergeser dalam
ruang kompromi dari satu tahun ke tahun berikutnya. Hasil lengkap disajikan pada
Tabel 4.22 berikut.
Tabel 4.22: Pergeseran Centroid Tahunan — STATIS (Skor Intra Fk)
Tahun Centroid PC1 Centroid PC2 D(ct, ct+1)
2020 -0,5932 -0,2966 0,0122
2021 -0,5999 -0,2865 0,0715
2022 -0,5329 -0,2614 0,0992
2023 -0,5914 -0,1813 0,2657
2024 -0,5601 -0,4452 –
Berdasarkan Tabel 4.22, pergeseran centroid terbesar terjadi pada transisi 2023–2024
dengan D(c2023, c2024) = 0,2657, yaitu sekitar 22 kali lebih besar dibandingkan transisi terkecil 2020–2021 (D = 0,0122). Lonjakan ini mengindikasikan adanya dinamika
struktural kolektif yang lebih intens menjelang akhir periode pengamatan, kemungkinan
mencerminkan percepatan pemulihan ekonomi yang tidak merata pasca-pandemi. Sebaliknya, kecilnya pergeseran pada 2020–2021 menunjukkan bahwa guncangan pandemi
bersifat seragam (uniform shock) sehingga tidak mengubah posisi relatif antar-wilayah
secara signifikan dalam ruang kompromi.
4.6.5 Jarak Pergeseran Per-Entitas Tahunan
Di samping pergeseran centroid kolektif, jarak pergeseran per-entitas mengukur seberapa jauh setiap wilayah atau indikator secara individual bergerak dalam ruang kompromi
dari satu tahun ke tahun berikutnya, dihitung menggunakan Persamaan 2.68. Penurunan
contoh komputasinya disajikan pada Paragraf 3.7.1.8 Bab 3.
Berdasarkan tabel jarak pergeseran per-entitas (tersedia di dashboard), wilayah dengan total jarak pergeseran kumulatif terbesar sepanjang 2020–2024 adalah Kabupaten
Universitas Indonesia

110
Pandeglang (PD = 6,641), Kabupaten Lebak (4,957), dan Kabupaten Serang (4,486) —
ketiganya merupakan wilayah pedalaman Banten yang berkarakter volatile. Sebaliknya,
wilayah dengan lintasan paling stabil adalah Kabupaten Tasikmalaya (PD = 0,912), Kota
Cirebon (1,014), dan Kabupaten Subang (1,037), yang menunjukkan profil pembangunan
yang relatif konsisten sepanjang periode. Pada dimensi indikator, Tingkat Partisipasi
Angkatan Kerja (TPAK, PD = 0,773) dan Gini Ratio (GR, 0,635) merupakan indikator
dengan pergeseran loading terbesar, mencerminkan dinamika distribusi tenaga kerja dan
ketimpangan yang fluktuatif; sebaliknya PDRB ADHK (0,006), RLS (0,020), dan IPG
(0,020) bersifat sebagai “jangkar” (anchor) stabilitas struktural pada sumbu kompromi.
4.6.6 Dinamika Struktur Indikator Tahunan
Selain untuk melihat pergeseran struktur wilayah tahunan, metode STATIS juga memungkinkan analisis terhadap perubahan hubungan antar-indikator itu sendiri pada setiap
tahun melalui matriks intra-loading (Gk) dan matriks korelasi terbobot (Sk). Tabel 4.24
hingga 4.29 menyajikan koordinat indikator pada tiap situasi tahunan.
Tabel 4.23: Loading Intra-Struktur Indikator Tahun 2020 hingga 2022
Tabel 4.24: Loading Intra
Tahun 2020 (G2020)
Indikator PC1 (Asli) PC2 (Asli) PC1 (Visual) PC2 (Visual)
IPG 0,8206 -0,2478 2,1230 -0,2729
PPM -0,8716 0,0502 -2,2548 0,0553
PDRB ADHK 0,7388 -0,1604 1,9113 -0,1766 RLS 0,9515 -0,0720 2,4616 -0,0793 AHH L 0,8936 0,0974 2,3119 0,1072 AHH P 0,8982 0,0893 2,3239 0,0984 GK 0,8798 -0,1781 2,2763 -0,1962
PKKP -0,8958 0,0309 -2,3176 0,0341 TPT 0,2634 0,7288 0,6814 0,8028 TPAK -0,2525 -0,7401 -0,6534 -0,8151 GR 0,5035 0,1117 1,3026 0,1231
Tabel 4.25: Loading Intra
Tahun 2021 (G2021)
Indikator PC1 (Asli) PC2 (Asli) PC1 (Visual) PC2 (Visual)
IPG 0,8216 -0,2493 2,1257 -0,2746
PPM -0,8859 0,0571 -2,2920 0,0628
PDRB ADHK 0,7391 -0,1574 1,9122 -0,1734 RLS 0,9534 -0,0714 2,4665 -0,0787 AHH L 0,8817 0,1121 2,2811 0,1235 AHH P 0,8902 0,1001 2,3031 0,1103 GK 0,8802 -0,1575 2,2772 -0,1734
PKKP -0,8768 0,0167 -2,2683 0,0184 TPT 0,1231 0,9153 0,3185 1,0082 TPAK -0,4197 -0,6099 -1,0857 -0,6718 GR 0,4518 0,1168 1,1690 0,1286
Tabel 4.26: Loading Intra
Tahun 2022 (G2022)
Indikator PC1 (Asli) PC2 (Asli) PC1 (Visual) PC2 (Visual)
IPG 0,8293 -0,2469 2,1456 -0,2720
PPM -0,8764 0,0488 -2,2674 0,0537
PDRB ADHK 0,7391 -0,1566 1,9122 -0,1725 RLS 0,9546 -0,0809 2,4696 -0,0891 AHH L 0,8844 0,1068 2,2880 0,1176 AHH P 0,8898 0,0945 2,3020 0,1041 GK 0,8855 -0,1501 2,2909 -0,1653
PKKP -0,9002 0,0413 -2,3290 0,0454 TPT 0,1752 0,9133 0,4532 1,0060 TPAK -0,5242 -0,3971 -1,3563 -0,4374 GR 0,5893 0,1459 1,5245 0,1607
Tabel 4.27: Loading Intra-Struktur Indikator Tahun 2023 hingga 2024
Tabel 4.28: Loading Intra Tahun 2023
(G2023)
Indikator PC1 (Asli) PC2 (Asli) PC1 (Visual) PC2 (Visual)
IPG 0,8246 -0,2522 2,1333 -0,2778
PPM -0,8852 0,0490 -2,2901 0,0540
PDRB ADHK 0,7390 -0,1570 1,9118 -0,1729
RLS 0,9544 -0,0866 2,4691 -0,0954
AHH L 0,8908 0,0955 2,3047 0,1052
AHH P 0,8730 0,1159 2,2585 0,1276
GK 0,8864 -0,1399 2,2932 -0,1540
PKKP -0,9286 0,0586 -2,4024 0,0645
TPT 0,1621 0,9065 0,4194 0,9985
TPAK -0,4860 -0,4521 -1,2573 -0,4979
GR 0,7771 0,0456 2,0105 0,0502
Tabel 4.29: Loading Intra Tahun 2024
(G2024)
Indikator PC1 (Asli) PC2 (Asli) PC1 (Visual) PC2 (Visual)
IPG 0,8278 -0,2532 2,1416 -0,2789
PPM -0,8888 0,0496 -2,2994 0,0547
PDRB ADHK 0,7392 -0,1568 1,9124 -0,1727
RLS 0,9552 -0,0889 2,4713 -0,0979
AHH L 0,8948 0,0799 2,3150 0,0880
AHH P 0,8656 0,1265 2,2394 0,1393
GK 0,8830 -0,1267 2,2844 -0,1396
PKKP -0,6501 0,3753 -1,6819 0,4133
TPT 0,2793 0,8422 0,7227 0,9277
TPAK -0,6862 -0,2910 -1,7753 -0,3206
GR 0,7751 0,2748 2,0054 0,3026
Stabilitas orientasi indikator dalam tabel di atas mengonfirmasi hasil uji RV yang
Universitas Indonesia

111
tinggi. Indikator IPM dan PDRB secara konsisten memiliki loading positif yang kuat pada
dimensi pertama di seluruh tahun, yang menunjukkan bahwa sumbu horizontal kompromi
memang merupakan representasi yang valid untuk kemajuan ekonomi dan kualitas manusia
secara stabil di seluruh periode pengamatan.
Matriks korelasi terbobot indikator (Sk) yang tersaji pada lampiran (atau file output)
memperlihatkan bahwa hubungan inter-indikator tidak mengalami patahan struktural
meskipun terjadi pandemi, yang memperkuat validitas penggunaan satu struktur kompromi
tunggal untuk keseluruhan periode pengamatan.
4.7 Analisis Klaster Spasio-Temporal
Analisis klaster dilakukan dalam tiga tahap: (1) penentuan jumlah klaster optimal, (2)
interpretasi pola pengelompokan dan matriks transisi, serta (3) evaluasi konsistensi klaster
terhadap standar kategori IPM UNDP.
4.7.1 Penentuan Jumlah Klaster Optimal
Dua metode digunakan secara komplementer untuk menentukan jumlah klaster k yang
optimal: Metode Elbow yang mengamati titik patahan kurva inersia (Within-Klaster Sum of
Squares/SSE) dan Agglomerative Hierarchical Clustering (Dendrogram) yang mengamati
jarak penggabungan antar-kelompok. Gambar 4.31 hingga Gambar 4.35 menyajikan hasil
kedua metode tersebut untuk setiap tahun pengamatan (2020–2024).
Gambar 4.31: Kurva Elbow dan Dendrogram Tahun 2020
Universitas Indonesia

112
Gambar 4.32: Kurva Elbow dan Dendrogram Tahun 2021
Gambar 4.33: Kurva Elbow dan Dendrogram Tahun 2022
Gambar 4.34: Kurva Elbow dan Dendrogram Tahun 2023
Universitas Indonesia

113
Gambar 4.35: Kurva Elbow dan Dendrogram Tahun 2024
Berdasarkan evaluasi kurva SSE pada Gambar 4.31 hingga 4.35, titik patahan yang
konsisten muncul pada k = 3 di seluruh tahun pengamatan. Hal ini dikonfirmasi oleh
dendrogram yang menunjukkan tiga cabang utama dengan jarak penggabungan yang
relatif besar. Sebagai konfirmasi kuantitatif tambahan, koefisien siluet rata-rata pada skor
kompromi aktual untuk k = 3 adalah SC = 0, 3883 (lihat prosedur perhitungan pada
Subbab 3.7.1.9 dan Persamaan 2.65), yang menunjukkan pemisahan klaster yang memadai.
Oleh karena itu, k = 3 ditetapkan sebagai jumlah klaster optimal yang digunakan pada
seluruh analisis selanjutnya.
4.7.2 Pola Pengelompokan Tahunan
Dengan k = 3, algoritma K-Means diterapkan pada skor kompromi setiap tahun.
Gambar 4.36 hingga 4.40 menunjukkan hasil pengelompokan tersebut dalam ruang biplot.
Gambar 4.36: Pengelompokan K-Means (k = 3) Tahun 2020: Biplot (kiri) dan Peta
(kanan)
Universitas Indonesia

114
Gambar 4.37: Pengelompokan K-Means (k = 3) Tahun 2021: Biplot (kiri) dan Peta
(kanan)
Gambar 4.38: Pengelompokan K-Means (k = 3) Tahun 2022: Biplot (kiri) dan Peta
(kanan)
Universitas Indonesia

115
Gambar 4.39: Pengelompokan K-Means (k = 3) Tahun 2023: Biplot (kiri) dan Peta
(kanan)
Gambar 4.40: Pengelompokan K-Means (k = 3) Tahun 2024: Biplot (kiri) dan Peta
(kanan)
Berdasarkan posisi sentroid klaster pada Dimensi 1 kompromi (skor rata-rata PC1
terendah ke tertinggi) dan wilayah yang termasuk di dalamnya, ketiga klaster dapat diidentifikasi secara ordinal sebagai berikut. Nama deskriptif setiap klaster ditetapkan
berdasarkan interpretasi terhadap profil sentroid dan keanggotaan wilayah pada Tabel 4.34
dan Tabel 4.35.
• Klaster 1 (Tertinggal): Klaster dengan sentroid Dimensi 1 terendah, mencerminkan
wilayah dengan nilai indikator pembangunan secara rata-rata paling rendah di antara
Universitas Indonesia

116
ketiga kelompok. Klaster ini menghimpun 12 wilayah: Kabupaten Pandeglang,
Lebak, Serang (Banten Selatan), Garut, Tasikmalaya, Cianjur, Majalengka, Indramayu, Ciamis, Pangandaran, Kota Banjar (Jawa Barat pedalaman selatan), serta
Kepulauan Seribu. Profil kelompok ini dicirikan oleh nilai PPM tinggi, RLS dan
AHH rendah, dengan 9 dari 12 anggota berkategori IPM “Sedang” menurut standar UNDP yang mengonfirmasi ketertinggalan pembangunan manusia yang terus
berlangsung selama periode pengamatan.
• Klaster 2 (Berkembang): Klaster dengan sentroid Dimensi 1 menengah, mengelompokkan wilayah-wilayah transisi yang menunjukkan profil pembangunan moderat.
Klaster ini terdiri dari 16 wilayah yang membentang di koridor industri Jawa Barat
bagian utara dan wilayah penyangga Banten: Bogor, Sukabumi, Bandung, Kuningan,
Cirebon, Sumedang, Subang, Purwakarta, Karawang, Bandung Barat, Kota Sukabumi, Kota Cirebon, Kota Tasikmalaya, Kabupaten Tangerang, Kota Cilegon, dan
Kota Serang. Sebanyak 15 dari 16 anggotanya berkategori IPM “Tinggi” (70–79,99),
mencerminkan wilayah yang sedang dalam proses transisi pembangunan menuju
tahap yang lebih maju.
• Klaster 3 (Maju): Klaster dengan sentroid Dimensi 1 tertinggi, dicirikan oleh
nilai skor kompromi positif yang kuat yang mencerminkan keunggulan relatif pada
indikator-indikator utama. Klaster ini mengelompokkan 13 wilayah metropolitan:
kelima kota administrasi DKI Jakarta, serta Kabupaten Bekasi, Kota Bogor, Kota
Bandung, Kota Bekasi, Kota Depok, Kota Cimahi, Kota Tangerang, dan Kota
Tangerang Selatan. Profil sentroid klaster ini unggul pada RLS, AHH, dan PDRB
ADHK tertinggi, dengan PPM terendah di antara ketiga klaster. Sebelas dari 13
anggota (84,6%) berkategori IPM “Sangat Tinggi” (≥ 80) menurut standar UNDP.
4.7.3 Matriks Transisi Klaster STATIS
pergeseran wilayah antar-klaster dari satu periode ke periode berikutnya dianalisis
melalui matriks transisi klaster berbasis koordinat STATIS tahunan (Fk) (lihat persamaan
2.40). Berbeda dengan matriks transisi biplot yang tidak mempertimbangkan bobot
temporal, matriks ini dibangun dari hasil K-Means pada ruang yang telah dipengaruhi
oleh nilai αk , sehingga perpindahan klaster yang terekam merupakan perpindahan yang
bermakna secara struktural. Sebagai periode representatif, disajikan matriks transisi 2023
hingga 2024 untuk kedua skenario pembobotan.
Universitas Indonesia

117
Skenario 1 – ACT Klasik, Host = Wilayah (Periode 2023 – 2024)
Tabel 4.30: Matriks Transisi Klaster 2023 → 2024
2023
2024 Cluster 1 (Rendah) (2024) Cluster 2 (Sedang) (2024) Cluster 3 (Tinggi) (2024)
Cluster 1 (Rendah) (2023) 7 0 0
Cluster 2 (Sedang) (2023) 14 6 0
Cluster 3 (Tinggi) (2023) 0 4 10
Tabel 4.31: Nama Daerah Perpindahan Klaster 2023 → 2024
Kategori Daftar Wilayah
Cluster 1 (Rendah)
(2023) → Cluster 1
(Rendah) (2024)
Cianjur, Tasikmalaya, Majalengka, Pangandaran, Pandeglang, Lebak, Serang
Cluster 2 (Sedang)
(2023) → Cluster 1
(Rendah) (2024)
Kep. Seribu, Sukabumi, Garut, Ciamis, Cirebon, Sumedang,
Indramayu, Subang, Bandung Barat, Kota Tasikmalaya, Kota
Banjar, Tangerang, Kota Cilegon, Kota Serang
Cluster 2 (Sedang)
(2023) → Cluster 2
(Sedang) (2024)
Bogor, Bandung, Kuningan, Purwakarta, Karawang, Kota
Cirebon
Cluster 3 (Tinggi)
(2023) → Cluster 2
(Sedang) (2024)
Bekasi, Kota Bogor, Kota Sukabumi, Kota Cimahi
Cluster 3 (Tinggi)
(2023) → Cluster 3
(Tinggi) (2024)
Kota Jakarta Selatan, Kota Jakarta Timur, Kota Jakarta Pusat,
Kota Jakarta Barat, Kota Jakarta Utara, Kota Bandung, Kota
Bekasi, Kota Depok, Kota Tangerang, Kota Tangerang Selatan
Skenario 2 – Dual STATIS, Host = Indikator (Periode 2023 – 2024)
Tabel 4.32: Matriks Transisi Klaster 2023 → 2024
2023
2024 Cluster 1 (Rendah) (2024) Cluster 2 (Sedang) (2024) Cluster 3 (Tinggi) (2024)
Cluster 1 (Rendah) (2023) 6 1 0
Cluster 2 (Sedang) (2023) 1 19 0
Cluster 3 (Tinggi) (2023) 0 0 14
Universitas Indonesia

118
Tabel 4.33: Nama Daerah Perpindahan Klaster 2023 → 2024
Kategori Daftar Wilayah
Cluster 1 (Rendah)
(2023) → Cluster 1
(Rendah) (2024)
Cianjur, Tasikmalaya, Majalengka, Pangandaran, Pandeglang, Lebak
Cluster 1 (Rendah)
(2023) → Cluster 2
(Sedang) (2024)
Serang
Cluster 2 (Sedang)
(2023) → Cluster 1
(Rendah) (2024)
Ciamis
Cluster 2 (Sedang)
(2023) → Cluster 2
(Sedang) (2024)
Kep. Seribu, Bogor, Sukabumi, Bandung, Garut, Kuningan,
Cirebon, Sumedang, Indramayu, Subang, Purwakarta,
Karawang, Bandung Barat, Kota Cirebon, Kota Tasikmalaya,
Kota Banjar, Tangerang, Kota Cilegon, Kota Serang
Cluster 3 (Tinggi)
(2023) → Cluster 3
(Tinggi) (2024)
Kota Jakarta Selatan, Kota Jakarta Timur, Kota Jakarta Pusat,
Kota Jakarta Barat, Kota Jakarta Utara, Bekasi, Kota Bogor,
Kota Sukabumi, Kota Bandung, Kota Bekasi, Kota Depok,
Kota Cimahi, Kota Tangerang, Kota Tangerang Selatan
Gambar 4.41 memperlihatkan perbandingan peta persebaran klaster STATIS antara
tahun 2023 dan 2024 untuk kedua skenario pembobotan.
Skenario 1 Skenario 2
Gambar 4.41: Perbandingan Peta Klaster STATIS Periode 2023 vs 2024 (Kiri: Tahun
2023; Kanan: Tahun 2024)
Elemen-elemen diagonal pada kedua matriks transisi memiliki nilai dominan, yang
mengonfirmasi stabilitas keanggotaan klaster yang tinggi sebagaimana diprediksikan oleh
nilai koefisien RV yang besar (> 0, 9543). Nilai koefisien RV yang tinggi antar-tahun
(Tabel 4.3, Subbab 4.3.1) mengimplikasikan bahwa perubahan konfigurasi multivariate
antar-wilayah bersifat inkremental, sehingga sebagian besar wilayah tetap mempertahankan posisi klasternya dari tahun ke tahun. Periode 2023–2024 mencatat restrukturisasi
klaster terbesar: 14 wilayah bergeser dari Klaster 2 ke Klaster 1 (Kepulauan Seribu,
Sukabumi, Garut, Ciamis, Cirebon, Sumedang, Indramayu, Subang, Bandung Barat, Kota
Tasikmalaya, Kota Banjar, Tangerang, Kota Cilegon, dan Kota Serang), serta 4 wilayah
Universitas Indonesia

119
bergeser dari Klaster 3 ke Klaster 2 (Bekasi, Kota Bogor, Kota Sukabumi, dan Kota
Cimahi). Sementara itu, 10 wilayah metropolitan inti (kelima kota DKI, Kota Bandung,
Kota Bekasi, Kota Depok, Kota Tangerang, dan Kota Tangerang Selatan) tetap bertahan
di Klaster 3, mengonfirmasi ketahanan struktural yang kokoh. Perpindahan ini mencerminkan pergeseran posisi relatif antar-wilayah. Hal ini bukan penurunan nilai indikator
yang sebenarnya, karena K-Means dijalankan ulang setiap tahun: sentroid setiap klaster
dihitung ulang dari data tahun berjalan, sehingga wilayah yang capaiannya stagnan atau
tumbuh lebih lambat dibanding wilayah-wilayah lain akan tergeser ke klaster yang lebih
rendah meskipun kondisi aslinya tidak memburuk.
4.7.4 Klaster Trajektori Keseluruhan
Selain pengelompokan tahunan, dilakukan pengelompokan berbasis trajektori keseluruhan selama 2020 hingga 2024 dengan menggunakan fitur koordinat gabungan dari semua
tahun pada ruang kompromi. Pendekatan ini menghasilkan label klaster yang bersifat
longitudinal, yaitu mencerminkan posisi rata-rata dan pola pergeseran wilayah sepanjang
periode pengamatan, bukan posisi pada satu tahun tertentu. Jumlah klaster optimal k = 3
mengacu pada evaluasi kurva SSE dan dendrogram pada Subbab 4.7.1, yang menunjukkan
titik patahan konsisten di k = 3 pada seluruh tahun pengamatan. Dengan demikian, klaster
trajektori keseluruhan dapat diinterpretasikan sebagai representasi yang lebih holistik dari
pola pembangunan wilayah selama lima tahun, yang mengintegrasikan informasi temporal
dan spasial secara simultan. Gambar 4.42 menyajikan scatter plot klaster trajektori pada
ruang kompromi untuk kedua skenario.
Universitas Indonesia

120
Skenario 1 Skenario 2
Gambar 4.42: Scatter Plot Klaster Trajektori pada Ruang Kompromi
Dari Gambar 4.42 terlihat bahwa klaster trajektori keseluruhan memiliki pola yang
sangat mirip dengan klaster tahunan, namun dengan beberapa perbedaan penting. Klaster
trajektori cenderung lebih homogen dalam hal pola pergeseran wilayah, sehingga wilayah
yang memiliki lintasan yang serupa selama lima tahun akan dikelompokkan bersama,
meskipun posisi mereka pada satu tahun tertentu mungkin berbeda. Selanjutnya, Gambar 4.43 memperlihatkan persebaran geografis klaster trajektori keseluruhan untuk kedua
skenario.
Skenario 1 Skenario 2
Gambar 4.43: Peta Persebaran Klaster Trajektori Keseluruhan
Kemudian dari Gambar 4.43 terlihat bahwa klaster trajektori keseluruhan memiliki
distribusi geografis yang cukup konsisten dengan klaster tahunan, namun dengan beberapa
perbedaan penting. Terdapat wilayah dengan klaster kompromi (trajektori) yang berbeda
Universitas Indonesia

121
dari klaster STATIS tahunan 2024: Kabupaten Bekasi, Kota Bogor, dan Kota Cimahi
termasuk Klaster 3 kompromi (berdasarkan Tabel 4.34) namun bergeser ke Klaster 2
pada tahun 2024, mencerminkan moderasi kinerja pada tahun terakhir yang tidak cukup
kuat untuk mengubah penilaian jangka-menengah. Kota Sukabumi berada di Klaster 2
kompromi meskipun sempat naik ke Klaster 3 pada 2023 sebelum kembali ke Klaster 2 di
2024, menunjukkan mobilitas naik yang bersifat temporer (Lihat Tabel 4.31). Perbedaan
ini mengonfirmasi nilai klaster kompromi sebagai representasi holistik jangka menengah
yang lebih stabil dan representatif dibandingkan klaster tahunan yang dapat terpengaruh
fluktuasi satu periode. Distribusi proporsi wilayah pada setiap klaster divisualisasikan pada
Gambar 4.44.
Skenario 1 Skenario 2
Gambar 4.44: Distribusi Proporsi Klaster Trajektori Keseluruhan
jika dilihat dari distribusi proporsi klaster pada Gambar 4.44, terlihat bahwa klaster
trajektori keseluruhan memiliki proporsi wilayah yang cukup seimbang, dengan sedikit
dominasi oleh Klaster 2 (Berkembang). Hal ini menunjukkan bahwa sebagian besar
wilayah mengalami pola pergeseran yang serupa selama lima tahun, yang mencerminkan
tren pembangunan yang relatif homogen di tingkat regional.
Tabel 4.34 dan Tabel 4.35 merinci anggota setiap klaster trajektori untuk masingmasing skenario.
Universitas Indonesia

122
Tabel 4.34: Ringkasan Anggota Klaster Kompromi STATIS (Skenario 1)
Kelompok (Cluster)Jumlah Wilayah Daftar Kabupaten/Kota
Cluster 1 (Rendah)12 Kep. Seribu, Cianjur, Garut, Tasikmalaya, Ciamis, Majalengka, Indramayu, Pangandaran, Kota Banjar, Pandeglang,
Lebak, Serang
Cluster 2 (Sedang)16 Bogor, Sukabumi, Bandung, Kuningan, Cirebon, Sumedang,
Subang, Purwakarta, Karawang, Bandung Barat, Kota Sukabumi, Kota Cirebon, Kota Tasikmalaya, Tangerang, Kota
Cilegon, Kota Serang
Cluster 3 (Tinggi) 13 Kota Jakarta Selatan, Kota Jakarta Timur, Kota Jakarta Pusat,
Kota Jakarta Barat, Kota Jakarta Utara, Bekasi, Kota Bogor,
Kota Bandung, Kota Bekasi, Kota Depok, Kota Cimahi, Kota
Tangerang, Kota Tangerang Selatan
Tabel 4.35: Ringkasan Anggota Klaster Kompromi STATIS (Skenario 2)
Kelompok (Cluster)Jumlah Wilayah Daftar Kabupaten/Kota
Cluster 1 (Rendah)12 Kep. Seribu, Cianjur, Garut, Tasikmalaya, Ciamis, Majalengka, Indramayu, Pangandaran, Kota Banjar, Pandeglang,
Lebak, Serang
Cluster 2 (Sedang)16 Bogor, Sukabumi, Bandung, Kuningan, Cirebon, Sumedang,
Subang, Purwakarta, Karawang, Bandung Barat, Kota Sukabumi, Kota Cirebon, Kota Tasikmalaya, Tangerang, Kota
Cilegon, Kota Serang
Cluster 3 (Tinggi) 13 Kota Jakarta Selatan, Kota Jakarta Timur, Kota Jakarta Pusat,
Kota Jakarta Barat, Kota Jakarta Utara, Bekasi, Kota Bogor,
Kota Bandung, Kota Bekasi, Kota Depok, Kota Cimahi, Kota
Tangerang, Kota Tangerang Selatan
4.8 Analisis Komparatif Sembilan Jenis Perbandingan Klaster
Evaluasi komprehensif terhadap hasil pengelompokan dilakukan melalui sembilan
jenis analisis perbandingan (comparative clustering analysis) yang diorganisasikan dalam
empat kategori: (i) pergeseran temporal berbasis metode yang berbeda (Subbab 4.8.1–
4.7.3), (ii) kongruensi klaster antar-metode pada tahun yang sama (Subbab 4.8.2), (iii)
validasi eksternal terhadap klasifikasi normatif IPM UNDP (Subbab 4.8.3–4.8.5), serta (iv)
evaluasi konsistensi klaster kompromi terhadap klaster tahunan dan standar internasional
(Subbab 4.8.6–4.8.8). Untuk setiap jenis perbandingan, disajikan satu periode representatif
dari masing-masing skenario pembobotan.
Sebelum menyusun matriks transisi, label klaster antar-tahun diselaraskan menggunakan Hungarian Algorithm (Kuhn, 1955) berbasis matriks jarak centroid (Persamaan 2.68): centroid setiap klaster tahun dasar c1i dibandingkan dengan centroid
Universitas Indonesia

123
klaster tahun uji c2j, dan pemetaan bipartit berbobot minimum dipilih sehingga label
yang diberikan konsisten antar-tahun. Penurunan contoh komputasi matriks jarak 3 × 3
beserta pemilihan penugasan optimal disajikan pada Subbab 3.7.1.11 Bab 3.
4.8.1 Perbandingan 1 — Matriks Transisi Klaster Biplot Antar-Tahun
Perbandingan pertama menelaah pergeseran wilayah pada ruang Biplot PCA tanpa
pembobotan temporal (unweighted), yakni berdasarkan dekomposisi XT
k Xk per tahun
tanpa pengaruh bobot αk dari analisis inter-struktur. Karena koordinat Biplot PCA bersifat
skenario-bebas (tidak bergantung pada skema pembobotan wilayah atau indikator), matriks
transisi ini merupakan referensi tunggal (baseline) yang berlaku untuk kedua skenario.
Hasil periode 2023–2024 disajikan sebagai representasi akhir periode pengamatan.
Tabel 4.36: Matriks Transisi Klaster Biplot (Unweighted) 2023 → 2024
2023
2024 Cluster 1 (Rendah) (2024) Cluster 2 (Sedang) (2024) Cluster 3 (Tinggi) (2024)
Cluster 1 (Rendah) (2023) 6 0 0
Cluster 2 (Sedang) (2023) 0 21 0
Cluster 3 (Tinggi) (2023) 0 0 14
Tabel 4.37: Nama Daerah Transisi Klaster Biplot 2023 → 2024
Kategori Daftar Wilayah
Cluster 1 (Rendah)
(2023) → Cluster 1
(Rendah) (2024)
Pandeglang, Lebak, Tangerang, Serang, Kota Cilegon, Kota
Serang
Cluster 2 (Sedang)
(2023) → Cluster 2
(Sedang) (2024)
Kep. Seribu, Bogor, Sukabumi, Cianjur, Bandung, Garut,
Tasikmalaya, Ciamis, Kuningan, Cirebon, Majalengka,
Sumedang, Indramayu, Subang, Purwakarta, Karawang,
Bandung Barat, Pangandaran, Kota Cirebon, Kota Tasikmalaya, Kota Banjar
Cluster 3 (Tinggi)
(2023) → Cluster 3
(Tinggi) (2024)
Kota Jakarta Selatan, Kota Jakarta Timur, Kota Jakarta Pusat,
Kota Jakarta Barat, Kota Jakarta Utara, Bekasi, Kota Bogor,
Kota Sukabumi, Kota Bandung, Kota Bekasi, Kota Depok,
Kota Cimahi, Kota Tangerang, Kota Tangerang Selatan
Gambar 4.45 memperlihatkan perbandingan peta klaster Biplot antara tahun 2023 dan
2024.
Universitas Indonesia

124
Gambar 4.45: Perbandingan Peta Klaster Biplot (Unweighted): Tahun 2023 (kiri) vs 2024
(kanan)
Tingginya nilai elemen diagonal pada matriks transisi Biplot mengindikasikan bahwa
posisi relatif wilayah dalam ruang PCA tahunan juga cukup stabil, meskipun tidak memperhitungkan struktur konsensus lintas tahun. Perbedaan antara stabilitas klaster Biplot
dan klaster STATIS (Subbab 4.7.3) akan dikuantifikasi lebih lanjut pada Subbab 4.8.2.
4.8.2 Perbandingan 3 — Kongruensi Klaster Biplot vs STATIS per Tahun
Perbandingan ketiga mengukur sejauh mana label klaster yang dihasilkan oleh KMeans Biplot PCA (unweighted) berbeda dari label klaster STATIS berbobot (Fk) pada
tahun yang sama. Perbedaan ini mencerminkan dampak kuantitatif penggabungan informasi inter-struktur (α) terhadap konfigurasi klaster. Disajikan tabulasi silang untuk tahun
2024 dari masing-masing skenario.
Skenario 1 – ACT Klasik (Tahun 2024)
Tabel 4.38: Matriks Perbandingan Klaster Biplot vs STATIS (Tahunan) (2024)
Biplot (Unweighted) Statis 1 (Rendah) Statis 2 (Sedang) Statis 3 (Tinggi)
Biplot 1 (Rendah) 6 0 0
Biplot 2 (Sedang) 15 6 0
Biplot 3 (Tinggi) 0 4 10
Universitas Indonesia

125
Tabel 4.39: Nama Daerah Perbandingan Klaster Biplot vs STATIS (Tahunan) (2024)
Kategori Daftar Wilayah
Biplot 1 (Rendah) &
Statis 1 (Rendah)
Pandeglang, Lebak, Tangerang, Serang, Kota Cilegon, Kota
Serang
Biplot 2 (Sedang) &
Statis 1 (Rendah)
Kep. Seribu, Sukabumi, Cianjur, Garut, Tasikmalaya,
Ciamis, Cirebon, Majalengka, Sumedang, Indramayu, Subang, Bandung Barat, Pangandaran, Kota Tasikmalaya, Kota
Banjar
Biplot 2 (Sedang) &
Statis 2 (Sedang)
Bogor, Bandung, Kuningan, Purwakarta, Karawang, Kota
Cirebon
Biplot 3 (Tinggi) &
Statis 2 (Sedang)
Bekasi, Kota Bogor, Kota Sukabumi, Kota Cimahi
Biplot 3 (Tinggi) &
Statis 3 (Tinggi)
Kota Jakarta Selatan, Kota Jakarta Timur, Kota Jakarta Pusat,
Kota Jakarta Barat, Kota Jakarta Utara, Kota Bandung, Kota
Bekasi, Kota Depok, Kota Tangerang, Kota Tangerang Selatan
Tabel 4.40: Matriks Perbandingan Klaster Biplot vs STATIS (vs Kompromi) (2024)
Statis (Compromise) Biplot 1 (Rendah) (2024) Biplot 2 (Sedang) (2024) Biplot 3 (Tinggi) (2024)
Statis 1 (Rendah) 12 0 0
Statis 2 (Sedang) 9 7 0
Statis 3 (Tinggi) 0 3 10
Tabel 4.41: Nama Daerah Perbandingan Klaster Biplot vs STATIS (vs Kompromi) (2024)
Kategori Daftar Wilayah
Statis 1 (Rendah) &
Biplot 1 (Rendah)
Kep. Seribu, Cianjur, Garut, Tasikmalaya, Ciamis, Majalengka, Indramayu, Pangandaran, Kota Banjar, Pandeglang,
Lebak, Serang
Statis 2 (Sedang) &
Biplot 1 (Rendah)
Sukabumi, Cirebon, Sumedang, Subang, Bandung Barat,
Kota Tasikmalaya, Tangerang, Kota Cilegon, Kota Serang
Statis 2 (Sedang) &
Biplot 2 (Sedang)
Bogor, Bandung, Kuningan, Purwakarta, Karawang, Kota
Sukabumi, Kota Cirebon
Statis 3 (Tinggi) &
Biplot 2 (Sedang)
Bekasi, Kota Bogor, Kota Cimahi
Statis 3 (Tinggi) &
Biplot 3 (Tinggi)
Kota Jakarta Selatan, Kota Jakarta Timur, Kota Jakarta Pusat,
Kota Jakarta Barat, Kota Jakarta Utara, Kota Bandung, Kota
Bekasi, Kota Depok, Kota Tangerang, Kota Tangerang Selatan
Universitas Indonesia

126
Skenario 2 – Dual STATIS (Tahun 2024)
Tabel 4.42: Matriks Perbandingan Klaster Biplot vs STATIS (Tahunan) (2024)
Biplot (Unweighted) Statis 1 (Rendah) Statis 2 (Sedang) Statis 3 (Tinggi)
Biplot 1 (Rendah) 2 4 0
Biplot 2 (Sedang) 5 16 0
Biplot 3 (Tinggi) 0 0 14
Tabel 4.43: Nama Daerah Perbandingan Klaster Biplot vs STATIS (Tahunan) (2024)
Kategori Daftar Wilayah
Biplot 1 (Rendah) &
Statis 1 (Rendah)
Pandeglang, Lebak
Biplot 1 (Rendah) &
Statis 2 (Sedang)
Tangerang, Serang, Kota Cilegon, Kota Serang
Biplot 2 (Sedang) &
Statis 1 (Rendah)
Cianjur, Tasikmalaya, Ciamis, Majalengka, Pangandaran
Biplot 2 (Sedang) &
Statis 2 (Sedang)
Kep. Seribu, Bogor, Sukabumi, Bandung, Garut, Kuningan,
Cirebon, Sumedang, Indramayu, Subang, Purwakarta,
Karawang, Bandung Barat, Kota Cirebon, Kota Tasikmalaya,
Kota Banjar
Biplot 3 (Tinggi) &
Statis 3 (Tinggi)
Kota Jakarta Selatan, Kota Jakarta Timur, Kota Jakarta Pusat,
Kota Jakarta Barat, Kota Jakarta Utara, Bekasi, Kota Bogor,
Kota Sukabumi, Kota Bandung, Kota Bekasi, Kota Depok,
Kota Cimahi, Kota Tangerang, Kota Tangerang Selatan
Tabel 4.44: Matriks Perbandingan Klaster Biplot vs STATIS (vs Kompromi) (2024)
Statis (Compromise) Biplot 1 (Rendah) (2024) Biplot 2 (Sedang) (2024) Biplot 3 (Tinggi) (2024)
Statis 1 (Rendah) 7 5 0
Statis 2 (Sedang) 0 15 1
Statis 3 (Tinggi) 0 0 13
Universitas Indonesia

127
Tabel 4.45: Nama Daerah Perbandingan Klaster Biplot vs STATIS (vs Kompromi) (2024)
Kategori Daftar Wilayah
Statis 1 (Rendah) &
Biplot 1 (Rendah)
Cianjur, Tasikmalaya, Ciamis, Majalengka, Pangandaran,
Pandeglang, Lebak
Statis 1 (Rendah) &
Biplot 2 (Sedang)
Kep. Seribu, Garut, Indramayu, Kota Banjar, Serang
Statis 2 (Sedang) &
Biplot 2 (Sedang)
Bogor, Sukabumi, Bandung, Kuningan, Cirebon, Sumedang,
Subang, Purwakarta, Karawang, Bandung Barat, Kota Cirebon, Kota Tasikmalaya, Tangerang, Kota Cilegon, Kota
Serang
Statis 2 (Sedang) &
Biplot 3 (Tinggi)
Kota Sukabumi
Statis 3 (Tinggi) &
Biplot 3 (Tinggi)
Kota Jakarta Selatan, Kota Jakarta Timur, Kota Jakarta Pusat,
Kota Jakarta Barat, Kota Jakarta Utara, Bekasi, Kota Bogor,
Kota Bandung, Kota Bekasi, Kota Depok, Kota Cimahi, Kota
Tangerang, Kota Tangerang Selatan
Gambar 4.46 memperlihatkan perbandingan distribusi geografis klaster Biplot dan
klaster STATIS tahun 2024 untuk kedua skenario.
Skenario 1 Skenario 2
Gambar 4.46: Perbandingan Peta Klaster Biplot (kiri) vs STATIS (kanan) Tahun 2024
Elemen-elemen diagonal yang dominan pada tabulasi silang mengindikasikan tingkat
kongruensi yang tinggi antara kedua metode, mengonfirmasi bahwa penggabungan bobot
temporal oleh STATIS tidak mengubah struktur klaster secara fundamental. Perbedaan
yang teramati pada elemen di luar diagonal utama merepresentasikan wilayah-wilayah yang
berada di batas pemisah antar-klaster, di mana pembobotan temporal STATIS berperan
sebagai penyeimbang (stabilizer) terhadap fluktuasi lokal tahunan.
4.8.3 Perbandingan 4 — Matriks Transisi Kategori IPM UNDP
Perbandingan keempat merupakan baseline normatif yang mengukur pergeseran
wilayah berdasarkan klasifikasi baku UNDP: Sangat Tinggi (≥ 80), Tinggi (70–79,99),
Sedang (55–69,99), dan Rendah (< 55). Matriks ini tidak bergantung pada metode klaster
apapun, melainkan langsung bersumber dari nilai IPM tahunan, sehingga berfungsi sebagai
Universitas Indonesia

128
acuan standar internasional. Karena nilai IPM merupakan besaran skalar yang sama untuk
kedua skenario, matriks IPM untuk skenario 1 dan skenario 2 adalah ekuivalen; perbedaan
antara kedua skenario hanya terletak pada format penyimpanan berkas.
Skenario 1 (Periode 2023 – 2024)
Tabel 4.46: Matriks Transisi Kategori IPM (UNDP) Tahun 2023 ke 2024
2023
2024 Sangat Tinggi Tinggi Sedang Rendah
Sangat Tinggi 9 0 0 0
Tinggi 2 19 0 0
Sedang 0 1 10 0
Rendah 0 0 0 0
Tabel 4.47: Nama Daerah Matriks Transisi Kategori IPM (UNDP) 2023 → 2024
Kategori Daftar Wilayah
Sangat Tinggi (2023)
→ Sangat Tinggi
(2024)
Kota Jakarta Selatan, Kota Jakarta Timur, Kota Jakarta Pusat,
Kota Jakarta Barat, Kota Jakarta Utara, Kota Bandung, Kota
Bekasi, Kota Depok, Kota Tangerang Selatan
Tinggi (2023) →
Sangat Tinggi (2024)
Kota Cimahi, Kota Tangerang
Tinggi (2023) →
Tinggi (2024)
Kep. Seribu, Bogor, Bandung, Ciamis, Kuningan, Cirebon,
Sumedang, Subang, Purwakarta, Karawang, Bekasi, Kota
Bogor, Kota Sukabumi, Kota Cirebon, Kota Tasikmalaya,
Kota Banjar, Tangerang, Kota Cilegon, Kota Serang
Sedang (2023) →
Tinggi (2024)
Bandung Barat
Sedang (2023) →
Sedang (2024)
Sukabumi, Cianjur, Garut, Tasikmalaya, Majalengka, Indramayu, Pangandaran, Pandeglang, Lebak, Serang
Skenario 2 (Periode 2023 – 2024)
Tabel 4.48: Matriks Transisi Kategori IPM (UNDP) Tahun 2023 ke 2024
2023
2024 Sangat Tinggi Tinggi Sedang Rendah
Sangat Tinggi 9 0 0 0
Tinggi 2 19 0 0
Sedang 0 1 10 0
Rendah 0 0 0 0
Universitas Indonesia

129
Tabel 4.49: Nama Daerah Matriks Transisi Kategori IPM (UNDP) 2023 → 2024
Kategori Daftar Wilayah
Sangat Tinggi (2023)
→ Sangat Tinggi
(2024)
Kota Jakarta Selatan, Kota Jakarta Timur, Kota Jakarta Pusat,
Kota Jakarta Barat, Kota Jakarta Utara, Kota Bandung, Kota
Bekasi, Kota Depok, Kota Tangerang Selatan
Tinggi (2023) →
Sangat Tinggi (2024)
Kota Cimahi, Kota Tangerang
Tinggi (2023) →
Tinggi (2024)
Kep. Seribu, Bogor, Bandung, Ciamis, Kuningan, Cirebon,
Sumedang, Subang, Purwakarta, Karawang, Bekasi, Kota
Bogor, Kota Sukabumi, Kota Cirebon, Kota Tasikmalaya,
Kota Banjar, Tangerang, Kota Cilegon, Kota Serang
Sedang (2023) →
Tinggi (2024)
Bandung Barat
Sedang (2023) →
Sedang (2024)
Sukabumi, Cianjur, Garut, Tasikmalaya, Majalengka, Indramayu, Pangandaran, Pandeglang, Lebak, Serang
Gambar 4.47 memperlihatkan perbandingan distribusi geografis kategori IPM UNDP
antara tahun 2023 dan 2024 untuk kedua skenario.
Skenario 1 Skenario 2
Gambar 4.47: Perbandingan Peta Kategori IPM UNDP: Tahun 2023 (kiri) vs 2024 (kanan)
Dominasi elemen diagonal pada kedua matriks mengonfirmasi tingginya stabilitas
kategori IPM, yaitu wilayah-wilayah cenderung bertahan pada kategori yang sama dari
satu tahun ke tahun berikutnya. pergeseran ke arah kategori yang lebih tinggi (elemen
sub-diagonal) mendominasi dibandingkan pergeseran ke arah kategori yang lebih rendah
(elemen super-diagonal), yang mengindikasikan tren pembangunan manusia yang secara
keseluruhan bersifat progresif selama 2020–2024.
4.8.4 Perbandingan 5 — Validasi Eksternal Klaster Biplot terhadap Kategori IPM
Perbandingan kelima melaksanakan validasi eksternal (external validation) terhadap
klaster Biplot PCA dengan membandingkannya terhadap klasifikasi normatif IPM UNDP
pada tahun yang sama. Tingginya kesesuaian antara klaster data-driven dengan kategori
institusional merupakan indikator validitas konstruk yang kuat. Disajikan tabulasi silang
untuk tahun 2024.
Universitas Indonesia

130
Skenario 1 (Tahun 2024)
Tabel 4.50: Perbandingan Klaster Biplot vs Kategori IPM (UNDP) (2024)
Klaster K-Means Sangat Tinggi Tinggi Sedang Rendah
Cluster 1 (Rendah) 0 11 10 0
Cluster 2 (Sedang) 1 9 0 0
Cluster 3 (Tinggi) 10 0 0 0
Tabel 4.51: Nama Daerah Perbandingan Klaster Biplot vs Kategori IPM (UNDP) (2024)
Kategori Daftar Wilayah
Cl. 1 (Rendah) &
Tinggi
Kep. Seribu, Ciamis, Cirebon, Sumedang, Subang, Bandung
Barat, Kota Tasikmalaya, Kota Banjar, Tangerang, Kota
Cilegon, Kota Serang
Cl. 1 (Rendah) &
Sedang
Sukabumi, Cianjur, Garut, Tasikmalaya, Majalengka, Indramayu, Pangandaran, Pandeglang, Lebak, Serang
Cl. 2 (Sedang) &
Sangat Tinggi
Kota Cimahi
Cl. 2 (Sedang) &
Tinggi
Bogor, Bandung, Kuningan, Purwakarta, Karawang, Bekasi,
Kota Bogor, Kota Sukabumi, Kota Cirebon
Cl. 3 (Tinggi) & Sangat
Tinggi
Kota Jakarta Selatan, Kota Jakarta Timur, Kota Jakarta Pusat,
Kota Jakarta Barat, Kota Jakarta Utara, Kota Bandung, Kota
Bekasi, Kota Depok, Kota Tangerang, Kota Tangerang Selatan
Skenario 2 (Tahun 2024)
Tabel 4.52: Perbandingan Klaster Biplot vs Kategori IPM (UNDP) (2024)
Klaster K-Means Sangat Tinggi Tinggi Sedang Rendah
Cluster 1 (Rendah) 0 1 6 0
Cluster 2 (Sedang) 0 16 4 0
Cluster 3 (Tinggi) 11 3 0 0
Universitas Indonesia

131
Tabel 4.53: Nama Daerah Perbandingan Klaster Biplot vs Kategori IPM (UNDP) (2024)
Kategori Daftar Wilayah
Cl. 1 (Rendah) &
Tinggi
Ciamis
Cl. 1 (Rendah) &
Sedang
Cianjur, Tasikmalaya, Majalengka, Pangandaran, Pandeglang, Lebak
Cl. 2 (Sedang) &
Tinggi
Kep. Seribu, Bogor, Bandung, Kuningan, Cirebon,
Sumedang, Subang, Purwakarta, Karawang, Bandung Barat,
Kota Cirebon, Kota Tasikmalaya, Kota Banjar, Tangerang,
Kota Cilegon, Kota Serang
Cl. 2 (Sedang) &
Sedang
Sukabumi, Garut, Indramayu, Serang
Cl. 3 (Tinggi) & Sangat
Tinggi
Kota Jakarta Selatan, Kota Jakarta Timur, Kota Jakarta Pusat,
Kota Jakarta Barat, Kota Jakarta Utara, Kota Bandung, Kota
Bekasi, Kota Depok, Kota Cimahi, Kota Tangerang, Kota
Tangerang Selatan
Cl. 3 (Tinggi) & Tinggi Bekasi, Kota Bogor, Kota Sukabumi
Gambar 4.48 memperlihatkan perbandingan distribusi geografis klaster Biplot dan
kategori IPM UNDP tahun 2024 untuk kedua skenario.
Skenario 1 Skenario 2
Gambar 4.48: Perbandingan Peta Klaster Biplot (kiri) vs Kategori IPM UNDP (kanan)
Tahun 2024
Korespondensi antara klaster Biplot dengan kategori IPM UNDP merupakan bukti
validitas konstruk yang mengonfirmasi bahwa pemisahan klaster oleh PCA mencerminkan
perbedaan pembangunan yang bermakna, bukan sekadar konstruk semu dari proses reduksi dimensi. Klaster 3 Biplot menunjukkan kesesuaian sempurna dengan IPM “Sangat
Tinggi”: seluruh 10 anggotanya (Kota Jakarta Selatan, Timur, Pusat, Barat, Utara, Kota
Bandung, Kota Bekasi, Kota Depok, Kota Tangerang, dan Kota Tangerang Selatan) berkategori IPM ≥ 80, tanpa pengecualian. Klaster 1 Biplot sepenuhnya bebas dari kategori
“Sangat Tinggi”: 10 dari 21 anggotanya berkategori “Sedang” (Sukabumi, Cianjur, Garut,
Tasikmalaya, Majalengka, Indramayu, Pangandaran, Pandeglang, Lebak, Serang) dan
11 lainnya berkategori “Tinggi”. Korespondesi yang kuat di kedua kutub ini memvaliUniversitas Indonesia

132
dasi secara empiris bahwa pemisahan klaster Biplot mencerminkan perbedaan kapasitas
pembangunan manusia yang nyata (Lihat Tabel 4.51).
4.8.5 Perbandingan 6 — Validasi Eksternal Klaster STATIS Tahunan terhadap
Kategori IPM
Perbandingan keenam melaksanakan validasi eksternal yang ekuivalen untuk klaster
STATIS tahunan (Fk), sehingga memungkinkan perbandingan langsung antara kualitas
validasi eksternal metode Biplot (Subbab 4.8.4) dengan metode STATIS.
Skenario 1 (Tahun 2024)
Tabel 4.54: Perbandingan Klaster STATIS Tahunan vs Kategori IPM (UNDP) (2024)
Klaster STATIS Sangat Tinggi Tinggi Sedang Rendah
Cluster 1 (Rendah) 0 11 10 0
Cluster 2 (Sedang) 1 9 0 0
Cluster 3 (Tinggi) 10 0 0 0
Tabel 4.55: Nama Daerah Perbandingan Klaster STATIS Tahunan vs IPM (UNDP) (2024)
Kategori Daftar Wilayah
Cl. 1 (Rendah) &
Tinggi
Kep. Seribu, Ciamis, Cirebon, Sumedang, Subang, Bandung
Barat, Kota Tasikmalaya, Kota Banjar, Tangerang, Kota
Cilegon, Kota Serang
Cl. 1 (Rendah) &
Sedang
Sukabumi, Cianjur, Garut, Tasikmalaya, Majalengka, Indramayu, Pangandaran, Pandeglang, Lebak, Serang
Cl. 2 (Sedang) &
Sangat Tinggi
Kota Cimahi
Cl. 2 (Sedang) &
Tinggi
Bogor, Bandung, Kuningan, Purwakarta, Karawang, Bekasi,
Kota Bogor, Kota Sukabumi, Kota Cirebon
Cl. 3 (Tinggi) & Sangat
Tinggi
Kota Jakarta Selatan, Kota Jakarta Timur, Kota Jakarta Pusat,
Kota Jakarta Barat, Kota Jakarta Utara, Kota Bandung, Kota
Bekasi, Kota Depok, Kota Tangerang, Kota Tangerang Selatan
Universitas Indonesia

133
Skenario 2 (Tahun 2024)
Tabel 4.56: Perbandingan Klaster STATIS Tahunan vs Kategori IPM (UNDP) (2024)
Klaster STATIS Sangat Tinggi Tinggi Sedang Rendah
Cluster 1 (Rendah) 0 1 6 0
Cluster 2 (Sedang) 0 16 4 0
Cluster 3 (Tinggi) 11 3 0 0
Tabel 4.57: Nama Daerah Perbandingan Klaster STATIS Tahunan vs IPM (UNDP) (2024)
Kategori Daftar Wilayah
Cl. 1 (Rendah) &
Tinggi
Ciamis
Cl. 1 (Rendah) &
Sedang
Cianjur, Tasikmalaya, Majalengka, Pangandaran, Pandeglang, Lebak
Cl. 2 (Sedang) &
Tinggi
Kep. Seribu, Bogor, Bandung, Kuningan, Cirebon,
Sumedang, Subang, Purwakarta, Karawang, Bandung Barat,
Kota Cirebon, Kota Tasikmalaya, Kota Banjar, Tangerang,
Kota Cilegon, Kota Serang
Cl. 2 (Sedang) &
Sedang
Sukabumi, Garut, Indramayu, Serang
Cl. 3 (Tinggi) & Sangat
Tinggi
Kota Jakarta Selatan, Kota Jakarta Timur, Kota Jakarta Pusat,
Kota Jakarta Barat, Kota Jakarta Utara, Kota Bandung, Kota
Bekasi, Kota Depok, Kota Cimahi, Kota Tangerang, Kota
Tangerang Selatan
Cl. 3 (Tinggi) & Tinggi Bekasi, Kota Bogor, Kota Sukabumi
Gambar 4.49 memperlihatkan perbandingan distribusi geografis klaster STATIS tahunan dan kategori IPM UNDP tahun 2024 untuk kedua skenario.
Skenario 1 Skenario 2
Gambar 4.49: Perbandingan Peta Klaster STATIS Tahunan (kiri) vs Kategori IPM UNDP
(kanan) Tahun 2024
Perbandingan antara hasil pada Subbab 4.8.4 dan 4.8.5 mengindikasikan bahwa klaster
STATIS, yang mengintegrasikan bobot temporal αk, cenderung menghasilkan kongruensi
Universitas Indonesia

134
yang sedikit lebih tinggi terhadap kategori IPM dibandingkan klaster Biplot PCA. Hal
ini dapat diinterpretasikan sebagai konsekuensi dari pembobotan STATIS yang meredam
anomali tahun tertentu (misalnya gangguan pandemi 2020–2021), sehingga klaster yang
terbentuk lebih merepresentasikan kondisi struktural jangka menengah yang lebih berkorespondensi dengan kategorisasi IPM.
4.8.6 Perbandingan 7 — Konsistensi Klaster Kompromi terhadap Klaster Biplot
per Tahun
Perbandingan ketujuh mengevaluasi seberapa konsisten klaster berbasis trajektori
keseluruhan (Kompromi STATIS, ktraj) dengan klaster Biplot PCA tahunan (kbiplot,t).
Tingginya kongruensi mengonfirmasi bahwa klaster kompromi merupakan representasi
yang stabil dari keanggotaan klaster sepanjang periode pengamatan, bukan hanya kondisi
satu tahun tertentu.
Skenario 1 (Tahun 2024)
Tabel 4.58: Matriks Perbandingan Klaster Biplot vs Kompromi STATIS (2024)
Biplot (Unweighted) Kompromi 1 (Rendah) Kompromi 2 (Sedang) Kompromi 3 (Tinggi)
Biplot 1 (Rendah) 3 3 0
Biplot 2 (Sedang) 9 12 0
Biplot 3 (Tinggi) 0 1 13
Tabel 4.59: Nama Daerah Perbandingan Klaster Biplot vs Kompromi STATIS (2024)
Kategori Daftar Wilayah
Biplot 1 (Rendah) &
Kompromi 1 (Rendah)
Pandeglang, Lebak, Serang
Biplot 1 (Rendah) &
Kompromi 2 (Sedang)
Tangerang, Kota Cilegon, Kota Serang
Biplot 2 (Sedang) &
Kompromi 1 (Rendah)
Kep. Seribu, Cianjur, Garut, Tasikmalaya, Ciamis, Majalengka, Indramayu, Pangandaran, Kota Banjar
Biplot 2 (Sedang) &
Kompromi 2 (Sedang)
Bogor, Sukabumi, Bandung, Kuningan, Cirebon, Sumedang,
Subang, Purwakarta, Karawang, Bandung Barat, Kota Cirebon, Kota Tasikmalaya
Biplot 3 (Tinggi) &
Kompromi 2 (Sedang)
Kota Sukabumi
Biplot 3 (Tinggi) &
Kompromi 3 (Tinggi)
Kota Jakarta Selatan, Kota Jakarta Timur, Kota Jakarta Pusat,
Kota Jakarta Barat, Kota Jakarta Utara, Bekasi, Kota Bogor,
Kota Bandung, Kota Bekasi, Kota Depok, Kota Cimahi, Kota
Tangerang, Kota Tangerang Selatan
Universitas Indonesia

135
Skenario 2 (Tahun 2024)
Tabel 4.60: Matriks Perbandingan Klaster Biplot vs Kompromi STATIS (2024)
Biplot (Unweighted) Kompromi 1 (Rendah) Kompromi 2 (Sedang) Kompromi 3 (Tinggi)
Biplot 1 (Rendah) 3 3 0
Biplot 2 (Sedang) 9 12 0
Biplot 3 (Tinggi) 0 1 13
Tabel 4.61: Nama Daerah Perbandingan Klaster Biplot vs Kompromi STATIS (2024)
Kategori Daftar Wilayah
Biplot 1 (Rendah) &
Kompromi 1 (Rendah)
Pandeglang, Lebak, Serang
Biplot 1 (Rendah) &
Kompromi 2 (Sedang)
Tangerang, Kota Cilegon, Kota Serang
Biplot 2 (Sedang) &
Kompromi 1 (Rendah)
Kep. Seribu, Cianjur, Garut, Tasikmalaya, Ciamis, Majalengka, Indramayu, Pangandaran, Kota Banjar
Biplot 2 (Sedang) &
Kompromi 2 (Sedang)
Bogor, Sukabumi, Bandung, Kuningan, Cirebon, Sumedang,
Subang, Purwakarta, Karawang, Bandung Barat, Kota Cirebon, Kota Tasikmalaya
Biplot 3 (Tinggi) &
Kompromi 2 (Sedang)
Kota Sukabumi
Biplot 3 (Tinggi) &
Kompromi 3 (Tinggi)
Kota Jakarta Selatan, Kota Jakarta Timur, Kota Jakarta Pusat,
Kota Jakarta Barat, Kota Jakarta Utara, Bekasi, Kota Bogor,
Kota Bandung, Kota Bekasi, Kota Depok, Kota Cimahi, Kota
Tangerang, Kota Tangerang Selatan
Gambar 4.50 memperlihatkan perbandingan distribusi geografis klaster kompromi dan
klaster Biplot tahun 2024 untuk kedua skenario.
Skenario 1 Skenario 2
Gambar 4.50: Perbandingan Peta Klaster Kompromi (kiri) vs Klaster Biplot Tahun 2024
(kanan)
Universitas Indonesia

136
4.8.7 Perbandingan 8 — Konsistensi Klaster Kompromi terhadap Klaster STATIS
Tahunan
Perbandingan kedelapan merupakan analisis konsistensi yang paling fundamental
dalam kerangka STATIS, yaitu evaluasi kesesuaian antara klaster kompromi (yang dibangun
dari seluruh periode) dengan klaster STATIS tahunan. Nilai kongruensi yang tinggi
merupakan bukti empiris bahwa struktur kompromi Wc =
P
k αkWk berhasil merangkum
informasi temporal secara representatif.
Skenario 1 (Tahun 2024)
Tabel 4.62: Perbandingan Klaster Kompromi STATIS vs Klaster STATIS Tahunan (2024)
Kompromi STATIS Statis 1 (Rendah) (2024) Statis 2 (Sedang) (2024) Statis 3 (Tinggi) (2024)
Kompromi 1 (Rendah) 12 0 0
Kompromi 2 (Sedang) 9 7 0
Kompromi 3 (Tinggi) 0 3 10
Tabel 4.63: Nama Daerah Perbandingan Klaster Kompromi STATIS vs STATIS Tahunan
(2024)
Kategori Daftar Wilayah
Kompromi 1 (Rendah)
& Statis 1 (Rendah)
Kep. Seribu, Cianjur, Garut, Tasikmalaya, Ciamis, Majalengka, Indramayu, Pangandaran, Kota Banjar, Pandeglang,
Lebak, Serang
Kompromi 2 (Sedang)
& Statis 1 (Rendah)
Sukabumi, Cirebon, Sumedang, Subang, Bandung Barat,
Kota Tasikmalaya, Tangerang, Kota Cilegon, Kota Serang
Kompromi 2 (Sedang)
& Statis 2 (Sedang)
Bogor, Bandung, Kuningan, Purwakarta, Karawang, Kota
Sukabumi, Kota Cirebon
Kompromi 3 (Tinggi) &
Statis 2 (Sedang)
Bekasi, Kota Bogor, Kota Cimahi
Kompromi 3 (Tinggi) &
Statis 3 (Tinggi)
Kota Jakarta Selatan, Kota Jakarta Timur, Kota Jakarta Pusat,
Kota Jakarta Barat, Kota Jakarta Utara, Kota Bandung, Kota
Bekasi, Kota Depok, Kota Tangerang, Kota Tangerang Selatan
Skenario 2 (Tahun 2024)
Tabel 4.64: Perbandingan Klaster Kompromi STATIS vs Klaster STATIS Tahunan (2024)
Kompromi STATIS Statis 1 (Rendah) (2024) Statis 2 (Sedang) (2024) Statis 3 (Tinggi) (2024)
Kompromi 1 (Rendah) 7 5 0
Kompromi 2 (Sedang) 0 15 1
Kompromi 3 (Tinggi) 0 0 13
Universitas Indonesia

137
Tabel 4.65: Nama Daerah Perbandingan Klaster Kompromi STATIS vs STATIS Tahunan
(2024)
Kategori Daftar Wilayah
Kompromi 1 (Rendah)
& Statis 1 (Rendah)
Cianjur, Tasikmalaya, Ciamis, Majalengka, Pangandaran,
Pandeglang, Lebak
Kompromi 1 (Rendah)
& Statis 2 (Sedang)
Kep. Seribu, Garut, Indramayu, Kota Banjar, Serang
Kompromi 2 (Sedang)
& Statis 2 (Sedang)
Bogor, Sukabumi, Bandung, Kuningan, Cirebon, Sumedang,
Subang, Purwakarta, Karawang, Bandung Barat, Kota Cirebon, Kota Tasikmalaya, Tangerang, Kota Cilegon, Kota
Serang
Kompromi 2 (Sedang)
& Statis 3 (Tinggi)
Kota Sukabumi
Kompromi 3 (Tinggi) &
Statis 3 (Tinggi)
Kota Jakarta Selatan, Kota Jakarta Timur, Kota Jakarta Pusat,
Kota Jakarta Barat, Kota Jakarta Utara, Bekasi, Kota Bogor,
Kota Bandung, Kota Bekasi, Kota Depok, Kota Cimahi, Kota
Tangerang, Kota Tangerang Selatan
Gambar 4.51 memperlihatkan perbandingan distribusi geografis klaster kompromi dan
klaster STATIS tahun 2024 untuk kedua skenario.
Skenario 1 Skenario 2
Gambar 4.51: Perbandingan Peta Klaster Kompromi (kiri) vs Klaster STATIS Tahun 2024
(kanan)
Dominasi elemen diagonal memvalidasi bahwa klaster kompromi STATIS bukanlah
sekadar rata-rata aritmetika temporal, melainkan representasi yang secara struktural konsisten terhadap kondisi tahunan. Ketidakkonsistenan minor pada elemen di luar diagonal
utama umumnya berasal dari wilayah-wilayah yang mengalami perubahan status klaster
pada tahun tertentu, khususnya pada 2020–2021 akibat guncangan pandemi yang sempat
menggeser posisi beberapa wilayah perbatasan klaster.
Universitas Indonesia

138
4.8.8 Perbandingan 9 — Validasi Eksternal Final Klaster Kompromi terhadap
Kategori IPM
Perbandingan kesembilan merupakan validasi eksternal final yang membandingkan
klaster kompromi STATIS, sebagai representasi kondisi pembangunan manusia selama
2020–2024 secara keseluruhan, dengan kategori IPM UNDP per tahun. Korespondensi
yang tinggi antara klaster longitudinal data-driven dengan kategori institusional normatif
merupakan validasi eksternal paling kuat atas relevansi substantif metode STATIS dalam
konteks kebijakan pembangunan manusia.
Skenario 1 (Tahun 2024)
Tabel 4.66: Perbandingan Klaster Kompromi STATIS vs Kategori IPM (UNDP) (2024)
Kompromi STATIS Sangat Tinggi Tinggi Sedang Rendah
Kompromi 1 (Rendah) 0 3 9 0
Kompromi 2 (Sedang) 0 15 1 0
Kompromi 3 (Tinggi) 11 2 0 0
Tabel 4.67: Nama Daerah Perbandingan Klaster Kompromi STATIS vs IPM (UNDP)
(2024)
Kategori Daftar Wilayah
Kompromi 1 (Rendah)
& Tinggi
Kep. Seribu, Ciamis, Kota Banjar
Kompromi 1 (Rendah)
& Sedang
Cianjur, Garut, Tasikmalaya, Majalengka, Indramayu, Pangandaran, Pandeglang, Lebak, Serang
Kompromi 2 (Sedang)
& Tinggi
Bogor, Bandung, Kuningan, Cirebon, Sumedang, Subang,
Purwakarta, Karawang, Bandung Barat, Kota Sukabumi,
Kota Cirebon, Kota Tasikmalaya, Tangerang, Kota Cilegon,
Kota Serang
Kompromi 2 (Sedang)
& Sedang
Sukabumi
Kompromi 3 (Tinggi) &
Sangat Tinggi
Kota Jakarta Selatan, Kota Jakarta Timur, Kota Jakarta Pusat,
Kota Jakarta Barat, Kota Jakarta Utara, Kota Bandung, Kota
Bekasi, Kota Depok, Kota Cimahi, Kota Tangerang, Kota
Tangerang Selatan
Kompromi 3 (Tinggi) &
Tinggi
Bekasi, Kota Bogor
Universitas Indonesia

139
Skenario 2 (Tahun 2024)
Tabel 4.68: Perbandingan Klaster Kompromi STATIS vs Kategori IPM (UNDP) (2024)
Kompromi STATIS Sangat Tinggi Tinggi Sedang Rendah
Kompromi 1 (Rendah) 0 3 9 0
Kompromi 2 (Sedang) 0 15 1 0
Kompromi 3 (Tinggi) 11 2 0 0
Tabel 4.69: Nama Daerah Perbandingan Klaster Kompromi STATIS vs IPM (UNDP)
(2024)
Kategori Daftar Wilayah
Kompromi 1 (Rendah)
& Tinggi
Kep. Seribu, Ciamis, Kota Banjar
Kompromi 1 (Rendah)
& Sedang
Cianjur, Garut, Tasikmalaya, Majalengka, Indramayu, Pangandaran, Pandeglang, Lebak, Serang
Kompromi 2 (Sedang)
& Tinggi
Bogor, Bandung, Kuningan, Cirebon, Sumedang, Subang,
Purwakarta, Karawang, Bandung Barat, Kota Sukabumi,
Kota Cirebon, Kota Tasikmalaya, Tangerang, Kota Cilegon,
Kota Serang
Kompromi 2 (Sedang)
& Sedang
Sukabumi
Kompromi 3 (Tinggi) &
Sangat Tinggi
Kota Jakarta Selatan, Kota Jakarta Timur, Kota Jakarta Pusat,
Kota Jakarta Barat, Kota Jakarta Utara, Kota Bandung, Kota
Bekasi, Kota Depok, Kota Cimahi, Kota Tangerang, Kota
Tangerang Selatan
Kompromi 3 (Tinggi) &
Tinggi
Bekasi, Kota Bogor
Gambar 4.52 memperlihatkan perbandingan distribusi geografis klaster kompromi dan
kategori IPM UNDP tahun 2024 untuk kedua skenario.
Skenario 1 Skenario 2
Gambar 4.52: Perbandingan Peta Klaster Kompromi (kiri) vs Kategori IPM UNDP Tahun
2024 (kanan)
Berdasarkan Tabel 4.67, pola kesesuaian antar-kelompok adalah sebagai berikut:
Universitas Indonesia

140
• Klaster 3 (Maju) dan Kategori IPM Atas: Sebelas dari 13 anggota Klaster 3
kompromi berkategori IPM “Sangat Tinggi” (≥ 80): kelima kota DKI Jakarta,
Kota Bandung, Kota Bekasi, Kota Depok, Kota Cimahi, Kota Tangerang, dan Kota
Tangerang Selatan. Dua anggota lainnya yakni Kabupaten Bekasi dan Kota Bogor
yang berkategori “Tinggi”, ini mencerminkan wilayah urban-pinggiran dengan
capaian tinggi namun belum mencapai ambang batas ≥ 80. Presisi Klaster 3
kompromi terhadap IPM “Sangat Tinggi” sebesar 84,6% (11 dari 13) merupakan
bukti validitas konstruk yang sangat kuat bahwa pemisahan dimensi kompromi (Fc)
mencerminkan hierarki pembangunan yang konsisten dengan standar internasional.
• Klaster 2 (Berkembang) yang Homogen: Klaster 2 kompromi justru bersifat cukup
homogen dari sisi IPM: 15 dari 16 anggotanya berkategori IPM “Tinggi” (70–79,99),
mewakili hampir seluruh koridor industri dan wilayah penyangga. Hanya Kabupaten Sukabumi yang berkategori “Sedang”, namun profil multivariat gabungannya
tetap menempatkannya lebih dekat ke kelompok berkembang daripada kelompok
tertinggal. Homogenitas ini menegaskan bahwa Klaster 2 secara substantif merepresentasikan wilayah-wilayah dalam koridor IPM “Tinggi” yang secara struktural
belum mencapai capaian wilayah metropolitan.
• Klaster 1 (Tertinggal) dan Kategori IPM Bawah: Klaster 1 kompromi didominasi
oleh 9 wilayah berkategori IPM “Sedang” (55–69,99): Cianjur, Garut, Tasikmalaya,
Majalengka, Indramayu, Pangandaran, Pandeglang, Lebak, dan Serang, seluruhnya merupakan kabupaten pedalaman atau rural di Banten Selatan dan Jawa Barat
Selatan. Tiga anggota lainnya, yaitu Kepulauan Seribu, Ciamis, dan Kota Banjar memiliki kategori “Tinggi” namun tetap terklasifikasikan ke Klaster 1 karena
profil multivariat gabungan 11 indikator yang menempatkan mereka lebih dekat
ke kelompok tertinggal. Kesesuaian ini menunjukkan bahwa posisi klaster dalam
ruang kompromi merupakan representasi yang valid secara substantif dari kondisi
pembangunan struktural jangka menengah.
Korespondensi tinggi antara klaster kompromi data-driven dengan kategori normatif
UNDP merupakan validasi eksternal yang memperkuat bahwa metode STATIS mampu
mengekstrak struktur laten pembangunan manusia yang bermakna secara kebijakan dari
data longitudinal multivariat.
Universitas Indonesia

141
4.9 Uji Autokorelasi Spasial (Global Moran’s I)
Sebagai tahap lanjutan, dilakukan uji autokorelasi spasial global menggunakan statistik Moran’s I untuk mendeteksi apakah terdapat pola pengelompokan geografis yang
signifikan pada skor kompromi. Penghitungan mengacu pada Persamaan 2.73 di Bab 2,
dengan prosedur inferensi permutasi yang diuraikan pada Bab 3.
Skenario 1: ACT Klasik
Skenario 2: Dual STATIS
Hasil penghitungan menunjukkan:
• Skenario 1 (ACT Klasik): Moran’s I = 0, 6903, p-value = 0, 3151, Tidak Signifikan (Pola Spasial Acak)
• Skenario 2 (Dual STATIS): Moran’s I = 0, 6930, p-value = 0, 3133, Tidak Signifikan (Pola Spasial Acak)
Kedua skenario menghasilkan p-value > 0,05, sehingga tidak terdapat autokorelasi
spasial yang signifikan pada skor kompromi PC1. Hasil ini konsisten dengan karakteristik
data yang digunakan: skor kompromi mencerminkan posisi relatif kabupaten/kota dalam
ruang multivariat pembangunan manusia, bukan besaran nilai geografis yang secara inheren
bersifat spasial. Ketiadaan autokorelasi spasial menunjukkan bahwa kinerja pembangunan
manusia di kawasan Banten–DKI Jakarta–Jawa Barat lebih ditentukan oleh kapabilitas
internal dan kebijakan sektoral masing-masing wilayah daripada oleh efek kedekatan
geografis (spillover spasial). Pola pengelompokan yang ditemukan (Klaster Tertinggal–
Berkembang–Maju) adalah kluster statistik (berbasis kemiripan profil indikator), bukan
kluster geografis yang tersebar membentuk hot spot atau cold spot dalam pengertian
autokorelasi spasial klasik (Anselin, 1988).
4.10 Analisis Hotspot (Local Moran’s I)
Apabila nilai Moran’s I positif dan signifikan (I > E(I), p < 0,05), hal ini
mengindikasikan terdapat pola pengelompokan spasial positif (positive spatial clustering): wilayah dengan skor kompromi tinggi cenderung bertetangga dengan wilayah
berskor tinggi, dan demikian sebaliknya. Namun, karena p-value > 0,05 pada kedua
skenario, tidak terdapat bukti statistik yang cukup untuk mengklaim adanya hotspot atau
coldspot geografis yang bermakna pada data ini. Klaster yang ditemukan (Tertinggal,
Universitas Indonesia

142
Berkembang, Maju) adalah kluster statistik berbasis kesamaan profil multivariat—bukan
hasil autokorelasi geografis. Oleh karena itu, analisis Local Moran’s I untuk mengidentifikasi hotspot lokal tidak relevan secara inferensial, dan kebijakan intervensi sebaiknya
diarahkan berdasarkan profil kapabilitas internal wilayah, bukan semata-mata kedekatan
geografisnya.
4.11 Ringkasan Temuan dan Implikasi Kebijakan
Berdasarkan keseluruhan analisis yang telah dilakukan, berikut adalah rangkuman
temuan utama penelitian ini yang disusun untuk menjawab tujuan penelitian serta memberikan manfaat bagi akademik dan praktisi:
4.11.1 Pembahasan Berdasarkan Tujuan Penelitian
Penelitian ini telah berhasil mencapai lima tujuan utama yang dirumuskan pada Bab 1
sebagai berikut:
1. Struktur Kompromi (Tujuan 1): Melalui metode STATIS, telah berhasil dibentuk
sebuah struktur kompromi yang merepresentasikan pola hubungan multivariat sebelas indikator pembangunan. Dimensi 1 kompromi berhasil merangkum 60, 85%
variansi total, yang menunjukkan kontribusi dominan dari indikator ekonomi (PDRB)
dan kualitas manusia (IPM). Hal ini memberikan gambaran umum yang stabil mengenai hierarki pembangunan di wilayah Banten, DKI Jakarta, dan Jawa Barat.
2. Keabsahan Dinamika Temporal (Tujuan 2): Uji stabilitas melalui permutasi
menunjukkan hasil yang signifikan (p = 0, 0010), membuktikan bahwa meskipun
terdapat fluktuasi tahunan (terutama akibat pandemi), struktur hubungan antarindikator yang ditemukan benar-benar ada secara statistik dan tetap kokoh sebagai
landasan analisis spasio-temporal.
3. Analisis Lintasan/Trajektori (Tujuan 3): Visualisasi Dynamic Biplot secara efektif
menggambarkan lintasan pergeseran 41 kabupaten/kota selama 2020–2024 pada
Gambar 4.25. Tiga pola dominan teridentifikasi: (i) Pemulihan V-shape pada
wilayah Klaster 3 metropolitan seperti Kota Tangerang Selatan mencatat peningkatan
PC1 terbesar (+0,85 poin antara 2020 dan 2024), diikuti Kota Bekasi (+0,39)
dan Kota Jakarta Pusat (+0,26); (ii) Progresif bertahap pada beberapa wilayah
penyangga seperti Kota Sukabumi yang berhasil naik dari Klaster 2 ke Klaster 3
pada 2023; dan (iii) Stagnasi dan penurunan relatif pada wilayah Klaster 1 seperti
Universitas Indonesia

143
Kabupaten Pangandaran (PC1 turun −1,37 poin) dan Garut (−0,61), yang skor
PC1-nya justru semakin jauh dari skor PC1 wilayah-wilayah Klaster 3. Hal ini
menandakan kesenjangan pembangunan yang melebar, bukan menyempit, selama
lima tahun pengamatan.
4. Pengelompokan dan Migrasi Wilayah (Tujuan 4): Pengelompokan berbasis KMeans (k = 3) mengungkap tiga kelompok wilayah yang berbeda secara struktural
(Tertinggal, Berkembang, Maju). Matriks transisi menunjukkan: nol perpindahan
pada 2020–2021 (stabilitas sempurna di tahun pertama pandemi); pemulihan selektif
2022–2023 dengan 6 wilayah naik K1→K2 (antara lain Garut, Ciamis, Kota Cilegon)
dan 1 wilayah naik K2→K3 (Kota Sukabumi); serta restrukturisasi terbesar pada
2023–2024 dengan 14 wilayah turun K2→K1 (termasuk Tangerang, Kota Serang,
Cirebon, Bandung Barat) dan 4 wilayah turun K3→K2 (Bekasi, Kota Bogor, Kota
Sukabumi, Kota Cimahi). Sementara itu, 10 wilayah metropolitan inti (5 kota DKI
beserta Kota Bandung, Kota Bekasi, Kota Depok, Kota Tangerang, Kota Tangerang
Selatan) tetap bertahan di Klaster 3 sepanjang 2020–2024, mengonfirmasi ketahanan
struktural yang kokoh pada wilayah-wilayah tersebut.
5. Platform Eksplorasi Interaktif (Tujuan 5): Seluruh hasil penghitungan dan visualisasi ini telah diintegrasikan ke dalam naskah bab ini, yang merupakan representasi
statis dari aplikasi berbasis web (Dash/Plotly) yang dikembangkan. Dashboard
tersebut memungkinkan pengguna melakukan pemantauan indikator secara mandiri
dan intuitif.
4.11.2 Kontribusi terhadap Manfaat Penelitian
Sesuai dengan manfaat penelitian yang diharapkan:
1. Manfaat Akademis: Penelitian ini memperkaya literatur statistik multivariat di
Indonesia dengan menerapkan metode STATIS dan Dynamic Biplot yang masih
jarang digunakan pada data regional. Penggabungan analisis inter-struktur, struktur kompromi, dan analisis klaster memberikan metodologi yang lebih dinamis
dibandingkan analisis dua-arah konvensional.
2. Manfaat Praktis: Keberadaan visualisasi peta dan matriks transisi dalam aplikasi dashboard menyediakan alat bantu yang kuat bagi pemerintah daerah untuk
mengevaluasi efektivitas kebijakan pembangunan secara komprehensif, tidak hanya
melihat satu indikator secara terisolasi.
Universitas Indonesia

144
4.11.3 Implikasi Kebijakan
Temuan bahwa struktur kesenjangan pembangunan bersifat struktural dan terpola secara spasial menunjukkan bahwa kebijakan pembangunan yang bersifat “merata” (uniform)
tidak akan efektif. Diperlukan kebijakan yang targeted, yang memfokuskan intervensi pada
wilayah-wilayah Klaster 1 yang secara struktural tertinggal, terutama di bidang pendidikan
(peningkatan RLS), pengurangan kemiskinan, dan penciptaan lapangan kerja. Strategi
pembangunan berbasis regional spillover dari pusat pertumbuhan (DKI Jakarta) ke wilayah
penyangga perlu ditingkatkan untuk memperpendek kesenjangan spasio-temporal yang
teramati dalam lima tahun terakhir.
Universitas Indonesia

BAB 5
KESIMPULAN DAN SARAN
Bab ini menyimpulkan hasil temuan utama dari analisis STATIS dan Dynamic Biplot
mengenai dinamika pembangunan ekonomi regional di Provinsi Banten, DKI Jakarta,
dan Jawa Barat, serta memberikan saran rekomendasi kebijakan maupun pengembangan
penelitian selanjutnya.
5.1 Kesimpulan
Berdasarkan hasil penelusuran dan pembahasan yang telah dilakukan pada bab sebelumnya, dapat diperoleh beberapa kesimpulan analisis sebagai berikut:
1. Karakteristik Struktur Kompromi (Menjawab Tujuan 1): Melalui sintesis matriks kompromi STATIS (Skenario 1 Wc dan Skenario 2 Sc), diperoleh representasi
rata-rata spasio-temporal yang stabil selama periode 2020–2024. Dua dimensi utama
berhasil diidentifikasi: Dimensi 1 merepresentasikan “Taraf Kemajuan Pembangunan Manusia dan Ekonomi” (didominasi oleh RLS, Angka Harapan Hidup L/P, GK,
IPG, dan PDRB ADHK dengan loading > 0,70) serta menjelaskan bagian terbesar
variansi total (60,85%). Dimensi 2 merepresentasikan “Tekanan Pasar Tenaga Kerja”
(didominasi oleh Tingkat Pengangguran Terbuka dengan loading +0,86 dan Tingkat
Partisipasi Angkatan Kerja dengan loading −0,49). Vektor bobot interstruktur (αW
dan αS) terdistribusi secara seimbang di antara kelima tahun pengamatan, yang
membuktikan representativitas yang adil tanpa distorsi tahun ekstrem.
2. Signifikansi Stabilitas Temporal dan Vektor Struktur (Menjawab Tujuan 2):
Analisis interstruktur STATIS menghasilkan koefisien RV yang sangat tinggi (seluruhnya di atas 0, 9543), membuktikan bahwa hierarki kemiripan antar-wilayah bersifat persisten lintas waktu. Uji permutasi Monte Carlo (B = 1000) membuktikan
bahwa struktur kompromi tersebut signifikan secara statistik (p-value = 0, 0010 <
0,05) dan bukan merupakan hasil fluktuasi acak. Kompresi interstruktur ke dalam
vektor struktur β˜ = λ1α1 secara longitudinal merekam dampak krisis pandemi
(2020–2021) dan pemulihan (2022–2024) yang nyata, sementara analisis transversal
membuktikan adanya heterogenitas respons temporal yang unik antar-provinsi.
3. Pola Lintasan Trajektori Spasio-Temporal (Menjawab Tujuan 3): Visualisasi
Dynamic Biplot berhasil menggambarkan lintasan (trajectory) pergerakan 41 kabupaten/kota secara dinamis dari masa krisis (2020) hingga masa pemulihan (2024).
145 Universitas Indonesia

146
Tiga pola pergerakan diidentifikasi: (1) pola pemulihan berbentuk huruf V (V-shape
rebound) pada kota-kota metropolitan berbasis jasa—Kota Tangerang Selatan mencatat peningkatan PC1 terbesar (+0,85 poin antara 2020 dan 2024), diikuti Kota
Bekasi (+0,39) dan seluruh kota DKI Jakarta; (2) pola progresif bertahap (monotone
improvement) pada wilayah penyangga seperti Kota Sukabumi yang berhasil naik
ke Klaster 3 pada 2023; dan (3) pola stagnasi dan penurunan relatif (inertia) pada
kabupaten pedalaman selatan—Kabupaten Pangandaran mengalami penurunan PC1
terbesar (−1,37 poin), diikuti Garut (−0,61) dan Cianjur (−0,51), menunjukkan
bahwa skor PC1 wilayah-wilayah ini justru semakin jauh dari skor PC1 wilayah
metropolitan Klaster 3—kesenjangan pembangunan melebar, bukan menyempit,
selama lima tahun pengamatan.
4. Pengelompokan Spasial dan Dinamika Transisi (Menjawab Tujuan 4): Jumlah klaster optimal k = 3 dipilih berdasarkan kombinasi kurva Elbow dan koefisien siluet rata-rata (SC = 0, 3883), yang mengonfirmasi pemisahan antar-klaster
yang memadai. Penerapan algoritma K-Means (k = 3) pada skor kompromi
berhasil mengklasifikasikan wilayah menjadi tiga kelompok: Klaster 1 (Tertinggal, 12 wilayah—Banten Selatan dan Jawa Barat pedalaman selatan), Klaster 2
(Berkembang, 16 wilayah—koridor industri dan penyangga), dan Klaster 3 (Maju,
13 wilayah—metropolitan). Matriks transisi temporal mengonfirmasi stabilitas sempurna pada 2020–2021 (nol perpindahan), diikuti pemulihan selektif 2022–2023
(6 wilayah naik K1→K2, termasuk Garut dan Kepulauan Seribu; 1 wilayah naik
K2→K3 yaitu Kota Sukabumi), dan restrukturisasi terbesar 2023–2024 (14 wilayah
turun K2→K1 dan 4 wilayah turun K3→K2). Sepuluh wilayah metropolitan inti
(5 kota DKI beserta Kota Bandung, Kota Bekasi, Kota Depok, Kota Tangerang,
Kota Tangerang Selatan) tetap bertahan di Klaster 3 sepanjang seluruh periode,
mengonfirmasi ketahanan struktural yang kokoh.
5. Pengembangan Dashboard Eksplorasi Interaktif (Menjawab Tujuan 5): Aplikasi berbasis web dikembangkan menggunakan framework Dash dan Plotly sebagai
instrumen eksplorasi visual interaktif yang andal. Fitur-fitur seperti pemetaan spasial
dinamis, visualisasi trajektori, seleksi tahun, dan rotasi sumbu baseline berhasil
memfasilitasi pemantauan indikator pembangunan secara interaktif oleh peneliti
dan pemangku kebijakan. Keandalan model terkonfirmasi melalui validasi eksternal
(tabulasi silang) yang menunjukkan kesesuaian yang sangat tinggi antara klaster
STATIS dengan klasifikasi normatif kategori IPM UNDP.
Universitas Indonesia

147
5.2 Saran
Berdasarkan kesimpulan di atas, diajukan beberapa saran sebagai berikut:
1. Implikasi Kebijakan Targeted Regional: Mengingat struktur ketimpangan yang
persisten dan terpola secara spasial, kebijakan pemerataan pembangunan tidak
cukup hanya mengandalkan pertumbuhan ekonomi nasional. Diperlukan intervensi
struktural yang targeted pada wilayah Klaster 1 (Tertinggal), khususnya Kabupaten
Lebak, Pandeglang, dan Garut, dengan fokus pada peningkatan kapabilitas dasar
(RLS, AHH) serta pengurangan kemiskinan melalui strategi regional spillover dari
pusat pertumbuhan DKI Jakarta.
2. Pengembangan Penelitian Lanjutan: Penelitian ini dibatasi pada indikator
makroekonomi tahunan dari 3 provinsi. Penelitian selanjutnya disarankan untuk
menggunakan data dengan resolusi waktu yang lebih tinggi (triwulanan) apabila
tersedia, memperluas cakupan wilayah ke seluruh Pulau Jawa atau Indonesia dengan mengintegrasikan metode Dual STATIS untuk analisis korelasi indikator, serta
menambahkan indikator sosial-infrastruktur guna memperkaya dimensi analisis.
3. Pengembangan Model Prediktif Spasio-Temporal (Forecasting): Karena penelitian ini terbatas pada analisis deskriptif-eksploratif data historis, penelitian selanjutnya sangat disarankan untuk mengintegrasikan model prediktif (peramalan) guna
memproyeksikan skor wilayah (Fk) dan loading indikator (Gk) ke tahun berikutnya
(k = K + 1). Hal ini dapat diwujudkan dengan menerapkan pemodelan deret waktu
(seperti ARIMA, Vector Autoregression (VAR), atau model Grey Systems GM(1,1)
yang sangat handal untuk deret data pendek K = 5) langsung pada barisan koordinat kompromi (fi,k dan gj,k) dan lintasan komponen vektor struktur (β˜
k). Dengan
memproyeksikan lintasan koordinat ini, pemangku kebijakan dapat mengantisipasi
pergeseran status kesejahteraan wilayah dan dinamika hubungan antar-indikator
secara prospektif.
4. Pengembangan Platform Dashboard Interaktif: Dashboard berbasis web
(Dash/Plotly) yang dikembangkan dalam penelitian ini dapat dikembangkan lebih
lanjut dengan menambahkan fitur pembaruan data secara otomatis dari API BPS
Open Data, sehingga pemantauan dinamika pembangunan regional dapat dilakukan
secara near real-time oleh pemangku kepentingan kebijakan.
Universitas Indonesia

DAFTAR PUSTAKA
Abdi, H., Williams, L. J., Valentin, D., & Bennani-Dosse, M. (2012). Statis and distatis:
optimum multitable principal component analysis and three way metric multidimensional
scaling. Wiley Interdisciplinary Reviews: Computational Statistics, 4(2), 124–167.
https://doi.org/10.1002/wics.198
Ahmad, A. (2024). Analisis regresi panel spasiotemporal menggunakan model spasial
durbin dinamis pada prevalensi kusta di pulau jawa. Institut Pertanian Bogor, Bogor.
Diakses: 20 November 2025. https://repository.ipb.ac.id/handle/123456789/154545
Anselin, L. (1988). Spatial Econometrics: Methods and Models. Kluwer Academic
Publishers. https://link.springer.com/book/10.1007/978-94-015-7799-1
Anton, H. & Rorres, C. (2013). Elementary Linear Algebra: Applications Version (11th
ed.). John Wiley & Sons. https://studylib.net/doc/27681515/linear-algebra-11th-edihoward-anton
Areia, A., Oliveira, M., & Mexia, J. T. (2008). Models for a series of studies based on
geometrical representation. Statistical Methodology, 5(3), 277–288. https://doi.org/10.
1016/j.stamet.2007.09.001
Badan Pusat Statistik (2024). Prevalensi ketidakcukupan konsumsi pangan (persen)
per kabupaten/kota (persen), 2024. Diakses dari situs resmi Badan Pusat Statistik. https://www.bps.go.id/id/statistics-table/2/MjI2OSMy/prevalensi-ketidakcukupankonsumsi-pangan--persen--per-kabupaten-kota--persen-.html
Badan Pusat Statistik (2025a). Indikator kesejahteraan rakyat 2025. Katalog 4102004, ISSN 0215-4641. https://www.bps.go.id/id/publication/2025/12/30/
bcc229ac31403d023f2d26a3/indikator-kesejahteraan-rakyat-2025.html
Badan Pusat Statistik (2025b). Tinjauan regional berdasarkan pdrb kabupaten/kota 2020–
2024. Diakses dari situs resmi Badan Pusat Statistik. https://www.bps.go.id/id/statisticstable/3/WkdVMWRYVnBkMnBvVEhKSVkyWXhNblZtTjJSbmR6MDkjMw==/
produk-domestik-regional-bruto-atas-dasar-harga-berlaku-menurut-provinsi--miliarrupiah---2022.html?year=2024
Badan Pusat Statistik Provinsi Banten (2020). Provinsi banten dalam angka
2020. https://banten.bps.go.id/id/publication/2020/04/27/41c9192acf8b24ecbaa2d8e5/
provinsi-banten-dalam-angka-2020.html
148 Universitas Indonesia

149
Badan Pusat Statistik Provinsi Banten (2021). Provinsi banten dalam angka
2021. https://banten.bps.go.id/id/publication/2021/02/26/344951df21aeedca2572f336/
provinsi-banten-dalam-angka-2021.html
Badan Pusat Statistik Provinsi Banten (2022). Provinsi banten dalam angka
2022. https://banten.bps.go.id/id/publication/2022/02/25/19658ae14140f17fc6ae9e3a/
provinsi-banten-dalam-angka-2022.html
Badan Pusat Statistik Provinsi Banten (2023). Provinsi banten dalam angka
2023. https://banten.bps.go.id/id/publication/2023/02/28/482ee839483674f34dd96faf/
provinsi-banten-dalam-angka-2023.html
Badan Pusat Statistik Provinsi Banten (2024). Provinsi banten dalam angka
2024. https://banten.bps.go.id/en/publication/2024/02/28/e9b15e76302d4876c37f60be/
provinsi-banten-dalam-angka-2024.html
Badan Pusat Statistik Provinsi DKI Jakarta (2020). Provinsi dki jakarta dalam angka
2020. https://jakarta.bps.go.id/id/publication/2020/04/27/20f5a58abcb80a0ad2a88725/
provinsi-dki-jakarta-dalam-angka-2020.html
Badan Pusat Statistik Provinsi DKI Jakarta (2021). Provinsi dki jakarta dalam angka
2021. https://jakarta.bps.go.id/id/publication/2021/02/26/bb7fa6dd5e90b534e3fa6984/
provinsi-dki-jakarta-dalam-angka-2021.html
Badan Pusat Statistik Provinsi DKI Jakarta (2022). Provinsi dki jakarta dalam angka
2022. https://jakarta.bps.go.id/id/publication/2022/02/25/5979600247867d861a1f334c/
provinsi-dki-jakarta-dalam-angka-2022.html
Badan Pusat Statistik Provinsi DKI Jakarta (2023). Provinsi dki jakarta dalam angka
2023. https://jakarta.bps.go.id/id/publication/2023/02/28/fd35fcb5d10a1e03f0d71348/
provinsi-dki-jakarta-dalam-angka-2023.html
Badan Pusat Statistik Provinsi DKI Jakarta (2024). Provinsi dki jakarta dalam angka
2024. https://jakarta.bps.go.id/id/publication/2024/02/28/baae7b80d16101c7bef30cc0/
provinsi-dki-jakarta-dalam-angka-2024.html
Badan Pusat Statistik Provinsi Jawa Barat (2020). Provinsi jawa barat dalam angka
2020. https://jabar.bps.go.id/id/publication/2020/04/27/cfab9a400cf304f800182a5f/
provinsi-jawa-barat-dalam-angka-2020.html
Universitas Indonesia

150
Badan Pusat Statistik Provinsi Jawa Barat (2021). Provinsi jawa barat dalam angka
2021. https://jabar.bps.go.id/id/publication/2021/02/26/4d3f7ec6c519dda0b9785d45/
provinsi-jawa-barat-dalam-angka-2021.html
Badan Pusat Statistik Provinsi Jawa Barat (2022). Provinsi jawa barat dalam angka
2022. https://jabar.bps.go.id/id/publication/2022/02/25/0d261f828b581d8082bbc6c1/
provinsi-jawa-barat-dalam-angka-2022.html
Badan Pusat Statistik Provinsi Jawa Barat (2023). Provinsi jawa barat dalam angka
2023. https://jabar.bps.go.id/id/publication/2023/02/28/57231a828abbfdd50a21fe31/
provinsi-jawa-barat-dalam-angka-2023.html
Badan Pusat Statistik Provinsi Jawa Barat (2024). Provinsi jawa barat dalam angka
2024. https://jabar.bps.go.id/id/publication/2024/02/28/35ffe2d35104b39feb577e8f/
provinsi-jawa-barat-dalam-angka-2024.html
Baltagi, B. (2021). Econometric Analysis of Panel Data. https://doi.org/10.1007/978-3-
030-53953-5
Bartlett, M. S. (1954). A note on the multiplying factors for various χ
2
approximations.
Journal of the Royal Statistical Society: Series B (Methodological), 16(2), 296–298.
https://doi.org/10.1111/j.2517-6161.1954.tb00174.x
Bholowalia, P. & Kumar, A. (2014). Ebk-means: A clustering technique based on elbow
method and k-means in wsn. International Journal of Computer Applications, 105(9).
https://doi.org/10.5120/18405-9674
Biegert, T., Ozcan, B., & Rossetti-Youlton, M. (2023). Household joblessness in u.s. ¨
metropolitan areas during the covid-19 pandemic: Polarization and the role of educational profiles. Socius: Sociological Research for a Dynamic World, 9, 1–18.
https://doi.org/10.1177/23780231231158087
Bulmer, M. G. (1979). Principles of Statistics (dover edition ed.). Dover Publications.
Klasifikasi standar skewness: |γ1| < 0.5 simetris; 0.5 ≤ |γ1| < 1 moderat; |γ1| ≥ 1
sangat menceng (hlm. 110-112).
Departemen Pertanian (Deptan) (2003). Analisis produktivitas dan stabilitas komoditas
strategis pangan propinsi di yogyakarta. Technical report, Departemen Pertanian RI.
Metodologi standar klasifikasi CV untuk stabilitas pertanian. https://jiip.polbangtanyoma.
ac.id/index.php/jiip/article/download/165/159/531
Universitas Indonesia

151
Dziuban, C. D. & Shirkey, E. C. (1974). When is a correlation matrix appropriate
for factor analysis? some decision rules. Psychological Bulletin, 81(6), 358–361.
https://doi.org/10.1037/h0036316
Egido, J. & Galindo, P. (2015). Dynamic biplot. evolution of the economic freedom in
the european union. British Journal of Applied Science & Technology, 11(3), 1–13.
https://doi.org/10.9734/BJAST/2015/20289
Escoufier, Y. (1973). Le traitement des variables vectorielles. Biometrics, 29(4), 751–760.
https://doi.org/10.2307/2529140
Gabriel, K. R. (1971). The biplot graphic display of matrices with application to principal
component analysis. Biometrika, 58(3), 453–467. https://doi.org/10.1093/biomet/58.3.
453
Hair, J. F., Black, W. C., Babin, B. J., & Anderson, R. E. (2019). Multivariate Data Analysis
(8 ed.). Cengage Learning. https://www.cengage.uk/c/multivariate-data-analysis-8ehair-black-babin-anderson/9781473756540/
Hotelling, H. (1933). Analysis of a complex of statistical variables into principal components. Journal of Educational Psychology, 24(6), 417–441. https://doi.org/10.1037/
h0071325
Johnson, R. A. & Wichern, D. W. (2007). Applied Multivariate Statistical Analysis (6
ed.). Pearson Prentice Hall. https://www.pearson.com/en-us/subject-catalog/p/appliedmultivariate-statistical-analysis/P200000003290
Kaiser, H. F. (1970). A second generation little jiffy. Psychometrika, 35(4), 401–415.
https://doi.org/10.1007/BF02291817
Kaiser, H. F. (1974). An index of factorial simplicity. Psychometrika, 39(1), 31–36.
https://doi.org/10.1007/BF02291575
Kazi-Aoual, F., Hitier, S., Sabatier, R., & Lebreton, J. D. (1995). Refined approximations to
permutation tests for multivariate inference. Computational Statistics & Data Analysis,
20(4), 385–394. https://doi.org/10.1016/0167-9473(94)00064-2
Kroonenberg, P. M. (2008). Applied Multiway Data Analysis. John Wiley & Sons.
https://doi.org/10.1002/9780470238004
Kuhn, H. W. (1955). The hungarian method for the assignment problem. Naval Research
Logistics Quarterly, 2(1–2), 83–97. https://doi.org/10.1002/nav.3800020109
Universitas Indonesia

152
Lavit, C., Escoufier, Y., Sabatier, R., & Traissac, P. (1994). The act (statis method).
Computational Statistics & Data Analysis, 18(1), 97–119. https://doi.org/10.1016/0167-
9473(94)90134-1
Legendre, P. & Legendre, L. (2012). Numerical ecology. Elsevier. https://www.
sciencedirect.com/book/9780444538680/numerical-ecology
MacQueen, J. (1967). Some methods for classification and analysis of multivariate
observations. Proceedings of the fifth Berkeley symposium on mathematical statistics and
probability, volume 1, 281–297. https://digicoll.lib.berkeley.edu/record/113015?v=pdf
Mexia, J. T. & Oliveira, M. (2007). Non-observable regressors. Biometrical Letters, 44(2),
97–103. https://bl.up.poznan.pl/full/BL%2044%202%202.pdf
Nuswantari, A. (2021). Pemodelan kemiskinan di pulau jawa tahun 2012 sampai 2018
dengan panel spasial dinamis. Institut Pertanian Bogor, Bogor. Diakses: 20 November
2025. https://repository.ipb.ac.id/handle/123456789/107909
Oliveira, M., Garc¸ao, E., Grilo, L. M., Alexandre, A., & Mexia, J. T. (2024). Statis ˜
methodology: An overview. AIP Conference Proceedings, volume 3094, 450006.
https://doi.org/10.1063/5.0212208
Pearson, K. (1901). Liii. on lines and planes of closest fit to systems of points in space.
The London, Edinburgh, and Dublin Philosophical Magazine and Journal of Science,
2(11), 559–572. https://doi.org/10.1080/14786440109462720
Phipson, B. & Smyth, G. K. (2010). Permutation p-values should never be zero: Calculating
exact p-values when permutations are randomly drawn. Statistical Applications in
Genetics and Molecular Biology, 9(1), Article 39. https://doi.org/10.48550/arXiv.1603.
05766
Plotly Technologies Inc. (2015). Collaborative data science. Plotly. https://plotly.com
Plotly Technologies Inc. (2024). Dash Python User Guide. Plotly Technologies Inc.
Dokumentasi resmi kerangka kerja Dash. https://dash.plotly.com/
Rousseeuw, P. J. (1987). Silhouettes: A graphical aid to the interpretation and validation
of cluster analysis. Journal of Computational and Applied Mathematics, 20, 53–65.
https://doi.org/10.1016/0377-0427(87)90125-7
Universitas Indonesia

153
Syakur, M., Khotimah, B., Rochman, E., & Satoto, B. (2018). Integration k-means
clustering method and elbow method for identification of the best customer profile
cluster. IOP conference series: materials science and engineering, volume 336, 012017.
https://doi.org/10.1088/1757-899X/336/1/012017
Ward Jr, J. H. (1963). Hierarchical grouping to optimize an objective function. Journal
of the American statistical association, 58(301), 236–244. https://doi.org/10.1080/
01621459.1963.10500845
Weisstein, Eric W. (2024). Kurtosis excess. https://mathworld.wolfram.com/
KurtosisExcess.html. Rumus lengkap gamma2 dan referensi historis.
Widyaningsih, Y. (2024). Perbandingan antara hasil biclustering dan hasil clusterbiplot pada data 119 kabupaten/kota di pulau jawa tahun 2022 berdasarkan indikator kesejahteraan rakyat. Universitas Indonesia, Depok. https://doi.org/10.30598/
barekengvol19iss2pp1009-1022
Universitas Indonesia

Lampiran 1 Data Penelitian
Berikut adalah data penelitian untuk masing-masing indikator dari tahun 2020 hingga
2024.
Data Indikator IPM
No Provinsi Kabupaten/Kota 2020 2021 2022 2023 2024
1 DKI JAKARTA Kep. Seribu 71,63 72,10 72,79 73,39 74,17
2 DKI JAKARTA Kota Jakarta Selatan 84,72 84,90 85,21 86,07 86,94
3 DKI JAKARTA Kota Jakarta Timur 82,66 82,97 83,45 84,13 84,65
4 DKI JAKARTA Kota Jakarta Pusat 81,39 81,56 82,11 82,53 83,01
5 DKI JAKARTA Kota Jakarta Barat 81,38 81,76 82,51 83,02 83,64
6 DKI JAKARTA Kota Jakarta Utara 80,29 80,51 80,81 81,11 81,39
7 JAWA BARAT Bogor 70,40 70,60 71,20 71,78 72,39
8 JAWA BARAT Sukabumi 66,88 67,07 67,64 68,49 68,97
9 JAWA BARAT Cianjur 65,36 65,56 65,94 66,55 67,24
10 JAWA BARAT Bandung 72,39 72,73 73,16 73,74 74,27
11 JAWA BARAT Garut 66,12 66,45 67,41 68,11 68,79
12 JAWA BARAT Tasikmalaya 65,67 65,90 66,84 67,76 68,36
13 JAWA BARAT Ciamis 70,49 70,93 71,45 72,05 72,56
14 JAWA BARAT Kuningan 69,38 69,71 70,16 70,72 71,26
15 JAWA BARAT Cirebon 68,75 69,12 70,06 70,95 71,44
16 JAWA BARAT Majalengka 67,59 67,81 68,56 69,13 69,74
17 JAWA BARAT Sumedang 71,64 71,80 72,69 73,18 73,73
18 JAWA BARAT Indramayu 67,29 67,64 68,55 69,25 69,83
19 JAWA BARAT Subang 68,95 69,13 69,87 70,70 71,36
20 JAWA BARAT Purwakarta 70,82 70,98 71,56 72,09 72,65
21 JAWA BARAT Karawang 70,66 70,94 71,74 72,35 72,90
22 JAWA BARAT Bekasi 74,07 74,45 75,22 75,76 76,44
23 JAWA BARAT Bandung Barat 68,08 68,29 69,04 69,61 70,03
24 JAWA BARAT Pangandaran 68,06 68,28 69,03 69,38 69,85
25 JAWA BARAT Kota Bogor 76,11 76,59 77,17 77,85 78,52
26 JAWA BARAT Kota Sukabumi 74,21 74,60 75,40 76,32 76,85
27 JAWA BARAT Kota Bandung 81,51 81,96 82,50 83,04 83,52
28 JAWA BARAT Kota Cirebon 74,89 75,25 75,89 76,46 77,08
29 JAWA BARAT Kota Bekasi 81,50 81,95 82,46 83,03 83,54
30 JAWA BARAT Kota Depok 80,97 81,37 81,86 82,38 82,91
31 JAWA BARAT Kota Cimahi 77,83 78,06 78,77 79,46 80,08
32 JAWA BARAT Kota Tasikmalaya 73,04 73,31 73,83 74,47 75,04
33 JAWA BARAT Kota Banjar 71,70 71,92 72,55 73,08 73,64
34 BANTEN Pandeglang 65,00 65,17 65,84 66,42 66,99
35 BANTEN Lebak 63,91 64,03 64,71 65,21 65,86
36 BANTEN Tangerang 71,92 72,29 72,97 73,43 74,09
37 BANTEN Serang 66,70 66,82 67,75 68,39 69,03
38 BANTEN Kota Tangerang 78,25 78,50 78,90 79,46 80,02
39 BANTEN Kota Cilegon 73,05 73,35 73,95 74,54 75,14
40 BANTEN Kota Serang 72,16 72,44 72,98 73,48 73,99
41 BANTEN Kota Tangerang Selatan 81,36 81,60 81,95 82,28 82,89
154 Universitas Indonesia

155
Data Indikator IPG
No Provinsi Kabupaten/Kota 2020 2021 2022 2023 2024
1 DKI JAKARTA Kep. Seribu 93,39 93,39 93,69 93,60 93,62
2 DKI JAKARTA Kota Jakarta Selatan 96,14 96,28 96,48 97,16 97,30
3 DKI JAKARTA Kota Jakarta Timur 95,48 95,44 95,92 95,98 95,90
4 DKI JAKARTA Kota Jakarta Pusat 96,00 96,11 96,41 96,28 96,24
5 DKI JAKARTA Kota Jakarta Barat 95,16 95,37 95,39 95,39 95,66
6 DKI JAKARTA Kota Jakarta Utara 93,53 93,60 93,84 94,03 94,43
7 JAWA BARAT Bogor 89,23 89,40 89,95 90,20 90,17
8 JAWA BARAT Sukabumi 87,37 87,44 87,64 88,30 88,59
9 JAWA BARAT Cianjur 84,14 84,30 84,90 85,17 85,45
10 JAWA BARAT Bandung 93,95 93,85 93,97 94,03 94,23
11 JAWA BARAT Garut 82,29 82,13 82,78 83,37 83,50
12 JAWA BARAT Tasikmalaya 86,02 86,36 86,64 87,01 87,10
13 JAWA BARAT Ciamis 86,61 86,55 87,04 87,25 87,48
14 JAWA BARAT Kuningan 86,81 86,97 87,17 87,31 87,55
15 JAWA BARAT Cirebon 83,47 83,49 84,01 84,28 84,74
16 JAWA BARAT Majalengka 85,91 86,24 86,69 87,06 87,16
17 JAWA BARAT Sumedang 95,18 95,18 95,09 95,41 95,55
18 JAWA BARAT Indramayu 88,15 87,92 88,17 88,81 89,03
19 JAWA BARAT Subang 90,53 90,51 90,95 91,29 91,55
20 JAWA BARAT Purwakarta 86,56 86,70 86,94 87,08 87,18
21 JAWA BARAT Karawang 90,12 90,29 90,50 91,07 91,23
22 JAWA BARAT Bekasi 88,58 88,88 89,37 89,53 89,72
23 JAWA BARAT Bandung Barat 79,06 79,07 79,69 80,37 81,02
24 JAWA BARAT Pangandaran 89,78 89,92 90,20 90,04 90,38
25 JAWA BARAT Kota Bogor 91,02 90,87 90,97 91,11 91,26
26 JAWA BARAT Kota Sukabumi 91,36 91,58 91,83 92,35 92,59
27 JAWA BARAT Kota Bandung 94,68 94,72 94,81 94,71 94,87
28 JAWA BARAT Kota Cirebon 94,39 94,46 94,91 95,55 95,83
29 JAWA BARAT Kota Bekasi 92,95 93,09 93,49 93,52 93,62
30 JAWA BARAT Kota Depok 93,01 92,96 93,47 93,44 93,70
31 JAWA BARAT Kota Cimahi 92,01 91,95 91,97 92,00 92,28
32 JAWA BARAT Kota Tasikmalaya 91,64 91,74 91,87 92,13 92,22
33 JAWA BARAT Kota Banjar 87,17 87,01 87,67 87,85 88,12
34 BANTEN Pandeglang 86,71 86,83 86,93 87,09 87,23
35 BANTEN Lebak 79,81 79,87 80,20 80,75 81,22
36 BANTEN Tangerang 90,98 91,08 91,24 91,67 91,69
37 BANTEN Serang 92,14 91,98 92,32 92,49 92,54
38 BANTEN Kota Tangerang 95,06 94,72 94,73 94,89 94,95
39 BANTEN Kota Cilegon 87,22 87,25 87,57 88,05 88,29
40 BANTEN Kota Serang 92,28 92,47 92,63 92,49 92,61
41 BANTEN Kota Tangerang Selatan 93,58 93,80 93,96 94,23 94,47
Universitas Indonesia

156
Data Indikator PPM
No Provinsi Kabupaten/Kota 2020 2021 2022 2023 2024
1 DKI JAKARTA Kep. Seribu 14,87 15,06 14,11 13,13 13,03
2 DKI JAKARTA Kota Jakarta Selatan 3,43 3,56 3,52 3,10 3,03
3 DKI JAKARTA Kota Jakarta Timur 4,16 4,28 4,30 4,20 4,09
4 DKI JAKARTA Kota Jakarta Pusat 4,51 4,94 4,90 4,68 4,63
5 DKI JAKARTA Kota Jakarta Barat 4,25 4,31 4,22 4,09 3,94
6 DKI JAKARTA Kota Jakarta Utara 6,78 7,24 7,24 6,78 6,44
7 JAWA BARAT Bogor 7,69 8,13 7,73 7,27 7,05
8 JAWA BARAT Sukabumi 7,09 7,70 7,34 7,01 6,87
9 JAWA BARAT Cianjur 10,36 11,18 10,55 10,22 10,14
10 JAWA BARAT Bandung 6,91 7,15 6,80 6,40 6,19
11 JAWA BARAT Garut 9,98 10,65 10,42 9,77 9,68
12 JAWA BARAT Tasikmalaya 10,34 11,15 10,73 10,28 10,23
13 JAWA BARAT Ciamis 7,62 7,97 7,72 7,42 7,39
14 JAWA BARAT Kuningan 12,82 13,10 12,76 12,12 11,88
15 JAWA BARAT Cirebon 11,24 12,30 12,01 11,20 11,00
16 JAWA BARAT Majalengka 11,43 12,33 11,94 11,21 10,82
17 JAWA BARAT Sumedang 10,26 10,71 10,14 9,36 9,10
18 JAWA BARAT Indramayu 12,70 13,04 12,77 12,13 11,93
19 JAWA BARAT Subang 9,31 10,03 9,75 9,52 9,49
20 JAWA BARAT Purwakarta 8,27 8,83 8,70 8,46 8,41
21 JAWA BARAT Karawang 8,26 8,95 8,44 7,87 7,86
22 JAWA BARAT Bekasi 4,82 5,21 5,01 4,93 4,80
23 JAWA BARAT Bandung Barat 10,49 11,30 10,82 10,52 10,49
24 JAWA BARAT Pangandaran 8,99 9,65 9,32 8,98 8,75
25 JAWA BARAT Kota Bogor 6,68 7,24 7,10 6,67 6,53
26 JAWA BARAT Kota Sukabumi 7,70 8,25 8,02 7,50 7,20
27 JAWA BARAT Kota Bandung 3,99 4,37 4,25 3,96 3,87
28 JAWA BARAT Kota Cirebon 9,52 10,03 9,82 9,16 9,02
29 JAWA BARAT Kota Bekasi 4,38 4,74 4,43 4,10 4,01
30 JAWA BARAT Kota Depok 2,45 2,58 2,53 2,38 2,34
31 JAWA BARAT Kota Cimahi 5,11 5,35 5,11 4,66 4,39
32 JAWA BARAT Kota Tasikmalaya 12,97 13,13 12,72 11,53 11,10
33 JAWA BARAT Kota Banjar 6,09 7,11 6,73 6,14 5,85
34 BANTEN Pandeglang 9,92 10,72 9,32 9,27 9,18
35 BANTEN Lebak 9,24 10,29 8,91 8,68 8,44
36 BANTEN Tangerang 6,23 7,12 6,92 6,93 6,55
37 BANTEN Serang 4,94 5,49 4,96 4,85 4,51
38 BANTEN Kota Tangerang 5,22 5,93 5,77 5,89 5,43
39 BANTEN Kota Cilegon 3,69 4,24 3,64 3,98 3,75
40 BANTEN Kota Serang 6,06 6,79 5,94 6,20 5,65
41 BANTEN Kota Tangerang Selatan 2,29 2,57 2,50 2,57 2,36
Universitas Indonesia

157
Data Indikator PDRB ADHK
No Provinsi Kabupaten/Kota 2020 2021 2022 2023 2024
1 DKI JAKARTA Kep. Seribu 3.706,25 3.648,80 3.586,89 3.305,13 3.268,50
2 DKI JAKARTA Kota Jakarta Selatan 419.257,19 429.259,25 451.772,69 475.806,20 500.422,68
3 DKI JAKARTA Kota Jakarta Timur 299.515,16 313.808,70 329.546,33 346.508,69 362.068,12
4 DKI JAKARTA Kota Jakarta Pusat 449.518,14 459.971,27 481.954,45 506.560,74 532.606,87
5 DKI JAKARTA Kota Jakarta Barat 316.172,61 327.599,16 345.619,56 363.931,23 383.113,08
6 DKI JAKARTA Kota Jakarta Utara 312.100,58 330.890,18 350.233,72 364.823,88 380.461,09
7 JAWA BARAT Bogor 154.113,60 159.589,55 167.966,18 176.683,58 185.882,34
8 JAWA BARAT Sukabumi 46.205,28 47.935,63 50.388,72 52.993,98 55.723,92
9 JAWA BARAT Cianjur 31.792,32 32.897,53 34.556,40 36.339,55 38.208,18
10 JAWA BARAT Bandung 81.060,97 83.949,37 88.437,96 92.830,17 97.509,91
11 JAWA BARAT Garut 38.598,43 39.981,19 42.012,96 44.087,22 46.270,92
12 JAWA BARAT Tasikmalaya 24.344,86 25.179,48 26.361,66 27.599,16 28.879,68
13 JAWA BARAT Ciamis 21.970,41 22.774,93 23.918,89 25.111,71 26.360,94
14 JAWA BARAT Kuningan 16.882,76 17.483,56 18.450,34 19.418,73 20.507,27
15 JAWA BARAT Cirebon 33.304,05 34.127,52 35.523,78 37.246,57 39.418,27
16 JAWA BARAT Majalengka 21.754,54 22.787,90 24.298,10 25.793,24 27.438,93
17 JAWA BARAT Sumedang 23.665,01 24.368,73 25.595,02 26.876,86 27.965,76
18 JAWA BARAT Indramayu 59.200,00 59.544,87 61.259,61 67.237,69 68.701,96
19 JAWA BARAT Subang 28.344,32 28.960,89 30.202,58 31.604,98 33.126,46
20 JAWA BARAT Purwakarta 45.293,24 46.840,15 49.292,98 51.740,50 54.142,29
21 JAWA BARAT Karawang 157.710,59 166.941,49 177.470,89 187.051,65 194.903,10
22 JAWA BARAT Bekasi 242.971,39 251.778,52 265.120,49 279.224,90 293.664,79
23 JAWA BARAT Bandung Barat 30.640,41 31.701,79 33.393,29 35.062,19 36.793,68
24 JAWA BARAT Pangandaran 7.738,97 8.022,78 8.426,40 8.869,86 9.324,31
25 JAWA BARAT Kota Bogor 32.162,74 33.372,48 35.258,87 37.055,36 38.964,67
26 JAWA BARAT Kota Sukabumi 8.534,72 8.851,05 9.324,16 9.801,88 10.303,16
27 JAWA BARAT Kota Bandung 193.144,95 200.414,03 211.250,67 221.969,13 233.053,79
28 JAWA BARAT Kota Cirebon 16.648,21 17.155,75 18.030,25 18.934,40 19.884,06
29 JAWA BARAT Kota Bekasi 67.619,24 69.796,94 73.260,65 77.241,79 81.251,61
30 JAWA BARAT Kota Depok 48.135,93 49.947,24 52.564,98 55.221,82 58.241,86
31 JAWA BARAT Kota Cimahi 22.340,56 23.275,78 24.652,73 25.931,98 27.296,93
32 JAWA BARAT Kota Tasikmalaya 15.430,02 15.980,75 16.781,04 17.781,93 18.709,87
33 JAWA BARAT Kota Banjar 3.251,76 3.371,78 3.516,57 3.682,30 3.883,98
34 BANTEN Pandeglang 19.511,68 20.116,90 20.814,77 21.340,73 21.928,59
35 BANTEN Lebak 20.622,04 21.278,48 22.087,37 22.706,47 23.377,24
36 BANTEN Tangerang 93.480,39 97.869,38 103.219,44 108.569,83 113.998,21
37 BANTEN Serang 52.866,43 54.844,65 57.605,99 60.369,96 63.255,33
38 BANTEN Kota Tangerang 102.415,68 106.413,71 112.778,10 119.059,85 125.060,32
39 BANTEN Kota Cilegon 73.319,12 77.163,94 80.646,38 84.536,65 88.628,11
40 BANTEN Kota Serang 22.518,66 23.392,75 24.495,02 25.655,25 26.863,08
41 BANTEN Kota Tangerang Selatan 59.531,08 62.393,12 66.021,24 69.561,56 73.052,84
Universitas Indonesia

158
Data Indikator RLS
No Provinsi Kabupaten/Kota 2020 2021 2022 2023 2024
1 DKI JAKARTA Kep. Seribu 8,68 8,81 9,02 9,03 9,26
2 DKI JAKARTA Kota Jakarta Selatan 11,63 11,64 11,66 11,75 11,95
3 DKI JAKARTA Kota Jakarta Timur 11,66 11,67 11,73 11,88 11,99
4 DKI JAKARTA Kota Jakarta Pusat 11,38 11,39 11,53 11,54 11,61
5 DKI JAKARTA Kota Jakarta Barat 10,63 10,78 11,13 11,23 11,24
6 DKI JAKARTA Kota Jakarta Utara 10,80 10,81 10,82 10,84 10,85
7 JAWA BARAT Bogor 8,30 8,31 8,34 8,37 8,39
8 JAWA BARAT Sukabumi 7,07 7,10 7,11 7,33 7,34
9 JAWA BARAT Cianjur 7,18 7,19 7,20 7,22 7,33
10 JAWA BARAT Bandung 8,96 9,07 9,08 9,10 9,15
11 JAWA BARAT Garut 7,52 7,53 7,83 7,84 7,85
12 JAWA BARAT Tasikmalaya 7,35 7,48 7,73 7,96 7,97
13 JAWA BARAT Ciamis 7,70 7,90 8,00 8,09 8,10
14 JAWA BARAT Kuningan 7,57 7,80 7,88 7,89 7,90
15 JAWA BARAT Cirebon 6,92 7,10 7,40 7,64 7,65
16 JAWA BARAT Majalengka 7,27 7,31 7,49 7,52 7,53
17 JAWA BARAT Sumedang 8,51 8,52 8,72 8,73 8,74
18 JAWA BARAT Indramayu 6,30 6,52 6,83 6,94 6,95
19 JAWA BARAT Subang 7,10 7,11 7,20 7,45 7,46
20 JAWA BARAT Purwakarta 8,09 8,10 8,11 8,13 8,14
21 JAWA BARAT Karawang 7,77 7,78 7,96 8,04 8,05
22 JAWA BARAT Bekasi 9,12 9,30 9,53 9,57 9,76
23 JAWA BARAT Bandung Barat 8,19 8,20 8,22 8,23 8,24
24 JAWA BARAT Pangandaran 7,74 7,85 8,03 8,04 8,10
25 JAWA BARAT Kota Bogor 10,33 10,53 10,63 10,64 10,71
26 JAWA BARAT Kota Sukabumi 9,59 9,81 10,14 10,37 10,38
27 JAWA BARAT Kota Bandung 10,75 10,99 11,00 11,06 11,07
28 JAWA BARAT Kota Cirebon 9,91 10,12 10,33 10,37 10,53
29 JAWA BARAT Kota Bekasi 11,16 11,31 11,44 11,66 11,79
30 JAWA BARAT Kota Depok 11,28 11,46 11,47 11,58 11,59
31 JAWA BARAT Kota Cimahi 10,96 11,08 11,21 11,39 11,52
32 JAWA BARAT Kota Tasikmalaya 9,33 9,52 9,53 9,54 9,63
33 JAWA BARAT Kota Banjar 8,63 8,77 8,78 8,79 8,83
34 BANTEN Pandeglang 7,10 7,11 7,13 7,15 7,16
35 BANTEN Lebak 6,40 6,41 6,59 6,60 6,61
36 BANTEN Tangerang 8,39 8,61 8,92 8,93 9,06
37 BANTEN Serang 7,50 7,51 7,78 7,79 7,88
38 BANTEN Kota Tangerang 10,69 10,83 10,84 10,91 11,14
39 BANTEN Kota Cilegon 9,87 10,08 10,34 10,38 10,39
40 BANTEN Kota Serang 8,76 8,89 8,90 8,91 8,93
41 BANTEN Kota Tangerang Selatan 11,81 11,82 11,84 11,85 11,86
Universitas Indonesia

159
Data Indikator AHH L
No Provinsi Kabupaten/Kota 2020 2021 2022 2023 2024
1 DKI JAKARTA Kep. Seribu 66,67 66,84 67,02 67,68 67,85
2 DKI JAKARTA Kota Jakarta Selatan 72,23 72,33 72,60 73,03 73,38
3 DKI JAKARTA Kota Jakarta Timur 72,60 72,70 72,98 73,38 73,73
4 DKI JAKARTA Kota Jakarta Pusat 72,22 72,30 72,55 73,06 73,65
5 DKI JAKARTA Kota Jakarta Barat 71,72 71,82 71,97 72,39 72,55
6 DKI JAKARTA Kota Jakarta Utara 71,36 71,44 71,58 71,80 71,97
7 JAWA BARAT Bogor 69,24 69,52 69,72 69,97 70,20
8 JAWA BARAT Sukabumi 69,03 69,27 69,53 69,89 70,11
9 JAWA BARAT Cianjur 68,17 68,45 68,62 68,87 69,03
10 JAWA BARAT Bandung 71,67 71,94 72,14 72,50 72,73
11 JAWA BARAT Garut 69,49 69,74 69,74 70,12 70,32
12 JAWA BARAT Tasikmalaya 67,51 67,78 67,97 68,27 68,47
13 JAWA BARAT Ciamis 69,92 70,19 70,39 70,61 70,90
14 JAWA BARAT Kuningan 71,72 71,98 72,18 72,52 72,72
15 JAWA BARAT Cirebon 70,08 70,34 70,54 70,79 71,03
16 JAWA BARAT Majalengka 68,31 68,58 68,79 69,12 69,37
17 JAWA BARAT Sumedang 70,52 70,77 70,95 71,13 71,29
18 JAWA BARAT Indramayu 69,72 70,01 70,24 70,50 70,80
19 JAWA BARAT Subang 70,46 70,78 71,03 71,16 71,34
20 JAWA BARAT Purwakarta 69,05 69,31 69,50 69,79 70,04
21 JAWA BARAT Karawang 70,25 70,50 70,69 70,93 71,10
22 JAWA BARAT Bekasi 71,81 71,99 72,17 72,53 72,75
23 JAWA BARAT Bandung Barat 70,43 70,70 70,88 71,07 71,23
24 JAWA BARAT Pangandaran 69,49 69,75 69,94 70,22 70,45
25 JAWA BARAT Kota Bogor 71,74 72,02 72,25 72,69 72,89
26 JAWA BARAT Kota Sukabumi 70,52 70,75 70,92 71,08 71,23
27 JAWA BARAT Kota Bandung 72,45 72,71 72,92 73,30 73,50
28 JAWA BARAT Kota Cirebon 70,35 70,61 70,82 71,06 71,22
29 JAWA BARAT Kota Bekasi 73,22 73,48 73,70 74,10 74,29
30 JAWA BARAT Kota Depok 72,61 72,87 73,08 73,49 73,68
31 JAWA BARAT Kota Cimahi 72,19 72,45 72,66 73,05 73,23
32 JAWA BARAT Kota Tasikmalaya 70,25 70,52 70,73 70,95 71,11
33 JAWA BARAT Kota Banjar 69,04 69,31 69,53 69,85 70,08
34 BANTEN Pandeglang 62,69 62,82 63,17 63,70 63,90
35 BANTEN Lebak 65,19 65,33 65,69 66,19 66,40
36 BANTEN Tangerang 67,86 67,93 68,25 68,73 68,95
37 BANTEN Serang 62,68 62,82 63,20 63,72 64,00
38 BANTEN Kota Tangerang 69,61 69,61 69,86 70,28 70,50
39 BANTEN Kota Cilegon 64,66 64,68 64,95 65,46 65,72
40 BANTEN Kota Serang 65,87 65,91 66,19 67,05 67,33
41 BANTEN Kota Tangerang Selatan 70,51 70,55 70,85 71,07 71,21
Universitas Indonesia

160
Data Indikator AHH P
No Provinsi Kabupaten/Kota 2020 2021 2022 2023 2024
1 DKI JAKARTA Kep. Seribu 70,74 71,03 71,27 71,63 71,80
2 DKI JAKARTA Kota Jakarta Selatan 76,04 76,12 76,51 76,69 76,86
3 DKI JAKARTA Kota Jakarta Timur 76,22 76,27 76,63 77,01 77,11
4 DKI JAKARTA Kota Jakarta Pusat 75,92 76,00 76,39 76,71 76,94
5 DKI JAKARTA Kota Jakarta Barat 75,43 75,49 75,65 76,09 76,28
6 DKI JAKARTA Kota Jakarta Utara 75,13 75,22 75,41 75,55 76,04
7 JAWA BARAT Bogor 73,06 73,25 73,54 73,96 74,21
8 JAWA BARAT Sukabumi 72,87 73,13 73,47 73,88 74,11
9 JAWA BARAT Cianjur 72,03 72,23 72,50 72,81 72,98
10 JAWA BARAT Bandung 75,34 75,52 75,81 76,20 76,44
11 JAWA BARAT Garut 73,30 73,48 73,74 74,12 74,33
12 JAWA BARAT Tasikmalaya 71,39 71,59 71,88 72,21 72,41
13 JAWA BARAT Ciamis 73,72 73,91 74,20 74,63 74,94
14 JAWA BARAT Kuningan 75,42 75,61 75,87 76,22 76,42
15 JAWA BARAT Cirebon 73,87 74,07 74,37 74,83 75,15
16 JAWA BARAT Majalengka 72,18 72,37 72,68 73,08 73,34
17 JAWA BARAT Sumedang 74,30 74,49 74,78 75,35 75,70
18 JAWA BARAT Indramayu 73,52 73,72 74,04 74,52 74,83
19 JAWA BARAT Subang 74,22 74,44 74,79 75,43 75,80
20 JAWA BARAT Purwakarta 72,89 73,08 73,37 73,78 74,04
21 JAWA BARAT Karawang 74,03 74,21 74,50 74,97 75,29
22 JAWA BARAT Bekasi 75,50 75,63 75,87 76,23 76,55
23 JAWA BARAT Bandung Barat 74,20 74,39 74,67 75,23 75,56
24 JAWA BARAT Pangandaran 73,29 73,48 73,77 74,22 74,46
25 JAWA BARAT Kota Bogor 75,44 75,65 75,96 76,37 76,69
26 JAWA BARAT Kota Sukabumi 74,28 74,44 74,72 75,25 75,57
27 JAWA BARAT Kota Bandung 75,97 76,16 76,46 76,93 77,24
28 JAWA BARAT Kota Cirebon 74,13 74,31 74,61 75,20 75,55
29 JAWA BARAT Kota Bekasi 76,71 76,89 77,18 77,45 77,75
30 JAWA BARAT Kota Depok 76,42 76,60 76,91 77,11 77,41
31 JAWA BARAT Kota Cimahi 75,77 75,96 76,26 76,70 77,01
32 JAWA BARAT Kota Tasikmalaya 74,02 74,21 74,51 74,99 75,32
33 JAWA BARAT Kota Banjar 72,89 73,08 73,39 73,84 74,08
34 BANTEN Pandeglang 66,54 66,68 67,10 67,56 67,78
35 BANTEN Lebak 69,13 69,24 69,66 70,17 70,38
36 BANTEN Tangerang 71,81 71,84 72,19 72,67 72,90
37 BANTEN Serang 66,50 66,61 67,03 67,58 67,89
38 BANTEN Kota Tangerang 73,49 73,49 73,80 74,29 74,52
39 BANTEN Kota Cilegon 68,58 68,61 68,95 69,42 69,70
40 BANTEN Kota Serang 69,94 70,06 70,48 71,02 71,29
41 BANTEN Kota Tangerang Selatan 74,31 74,31 74,62 75,24 75,52
Universitas Indonesia

161
Data Indikator GK
No Provinsi Kabupaten/Kota 2020 2021 2022 2023 2024
1 DKI JAKARTA Kep. Seribu 661.852,00 670.785,00 701.088,00 743.820,00 775.909,00
2 DKI JAKARTA Kota Jakarta Selatan 782.731,00 792.684,00 842.630,00 891.580,00 928.278,00
3 DKI JAKARTA Kota Jakarta Timur 581.954,00 594.849,00 630.842,00 669.713,00 697.472,00
4 DKI JAKARTA Kota Jakarta Pusat 672.328,00 689.379,00 733.176,00 784.679,00 818.722,00
5 DKI JAKARTA Kota Jakarta Barat 557.803,00 570.330,00 599.514,00 657.247,00 685.071,00
6 DKI JAKARTA Kota Jakarta Utara 594.240,00 609.310,00 645.431,00 684.557,00 712.835,00
7 JAWA BARAT Bogor 402.877,00 418.483,00 443.787,00 482.174,00 513.512,00
8 JAWA BARAT Sukabumi 328.284,00 342.094,00 357.636,00 392.705,00 416.751,00
9 JAWA BARAT Cianjur 371.699,00 387.631,00 406.829,00 437.327,00 466.509,00
10 JAWA BARAT Bandung 367.403,00 378.819,00 398.884,00 426.069,00 455.636,00
11 JAWA BARAT Garut 310.437,00 320.050,00 335.134,00 367.681,00 393.464,00
12 JAWA BARAT Tasikmalaya 323.880,00 333.909,00 347.917,00 378.958,00 400.148,00
13 JAWA BARAT Ciamis 378.108,00 389.676,00 405.294,00 442.108,00 467.497,00
14 JAWA BARAT Kuningan 352.358,00 358.069,00 371.665,00 402.767,00 420.867,00
15 JAWA BARAT Cirebon 393.452,00 404.635,00 416.914,00 451.853,00 475.046,00
16 JAWA BARAT Majalengka 453.201,00 466.813,00 483.319,00 520.713,00 547.912,00
17 JAWA BARAT Sumedang 353.092,00 360.054,00 371.870,00 396.573,00 422.714,00
18 JAWA BARAT Indramayu 474.807,00 481.754,00 499.805,00 532.545,00 560.159,00
19 JAWA BARAT Subang 351.885,00 360.691,00 372.308,00 404.320,00 434.161,00
20 JAWA BARAT Purwakarta 378.287,00 387.754,00 403.663,00 434.187,00 467.948,00
21 JAWA BARAT Karawang 466.152,00 494.201,00 521.158,00 555.889,00 597.345,00
22 JAWA BARAT Bekasi 518.667,00 549.875,00 579.221,00 632.811,00 674.924,00
23 JAWA BARAT Bandung Barat 362.212,00 374.470,00 393.956,00 426.355,00 455.325,00
24 JAWA BARAT Pangandaran 382.264,00 394.101,00 411.174,00 447.242,00 470.068,00
25 JAWA BARAT Kota Bogor 547.399,00 571.425,00 608.949,00 661.384,00 699.861,00
26 JAWA BARAT Kota Sukabumi 547.415,00 567.734,00 594.118,00 644.276,00 678.258,00
27 JAWA BARAT Kota Bandung 500.452,00 515.396,00 545.675,00 591.124,00 614.707,00
28 JAWA BARAT Kota Cirebon 457.954,00 467.248,00 485.613,00 520.579,00 545.592,00
29 JAWA BARAT Kota Bekasi 657.953,00 692.885,00 731.392,00 795.965,00 842.927,00
30 JAWA BARAT Kota Depok 688.194,00 705.084,00 744.771,00 804.984,00 843.893,00
31 JAWA BARAT Kota Cimahi 511.375,00 522.281,00 549.450,00 590.197,00 612.567,00
32 JAWA BARAT Kota Tasikmalaya 470.150,00 480.341,00 498.711,00 537.497,00 565.377,00
33 JAWA BARAT Kota Banjar 344.363,00 357.210,00 373.510,00 408.710,00 430.590,00
34 BANTEN Pandeglang 361.713,00 390.754,00 416.248,00 455.297,00 479.886,00
35 BANTEN Lebak 334.509,00 358.479,00 380.681,00 419.066,00 440.705,00
36 BANTEN Tangerang 520.742,00 545.065,00 576.122,00 604.362,00 632.941,00
37 BANTEN Serang 341.074,00 362.102,00 385.864,00 425.751,00 449.291,00
38 BANTEN Kota Tangerang 632.835,00 655.061,00 695.225,00 746.468,00 785.113,00
39 BANTEN Kota Cilegon 504.571,00 540.032,00 572.780,00 632.703,00 663.533,00
40 BANTEN Kota Serang 411.595,00 431.210,00 460.613,00 511.575,00 540.786,00
41 BANTEN Kota Tangerang Selatan 643.083,00 665.610,00 712.717,00 782.543,00 828.278,00
Universitas Indonesia

162
Data Indikator PKKP
No Provinsi Kabupaten/Kota 2020 2021 2022 2023 2024
1 DKI JAKARTA Kep. Seribu 2,86 2,41 3,91 3,63 4,39
2 DKI JAKARTA Kota Jakarta Selatan 1,50 1,63 2,51 2,28 2,45
3 DKI JAKARTA Kota Jakarta Timur 2,18 2,24 3,71 3,08 4,38
4 DKI JAKARTA Kota Jakarta Pusat 1,97 3,08 3,82 2,91 3,63
5 DKI JAKARTA Kota Jakarta Barat 1,78 2,00 2,80 2,03 3,55
6 DKI JAKARTA Kota Jakarta Utara 2,18 2,63 4,78 2,62 3,39
7 JAWA BARAT Bogor 4,06 4,73 7,03 5,30 5,68
8 JAWA BARAT Sukabumi 4,89 4,81 7,42 5,98 5,82
9 JAWA BARAT Cianjur 4,75 6,20 9,54 7,02 5,37
10 JAWA BARAT Bandung 5,07 5,53 8,30 5,19 5,80
11 JAWA BARAT Garut 4,06 5,00 7,74 5,76 5,02
12 JAWA BARAT Tasikmalaya 5,64 7,12 10,45 8,65 5,95
13 JAWA BARAT Ciamis 5,11 5,87 6,74 6,10 5,78
14 JAWA BARAT Kuningan 3,51 3,73 6,63 6,43 5,83
15 JAWA BARAT Cirebon 5,41 6,19 9,04 7,28 5,65
16 JAWA BARAT Majalengka 4,17 4,94 9,88 7,90 5,95
17 JAWA BARAT Sumedang 4,01 5,09 7,43 5,55 5,43
18 JAWA BARAT Indramayu 4,28 3,97 5,99 5,40 5,46
19 JAWA BARAT Subang 4,81 4,16 7,06 6,23 6,19
20 JAWA BARAT Purwakarta 4,08 3,45 6,72 5,07 5,84
21 JAWA BARAT Karawang 3,79 4,80 6,73 5,93 5,10
22 JAWA BARAT Bekasi 2,26 2,38 4,30 4,36 4,60
23 JAWA BARAT Bandung Barat 5,07 6,04 7,81 7,48 6,25
24 JAWA BARAT Pangandaran 4,39 4,32 5,77 6,04 5,22
25 JAWA BARAT Kota Bogor 2,95 3,08 5,35 4,86 7,98
26 JAWA BARAT Kota Sukabumi 3,65 4,90 8,46 5,23 5,57
27 JAWA BARAT Kota Bandung 2,82 2,92 4,55 3,64 5,48
28 JAWA BARAT Kota Cirebon 4,52 5,71 7,97 6,85 6,56
29 JAWA BARAT Kota Bekasi 1,33 1,79 3,33 2,27 4,88
30 JAWA BARAT Kota Depok 1,26 1,55 2,15 1,99 4,52
31 JAWA BARAT Kota Cimahi 2,66 3,23 4,63 2,98 5,43
32 JAWA BARAT Kota Tasikmalaya 4,17 5,35 7,92 5,14 5,31
33 JAWA BARAT Kota Banjar 3,93 5,93 6,40 5,55 5,13
34 BANTEN Pandeglang 3,60 4,37 4,39 4,69 2,70
35 BANTEN Lebak 3,14 3,87 3,45 4,82 2,52
36 BANTEN Tangerang 1,84 2,65 2,28 2,59 2,22
37 BANTEN Serang 2,44 3,03 2,90 3,23 2,13
38 BANTEN Kota Tangerang 1,42 2,16 1,71 1,58 1,89
39 BANTEN Kota Cilegon 2,02 2,46 2,04 2,19 1,96
40 BANTEN Kota Serang 1,29 1,99 2,42 2,30 2,18
41 BANTEN Kota Tangerang Selatan 1,51 1,89 1,50 1,92 3,13
Universitas Indonesia

163
Data Indikator TPT
No Provinsi Kabupaten/Kota 2020 2021 2022 2023 2024
1 DKI JAKARTA Kep. Seribu 7,37 8,58 8,47 8,11 7,93
2 DKI JAKARTA Kota Jakarta Selatan 10,79 7,33 5,63 5,37 5,22
3 DKI JAKARTA Kota Jakarta Timur 9,29 8,23 8,39 7,24 6,95
4 DKI JAKARTA Kota Jakarta Pusat 10,97 7,75 5,88 6,42 6,24
5 DKI JAKARTA Kota Jakarta Barat 12,27 9,06 7,10 6,39 6,18
6 DKI JAKARTA Kota Jakarta Utara 11,79 9,84 8,04 7,05 6,18
7 JAWA BARAT Bogor 14,29 12,22 10,64 8,47 7,34
8 JAWA BARAT Sukabumi 9,60 9,51 7,77 7,32 7,11
9 JAWA BARAT Cianjur 11,05 9,32 8,41 7,71 5,99
10 JAWA BARAT Bandung 8,58 8,32 6,98 6,52 6,36
11 JAWA BARAT Garut 8,95 8,68 7,60 7,33 6,96
12 JAWA BARAT Tasikmalaya 7,12 6,16 4,17 3,89 3,74
13 JAWA BARAT Ciamis 5,66 5,06 3,75 3,52 3,37
14 JAWA BARAT Kuningan 11,22 11,68 9,81 9,49 7,78
15 JAWA BARAT Cirebon 11,52 10,38 8,11 7,65 6,74
16 JAWA BARAT Majalengka 5,84 5,71 4,16 4,12 4,01
17 JAWA BARAT Sumedang 9,89 9,18 7,72 6,94 6,16
18 JAWA BARAT Indramayu 9,21 8,30 6,49 6,46 6,25
19 JAWA BARAT Subang 9,48 9,77 7,77 7,65 6,73
20 JAWA BARAT Purwakarta 11,07 10,70 8,75 7,72 7,34
21 JAWA BARAT Karawang 11,52 11,83 9,87 8,95 8,04
22 JAWA BARAT Bekasi 11,54 10,09 10,31 8,87 8,82
23 JAWA BARAT Bandung Barat 12,25 11,65 9,63 8,11 6,70
24 JAWA BARAT Pangandaran 5,08 3,25 1,56 1,52 1,58
25 JAWA BARAT Kota Bogor 12,68 11,79 10,78 9,39 8,13
26 JAWA BARAT Kota Sukabumi 12,17 10,78 8,83 8,53 8,34
27 JAWA BARAT Kota Bandung 11,19 11,46 9,55 8,83 7,40
28 JAWA BARAT Kota Cirebon 10,97 10,53 8,42 7,66 6,29
29 JAWA BARAT Kota Bekasi 10,68 10,88 8,81 7,90 7,82
30 JAWA BARAT Kota Depok 9,87 9,76 7,82 6,97 6,27
31 JAWA BARAT Kota Cimahi 13,30 13,07 10,77 10,52 8,97
32 JAWA BARAT Kota Tasikmalaya 7,99 7,66 6,62 6,55 6,49
33 JAWA BARAT Kota Banjar 6,73 6,09 5,53 5,43 5,44
34 BANTEN Pandeglang 9,15 7,70 9,24 9,05 8,09
35 BANTEN Lebak 9,63 7,86 8,55 7,57 6,23
36 BANTEN Tangerang 13,06 9,06 7,88 6,94 6,06
37 BANTEN Serang 12,22 10,58 10,61 9,94 9,18
38 BANTEN Kota Tangerang 8,63 9,07 7,16 6,76 5,92
39 BANTEN Kota Cilegon 12,69 10,13 8,10 7,25 6,08
40 BANTEN Kota Serang 9,26 9,41 8,17 7,45 7,12
41 BANTEN Kota Tangerang Selatan 8,48 8,60 6,59 5,81 5,09
Universitas Indonesia

164
Data Indikator TPAK
No Provinsi Kabupaten/Kota 2020 2021 2022 2023 2024
1 DKI JAKARTA Kep. Seribu 60,67 65,45 67,72 73,72 72,66
2 DKI JAKARTA Kota Jakarta Selatan 63,98 61,64 63,28 66,05 64,84
3 DKI JAKARTA Kota Jakarta Timur 62,12 61,03 61,47 65,14 64,55
4 DKI JAKARTA Kota Jakarta Pusat 63,98 63,40 63,82 63,83 64,97
5 DKI JAKARTA Kota Jakarta Barat 64,40 63,29 64,12 64,77 64,10
6 DKI JAKARTA Kota Jakarta Utara 66,33 66,16 64,81 65,43 67,72
7 JAWA BARAT Bogor 62,19 62,08 63,20 63,65 65,78
8 JAWA BARAT Sukabumi 61,33 64,70 68,92 67,51 69,56
9 JAWA BARAT Cianjur 68,69 67,44 69,48 71,88 72,23
10 JAWA BARAT Bandung 61,98 64,91 63,44 66,81 66,98
11 JAWA BARAT Garut 61,05 62,72 68,81 69,75 71,01
12 JAWA BARAT Tasikmalaya 69,83 68,50 68,14 68,22 68,62
13 JAWA BARAT Ciamis 71,75 70,21 68,86 66,35 68,22
14 JAWA BARAT Kuningan 62,05 63,90 61,88 61,78 65,71
15 JAWA BARAT Cirebon 63,42 64,15 65,16 66,01 69,16
16 JAWA BARAT Majalengka 69,81 67,77 66,38 68,53 69,67
17 JAWA BARAT Sumedang 69,13 68,66 64,84 67,69 69,19
18 JAWA BARAT Indramayu 69,04 69,41 68,64 63,90 66,46
19 JAWA BARAT Subang 66,62 67,70 68,80 70,11 71,80
20 JAWA BARAT Purwakarta 60,55 60,84 64,92 66,15 66,41
21 JAWA BARAT Karawang 64,38 63,71 65,02 63,23 63,52
22 JAWA BARAT Bekasi 63,88 65,60 65,06 64,78 64,93
23 JAWA BARAT Bandung Barat 59,60 60,48 64,09 66,63 67,98
24 JAWA BARAT Pangandaran 76,97 74,92 80,03 80,22 80,21
25 JAWA BARAT Kota Bogor 60,61 61,58 64,03 64,64 67,20
26 JAWA BARAT Kota Sukabumi 55,57 56,72 62,33 62,59 62,61
27 JAWA BARAT Kota Bandung 64,61 65,22 69,33 66,97 67,95
28 JAWA BARAT Kota Cirebon 63,87 62,97 65,54 68,75 68,59
29 JAWA BARAT Kota Bekasi 64,67 64,69 65,25 64,64 65,59
30 JAWA BARAT Kota Depok 63,91 62,58 63,31 62,72 62,87
31 JAWA BARAT Kota Cimahi 62,60 61,74 67,16 68,36 66,45
32 JAWA BARAT Kota Tasikmalaya 66,53 65,75 65,98 65,28 68,78
33 JAWA BARAT Kota Banjar 67,61 64,34 64,13 67,47 67,25
34 BANTEN Pandeglang 59,91 61,98 61,32 59,85 66,47
35 BANTEN Lebak 69,58 66,39 65,81 67,40 72,40
36 BANTEN Tangerang 65,02 63,54 66,69 66,52 65,53
37 BANTEN Serang 63,29 63,73 64,14 62,83 67,45
38 BANTEN Kota Tangerang 64,62 64,21 65,73 63,31 65,96
39 BANTEN Kota Cilegon 61,38 60,98 66,99 62,49 60,66
40 BANTEN Kota Serang 63,33 61,75 64,45 68,11 68,80
41 BANTEN Kota Tangerang Selatan 62,42 62,51 58,62 60,66 58,30
Universitas Indonesia

165
Data Indikator GR
No Provinsi Kabupaten/Kota 2020 2021 2022 2023 2024
1 DKI JAKARTA Kep. Seribu 0,31 0,24 0,21 0,43 0,43
2 DKI JAKARTA Kota Jakarta Selatan 0,43 0,38 0,38 0,43 0,43
3 DKI JAKARTA Kota Jakarta Timur 0,36 0,37 0,37 0,43 0,43
4 DKI JAKARTA Kota Jakarta Pusat 0,39 0,40 0,34 0,43 0,43
5 DKI JAKARTA Kota Jakarta Barat 0,38 0,37 0,42 0,43 0,43
6 DKI JAKARTA Kota Jakarta Utara 0,41 0,52 0,54 0,43 0,43
7 JAWA BARAT Bogor 0,38 0,40 0,40 0,37 0,40
8 JAWA BARAT Sukabumi 0,33 0,34 0,31 0,35 0,34
9 JAWA BARAT Cianjur 0,35 0,37 0,31 0,32 0,34
10 JAWA BARAT Bandung 0,41 0,39 0,37 0,35 0,36
11 JAWA BARAT Garut 0,35 0,32 0,30 0,36 0,32
12 JAWA BARAT Tasikmalaya 0,34 0,38 0,38 0,40 0,36
13 JAWA BARAT Ciamis 0,34 0,34 0,34 0,34 0,33
14 JAWA BARAT Kuningan 0,36 0,35 0,34 0,35 0,37
15 JAWA BARAT Cirebon 0,34 0,34 0,35 0,36 0,37
16 JAWA BARAT Majalengka 0,34 0,35 0,37 0,34 0,35
17 JAWA BARAT Sumedang 0,38 0,36 0,39 0,37 0,34
18 JAWA BARAT Indramayu 0,32 0,32 0,30 0,34 0,31
19 JAWA BARAT Subang 0,35 0,33 0,36 0,33 0,32
20 JAWA BARAT Purwakarta 0,41 0,39 0,37 0,37 0,40
21 JAWA BARAT Karawang 0,33 0,37 0,35 0,39 0,38
22 JAWA BARAT Bekasi 0,37 0,34 0,37 0,40 0,37
23 JAWA BARAT Bandung Barat 0,38 0,36 0,37 0,37 0,40
24 JAWA BARAT Pangandaran 0,34 0,32 0,29 0,32 0,29
25 JAWA BARAT Kota Bogor 0,41 0,44 0,43 0,46 0,48
26 JAWA BARAT Kota Sukabumi 0,40 0,44 0,48 0,46 0,44
27 JAWA BARAT Kota Bandung 0,40 0,43 0,46 0,46 0,44
28 JAWA BARAT Kota Cirebon 0,42 0,49 0,47 0,47 0,42
29 JAWA BARAT Kota Bekasi 0,35 0,40 0,44 0,41 0,43
30 JAWA BARAT Kota Depok 0,34 0,37 0,39 0,40 0,40
31 JAWA BARAT Kota Cimahi 0,38 0,41 0,41 0,39 0,40
32 JAWA BARAT Kota Tasikmalaya 0,37 0,41 0,41 0,43 0,36
33 JAWA BARAT Kota Banjar 0,31 0,34 0,36 0,34 0,32
34 BANTEN Pandeglang 0,30 0,27 0,30 0,24 0,24
35 BANTEN Lebak 0,31 0,33 0,31 0,27 0,26
36 BANTEN Tangerang 0,28 0,29 0,28 0,29 0,27
37 BANTEN Serang 0,30 0,26 0,28 0,27 0,27
38 BANTEN Kota Tangerang 0,34 0,34 0,38 0,38 0,33
39 BANTEN Kota Cilegon 0,37 0,36 0,32 0,36 0,36
40 BANTEN Kota Serang 0,35 0,34 0,35 0,31 0,32
41 BANTEN Kota Tangerang Selatan 0,37 0,37 0,36 0,36 0,40
Universitas Indonesia

166
Data Indikator JP
No Provinsi Kabupaten/Kota 2020 2021 2022 2023 2024
1 DKI JAKARTA Kep. Seribu 27.749,00 27.996,00 28.262,00 28.523,00 30.420,00
2 DKI JAKARTA Kota Jakarta Selatan 2.226.812,00 2.232.442,00 2.234.262,00 2.235.606,00 2.410.000,00
3 DKI JAKARTA Kota Jakarta Timur 3.037.139,00 3.051.866,00 3.066.074,00 3.079.618,00 3.310.000,00
4 DKI JAKARTA Kota Jakarta Pusat 1.056.896,00 1.057.465,00 1.053.482,00 1.049.314,00 1.100.000,00
5 DKI JAKARTA Kota Jakarta Barat 2.434.511,00 2.446.687,00 2.458.707,00 2.470.054,00 2.610.000,00
6 DKI JAKARTA Kota Jakarta Utara 1.778.981,00 1.788.981,00 1.799.220,00 1.808.985,00 1.870.000,00
7 JAWA BARAT Bogor 5.409.370,00 5.484.150,00 5.556.310,00 5.627.020,00 5.682.300,00
8 JAWA BARAT Sukabumi 2.718.740,00 2.747.450,00 2.775.310,00 2.802.400,00 2.828.020,00
9 JAWA BARAT Cianjur 2.470.560,00 2.500.640,00 2.529.810,00 2.558.140,00 2.584.990,00
10 JAWA BARAT Bandung 3.616.540,00 3.652.400,00 3.687.250,00 3.721.110,00 3.753.120,00
11 JAWA BARAT Garut 2.577.350,00 2.613.530,00 2.648.950,00 2.683.670,00 2.716.950,00
12 JAWA BARAT Tasikmalaya 1.861.060,00 1.876.890,00 1.892.220,00 1.907.050,00 1.920.920,00
13 JAWA BARAT Ciamis 1.226.050,00 1.234.830,00 1.243.320,00 1.251.540,00 1.259.230,00
14 JAWA BARAT Kuningan 1.162.590,00 1.175.950,00 1.189.010,00 1.201.760,00 1.213.930,00
15 JAWA BARAT Cirebon 2.270.350,00 2.301.330,00 2.331.360,00 2.360.440,00 2.387.960,00
16 JAWA BARAT Majalengka 1.301.570,00 1.315.010,00 1.328.010,00 1.340.620,00 1.352.540,00
17 JAWA BARAT Sumedang 1.149.470,00 1.159.260,00 1.168.840,00 1.178.240,00 1.187.130,00
18 JAWA BARAT Indramayu 1.829.230,00 1.851.730,00 1.873.400,00 1.894.330,00 1.914.040,00
19 JAWA BARAT Subang 1.605.250,00 1.620.700,00 1.635.560,00 1.649.820,00 1.663.160,00
20 JAWA BARAT Purwakarta 994.330,00 1.008.930,00 1.023.180,00 1.037.070,00 1.050.340,00
21 JAWA BARAT Karawang 2.434.090,00 2.465.570,00 2.496.190,00 2.526.000,00 2.554.380,00
22 JAWA BARAT Bekasi 3.101.730,00 3.148.740,00 3.193.840,00 3.237.420,00 3.273.870,00
23 JAWA BARAT Bandung Barat 1.782.190,00 1.808.420,00 1.834.230,00 1.859.640,00 1.884.190,00
24 JAWA BARAT Pangandaran 422.410,00 425.590,00 428.600,00 431.460,00 434.100,00
25 JAWA BARAT Kota Bogor 1.040.570,00 1.050.920,00 1.060.940,00 1.070.720,00 1.078.350,00
26 JAWA BARAT Kota Sukabumi 344.830,00 350.150,00 355.420,00 360.640,00 365.740,00
27 JAWA BARAT Kota Bandung 2.438.350,00 2.461.410,00 2.484.150,00 2.506.600,00 2.528.160,00
28 JAWA BARAT Kota Cirebon 332.580,00 335.810,00 338.940,00 341.980,00 344.850,00
29 JAWA BARAT Kota Bekasi 2.536.120,00 2.568.020,00 2.598.070,00 2.627.210,00 2.644.060,00
30 JAWA BARAT Kota Depok 2.046.830,00 2.081.130,00 2.113.620,00 2.145.400,00 2.163.640,00
31 JAWA BARAT Kota Cimahi 566.170,00 574.450,00 582.650,00 590.780,00 598.700,00
32 JAWA BARAT Kota Tasikmalaya 713.670,00 723.100,00 732.480,00 741.760,00 750.730,00
33 JAWA BARAT Kota Banjar 200.270,00 202.720,00 205.140,00 207.510,00 209.790,00
34 BANTEN Pandeglang 1.270,09 1.284,64 1.298,85 1.312,77 1.325,95
35 BANTEN Lebak 1.383,54 1.400,63 1.417,35 1.433,70 1.449,21
36 BANTEN Tangerang 3.236,70 3.279,71 3.321,65 3.362,61 3.400,49
37 BANTEN Serang 1.618,36 1.640,11 1.661,37 1.682,13 1.701,80
38 BANTEN Kota Tangerang 1.891,40 1.912,03 1.931,64 1.950,58 1.963,97
39 BANTEN Kota Cilegon 433,76 439,48 445,06 450,51 455,62
40 BANTEN Kota Serang 689,41 700,95 712,41 723,79 734,87
41 BANTEN Kota Tangerang Selatan 1.351,31 1.365,39 1.378,71 1.391,65 1.399,50
Universitas Indonesia

Lampiran 2 Kode Program Aplikasi Dashboard
Kode program aplikasi dashboard yang digunakan dalam penelitian ini dapat diakses
secara daring melalui tautan berikut:
https://huggingface.co/spaces/ammarhnfi/skripsi/blob/main/app.py
Berikut adalah tampilan antarmuka aplikasi dashboard:
Gambar 1: Tampilan Antarmuka Aplikasi Dashboard Analisis Spasio-Temporal
Sosio-Ekonomi
167 Universitas Indonesia
